% !TeX program = xelatex
% ============================================================================
%  《数学》 For OI  ——  面向信息学竞赛的数学教材（教师版）
%
%  全书按“部分（Part）— 章（Chapter）”组织：
%    第一部分  基础数学
%    第二部分  数论
%    第三部分  组合数学
%    第四部分  博弈论与概率论
%    第五部分  多项式与生成函数
%    附录 A--D
%
%  每个部分开头是一张“部分页”：页面右下角是由 TikZ 绘制的旋转叠加图案，
%  左下角是罗马数字与部分名称、部分标题。全部图案均为矢量图，无外部图片依赖。
%
%  编译方式：XeLaTeX 编译两遍（第一遍生成目录，第二遍稳定交叉引用）。
% ============================================================================
\documentclass[a4paper,twoside,openright,11pt,fontset=fandol]{ctexbook}
\usepackage{amsmath,amssymb,amsthm,mathtools}
\usepackage[twoside,top=2.3cm,bottom=2.3cm,inner=2.75cm,outer=1.85cm,headheight=16pt]{geometry}
\usepackage{enumitem,booktabs,longtable,array,multicol}
\usepackage[dvipsnames]{xcolor}
\usepackage[most]{tcolorbox}
\usepackage{fancyhdr,listings,tikz}
\usepackage{emptypage}
\usetikzlibrary{arrows.meta,positioning,trees}
\usepackage{hyperref}
\hypersetup{colorlinks=true,linkcolor=MidnightBlue,urlcolor=MidnightBlue,bookmarksnumbered=true,
  pdfauthor={Federico Prask},pdftitle={数学 For OI},pdfsubject={信息学竞赛数学教材}}

% ------------------------------------------------------------------ 配色 ----
\definecolor{ink}{HTML}{173A51}      % 主色：深蓝黑
\definecolor{accent}{HTML}{187C82}   % 辅色：青绿
\definecolor{soft}{HTML}{F1F6F8}     % 浅底
\definecolor{warn}{HTML}{AA8055}     % 警示：棕
\definecolor{tipbase}{HTML}{079C81}  % 提示：绿
\definecolor{chapterguide}{HTML}{DB8080} % 章节导读

\ctexset{chapter={format=\raggedright,
  nameformat=\Large\bfseries\color{ink},
  titleformat=\Huge\bfseries\color{ink},
  name={第,章},aftername=\par,number=\arabic{chapter}},
  section={format=\Large\bfseries\color{ink}},
  subsection={format=\large\bfseries\color{accent}}}
\setlength{\parskip}{0.25em}
\linespread{1.16}
\setlist{nosep,leftmargin=2em,topsep=0.25em}
\allowdisplaybreaks[2]
\emergencystretch=2em

% ------------------------------------------------------------ 页眉页脚 ----
\pagestyle{fancy}
\fancyhf{}
\fancyhead[LE]{\small\color{ink} 数学 \ For OI}          % 偶数页（左手页）外侧
\fancyhead[RO]{\small\nouppercase{\leftmark}}            % 奇数页（右手页）外侧
\fancyfoot[C]{\thepage}
\renewcommand{\headrulewidth}{0.3pt}
\fancypagestyle{plain}{\fancyhf{}\fancyfoot[C]{\thepage}\renewcommand{\headrulewidth}{0pt}}

% ------------------------------------------------------------ 定理环境 ----
\theoremstyle{definition}
\newtheorem{definition}{定义}[chapter]
\newtheorem{theorem}[definition]{定理}
\newtheorem{lemma}[definition]{引理}
\newtheorem{corollary}[definition]{推论}
\newtheorem{example}[definition]{例}
\renewcommand{\proofname}{证明}

% -------------------------------------------------------------- 色框 -------
\newtcolorbox{problem}[1]{enhanced,breakable,colback=soft,colframe=ink!60,boxrule=0.5pt,arc=1mm,
  title={#1},fonttitle=\bfseries,coltitle=white,colbacktitle=ink,
  left=2mm,right=2mm,top=1mm,bottom=1mm,before skip=8pt,after skip=6pt}
\newtcolorbox{teach}[1][教学提示]{enhanced,breakable,colback=accent!4,colframe=accent!60,boxrule=0.4pt,arc=2mm,
  title={#1},fonttitle=\bfseries,colbacktitle=accent,coltitle=white,left=2mm,right=2mm}
\newtcolorbox{pitfall}[1][易错点]{enhanced,breakable,colback=warn!4,colframe=warn!60,boxrule=0.4pt,arc=2mm,
  title={#1},fonttitle=\bfseries,colbacktitle=warn,coltitle=white,left=2mm,right=2mm}
\definecolor{tipbase}{HTML}{079C81}
\newtcolorbox{tip}[1][顺带说明]{enhanced,breakable,colback=tipbase!4,colframe=tipbase!60,boxrule=0.4pt,arc=2mm,
  title={#1},fonttitle=\bfseries,colbacktitle=tipbase,coltitle=white,
  left=2mm,right=2mm,top=1mm,bottom=1mm,before skip=6pt,after skip=6pt}
\newtcolorbox{chapterintro}[1][章节导读]{enhanced,breakable,colback=chapterguide!5,colframe=chapterguide,boxrule=0.4pt,arc=2mm,
  title={#1},fonttitle=\bfseries,colbacktitle=chapterguide,coltitle=white,
  left=2mm,right=2mm,top=1mm,bottom=1mm,before skip=6pt,after skip=8pt}

% ------------------------------------------------------------ 简写命令 ----
\newcommand{\Z}{\mathbb Z}
\newcommand{\N}{\mathbb N}
\newcommand{\Pp}{\mathbb P}
\newcommand{\Q}{\mathbb Q}
\newcommand{\R}{\mathbb R}
\newcommand{\eps}{\varepsilon}
\newcommand{\one}{\mathbf 1}
\newcommand{\id}{\operatorname{id}}
\newcommand{\ord}{\operatorname{ord}}
\newcommand{\lcm}{\operatorname{lcm}}
\newcommand{\rad}{\operatorname{rad}}
\newcommand{\iva}[1]{\left[#1\right]}
\newcommand{\floor}[1]{\left\lfloor #1\right\rfloor}
\newcommand{\ceil}[1]{\left\lceil #1\right\rceil}
\newcommand{\sol}{\par\noindent\textbf{解析.}\ }
\newcommand{\hint}{\par\noindent\textbf{解题方向.}\ }
\newcommand{\cost}{\par\noindent\textbf{复杂度与实现.}\ }
\newcommand{\src}[1]{\par\noindent{\footnotesize 题目／延伸资料：\url{#1}}\par}
% 博弈论
\newcommand{\mex}{\operatorname{mex}}
\newcommand{\SG}{\operatorname{SG}}
\newcommand{\Np}{N^+}
\newcommand{\topbit}{\operatorname{topbit}}
\newcommand{\dep}{\operatorname{dep}}
\newcommand{\Val}{V}
% 概率论
\newcommand{\E}{\mathbb E}
\newcommand{\Var}{\operatorname{Var}}
\newcommand{\Cov}{\operatorname{Cov}}
\newcommand{\Bern}{\operatorname{Bern}}
\newcommand{\Bin}{\operatorname{Bin}}
\newcommand{\Pois}{\operatorname{Pois}}
\newcommand{\Geom}{\operatorname{Geom}}
\newcommand{\Harm}{H}

% ------------------------------------------------------------ 代码样式 ----
\lstset{language=C++,basicstyle=\ttfamily\footnotesize,
  keywordstyle=\color{MidnightBlue}\bfseries,commentstyle=\color{gray},
  stringstyle=\color{accent},showstringspaces=false,breaklines=true,
  columns=fullflexible,keepspaces=true,frame=single,rulecolor=\color{ink!25},
  backgroundcolor=\color{soft},tabsize=2}

% ============================================================================
%                        部分页（Part Page）机制
% ============================================================================
% 正式书稿模式：每个部分页都从右页（奇数页）开始，中间留空白页。
% 若想合辑连续排版（不插空白偶数页），把下面一行改成 \comparisontrue。
\newif\ifcomparison
\comparisonfalse

% 定义：\DeclareMotif{编号}{沿局部 x 轴的平移量/mm}{单个元素的路径}
% 各版的副本大小统一，仅平移/旋转；P2 单元素外接圆半径为 78.4mm。
\newcommand{\DeclareMotif}[3]{%
  \expandafter\def\csname Offset#1\endcsname{#2}%
  \expandafter\def\csname Motif#1\endcsname{#3}%
}

% 01 正三角形。奇数边正多边形本身不是中心对称图形，
% 因此这类带平移的图案使用偶数 n，保证成对副本相差 180 度。
\DeclareMotif{1}{20}{
  (90:70mm) -- (210:70mm) -- (330:70mm) -- cycle
}
% 02 正方形。自身中心对称，n=5，旋转角度 i*72 度。
\DeclareMotif{2}{0}{
  (45:78.4mm) -- (135:78.4mm) -- (225:78.4mm) -- (315:78.4mm) -- cycle
}
% 03 正五边形。
\DeclareMotif{3}{18}{
  (90:65mm) -- (162:65mm) -- (234:65mm) -- (306:65mm) -- (378:65mm) -- cycle
}
% 04 四角凹星形。
\DeclareMotif{4}{0}{
  (0,86) -- (-30,30) -- (-86,0) -- (-30,-30)
  -- (0,-86) -- (30,-30) -- (86,0) -- (30,30) -- cycle
}
% 05 正七边形。
\DeclareMotif{5}{18}{
  (90:65mm) -- (141.428571:65mm) -- (192.857143:65mm)
  -- (244.285714:65mm) -- (295.714286:65mm)
  -- (347.142857:65mm) -- (398.571429:65mm) -- cycle
}
% 07 圆形。先偏心平移，再按等角度复制；否则同心圆旋转后会完全重合。
\DeclareMotif{7}{30}{(0,0) circle[radius=48mm]}
% 12 八角凹星。外半径 84mm，内半径 43mm。
\DeclareMotif{12}{0}{
  (90:84mm) -- (112.5:43mm) -- (135:84mm) -- (157.5:43mm)
  -- (180:84mm) -- (202.5:43mm) -- (225:84mm) -- (247.5:43mm)
  -- (270:84mm) -- (292.5:43mm) -- (315:84mm) -- (337.5:43mm)
  -- (360:84mm) -- (382.5:43mm) -- (405:84mm) -- (427.5:43mm) -- cycle
}

% 用法：\PartPage{形状编号}{副本数 n}{罗马数字}{部分名称}{标题}
% n 为 >=2 的整数；编号 1、3、5、7 请用偶数 n 保证整体中心对称。
% 为保留可见的叠加，避免 n 与单个元素的旋转对称性导致完全重合。
% 本宏绘制开篇页，不自动递增 part 计数器或写入目录。
\newcommand{\PartPage}[5]{%
  \ifcomparison\clearpage\else\cleardoublepage\fi
  \thispagestyle{empty}\null
  \begin{tikzpicture}[remember picture,overlay]
    \begin{scope}[shift={(current page.south west)},x=1mm,y=1mm]
      \clip (0,0) rectangle (210,297);
      % 默认中心位于纸张右下角；P2 向右、下各偏移2mm。
      % 该偏移只移动图案，不移动罗马数字或标题。
      \def\MotifX{210}\def\MotifY{0}
      \ifnum#1=2\relax\def\MotifX{212}\def\MotifY{-2}\fi
      % P1（编号1）向左上各偏移16mm；P6（编号7）向左上各偏移8mm。
      \ifnum#1=1\relax\def\MotifX{194}\def\MotifY{16}\fi
      \ifnum#1=7\relax\def\MotifX{202}\def\MotifY{8}\fi
      \begin{scope}[shift={(\MotifX,\MotifY)}]
        \pgfmathtruncatemacro{\LastCopy}{#2-1}
        \foreach \i in {0,...,\LastCopy}{%
          % 统一规则：theta_i = i * (360/n)，i=0,...,n-1。
          \pgfmathsetmacro{\Angle}{\i*360/#2}
          \begin{scope}[rotate=\Angle]
            \begin{scope}[shift={(\csname Offset#1\endcsname,0)}]
              % 单独填充透明度；不是对整体使用透明组。
              \fill[ink,fill opacity=0.20] \csname Motif#1\endcsname;
            \end{scope}
          \end{scope}
        }
      \end{scope}
      % 数字距纸张右侧、底部各约 12 mm；II、III 向左扩展。
      \node[anchor=south east,inner sep=0pt,outer sep=0pt,text=white,
        font=\rmfamily\fontsize{56}{60}\selectfont] at (198,12) {#3};
      \node[anchor=base west,inner sep=0pt,text=ink,
        font=\songti\fontsize{20}{26}\selectfont] at (28,162) {#4};
      \node[anchor=base west,inner sep=0pt,text=ink,
        font=\songti\fontsize{58}{68}\selectfont] at (28,135) {#5};
    \end{scope}
  \end{tikzpicture}%
  \clearpage
}

% 部分编号：只用于生成唯一的超链接锚点，不出现在页面上。
\newcounter{partno}
\renewcommand{\thepartno}{\Roman{partno}}
% 各部分的章号都从 1 重新开始；用 partno 限定锚点名，避免重复目标。
\renewcommand{\theHchapter}{\arabic{partno}.\arabic{chapter}}
\renewcommand{\theHsection}{\theHchapter.\arabic{section}}
\renewcommand{\theHsubsection}{\theHsection.\arabic{subsection}}
\renewcommand{\theHdefinition}{\theHchapter.\arabic{definition}}
\renewcommand{\theHequation}{\theHchapter.\arabic{equation}}
\renewcommand{\theHtable}{\theHchapter.\arabic{table}}
\renewcommand{\theHfigure}{\theHchapter.\arabic{figure}}

% 用法：\BeginPart{形状编号}{副本数 n}{罗马数字}{部分名称}{标题}
% 依次完成：递增部分编号、重置章号、绘制部分页、写入目录。
\newcommand{\BeginPart}[5]{%
  \stepcounter{partno}%
  \setcounter{chapter}{0}%
  \PartPage{#1}{#2}{#3}{#4}{#5}%
  \phantomsection
  \addcontentsline{toc}{part}{#4\quad #5}%
}

\begin{document}
\frontmatter

% ---------------------------------------------------------------- 封面 ----
\begin{titlepage}
\setlength{\parindent}{0pt}
\vspace*{1.2cm}
{\large\color{accent}MATHEMATICS FOR OLYMPIAD IN INFORMATICS}\par
\vspace{1cm}
{\fontsize{46}{65}\selectfont\bfseries\color{ink}数学 \huge{For OI}}\par
\vspace{0.5cm}
{\Large 教师版}\par
\vspace{0.5cm}
{\color{accent}\rule{\textwidth}{1.2pt}}\par
\vspace{0.3cm}
\begin{minipage}{0.9\textwidth}
\large Federico Prask · 2026\\[0.8em]
\normalsize 基础数学\quad 数论\quad 组合数学\quad 博弈论与概率论\quad 多项式与生成函数
\end{minipage}\par
\vfill
{\small 全书按“部分—章”组织。每一部分从一张部分页开始：页脚右下角的旋转叠加图案是该部分的标识，罗马数字标出部分的次序。}\par
{\small XeLaTeX 编译两遍；全部插图均为 TikZ 矢量图，无外部图片依赖。}\par
\end{titlepage}

% ------------------------------------------------------- 阅读说明 ----
\chapter*{阅读说明与课程地图}
\addcontentsline{toc}{chapter}{阅读说明与课程地图}
\markboth{阅读说明与课程地图}{}

\begin{chapterintro}[章节导读]
本书把信息学竞赛需要的数学分成五个部分，从最基础的集合、逻辑与代数，一路走到数论、组合数学、博弈论与概率论，最后以多项式与生成函数收尾。每一部分都是自足的：可以从任何一部分开始读，但前面的工具会在后面反复出现。
\end{chapterintro}

\begin{pitfall}[适用范围与阅读边界]
本书面向已经掌握初中代数、会写基本循环与递归、了解时间复杂度的信息学竞赛学生。第一部分的第 1--5 章定位为“把高中内容重新讲一遍并补上竞赛缺口”，默认读者已经见过这些概念，只是未必习惯于竞赛里的用法；第 6 章（复数与向量）对纯算法题帮助有限，可以晚一点读。第二部分数论的第 0 章是术语对照，遇到不认识的词先查那里。

同余、线性不定方程可以入门；二维积性函数、洲阁筛、超现实数属于选修提高内容。例题均给教学性摘要，不替代评测网站的完整输入输出说明；标为“研究拓展”的内容只给出可验证的数学路线，不冒充可直接提交的完整程序。实现模板不等于各原题的完整提交代码。
\end{pitfall}

\begin{tip}[本书的四种框]
正文里有四种带颜色的框，分工不同：
\begin{itemize}
\item \textbf{定义/定理/引理/推论/例}：黑色小标题，是必须记住的结论。先问自己“条件是什么、结论是什么”，再往下看。
\item \textbf{题目框}（深蓝色标题）：题目与它的解析。标“统一练习”的不是原题，而是把几道题共用的想法抽出来。
\item \textbf{易错框}（棕黄色标题，标注“易错点”）：这里写的是“大多数人会错在哪”，以及一个具体反例。
\item \textbf{提示框}（绿色标题）：只解释一个词、一步跳步或一个例子，可以单独读懂。
\end{itemize}
\begin{teach}[教学提示框]
青色标题的“教学提示”框是教师版特有的内容：讲法设计、课堂追问、易漏点检查表与实验任务。课堂追问建议先遮住答案再展示；例题解析不建议原样念完，应在关键换元前暂停。
\end{teach}
\end{tip}

\section*{全书地图}
\begin{longtable}{p{0.09\textwidth}p{0.20\textwidth}p{0.63\textwidth}}
\toprule
部分 & 主题 & 主线与可检查的学习产出\\
\midrule
\endfirsthead
\toprule
部分 & 主题 & 主线与可检查的学习产出\\
\midrule
\endhead
第一部分 & 基础数学 & 集合与逻辑、代数恒等变形、方程与不等式、函数、数列与归纳、复数与向量。产出：能把一句自然语言命题改写成带量词的数学命题，并写出完整的充要性证明。\\
第二部分 & 数论 & gcd 与 exgcd、同余与逆元、CRT、阶与原根、LTE、阶乘与组合数取模、反演、筛法、随机分解。产出：能写出条件、通解、最小可行解，并说明一次取模是否丢信息。\\
第三部分 & 组合数学 & 排列组合、二项式系数与恒等式、容斥、递推数列、抽屉原理、格路计数。产出：能把一个计数问题说清“数的是什么”，再选加法或乘法、容斥或递推。\\
第四部分 & 博弈论与概率论 & PN 态、模仿棋与策略窃取、Nim 与 SG 理论、纳什均衡；概率空间、期望与方差、常用分布与概率算法。产出：能把一个博弈画成状态图并倒推胜负；能把一个随机过程写成期望的线性组合。\\
第五部分 & 多项式与生成函数 & 形式幂级数、多项式乘法与快速傅里叶变换、生成函数计数。\\
附录 & A--D & 分层训练、核心模板、课堂设计与判定模板。\\
\bottomrule
\end{longtable}

\begin{teach}[建议的教学分层]
\begin{longtable}{p{0.12\textwidth}p{0.40\textwidth}p{0.38\textwidth}}
\toprule
层次 & 主线 & 可检查的学习产出\\
\midrule
A：基础 & 集合与逻辑、代数恒等变形、排列组合、PN 态递推 & 能完成充要性两段证明；能对一个小的状态图手工完成逆向归纳。\\
B：核心 & 同余与逆元、二项式系数、容斥、Nim 与 SG、期望线性性 & 能独立写出“两段式”证明；能把多堆游戏拆成和游戏并异或结算。\\
C：提高 & LTE、杜教筛、递推数列、反射原理、概率算法 & 能区分预处理和查询成本；能说明每个定理的适用条件缺一条会怎样。\\
D：专题 & 二维积性函数、洲阁筛、超现实数、混合策略 & 能识别稀疏支撑、状态等价与证明缺口，不只记最终公式。\\
\bottomrule
\end{longtable}
每讲按 $90$ 分钟安排：A 层约 $6$--$8$ 讲，B 层约 $8$--$10$ 讲，C 层约 $6$--$8$ 讲，D 层按选题增加。课时仅作备课建议，不应在一节课内从逆元跳到二维卷积。
\end{teach}

\section*{统一记号}
\begin{longtable}{p{0.23\textwidth}p{0.69\textwidth}}
\toprule
记号 & 含义与约定\\
\midrule
$\Pp,\Z,\N,\Q,\R$ & 分别为素数集、整数集、非负整数集、有理数集、实数集；数论函数的定义域为正整数。本书约定 $\N$ 含 $0$。\\
$m,M$ & 一般模数，默认为正整数；讨论单位群与逆元时通常 $m\ge2$。$m=1$ 的模运算结果一律为 $0$。\\
$p,q$ & 素数；$Q=p^\alpha$ 表示素数幂，避免把合数也一律记成 $p$。\\
$\iva{P}$ & Iverson 括号，命题成立为 $1$，否则为 $0$。\\
$\nu_p(x)$ & 非零整数 $x$ 中 $p$ 的指数，即 $\nu_p(|x|)$；需要时约定 $\nu_p(0)=+\infty$。\\
$\varphi,\mu,\tau$ & 欧拉函数、莫比乌斯函数、约数个数函数。原稿的 $d(n)$ 统一写作 $\tau(n)$。\\
$\one,\eps,\id_t$ & 常一函数、卷积单位函数、$\id_t(n)=n^t$。原稿 $I$ 统一写作 $\one$。\\
$f*g$ & 狄利克雷卷积；$fg$ 为逐点乘积。卷积除法不使用普通分数线混写。\\
$S_f(n)$ & $\sum_{i=1}^n f(i)$，并约定 $S_f(0)=0$。\\
$\ord_m(a)$ & $a$ 模 $m$ 的乘法阶；与原稿 $\delta_m(a)$ 相同。\\
$A,B$ & 初始状态下的先后手，一旦定下不再变化；具体问题中常用固定的人名。\\
先手／后手 & 当前局面下的下一步行动者／另一方，随局面变化；“把局面做成 P 态”即将先手让给对方。\\
P 态／N 态 & 公平组合游戏中后手必胜态（Previous player wins）／先手必胜态（Next player wins）。\\
$u\to v$，$\Np(u)$ & 状态图的一条合法转移；$u$ 的全部后继构成的集合。\\
$\oplus$ & 按位异或（Nim 和）；$S=a_1\oplus a_2\oplus\cdots\oplus a_n$。\\
$(x)_k$ & 非负整数 $x$ 二进制表示的第 $k$ 位；$\topbit(S)$ 为 $S$ 的最高置位。\\
$\mex(X)$ & 集合 $X$ 中未出现的最小非负整数。\\
$g(u)$，$\SG(u)$ & 状态 $u$ 的 SG 值（Grundy 值），两者完全同义；$g(u)=0$ 恰为 P 态。\\
$G=\{G^L\mid G^R\}$ & 游戏的双视角记号：$G^L,G^R$ 分别为 Left／Right 走一步能到达的局面集合。\\
$\ast$，$\ast_n$ & star：$\ast=\{0\mid 0\}$；一般地 $\ast_n$ 表示 SG 值为 $n$ 的公平游戏（大小为 $n$ 的 Nim 堆）。\\
$\omega$ & 最小无穷序数对应的超现实数 $\{0,1,2,\ldots\mid\}$。\\
$W$ & 双人零和博弈的收益矩阵，$W_{i,j}$ 为玩家 A 的收益；B 的收益为其相反数。\\
纯策略／混合策略 & 固定一种决策／为每种决策分配概率后随机决策。\\
misère（反规则） & “无法行动者胜”的规则族，如反 Nim 游戏；与正常规则结论可能截然不同。\\
$\floor{x},\ceil{x}$ & 下取整与上取整。\\
\bottomrule
\end{longtable}

\begin{tip}[速查表怎么用]
上面这张表是\textbf{速查表}，用术语解释术语，适合已经会的人。本书第 0 章（第二部分）与第一部分第 1 章把每一行都翻译成了大白话，并说明这个词为什么会出现。
\end{tip}

所有复杂度默认采用竞赛中的字长运算模型。高精度题目的“多项式时间”以输入位数 $\log n$ 为变量；$O(n^{1/4})$ 并不是关于输入位数的多项式时间。哈希表复杂度默认期望，随机算法的经验性能与严格最坏界会明确区分。

\begin{tip}[词义：字长模型]
\textbf{字长模型}= 假设 $64$ 位整数的一次加减乘比较都算一步。它不是真实机器模型（内存不连续时会慢），也不是高精度模型（那时乘法是 $O(n^2)$）。凡是遇到 $10^9$ 以上相乘、$10^{18}$ 以上相加，先想\textbf{会不会溢出 \texttt{long long}}。
\end{tip}

\tableofcontents
\mainmatter

% ============================================================================
%  第一部分  基础数学
% ============================================================================
\BeginPart{1}{6}{I}{第一部分}{基础数学}
% !TeX program = xelatex
% ============================================================================
%  第一部分  基础数学
%  第 1 章  集合、逻辑与证明方法
%  第 2 章  数系、运算与代数恒等变形
%  第 3 章  方程与不等式
%  第 4 章  函数
%  第 5 章  数列、求和与数学归纳法
%  第 6 章  复数与平面向量初步
% ============================================================================

\chapter{集合、逻辑与证明方法}
\begin{chapterintro}[章节导读]
从这一章开始，我们把“怎么把一件事说清楚”本身当作一门技术来练。信息学竞赛里大量的失分不是不会算，而是题意翻译错了、条件漏了一条、或者把“有解”和“有正整数解”混为一谈。集合与逻辑给出无歧义的表述工具，证明方法给出检验自己结论的工具。全书后面每一个定理的陈述与证明，用的都是本章的语言。
\end{chapterintro}

\begin{tip}[这一章解决什么问题]
本书第二部分的第 0 章是一张“术语速查表”，用术语解释术语；本章则是把这些词从零讲一遍。如果你读到一个词被卡住，先回第二部分第 0 章查表；如果你从来没有用过这些词，从本章读起。
\end{tip}

\section{集合与元素}
集合是“一些确定的对象放在一起”。对象叫作\textbf{元素}，$a\in A$ 读作“$a$ 是 $A$ 的元素”，$a\notin A$ 读作“不是”。$A\subseteq B$ 读作“$A$ 的每个元素都在 $B$ 里”（$A$ 含于 $B$）；$\varnothing$ 是空集，不含任何元素。两个集合相等，当且仅当它们互相含于：
\[
A=B\iff A\subseteq B\ \text{且}\ B\subseteq A.
\]
这句话是证明集合相等的\textbf{唯一标准动作}：要证 $A=B$，就分别证两个方向，不要试图“一步到位”。

本书最常用的几个数集：
\[
\Z=\{\ldots,-2,-1,0,1,2,\ldots\},\qquad
\N=\{0,1,2,\ldots\},\qquad
\Pp=\{2,3,5,7,11,\ldots\},
\]
分别是整数集、非负整数集、素数集；$\Q$、$\R$ 是有理数集与实数集。\textbf{注意本书约定 $\N$ 含 $0$}；如果你手上的教材不含 $0$，读本书时按本书的约定来，否则很多递推的边界会差一项。

描述法的写法 $\{x\in\Z:0\le x<m\}$ 表示“满足后面条件的那些 $x$”，冒号也可写成竖线。集合 $\{0,1,\ldots,m-1\}$ 有 $m$ 个元素——别把“从 $0$ 到 $m-1$”数成 $m-1$ 个，这是最常犯的低级错误之一。

\begin{definition}[交、并、差、补]
$A\cap B$ 是同时在两者中的元素；$A\cup B$ 是至少在一个中的元素；$A\setminus B$ 是在 $A$ 中但不在 $B$ 中的元素。在全集 $U$ 中，$\overline A=U\setminus A$ 称为 $A$ 的补集。
\end{definition}

\begin{theorem}[德摩根定律]
\[
\overline{A\cup B}=\overline A\cap\overline B,\qquad
\overline{A\cap B}=\overline A\cup\overline B.
\]
\end{theorem}
\begin{proof}
证第一个等式，用“两个方向”的标准动作。

充分性（$\overline{A\cup B}\subseteq\overline A\cap\overline B$）。设 $x\in\overline{A\cup B}$，则 $x\notin A\cup B$，于是 $x\notin A$ 且 $x\notin B$，即 $x\in\overline A$ 且 $x\in\overline B$，故 $x\in\overline A\cap\overline B$。

必要性（$\overline A\cap\overline B\subseteq\overline{A\cup B}$）。设 $x\in\overline A\cap\overline B$，则 $x\notin A$ 且 $x\notin B$，所以 $x$ 不属于 $A\cup B$（否则它至少在 $A$ 或 $B$ 之一中），故 $x\in\overline{A\cup B}$。

第二个等式完全同理。
\end{proof}

\begin{teach}[为什么要练“证集合相等”]
因为本书后面几乎所有“算法正确”的论证，最后都落在“这两个集合相同”上：欧拉定理里“乘以 $a$ 以后简化剩余系被重排”、中国剩余定理里“解集是一个同余类”、杜教筛里“递归状态落在原商集中”。让学生把“证集合相等=证两个包含”练成肌肉记忆，比多做十道题值钱。
\end{teach}

\subsection{一个必须分清的区别}
“$A$ 是 $B$ 的子集”（$A\subseteq B$，集合与集合的关系）和“$a$ 是 $A$ 的元素”（$a\in A$，元素与集合的关系）是两回事。写成 $\{\{1,2\}\}\subseteq\{\{1,2\},\{3\}\}$ 是对的，而 $\{1,2\}\in\{\{1,2\},\{3\}\}$ 也是对的——但 $\{1,2\}\subseteq\{\{1,2\},\{3\}\}$ 是错的。读题时看到“选出一个子集”和“选出一个元素”，先分清在说哪一层。

\section{逻辑联结词与量词}
\subsection{“或”“且”“非”“当且仅当”}
\begin{itemize}
\item “$P$ 且 $Q$”：两者同时成立，符号 $P\land Q$。
\item “$P$ 或 $Q$”：至少一个成立（可以两个都成立），符号 $P\lor Q$。\textbf{数学里的“或”不是“二选一”}。
\item “非 $P$”：$P$ 不成立，符号 $\lnot P$。
\item “$P\Rightarrow Q$”：$P$ 成立时 $Q$ 必须成立，读作“$P$ 推出 $Q$”。这时说“$P$ 是 $Q$ 的\textbf{充分}条件”“$Q$ 是 $P$ 的\textbf{必要}条件”。
\item “$P\iff Q$”：互相推出，读作“$P$ 当且仅当 $Q$”，也说“$P$ 是 $Q$ 的\textbf{充要条件}”。
\end{itemize}

\begin{pitfall}[把“或”读成“异或”]
题目说“$x$ 是偶数或 $x$ 是 $3$ 的倍数”，$x=6$ 当然满足。若你在草稿纸上写下“二选一”，就会把 $6$ 排除，后面的计数全错。组合数学的容斥原理之所以存在，就是因为在数“或”。
\end{pitfall}

\subsection{量词：任意与存在}
$\forall$ 读“对任意（所有）”，$\exists$ 读“存在（至少有一个）”，$\exists!$ 读“存在且唯一”。\textbf{量词的顺序改变含义}：
\[
\forall n\in\N,\ \exists p\in\Pp,\ p>n\qquad\text{（素数无穷多个）}.
\]
$\forall$ 在前表示“先给我 $n$，我再找 $p$”，$p$ 可以依赖 $n$。若写成 $\exists p,\ \forall n$，就变成“有一个固定素数比所有数都大”，荒谬。

否定规则要背熟：
\[
\lnot(\forall x,P(x))\iff \exists x,\lnot P(x),\qquad
\lnot(\exists x,P(x))\iff \forall x,\lnot P(x).
\]
于是要推翻“所有情况都好”，只需\textbf{举一个反例}；要反驳“存在坏情况”，必须证明“全都好”。本书第二部分的 Miller--Rabin、第三部分的抽屉原理，反复用到这两句话。

\begin{tip}[做题时的第一动作：把量词写出来]
拿到一道题，先在草稿纸上把结论的量词结构写出来：“对所有 $x$ 存在 $y$……”还是“存在 $x$ 对所有 $y$……”。很多题的难点就在于这个顺序，而题面用的是自然语言，容易含混。
\end{tip}

\section{证明的三种基本方法}
\subsection{直接证明}
从条件出发，沿着“定义—已知定理—定义”的链条推到结论。本书第二部分第 1 章裴蜀定理的证明就是标准样板：先用 $d=\gcd(a,b)$ 把 $a,b$ 写成 $da',db'$，再把目标式子提一个 $d$ 出来。

\subsection{反证法与最小反例}
要证“所有 $n$ 都满足 $P(n)$”，可设存在反例，取\textbf{最小的}那个 $n_0$，再由 $n_0$ 造出更小的反例 $n'<n_0$，得到矛盾。这就是\textbf{最小反例法}，它与数学归纳法等价。第二部分第 1 章证明“欧几里得算法正确”、第 15 章的韦达跳跃都用它，因为它们都需要“$n_0$ 的最小性”这一件武器。

\begin{pitfall}[“不停得到更小”需要一个停止条件]
反证法里“能无限变小”之所以是矛盾，靠的是\textbf{自然数的最小数原理}：非空正整数集有最小元。如果你递减的量不是整数（比如实数、字符串长度），要先说明它良基，否则“无限递减”不构成矛盾。第四部分博弈论里“状态图有环所以没有良基归纳”，正是这一条失效的例子。
\end{pitfall}

\subsection{数学归纳法}
写清三件事：命题 $P(n)$ 是什么；$P(1)$（或 $P(0)$）成立；“若对所有 $k<n$ 成立则 $P(n)$ 成立”。需要“所有更小的”而不是只有 $n-1$ 时，叫\textbf{第二归纳法}，本书也说“强归纳”。

\begin{theorem}[二项式定理的归纳证明]
对任意非负整数 $n$ 与实数 $x,y$，
\[
(x+y)^n=\sum_{k=0}^{n}\binom nk x^{n-k}y^{k}.
\]
\end{theorem}
\begin{proof}
对 $n$ 归纳。$n=0$ 时两边都是 $1$。设结论对 $n$ 成立，则
\[
(x+y)^{n+1}=(x+y)\sum_{k=0}^{n}\binom nk x^{n-k}y^{k}
=\sum_{k=0}^{n}\binom nk x^{n+1-k}y^{k}
 +\sum_{k=0}^{n}\binom nk x^{n-k}y^{k+1}.
\]
把第二个和式的指标换成 $j=k+1$：
\[
(x+y)^{n+1}=x^{n+1}+\sum_{k=1}^{n}\left[\binom nk+\binom n{k-1}\right]x^{n+1-k}y^{k}+y^{n+1}.
\]
用恒等式 $\dbinom nk+\dbinom n{k-1}=\dbinom{n+1}k$（它是第三部分第 1 章杨辉三角的直接写法），并把两端 $k=0$、$k=n+1$ 的项并回去，即得
\[
(x+y)^{n+1}=\sum_{k=0}^{n+1}\binom{n+1}{k}x^{n+1-k}y^{k}.
\]
归纳完成。
\end{proof}

\begin{teach}[归纳法的三个常见断点]
\begin{enumerate}
\item 忘了写 $P(1)$：归纳缺了起点，后面的链条挂在空中。
\item 归纳假设用错：证 $P(n)$ 时只能用 $P(1),\ldots,P(n-1)$，不能顺手用 $P(n)$。
\item 归纳步里悄悄用了“显然”：把一步代数留给读者，往往正是错的地方。要求学生在归纳步里把换指标的那一行显式写出来。
\end{enumerate}
\end{teach}

\section{双向证明的写法}
要证“$P$ 当且仅当 $Q$”，写两个方向、各起一段，并给每段一个开头词（“必要性.”“充分性.”）。若只证了一边，就把结论降格为“$P$ 推出 $Q$”。

\begin{tip}[“必要性来自……充分性来自……”是在说什么]
翻译成人话就是：“两个方向分开证”。\textbf{必要性}指“已知 $P$ 成立，去证 $Q$”（不做到的话 $P$ 就不成立，所以 $Q$ 是\emph{必须}的）；\textbf{充分性}指“已知 $Q$ 成立，去造出 $P$”（$Q$ 一条就\emph{足够}了）。一个方向没证，结论就还只是猜测。
\end{tip}

\section{分层例题与解析}
以下三题分别练“量词翻译”“集合相等”“归纳与递推”三件基本功。

\begin{problem}{例 1：把自然语言改写成带量词的命题}
设 $f:\N\to\N$。用 $\forall,\exists$ 与逻辑联结词写出下面两句话，并说明它们是否等价：
(1) “$f$ 有不动点”，即存在 $n$ 使 $f(n)=n$；
(2) “$f$ 不把任何两个不同的数映到同一个值”，即 $f$ 是单射。
\end{problem}
\hint (1) 只有一个存在量词。(2) 要对“任意两个不同的输入”发言，先固定它们。

\sol (1) 写成 $\exists n\in\N,\ f(n)=n$。(2) 写成
\[
\forall a,b\in\N,\ a\ne b\Rightarrow f(a)\ne f(b).
\]
两者不等价。$f(n)=n+1$ 是单射但没有不动点；$f(n)=0$（常函数）有不动点（$n=0$）但不是单射。\textbf{把“不变量”和“存在某个特殊值”分开}，是读题的基本功。

\begin{problem}{例 2：证明两个集合相等}
设 $A=\{3k+1:k\in\N\}$、$B=\{3k-2:k\in\N_{>0}\}$。证明 $A=B$。
\end{problem}
\hint 用“互相含于”的标准动作，每一边都把元素写成参数的形式。

\sol 设 $x\in A$，则 $x=3k+1=3(k+1)-2$，其中 $k+1\ge1$，故 $x\in B$。设 $x\in B$，则 $x=3k-2=3(k-1)+1$，其中 $k\ge1$ 故 $k-1\ge0$，于是 $x\in A$。两个方向都成立，$A=B$。

\begin{problem}{例 3：一个容易漏条件的递推}
在 $n\times n$ 的棋盘上放棋子，要求每行每列恰有一个。设 $n=1,2,3$ 时的放法数为 $f_1,f_2,f_3$。有人断言 $f_n=n!$，并用归纳法“证明”如下：假设 $f_n=n!$，则 $f_{n+1}=(n+1)f_n=(n+1)!$。请指出漏洞。
\end{problem}
\hint 归纳步里“$(n+1)f_n$”这一步凭什么成立？$n\times n$ 的解如何变成 $(n+1)\times(n+1)$ 的解？

\sol 结论 $f_n=n!$ 其实是对的（这就是排列数，$n$ 行各选一列且不重复），但\textbf{归纳步的论证不成立}：不是每一个 $(n+1)\times(n+1)$ 的解都能由唯一的 $n\times n$ 的解“加一行一列”得到，因为新增的那一行可能占用原来任意一列，从而把原来的解整体重排。正确的论证是直接的：第 $1$ 行有 $n$ 种选择，第 $2$ 行有 $n-1$ 种，……，共 $n!$ 种。这个例子用来说明：\textbf{递推关系必须先说清“每个对象被数了几次”}，不能只看数值碰巧对上。

\begin{teach}[这一章的离堂检查]
让学生在 $3$ 分钟内写出：(1) $\lnot(\forall x\exists y,P(x,y))$ 的等价形式；(2) 证 $A\subseteq B$ 的第一步该写什么；(3) 数学归纳法必须写清的三件事。三题全对即可进入第 2 章。
\end{teach}

\chapter{数系、运算与代数恒等变形}
\begin{chapterintro}[章节导读]
代数变形的能力决定了后面所有章节的速度：同余式要变形、不等式要变形、递推要变形、生成函数更要变形。本章不引入新算法，只把运算律、绝对值、乘法公式、指数对数这些“手里的工具”擦亮，并用竞赛里最常见的几种变形模式把它们串起来。
\end{chapterintro}

\section{数系与运算律}
自然数 $\N$、整数 $\Z$、有理数 $\Q$、实数 $\R$ 逐层扩张，每一次扩张都是为了“能解更多的方程”：减法要负数，除法要分数，开方与极限要实数。竞赛里绝大多数对象在 $\Z$ 或 $\Q$ 中，$\R$ 只在不等式与连续性讨论中出现。

四则运算满足交换律、结合律、分配律：
\[
a+b=b+a,\quad a(b+c)=ab+ac,\quad (ab)c=a(bc).
\]
\begin{pitfall}[分配律的常见笔误]
$(a+b)^2=a^2+2ab+b^2$，不是 $a^2+b^2$；$(a+b)(c+d)=ac+ad+bc+bd$，不是 $ac+bd$。第二部分数论里欧拉定理的证明中，原稿曾把 $(a+b)c$ 写成 $ac+ab$，一个符号错误就让整段证明失效。\textbf{凡是用分配律的地方，逐项写全，不要跳。}
\end{pitfall}

\subsection{整数除法与取整}
对整数 $a$ 与非零整数 $b$，带余除法给出唯一的 $q,r$ 使
\[
a=qb+r,\qquad 0\le r<|b|.
\]
$\floor{x}$ 是“不超过 $x$ 的最大整数”，$\ceil{x}$ 是“不小于 $x$ 的最小整数”。两条常用性质：
\[
\floor{\frac ab}=q\iff qb\le a<(q+1)b,\qquad
\ceil{\frac ab}=q\iff (q-1)b<a\le qb\quad(b>0).
\]
第二条也可写成 $\ceil{a/b}=\floor{(a+b-1)/b}$（$a\ge0$、$b>0$ 时）。

\begin{pitfall}[C++ 的除法不是向下取整]
C++ 的 \texttt{/} 是\textbf{向零取整}（截断）：\texttt{-7/2} 得 $-3$，而 $\floor{-7/2}=-4$。所以本书第二部分模板里所有“整除取商”都写成 \texttt{floor\_div}。对 $b>0$，设 $q_0$ 为 \texttt{a/b}、$r_0$ 为 \texttt{a\%b}，则
\[
\operatorname{floorDiv}(a,b)=q_0-\iva{r_0<0},\qquad
\operatorname{ceilDiv}(a,b)=q_0+\iva{r_0>0}.
\]
这样还避免了对最小负整数直接求 $-a$。测试 $a=-5,b=3$ 时应得 floor 为 $-2$、ceil 为 $-1$。
\end{pitfall}

\section{绝对值}
\[
|x|=\begin{cases}x,&x\ge0,\\-x,&x<0.\end{cases}
\]
它的代数意义是“到原点的距离”，因此一切距离性质都能翻译成绝对值不等式：
\[
|xy|=|x||y|,\qquad \left|\frac xy\right|=\frac{|x|}{|y|}\ (y\ne0),
\]
\[
|x|-|y|\le |x+y|\le |x|+|y|\qquad\text{（三角不等式）},
\]
\[
|x|=|y|\iff x^2=y^2,\qquad |x|<a\iff -a<x<a\quad(a>0).
\]
\begin{tip}[绝对值问题的标准动作]
遇到含绝对值的方程或不等式，标准动作是\textbf{分段去掉绝对值符号}：找到使每个绝对值内部变号的点，把数轴分成若干段，在每一段内绝对值可以放心展开。分段点是唯一需要小心的地方，边界上取等号与否要单独验证。
\end{tip}

\section{乘法公式与因式分解}
下面五个公式要能做到“闭着眼睛写出来”：
\begin{align}
(a\pm b)^2&=a^2\pm2ab+b^2,\\
(a+b+c)^2&=a^2+b^2+c^2+2ab+2bc+2ca,\\
a^2-b^2&=(a-b)(a+b),\\
a^3\pm b^3&=(a\pm b)(a^2\mp ab+b^2),\\
(a\pm b)^3&=a^3\pm3a^2b+3ab^2\pm b^3.
\end{align}
更一般地，$a^n-b^n$ 总能被 $a-b$ 整除；$n$ 为奇数时 $a^n+b^n$ 能被 $a+b$ 整除。第二部分第 5 章的升幂引理（LTE）研究的就是“$a^n-b^n$ 里到底含多少个 $p$”，那是这条路线的终点。

\subsection{换元的威力}
很多难题的难点不在公式，而在\textbf{看不出该换什么元}。三条高频换元：
\begin{itemize}
\item 对称式：遇到 $x+y$ 与 $xy$ 反复出现，令 $s=x+y$、$p=xy$，把两变量问题化成一变量问题。
\item 公共因子：遇到 $\gcd(a,b)$，令 $d=\gcd(a,b)$、$a=du$、$b=dv$、$\gcd(u,v)=1$，把“公约数”从“互质部分”中分离出来。这是第二部分第 1 章的起手式。
\item 缩放：遇到 $ax\equiv bx\pmod m$ 型的同余式，先提出 $\gcd(x,m)$，见第二部分第 2 章。
\end{itemize}

\begin{teach}[换元的教学顺序]
先给一个不用换元很难做的题（例如第二部分第 1 章的“数学游戏”），让学生先试 $5$ 分钟；再把换元写在黑板上，问“这个换元把什么分开了”。学生说不出“分开了公约数与互质部分”，就说明还没有学会，只是记住了形式。
\end{teach}

\section{分式、根式与指数对数}
\subsection{分式}
\[
\frac ab\pm\frac cd=\frac{ad\pm bc}{bd},\qquad
\frac ab\cdot\frac cd=\frac{ac}{bd},\qquad
\frac ab\div\frac cd=\frac{ad}{bc}\ (c\ne0).
\]
分式运算的第一条纪律：\textbf{先看定义域}。$\dfrac1{x-1}$ 在 $x=1$ 无定义，$\dfrac{x^2-1}{x-1}$ 在 $x=1$ 也无定义（虽然它在 $x\ne1$ 时等于 $x+1$）。竞赛里“约分前后定义域不同”是丢分的常见原因。

\subsection{根式}
$\sqrt{a}$（$a\ge0$）是 $a$ 的算术平方根。两条高频变形：
\[
\sqrt{a^2}=|a|,\qquad
\frac{1}{\sqrt a+\sqrt b}=\frac{\sqrt a-\sqrt b}{a-b}\quad\text{（分母有理化）}.
\]

\subsection{指数与对数}
$a^n$ 表示 $n$ 个 $a$ 相乘，约定 $a^0=1$（包括 $0^0=1$，按组合意义）。指数律：
\[
a^ma^n=a^{m+n},\quad (a^m)^n=a^{mn},\quad (ab)^n=a^nb^n.
\]
\textbf{注意 $(a^m)^n=a^{mn}$ 与 $a^{m^n}$ 完全不同}：幂塔 $a^{b^c}$ 按\textbf{右结合}读，等于 $a^{(b^c)}$。$3^{3^3}=3^{27}$ 已经约 $7.6\times10^{12}$，所以处理幂塔的核心矛盾是：指数巨大，而余数只有 $m$ 种可能。第二部分第 2 章的扩展欧拉定理就是为此准备的。

对数是指数的逆运算：$y=\log_a x$ 表示 $a^y=x$（$a>0,a\ne1,x>0$）。三条性质：
\[
\log_a(xy)=\log_a x+\log_a y,\quad
\log_a(x^y)=y\log_a x,\quad
\log_a b=\frac{\log_c b}{\log_c a}.
\]
\begin{tip}[换底公式的竞赛用法]
算法竞赛里 $\log$ 几乎只用来估计复杂度，此时\textbf{底数不影响渐近结论}，$\log_2 n$ 与 $\ln n$ 只差常数倍。但一旦进入精确计算（例如判断 $2^k\ge n$），底数必须写清楚。
\end{tip}

\section{分层例题与解析}
\begin{problem}{例 1：化简与定义域}
化简 $\dfrac{x^2-4}{x^2-5x+6}\cdot\dfrac{x-3}{x+2}$，并指出结果在哪些 $x$ 上成立。
\end{problem}
\hint 先因式分解，再约分；约掉因式以后要记住被约掉的因式不能为 $0$。

\sol 分子 $x^2-4=(x-2)(x+2)$，分母 $x^2-5x+6=(x-2)(x-3)$。于是
\[
\frac{(x-2)(x+2)}{(x-2)(x-3)}\cdot\frac{x-3}{x+2}=1,
\]
但\textbf{仅在 $x\ne\pm2$ 且 $x\ne3$ 时成立}。若题目问“对哪些 $x$ 等式成立”，答案是 $x\in\R\setminus\{-2,2,3\}$；只写 $1$ 不扣定义域是常见错误。

\begin{problem}{例 2：含绝对值的方程}
解方程 $|x-1|+|x-2|+|x-3|=4$。
\end{problem}
\hint 分段点是 $1,2,3$；在每一段内绝对值可以放心展开。

\sol 分四段。$x<1$ 时，原式为 $(1-x)+(2-x)+(3-x)=6-3x=4$，得 $x=\dfrac23$（符合 $x<1$）。$1\le x<2$ 时，原式为 $(x-1)+(2-x)+(3-x)=4-x=4$，得 $x=0$，\textbf{不在本段，舍去}。$2\le x<3$ 时，原式为 $(x-1)+(x-2)+(3-x)=x=4$，舍去。$x\ge3$ 时，原式为 $3x-6=4$，得 $x=\dfrac{10}{3}$。故解为 $x=\dfrac23$ 或 $x=\dfrac{10}{3}$。

\begin{teach}[“舍去”这一步不能省]
中间两段求出的解不在本段，必须显式写出“舍去”。学生常常算对了却忘了检查分段条件，写出 $x=0$ 这样的假解。要求每段都写“解得 $x=\cdots$，符合／不符合 $x\in[\cdots]$，故保留／舍去”。
\end{teach}

\begin{problem}{例 3：对称换元}
已知实数 $x,y$ 满足 $x+y=7$、$xy=12$，求 $x^3+y^3$。
\end{problem}
\hint 令 $s=x+y$、$p=xy$，把目标写成 $s,p$ 的多项式。

\sol 用恒等式 $x^3+y^3=(x+y)^3-3xy(x+y)$，代入得
\[
x^3+y^3=7^3-3\cdot12\cdot7=343-252=91.
\]
不需要解出 $x,y$（它们是 $3$ 与 $4$）——\textbf{能不解就不解}，是代数题的基本修养。

\begin{problem}{例 4：幂塔的结合性}
比较 $2^{3^2}$ 与 $(2^3)^2$，并计算 $2^{2^{2^2}}$。
\end{problem}
\hint 幂塔右结合；$(a^m)^n=a^{mn}$。

\sol $2^{3^2}=2^9=512$，而 $(2^3)^2=2^6=64$，两者不等。$2^{2^{2^2}}=2^{(2^{(2^2)})}=2^{16}=65536$；若误按左结合算成 $((2^2)^2)^2=2^8=256$，两者相差 $256$ 倍。第二部分第 4 章处理高阶箭号时，这个区别会反复出现。

\begin{pitfall}[指数运算的三条红线]
\begin{enumerate}
\item $(a^m)^n=a^{mn}$，但 $a^{m^n}=a^{(m^n)}$，两者一般不等。
\item $(ab)^n=a^nb^n$，但 $(a+b)^n\ne a^n+b^n$。
\item $a^{m+n}=a^ma^n$，但 $a^{mn}\ne a^ma^n$。
\end{enumerate}
\end{pitfall}

\chapter{方程与不等式}
\begin{chapterintro}[章节导读]
方程与不等式是“把条件写成式子再求解”的两种基本形态。竞赛里方程的难点常在“解不出来但不需要解出来”（韦达定理、对称性），不等式的难点常在“证明一个数非负”。本章把一元二次方程、韦达定理、基本不等式与柯西不等式讲透，它们是第三部分组合计数与第四部分概率期望的常备工具。
\end{chapterintro}

\section{一元二次方程}
$a x^2+bx+c=0$（$a\ne0$）的判别式 $\Delta=b^2-4ac$：
\begin{itemize}
\item $\Delta>0$：两个不等实根 $x=\dfrac{-b\pm\sqrt{\Delta}}{2a}$；
\item $\Delta=0$：重根 $x=-\dfrac b{2a}$；
\item $\Delta<0$：无实根。
\end{itemize}
\begin{theorem}[求根公式]
$a\ne0$ 时，$ax^2+bx+c=0$ 的根为
\[
x_{1,2}=\frac{-b\pm\sqrt{b^2-4ac}}{2a}.
\]
\end{theorem}
\begin{proof}
$a\ne0$，两边除以 $a$ 并配方：
\[
x^2+\frac ba x+\frac ca=0
\iff \left(x+\frac b{2a}\right)^2=\frac{b^2-4ac}{4a^2}.
\]
右边开方（注意 $\sqrt{4a^2}=2|a|$，把符号吸收进 $\pm$ 中即可），移项即得。
\end{proof}

\begin{pitfall}[数值稳定性]
用求根公式算根时，若 $b^2\gg4ac$，$\sqrt{b^2-4ac}\approx|b|$，两个根中有一个由“两个相近的数相减”得到，有效数字大量丢失。竞赛中遇到这种情形（例如解 $x^2-10^9x+1=0$），应改用
\[
x_1x_2=\frac ca,\qquad x_1+x_2=-\frac ba,
\]
先算绝对值大的根，再用乘积关系算另一个根。
\end{pitfall}

\section{韦达定理}
\begin{theorem}[韦达定理]
设 $ax^2+bx+c=0$（$a\ne0$）的两根为 $x_1,x_2$，则
\[
x_1+x_2=-\frac ba,\qquad x_1x_2=\frac ca.
\]
\end{theorem}
\begin{proof}
由求根公式直接相加、相乘：
\[
x_1+x_2=\frac{-b+\sqrt\Delta}{2a}+\frac{-b-\sqrt\Delta}{2a}=-\frac ba,
\]
\[
x_1x_2=\frac{(-b)^2-\Delta}{4a^2}=\frac{b^2-(b^2-4ac)}{4a^2}=\frac ca.
\]
\end{proof}

韦达定理的真正用处是\textbf{不解方程就能做关于根的对称式计算}。若 $f(x_1,x_2)$ 是对称多项式（交换 $x_1,x_2$ 不变），它一定能写成 $s=x_1+x_2$ 与 $p=x_1x_2$ 的多项式。常用：
\[
x_1^2+x_2^2=s^2-2p,\qquad
x_1^3+x_2^3=s^3-3sp,\qquad
\frac1{x_1}+\frac1{x_2}=\frac{s}{p}\ (p\ne0).
\]

\begin{tip}[韦达跳跃]
把韦达定理反过来用：固定其他变量，把方程看成关于某一个变量的二次方程，那么“另一个根”可以用 $s=\dfrac{-b}{a}$ 减出来。这个技巧在第二部分第 15 章（构造不定方程的解）里会变成一件主武器——从一个已知解造出一个更小的解，从而完成无穷递降。
\end{tip}

\section{不等式的基本性质}
对实数 $a,b,c,d$：
\begin{itemize}
\item $a>b\Rightarrow a\pm c>b\pm c$（加减不改变方向）；
\item $a>b$ 且 $c>0\Rightarrow ac>bc$；$a>b$ 且 $c<0\Rightarrow ac<bc$（\textbf{乘负数要变向}）；
\item $a>b>0\Rightarrow a^2>b^2$、$\dfrac1a<\dfrac1b$；
\item $a>b$ 且 $c>d\Rightarrow a+c>b+d$（同向可加）；$a>b>0$ 且 $c>d>0\Rightarrow ac>bd$（同正可乘）。
\end{itemize}

\begin{pitfall}[“同向可乘”要求两边都正]
$a>b$ 且 $c>d$ \textbf{不能}推出 $ac>bd$：取 $a=1,b=-2,c=-1,d=-3$，有 $ac=-1<bd=6$。乘法保序需要两边同号。
\end{pitfall}

\section{基本不等式与柯西不等式}
\begin{theorem}[均值不等式（AM--GM）]
对正实数 $a_1,\ldots,a_n$，
\[
\frac{a_1+\cdots+a_n}{n}\ge\sqrt[n]{a_1a_2\cdots a_n},
\]
等号成立当且仅当 $a_1=a_2=\cdots=a_n$。
\end{theorem}
\begin{proof}
$n=2$ 的情形：$(a-b)^2\ge0$ 展开得 $a^2-2ab+b^2\ge0$，即 $a^2+b^2\ge2ab$，令 $a=\sqrt x$、$b=\sqrt y$ 得 $\dfrac{x+y}{2}\ge\sqrt{xy}$。一般情形可用归纳或詹森不等式证明，此处从略；等号条件由 $(a-b)^2=0$ 直接看出。
\end{proof}

\begin{theorem}[柯西不等式]
对实数 $a_1,\ldots,a_n$ 与 $b_1,\ldots,b_n$，
\[
\left(\sum_{i=1}^n a_ib_i\right)^2\le\left(\sum_{i=1}^n a_i^2\right)\left(\sum_{i=1}^n b_i^2\right),
\]
等号成立当且仅当两组数成比例（$a_i=\lambda b_i$ 或 $b_i=\lambda a_i$）。
\end{theorem}
\begin{proof}
考虑关于 $t$ 的二次函数
\[
f(t)=\sum_{i=1}^n(a_it-b_i)^2
=t^2\sum a_i^2-2t\sum a_ib_i+\sum b_i^2.
\]
它是平方和，故对一切实数 $t$ 有 $f(t)\ge0$，即判别式非正：
\[
4\left(\sum a_ib_i\right)^2-4\left(\sum a_i^2\right)\left(\sum b_i^2\right)\le0.
\]
整理即得。等号成立要求 $f$ 有实根 $t_0$，即 $a_it_0-b_i=0$ 对所有 $i$ 成立，两组数成比例。
\end{proof}

\begin{teach}[配方法的示范价值]
柯西不等式的证明是“构造一个非负二次函数”的标准样板。这个套路在第三部分（证明组合数不等式）、第四部分（证明概率不超过 $1$）会一再出现。建议把“要证 $X\le Y$，就构造一个平方和”作为一条独立的方法教给学生。
\end{teach}

\section{分层例题与解析}
\begin{problem}{例 1：韦达定理求对称式}
设 $x^2-5x+3=0$ 的两根为 $x_1,x_2$，求 $x_1^2+x_2^2$ 与 $\dfrac1{x_1^2}+\dfrac1{x_2^2}$。
\end{problem}
\hint 用 $s=x_1+x_2=5$、$p=x_1x_2=3$。

\sol $x_1^2+x_2^2=s^2-2p=25-6=19$。又
\[
\frac1{x_1^2}+\frac1{x_2^2}=\frac{x_1^2+x_2^2}{x_1^2x_2^2}=\frac{19}{9}.
\]
注意 $p=3\ne0$，两个倒数都有定义。

\begin{problem}{例 2：均值不等式求最值}
求 $f(x)=x+\dfrac4x$ 在 $x>0$ 时的最小值。
\end{problem}
\hint 直接用两项的 AM--GM，注意等号条件。

\sol 由 AM--GM，$x+\dfrac4x\ge2\sqrt{x\cdot\dfrac4x}=4$，等号当且仅当 $x=\dfrac4x$，即 $x=2$。故最小值为 $4$，在 $x=2$ 处取得。\textbf{必须检查等号条件}：若定义域不含 $2$，就不能说最小值是 $4$。

\begin{pitfall}[忘了检查等号条件]
求 $x+\dfrac1x$ 在 $x\in[2,3]$ 的最小值：AM--GM 给出下界 $2$，但等号在 $x=1$ 处，不在 $[2,3]$ 内。实际最小值在 $x=2$ 处，为 $2.5$。凡是写“最小值为 $\cdots$”，都要跟一句“等号在 $x=\cdots$ 处成立，且该点在定义域内”。
\end{pitfall}

\begin{problem}{例 3：柯西不等式}
设 $a,b,c>0$ 且 $a+b+c=1$，证明 $\dfrac1a+\dfrac1b+\dfrac1c\ge9$。
\end{problem}
\hint 把左边看成 $\left(\dfrac1{\sqrt a}\right)^2+\cdots$ 与 $(\sqrt a)^2+\cdots$ 的乘积。

\sol 由柯西不等式，
\[
\left(\frac1a+\frac1b+\frac1c\right)(a+b+c)
=\left[\left(\frac1{\sqrt a}\right)^2+\left(\frac1{\sqrt b}\right)^2+\left(\frac1{\sqrt c}\right)^2\right]
 \left[(\sqrt a)^2+(\sqrt b)^2+(\sqrt c)^2\right]
\]
\[
\ge\left(\frac1{\sqrt a}\sqrt a+\frac1{\sqrt b}\sqrt b+\frac1{\sqrt c}\sqrt c\right)^2=3^2=9.
\]
左边又等于 $\dfrac1a+\dfrac1b+\dfrac1c$（因为 $a+b+c=1$），故得证。等号在 $a=b=c=\dfrac13$ 时成立。

\begin{problem}{例 4：判别式判定无解}
证明：不存在实数 $x$ 使 $x^2-4x+7=0$。
\end{problem}
\hint 算判别式；或者配方。

\sol $\Delta=16-28=-12<0$，故无实根。也可以配方：$x^2-4x+7=(x-2)^2+3\ge3>0$，左边恒正，不可能为 $0$。两种方法都要会；配方在“证明恒正”时更常用。

\chapter{函数}
\begin{chapterintro}[章节导读]
函数是“从一个集合到另一个集合的对应规则”。竞赛里的函数题分两类：一类是解析性质（单调、奇偶、周期、值域），一类是迭代与不动点。本章把初等函数的性质整理成可直接调用的工具，并在最后一节给出迭代问题的基本框架——它是第四部分概率与第五部分生成函数的共同基础。
\end{chapterintro}

\section{函数的定义与表示}
\begin{definition}[函数]
从集合 $A$ 到集合 $B$ 的函数 $f$ 是一个规则：对每个 $x\in A$，指定唯一的 $y\in B$ 与之对应，记作 $f(x)$。$A$ 叫定义域，$\{f(x):x\in A\}$ 叫值域。
\end{definition}
\textbf{“唯一”是函数的全部要点}。$y^2=x$ 不是 $x$ 的函数（一个 $x$ 对应两个 $y$），$y=|x|$ 是。竞赛里“判断是否为函数”的题，都是考这一条。

\begin{definition}[单射、满射、双射]
$f$ 是\textbf{单射}：$x_1\ne x_2\Rightarrow f(x_1)\ne f(x_2)$（不同的输入得到不同的输出）。$f$ 是\textbf{满射}：值域等于整个 $B$（每个输出都被取到）。$f$ 既是单射又是满射叫\textbf{双射}。
\end{definition}

\begin{tip}[有限集上单射自动是双射]
若 $A$、$B$ 都是有限集且元素个数相同，$f:A\to B$ 是单射就自动是双射——这就是抽屉原理。第二部分第 2 章证明欧拉定理时，“乘以 $a$ 是简化剩余系上的一一对应”只验证了单射那一半，靠的就是这条。
\end{tip}

\section{单调性、奇偶性、周期性}
\begin{itemize}
\item \textbf{单调递增}：$x_1<x_2\Rightarrow f(x_1)\le f(x_2)$；\textbf{严格递增}：$x_1<x_2\Rightarrow f(x_1)<f(x_2)$。单调递减同理。
\item \textbf{奇函数}：$f(-x)=-f(x)$，图像关于原点对称。\textbf{偶函数}：$f(-x)=f(x)$，图像关于 $y$ 轴对称。
\item \textbf{周期函数}：存在 $T>0$ 使 $f(x+T)=f(x)$ 对一切 $x$ 成立，最小的这样的 $T$ 叫最小正周期。
\end{itemize}
\begin{teach}[单调性怎么用]
单调性最常见的用法是“把函数值的不等式搬回自变量”：若 $f$ 严格递增，则 $f(a)<f(b)\iff a<b$。第二部分第 4 章求乘法阶、第四部分分析“最优策略随参数变化”时，都是这个套路。反函数存在也要求单调。
\end{teach}

\section{一次、二次与反比例函数}
\subsection{一次函数}
$f(x)=kx+b$（$k\ne0$）严格单调，$k>0$ 递增、$k<0$ 递减，与 $x$ 轴交于 $-\dfrac bk$。

\subsection{二次函数}
$f(x)=ax^2+bx+c$（$a\ne0$）配方得
\[
f(x)=a\left(x+\frac b{2a}\right)^2+\frac{4ac-b^2}{4a^2}.
\]
$a>0$ 时开口向上，在 $x=-\dfrac b{2a}$ 处取最小值 $\dfrac{4ac-b^2}{4a^2}$；$a<0$ 时相反。顶点、对称轴、判别式三件套必须熟练。

\subsection{反比例函数}
$f(x)=\dfrac kx$（$k\ne0$）在 $(-\infty,0)$ 与 $(0,+\infty)$ 上分别单调，但\textbf{在整个定义域上不单调}——这是分段讨论的典型场合。

\section{指数函数与对数函数}
\begin{itemize}
\item $a^x$（$a>0,a\ne1$）：$a>1$ 时严格递增，$0<a<1$ 时严格递减；值域 $(0,+\infty)$；恒过 $(0,1)$。
\item $\log_a x$：$a^x$ 的反函数；$a>1$ 时严格递增，值域 $\R$；恒过 $(1,0)$。
\end{itemize}
两条常用的比较：
\[
a>1,\ x>0\ \text{时}\ a^x>1;\qquad
a>1,\ x>1\ \text{时}\ \log_a x>0.
\]
\begin{tip}[算法竞赛里的对数]
估计复杂度时常用两条经验：$\log_2 10^9\approx30$，$\log_2(10^{18})\approx60$。所以“$10^9$ 以内的数二分查找约 $30$ 次”是可靠的第一近似。$\log\log n$ 增长极慢，$\log_2 30\approx5$。
\end{tip}

\section{函数的复合与迭代}
\subsection{复合}
$(f\circ g)(x)=f(g(x))$。复合不满足交换律：一般 $(f\circ g)(x)\ne(g\circ f)(x)$。若 $f$ 与 $g$ 都严格递增，则 $f\circ g$ 也严格递增——这个事实在第四部分证明“最优策略可比较”时有用。

\subsection{迭代与不动点}
从 $x_0$ 出发反复作用 $f$：$x_{n+1}=f(x_n)$。若存在 $x^*$ 使 $f(x^*)=x^*$，称 $x^*$ 为\textbf{不动点}。若 $f$ 连续且迭代收敛，则极限只能是不动点。

\begin{theorem}[压缩映射的收敛性]
设 $f$ 在区间 $I$ 上满足 $|f(x)-f(y)|\le q|x-y|$（$0<q<1$），且 $f(I)\subseteq I$。则从任意 $x_0\in I$ 出发，迭代 $x_{n+1}=f(x_n)$ 收敛到唯一的不动点。
\end{theorem}
\begin{proof}[证明思路]
先证 $\{x_n\}$ 是柯西列：$|x_{n+1}-x_n|\le q|x_n-x_{n-1}|\le\cdots\le q^n|x_1-x_0|$，于是对 $m>n$，
\[
|x_m-x_n|\le\sum_{i=n}^{m-1}|x_{i+1}-x_i|
\le(q^n+\cdots+q^{m-1})|x_1-x_0|\le\frac{q^n}{1-q}|x_1-x_0|\to0.
\]
由实数完备性，$x_n\to x^*$。再由 $f$ 连续，$x^*=f(x^*)$。唯一性：若 $x^*,y^*$ 都是不动点，则 $|x^*-y^*|=|f(x^*)-f(y^*)|\le q|x^*-y^*|$，因 $q<1$ 只能 $x^*=y^*$。
\end{proof}

\begin{tip}[迭代在竞赛里的样子]
真正的竞赛题很少让你算迭代极限，而是问“迭代多少步以后会怎样”。典型形态是第二部分第 4 章的幂塔：指数巨大，但余数只有 $m$ 种，所以关键在于\textbf{证明多少次迭代以后进入周期}，而不是求出极限。第五部分的形式幂级数会把这个想法推到“无穷次迭代的和”。
\end{tip}

\section{分层例题与解析}
\begin{problem}{例 1：单调性搬不等式}
设 $f(x)=x^3+x$。解不等式 $f(2x-1)>f(x+3)$。
\end{problem}
\hint $f$ 是否严格递增？递增就可以把 $f$ 去掉。

\sol $f$ 严格递增（$x_1<x_2$ 时 $x_2^3-x_1^3=(x_2-x_1)(x_2^2+x_1x_2+x_1^2)>0$，且 $x_2-x_1>0$）。于是
\[
2x-1>x+3\iff x>4.
\]
解集为 $(4,+\infty)$。

\begin{problem}{例 2：二次函数的最值}
求 $f(x)=x^2-6x+11$ 在 $[0,4]$ 上的最大值与最小值。
\end{problem}
\hint 配方找顶点，再看顶点是否在区间内。

\sol $f(x)=(x-3)^2+2$，顶点 $x=3\in[0,4]$，故最小值为 $f(3)=2$。最大值在离顶点最远的端点处取得：$|0-3|=3>|4-3|=1$，故最大值为 $f(0)=11$。

\begin{problem}{例 3：复合与定义域}
设 $f(x)=\dfrac1{x-1}$，$g(x)=x^2$。求 $(f\circ g)(x)$ 与 $(g\circ f)(x)$ 及其定义域。
\end{problem}
\hint 逐层代入；定义域要逐层检查。

\sol $(f\circ g)(x)=f(x^2)=\dfrac1{x^2-1}$，定义域 $x\ne\pm1$。$(g\circ f)(x)=g\!\left(\dfrac1{x-1}\right)=\dfrac1{(x-1)^2}$，定义域 $x\ne1$。两者不同，\textbf{复合不可交换}。

\begin{problem}{例 4：一个简单迭代}
设 $x_0=1$，$x_{n+1}=\dfrac{x_n}{2}+1$。证明 $x_n\to2$。
\end{problem}
\hint 算 $|x_{n+1}-2|$ 与 $|x_n-2|$ 的关系。

\sol $x_{n+1}-2=\dfrac{x_n}{2}-1=\dfrac{x_n-2}{2}$，故 $|x_{n+1}-2|=\dfrac12|x_n-2|$，于是 $|x_n-2|=2^{-n}|x_0-2|=2^{-n}\to0$，即 $x_n\to2$。这正是压缩映射的最简单例子。

\chapter{数列、求和与数学归纳法}
\begin{chapterintro}[章节导读]
数列是“定义在正整数上的函数”，也是算法分析的基本语言：一个循环跑多少轮、一次递归分几层，答案都是一个数列。本章整理等差数列、等比数列、常用求和方法与一阶线性递推，并把数学归纳法作为证明数列结论的标准工具。
\end{chapterintro}

\section{等差数列与等比数列}
\begin{definition}[等差数列]
若 $a_{n+1}-a_n=d$（常数），则 $\{a_n\}$ 是等差数列，
\[
a_n=a_1+(n-1)d,\qquad S_n=\frac{n(a_1+a_n)}2=\frac{n(2a_1+(n-1)d)}2.
\]
\end{definition}
\begin{definition}[等比数列]
若 $\dfrac{a_{n+1}}{a_n}=q$（常数，$q\ne0$），则 $\{a_n\}$ 是等比数列，
\[
a_n=a_1q^{n-1},\qquad
S_n=\begin{cases}
\dfrac{a_1(1-q^n)}{1-q},&q\ne1,\\[2mm]
na_1,&q=1.
\end{cases}
\]
\end{definition}
\begin{proof}[等比数列求和公式的推导]
设 $S_n=a_1+a_1q+\cdots+a_1q^{n-1}$，则
\[
qS_n=a_1q+a_1q^2+\cdots+a_1q^n=S_n+a_1q^n-a_1,
\]
移项得 $(q-1)S_n=a_1(q^n-1)$。$q\ne1$ 时即得公式；$q=1$ 时显然 $S_n=na_1$。
\end{proof}

\begin{pitfall}[公比为 1 要单独讨论]
公式 $\dfrac{a_1(1-q^n)}{1-q}$ 在 $q=1$ 时分母为零。写代码时若公比可能为 $1$（例如递推 $a_{n+1}=a_n$），必须特判。这个“分母可能为零”的模式在第二部分（求逆元要先查 $\gcd$）会以更严重的形式出现。
\end{pitfall}

\section{常用求和方法}
\subsection{分组求和}
把和式中的项按来源分组，每组分别用公式。例如 $\sum_{k=1}^n(k^2+2k)$ 分成 $\sum k^2$ 与 $2\sum k$。

\subsection{错位相减（差分抵消）}
形如 $\sum a_kb_k$，其中 $\{a_k\}$ 等差、$\{b_k\}$ 等比。推导：
\[
S_n=\sum_{k=1}^n a_kb_k,\qquad qS_n=\sum_{k=1}^n a_kb_{k+1},
\]
两式相减，右边的项两两错位相消。

\subsection{裂项相消}
若 $a_k=f(k)-f(k+1)$，则 $\sum_{k=1}^n a_k=f(1)-f(n+1)$。最常见的形式：
\[
\frac1{k(k+1)}=\frac1k-\frac1{k+1},\qquad
\frac1{\sqrt{k+1}+\sqrt k}=\sqrt{k+1}-\sqrt k.
\]

\begin{teach}[裂项的三个步骤]
\begin{enumerate}
\item 猜分母（或根式）的差分结构；
\item 待定系数：设 $\dfrac{1}{k(k+2)}=\dfrac{A}{k}+\dfrac{B}{k+2}$，解得 $A=\dfrac12$、$B=-\dfrac12$；
\item 写出来验证前两项与最后两项是否抵消。
\end{enumerate}
要求学生每步都写“抵消后剩下 $f(1)-f(n+1)$”，而不是直接写答案。
\end{teach}

\subsection{递推求和的识别}
看到 $\sum_{k=1}^n f(k)$ 而 $f$ 有递推关系（例如 $f(k)=f(k-1)+\cdots$），优先尝试\textbf{裂项}或\textbf{引出一个辅助和式}。第三部分第 4 章的经典数列大量使用这一招。

\section{一阶线性递推}
\begin{theorem}[一阶线性递推的通项]
设 $x_{n+1}=ax_n+b$（$a\ne0$），$x_0$ 给定。若 $a\ne1$，令 $c=\dfrac{b}{1-a}$，则
\[
x_n-c=a^n(x_0-c),
\]
即 $x_n=c+a^n(x_0-c)$。若 $a=1$，则 $x_n=x_0+nb$。
\end{theorem}
\begin{proof}
设 $y_n=x_n-c$，则
\[
y_{n+1}=x_{n+1}-c=ax_n+b-c=a(y_n+c)+b-c=ay_n+(ac+b-c).
\]
取 $c=\dfrac{b}{1-a}$ 即使 $ac+b-c=0$，于是 $y_{n+1}=ay_n$，即 $y_n=a^ny_0$。代回即得。
\end{proof}

\begin{tip}[不动点换元的通用性]
上面令 $c=\dfrac{b}{1-a}$，$c$ 正是不动点方程 $x=ax+b$ 的解。\textbf{“减去不动点把线性递推变成等比数列”是一个通用技巧}：凡是形如 $x_{n+1}=f(x_n)$ 且 $f$ 是线性分式变换的，都可以先找不动点再换元。第四部分概率里的很多期望递推也是这个形状。
\end{tip}

\section{数学归纳法的应用}
归纳法在数列里的标准用法是“猜通项—证通例”。操作流程：
\begin{enumerate}
\item 写出前若干项，观察规律，猜出 $a_n$ 的表达式；
\item 用归纳法证明：$n=1$ 验证；假设 $a_k$ 已知，从递推式推出 $a_{k+1}$ 的表达式。
\end{enumerate}

\begin{problem}{统一练习：斐波那契数列}
设 $F_1=F_2=1$，$F_{n+2}=F_{n+1}+F_n$。证明 $F_1+F_2+\cdots+F_n=F_{n+2}-1$。
\end{problem}
\hint 对 $n$ 归纳；归纳步只需要用递推式，不需要通项。

\sol $n=1$ 时 $F_1=1=F_3-1$，成立。设结论对 $n$ 成立，则
\[
F_1+\cdots+F_n+F_{n+1}=F_{n+2}-1+F_{n+1}=F_{n+3}-1,
\]
即结论对 $n+1$ 成立。归纳完成。注意这里完全不需要 $F_n$ 的通项公式——\textbf{能用递推就不要用通项}。

\section{分层例题与解析}
\begin{problem}{例 1：错位相减}
求 $S_n=1\cdot2+2\cdot2^2+3\cdot2^3+\cdots+n2^n$。
\end{problem}
\hint 写 $S_n$ 与 $2S_n$，相减。

\sol
\[
S_n=\sum_{k=1}^n k2^k,\qquad
2S_n=\sum_{k=1}^n k2^{k+1}=\sum_{k=2}^{n+1}(k-1)2^k.
\]
相减：
\[
S_n-2S_n=\sum_{k=1}^n k2^k-\sum_{k=2}^{n+1}(k-1)2^k
=2+\sum_{k=2}^{n}(k-(k-1))2^k-(n)2^{n+1}
=2+(2^2+\cdots+2^n)-n2^{n+1}.
\]
于是 $-S_n=2+2^{n+1}-4-n2^{n+1}=(1-n)2^{n+1}-2$，故
\[
S_n=(n-1)2^{n+1}+2.
\]
验算 $n=1$：$1\cdot2=2=(1-1)2^2+2$ ✓；$n=2$：$2+8=10=(2-1)2^3+2$ ✓。

\begin{problem}{例 2：裂项}
求 $\sum_{k=1}^n\dfrac1{k(k+2)}$。
\end{problem}
\hint 待定系数拆成两个分式之差。

\sol 设 $\dfrac1{k(k+2)}=\dfrac A{k}+\dfrac B{k+2}$，通分得 $1=A(k+2)+Bk=(A+B)k+2A$，故 $A=\dfrac12$、$B=-\dfrac12$。于是
\[
\sum_{k=1}^n\frac1{k(k+2)}
=\frac12\sum_{k=1}^n\left(\frac1k-\frac1{k+2}\right)
=\frac12\left(1+\frac12-\frac1{n+1}-\frac1{n+2}\right).
\]
注意\textbf{剩下的是前两项与后两项}，不是首尾各一项——这是裂项时最容易错的地方。

\begin{problem}{例 3：一阶线性递推}
设 $a_1=3$，$a_{n+1}=2a_n+1$。求 $a_n$。
\end{problem}
\hint 找不动点。

\sol 不动点方程 $x=2x+1$ 给出 $x=-1$。令 $b_n=a_n+1$，则
\[
b_{n+1}=a_{n+1}+1=2a_n+2=2(a_n+1)=2b_n,
\]
$b_1=4$，故 $b_n=4\cdot2^{n-1}=2^{n+1}$，$a_n=2^{n+1}-1$。验算 $a_1=2^2-1=3$ ✓，$a_2=2\cdot3+1=7=2^3-1$ ✓。

\begin{problem}{例 4：归纳法证明不等式}
证明：对一切正整数 $n$，$2^n\ge n+1$。
\end{problem}
\hint 对 $n$ 归纳；归纳步用 $2^{n+1}=2\cdot2^n$。

\sol $n=1$ 时 $2\ge2$ ✓。设 $2^n\ge n+1$，则
\[
2^{n+1}=2\cdot2^n\ge2(n+1)=2n+2\ge n+2,
\]
最后一步因 $n\ge1$。归纳完成。

\begin{teach}[数列章的离堂检查]
给一道“求和 $\sum_{k=1}^{n}\dfrac{k}{(k+1)!}$”的小题（提示：$\dfrac{k}{(k+1)!}=\dfrac1{k!}-\dfrac1{(k+1)!}$），要求写出裂项的两步与最终剩余项。做不出就回到裂项小节重做例 2。
\end{teach}

\chapter{复数与平面向量初步}
\begin{chapterintro}[章节导读]
复数在算法竞赛里出场不多，但它是两条重要路线的入口：一是旋转与几何（复数乘法就是旋转加缩放），二是单位根与分治（第五部分的快速傅里叶变换全靠 $n$ 次单位根）。本章只讲竞赛用得着的那一部分：平面向量、复数四则运算、单位根。
\end{chapterintro}

\section{平面向量}
\begin{definition}[平面向量]
平面向量是有大小和方向的量，记作 $\vec v=(x,y)$ 或粗体 $\mathbf v$。两个向量相等当且仅当它们的分量分别相等。
\end{definition}
运算：
\[
\mathbf u+\mathbf v=(u_1+v_1,u_2+v_2),\qquad
\lambda\mathbf v=(\lambda v_1,\lambda v_2),
\]
\[
\mathbf u\cdot\mathbf v=u_1v_1+u_2v_2=|\mathbf u||\mathbf v|\cos\theta,
\]
其中 $\theta$ 是两向量夹角。由第二式得两条常用结论：
\[
\mathbf u\perp\mathbf v\iff \mathbf u\cdot\mathbf v=0,\qquad
\cos\theta=\frac{\mathbf u\cdot\mathbf v}{|\mathbf u||\mathbf v|}.
\]
\textbf{点积为 0 是判定垂直的最快方法}，在计算几何题里比斜率比较可靠（斜率在竖直时会爆炸）。

\section{复数及其运算}
\begin{definition}[复数]
形如 $z=a+b\mathrm i$ 的数叫复数，其中 $a,b\in\R$，$\mathrm i^2=-1$。$a$ 叫实部，$b$ 叫虚部。$\overline z=a-b\mathrm i$ 叫共轭复数，$|z|=\sqrt{a^2+b^2}$ 叫模。
\end{definition}
四则运算按 $\mathrm i^2=-1$ 展开：
\[
(a+b\mathrm i)+(c+d\mathrm i)=(a+c)+(b+d)\mathrm i,
\]
\[
(a+b\mathrm i)(c+d\mathrm i)=(ac-bd)+(ad+bc)\mathrm i.
\]
共轭的两条性质：
\[
z\overline z=|z|^2,\qquad \overline{z+w}=\overline z+\overline w,\qquad \overline{zw}=\overline z\,\overline w.
\]
由第一条得到\textbf{除法的标准做法}——分子分母同乘分母的共轭：
\[
\frac{z_1}{z_2}=\frac{z_1\overline{z_2}}{|z_2|^2}\quad(z_2\ne0).
\]

\begin{pitfall}[$|z_1+z_2|$ 不等于 $|z_1|+|z_2|$]
复数也有三角不等式：$|z_1+z_2|\le|z_1|+|z_2|$，等号仅当两复数同向（辐角相同）时成立。把它当成等式是最常见的错误。
\end{pitfall}

\section{复数的几何意义与单位根}
把 $z=a+b\mathrm i$ 看成平面上的点 $(a,b)$，则
\[
z=r(\cos\theta+\mathrm i\sin\theta),\qquad r=|z|,\ \theta=\arg z,
\]
这叫复数的\textbf{三角形式}。两个复数相乘，模相乘、辐角相加：
\[
r_1(\cos\theta_1+\mathrm i\sin\theta_1)\cdot r_2(\cos\theta_2+\mathrm i\sin\theta_2)
=r_1r_2\bigl(\cos(\theta_1+\theta_2)+\mathrm i(\theta_1+\theta_2)\bigr).
\]
于是\textbf{乘以一个模为 $1$ 的复数就是一次旋转}——这是复数在几何题里最有用的性质。

\begin{definition}[$n$ 次单位根]
满足 $z^n=1$ 的复数叫 $n$ 次单位根。它们恰好是
\[
\omega_k=\mathrm e^{2\pi\mathrm ik/n}=\cos\frac{2\pi k}{n}+\mathrm i\sin\frac{2\pi k}{n},
\qquad k=0,1,\ldots,n-1.
\]
\end{definition}
\begin{theorem}[单位根的两条基本性质]
设 $\omega=\mathrm e^{2\pi\mathrm i/n}$，则
\[
1+\omega+\omega^2+\cdots+\omega^{n-1}=0\quad(n\ge2),
\]
且 $\omega^k\ne1$ 对 $1\le k<n$ 成立（即 $\omega$ 的阶恰为 $n$）。
\end{theorem}
\begin{proof}
第一条：$(\omega^n-1)=(\omega-1)(1+\omega+\cdots+\omega^{n-1})$，而 $\omega^n=1$ 且 $\omega\ne1$，故第二个因子为 $0$。这条证明正是第二部分第 2 章“等比数列求和”的直接应用。

第二条：若 $\omega^k=1$ 且 $1\le k<n$，则 $\cos\dfrac{2\pi k}{n}=1$，即 $\dfrac{2\pi k}{n}$ 是 $2\pi$ 的倍数，于是 $n\mid k$，与 $1\le k<n$ 矛盾。
\end{proof}

\begin{tip}[单位根为什么重要]
第五部分的多项式乘法需要“把两个次数为 $n$ 的多项式相乘从 $O(n^2)$ 降到 $O(n\log n)$”，核心就是取 $2n$ 次单位根作为求值点：单位根处求值有大量重复结构，可以分治。第二部分第 4 章讲“阶”的时候也会再用到 $\omega$ 的阶是 $n$ 这一事实。
\end{tip}

\section{分层例题与解析}
\begin{problem}{例 1：复数运算}
计算 $(1+\mathrm i)^2$ 与 $\dfrac{1+\mathrm i}{1-\mathrm i}$。
\end{problem}
\hint 直接展开；除法用共轭。

\sol $(1+\mathrm i)^2=1+2\mathrm i+\mathrm i^2=2\mathrm i$。又
\[
\frac{1+\mathrm i}{1-\mathrm i}=\frac{(1+\mathrm i)^2}{(1-\mathrm i)(1+\mathrm i)}=\frac{2\mathrm i}{2}=\mathrm i.
\]
注意 $\dfrac{1+\mathrm i}{1-\mathrm i}=\mathrm i$ 意味着“乘以 $\mathrm i$ 就是旋转 $90^\circ$”。

\begin{problem}{例 2：单位根求和}
设 $\omega=\mathrm e^{2\pi\mathrm i/5}$，求 $1+\omega+\omega^2+\omega^3+\omega^4$。
\end{problem}
\hint 用定理中的求和公式。

\sol 由定理，直接得 $0$。

\begin{problem}{例 3：旋转}
把点 $(1,0)$ 绕原点逆时针旋转 $60^\circ$，求新坐标。
\end{problem}
\hint 旋转 $60^\circ$ 就是乘以 $\cos60^\circ+\mathrm i\sin60^\circ$。

\sol 乘以 $\dfrac12+\dfrac{\sqrt3}{2}\mathrm i$，得
\[
1\cdot\left(\frac12+\frac{\sqrt3}{2}\mathrm i\right)=\frac12+\frac{\sqrt3}{2}\mathrm i,
\]
即新坐标为 $\left(\dfrac12,\dfrac{\sqrt3}{2}\right)$。

\begin{problem}{例 4：模长计算}
设 $|z|=2$，求 $|z+\overline z|$ 的最大值。
\end{problem}
\hint 设 $z=a+b\mathrm i$，则 $z+\overline z=2a$。

\sol $z+\overline z=2a=2\operatorname{Re}(z)$，而 $|a|\le|z|=2$，故 $|z+\overline z|=2|a|\le4$。最大值 $4$ 在 $z=\pm2$ 时取得。

\begin{teach}[复数章的取舍建议]
如果课时紧张，本章可以只讲第 1、2 节（向量点积与复数四则运算），把单位根留给第五部分讲快速傅里叶变换时再引入。但要注意：第二部分第 4 章讲“原根”时用的是完全不同的乘法结构（模 $m$ 的整数乘法），不要与这里的单位根混淆——一个是“旋转”，一个是“模幂周期”。
\end{teach}


% ============================================================================
%  第二部分  数论
% ============================================================================
\BeginPart{2}{5}{II}{第二部分}{数论}
% !TeX program = xelatex
% 本文件由 assemble.py 自 OI-Docs 原稿抽取生成，请勿手改（会被覆盖）。
% 源文件：Number Theory/number_theory_teacher.tex
% 源行号：169--3761（已删 3756-3761 的 \appendix 与页码重置）
% 说明：已删去独立成册所需的导言区与页码重置；章节正文与交叉引用原样保留。

% 第 0 章为预备章：把全书的术语与记号翻译成高中语言
\setcounter{chapter}{-1}
\chapter{数学用词与使用方法说明}
\begin{chapterintro}[章节导读]
后面的数论章节会默认使用约数、同余、可逆、赋值、卷积、双射等术语。先把这些词和记号统一起来，后面遇到新算法时就可以把注意力放在数学本身，而不是被术语卡住。
\end{chapterintro}

\begin{tip}[这一章解决什么问题]
作者在后面的章节默认你会说一种“数学黑话”：约数、模、可逆、赋值、卷积、双射、记忆化、阈值……这些词每一个似乎都能用一句话解释清楚，但正文一般不会专门说。本章把全书会用到、而高中课本不一定会说的词与记号一次交代清楚，并在末尾给出一张“高中叫法 $\leftrightarrow$ 本书叫法”的对照表和一段的证明示范。

本章不教新的数学知识。你以后在阅读正文遇到哪个词把你卡住了，就可以回到这里查它到底是什么意思。
\end{tip}

\section{怎么读这本书}
% 本书有三个版本，内容按“同一条主线、三种深浅”组织：

%\begin{itemize}
%\item \textbf{基础版}：定理、定义、例题与提示，不展开证明。适合第一遍上课用，或你还没学过某个定理时先建立整体印象。
%\item \textbf{完整版}：在基础版之上补上证明、解析、易错点与代码。适合自学与备赛。
%\item \textbf{教师版}：完整版 + 课堂设计、追问与评量。
%\end{itemize}

正文里有四种带颜色的框，它们分工不同：

\begin{itemize}
\item \textbf{定义/定理/引理/推论/例}：黑色小标题，是必须记住的结论。看书时先问自己“条件是什么、结论是什么”，再往下看。
\item \textbf{题目框}（深蓝色标题）：题目与它的解析。标“统一练习”的不是原题，而是把几道题共用的想法抽出来。
\item \textbf{易错框}（橙棕色标题，标注“易错点”）：这里写的是“大多数人会错在哪”，以及一个具体反例。跳过它最省时间，也最亏。
\item \textbf{本章新增的绿色 tip 框}：只解释一个词、一步跳步或一个例子。它可以单独读懂，不需要你先看懂上下文。
\end{itemize}

\begin{tip}[一句话学习方法]
每读完一节，合上书回答三个问题：这一节要解决什么问题？它的条件少一条会怎样（举出反例）？我能把结论写成几行代码吗？三个问题答不上来，就说明只是“看过了”，不是“会了”。
\end{tip}

\section{集合、元素与逻辑用语}
高中课本已经给了你绝大部分记号，这里只把本书用得多的补齐。

\subsection{集合与元素}
集合就是“一些确定的东西放在一起”，记 $a\in A$ 读作“$a$ 是 $A$ 的元素”，$a\notin A$ 读作“不是”；$A\subseteq B$ 读作“$A$ 的每个元素都在 $B$ 里”（$A$ 含于 $B$）；$\varnothing$ 是空集，即不含任何元素的集合。两个集合相等，当且仅当它们互相含于。

本书最常用的几个数集：
\[
\Z=\{\ldots,-2,-1,0,1,2,\ldots\},\qquad \N=\{0,1,2,\ldots\},\qquad \Pp=\{2,3,5,7,11,\ldots\},
\]
分别是整数集、非负整数集、素数集。$\mathbb Q$、$\mathbb R$ 是有理数集与实数集，本书只在说明“范围放宽了”时提到。注意本书约定 $\N$ 含 $0$；如果你所在的教材不含 $0$，读到 $\N$ 时就按本书的约定来。

描述法的写法 $\{x\in\Z:0\le x<m\}$ 表示“满足后面条件的那些 $x$”，它等于 $\{0,1,\ldots,m-1\}$。冒号也可写成竖线。$\{0,1,\ldots,m-1\}$ 有 $m$ 个元素，别把“从 $0$ 到 $m-1$”数成 $m-1$ 个。

\subsection{“或”“且”“当且仅当”}
\begin{itemize}
\item “$P$ 且 $Q$”：两者同时成立，符号 $P\land Q$。
\item “$P$ 或 $Q$”：至少一个成立（可以两个都成立），符号 $P\lor Q$。数学里的“或”不是“二选一”。
\item “$P\Rightarrow Q$”：$P$ 成立时 $Q$ 必须成立，读作“$P$ 推出 $Q$”。这时也说“$P$ 是 $Q$ 的\textbf{充分}条件”“$Q$ 是 $P$ 的\textbf{必要}条件”。
\item “$P\iff Q$”（$P\Leftrightarrow Q$）：互相推出，读作“$P$ 当且仅当 $Q$”，也读作“$P$ 是 $Q$ 的充要条件”。
\end{itemize}

\begin{tip}[“必要性来自……充分性来自……”是在说什么]
这是正文证明里最常见的句式，翻译成人话就是：“两个方向分开证”。写 $P\iff Q$ 时，
\textbf{必要性}指“已知 $P$ 成立，去证 $Q$”（不做到的话 $P$ 就不成立，所以 $Q$ 是\emph{必须}的）；
\textbf{充分性}指“已知 $Q$ 成立，去造出 $P$”（$Q$ 一条就\emph{足够}了）。
一个方向没证，结论就还只是猜测。本书第 0 章第 11 节会完整示范一次双向证明，第 1 章的裴蜀定理也给了不跳步的版本。
\end{tip}

\subsection{量词：任意与存在}
$\forall$ 读“对任意（所有）”，$\exists$ 读“存在（至少有一个）”，$\exists!$ 读“存在且唯一”。顺序改变含义：
\[
\forall n\in\N,\ \exists p\in\Pp,\ p>n\qquad\text{（素数无穷多个）},
\]
$\forall$ 在前表示“先给我 $n$，我再找 $p$”，$p$ 可以依赖 $n$。若写成 $\exists p,\ \forall n$，就变成“有一个固定素数比所有数都大”，荒谬。

否定规则要背熟：$\lnot(\forall x,P(x))$ 等价于“存在一个 $x$ 使 $P(x)$ 不成立”，即 $\exists x,\lnot P(x)$。因此要推翻“所有情况都好”，只需\textbf{举一个反例}；反之，要反驳“存在坏情况”，必须证明“全都好”。

\subsection{Iverson 括号与“示性函数”}
本书写 $\iva{\gcd(i,n)=1}$，读作“Iverson 括号”：方括号里是一句话（命题），成立时这个式子的值为 $1$，不成立时为 $0$。例如 $\iva{5>3}=1$，$\iva{4\text{ 是奇数}}=0$。

它的好处是把“只在某些情况下计入”写成算式，从而能自由交换求和顺序：
\[
\sum_{i=1}^{n}\iva{i\text{ 与 }n\text{ 互质}}\ =\ \text{“}1\text{ 到 }n\text{ 中与 }n\text{ 互质的数的个数”}.
\]
看到 $\iva{\cdot}$ 别慌，它只是一个“条件成立记 $1$、否则记 $0$”的开关。

\section{求和、取整、下标：把式子读成程序}
\subsection{求和与求积}
$\sum_{i=1}^{n}a_i=a_1+a_2+\cdots+a_n$，$\prod_{i=1}^{n}a_i=a_1a_2\cdots a_n$。下标从 $1$ 还是 $0$ 开始、到 $n$ 还是 $n-1$ 结束，务必按式子读；差一项在很多题里就是 $100$ 分与 $0$ 分的差距。空和（一项都没有）约定为 $0$，空积约定为 $1$，这是为了让递推式在边界也成立。

换序（“交换求和”）是本书用得最多的技术，合法条件是：项数有限。有限项相加，先算哪一项都行。示例：
\[
\sum_{d=1}^{n}\ \sum_{\substack{k\ge 1\\dk\le n}}\ 1\ =\ \sum_{m=1}^{n}\ \sum_{d\mid m}\ 1,
\]
左边是“对每个 $d$ 数它有多少个不超过 $n/d$ 的倍数 $k$”，右边是“对每个 $m$ 数它有多少个约数 $d$”；两边数的是同一批满足 $dk\le n$ 的整数对 $(d,k)$，所以相等。\textbf{遇到换序，先像这样把“数的是什么”说清楚}，比背公式可靠。

\subsection{取整}
$\floor{x}$ 是“不超过 $x$ 的最大整数”（向下取整，floor 是地板），$\ceil{x}$ 是“不小于 $x$ 的最小整数”（向上取整）。例如 $\floor{7/2}=3$，$\ceil{7/2}=4$，$\floor{-7/2}=-4$（注意不是 $-3$：$-4$ 才是不超过 $-3.5$ 的最大整数）。

两条最常用性质：对整数 $a$ 与正整数 $b$，
\[
\floor{\frac{a}{b}}=q\iff qb\le a<(q+1)b,\qquad \ceil{\frac{a}{b}}=q\iff (q-1)b<a\le qb.
\]
第二条也可以写成 $\ceil{a/b}=\floor{(a+b-1)/b}$（$a\ge0$ 时）。C++ 的 `/` 是\textbf{向零取整}（截断），不是向下取整：$-7/2=-3\ne\floor{-3.5}$。所以本书模板里所有“整除取商”都写成 \texttt{floor\_div}（见附录“核心模板与接口”与第 1 章）。

\subsection{下标、递推与“记号说明”}
$a_i$ 读“$a$ 下标 $i$”，是数列的第 $i$ 项（本书默认从 $a_1$ 开始，除非写了 $a_0$）。“递推式”是用前一项表示后一项，例如 $T_h=a^{T_{h-1}}$；要它可用，必须同时给\textbf{起点}（$T_0=1$）和\textbf{方向}（$h$ 递增）。看到“类似地”“同理”“不难验证”，但你自己验不出来，就把它当成一个待填的空。

\section{整数：约数、倍数、最大公约数与素因数分解}
\begin{itemize}
\item \textbf{整除}：$a\mid b$ 读“$a$ 整除 $b$”，意思是“存在整数 $k$ 使 $b=ka$”。此时 $a$ 是 $b$ 的\textbf{约数}（高中课本常叫“因数”），$b$ 是 $a$ 的\textbf{倍数}。规定 $0\mid 0$ 成立（取 $k$ 任意），而 $a\mid b$ 中若 $a=0,b\ne0$ 则不成立。符号 $\nmid$ 表示“不整除”。
\item \textbf{公约数与最大公约数}：同时整除 $a,b$ 的正整数叫公约数，其中最大的叫 $\gcd(a,b)$（greatest common divisor）。本书也写 $(a,b)$ 表示同一个东西。
\item \textbf{公倍数与最小公倍数}：同时被 $a,b$ 整除的正整数中最小的叫 $\lcm(a,b)$。有 $\gcd(a,b)\lcm(a,b)=ab$（$a,b>0$）。
\item \textbf{互质}：$\gcd(a,b)=1$，即“公约数只有 $1$”。注意：$1$ 与 $1$ 互质；$0$ 与 $1$ 互质；$0$ 与 $5$ 不互质（$\gcd(0,5)=5$）。两个互质的数不必都是素数（$8$ 与 $9$）。
\item \textbf{素数与合数}：大于 $1$ 且约数只有 $1$ 与自身的叫素数（本书与“质数”混用，同义）；大于 $1$ 又不是素数的叫合数。$1$ 既不是素数也不是合数。
\item \textbf{唯一分解定理}：每个大于 $1$ 的整数都能写成素数的乘积，且不计顺序时写法唯一，$n=p_1^{e_1}\cdots p_t^{e_t}$。于是“$n$ 的约数”与“指数向量 $(f_1,\ldots,f_t)$，$0\le f_j\le e_j$”一一对应，约数个数为
$\tau(n)=\prod_j(e_j+1)$。本书把原稿用的 $d(n)$ 统一记作 $\tau(n)$，避免与 $\gcd$ 的括号写法混淆。
\item \textbf{素因数（本书常写“素因子”）}：出现在分解式里的素数。$\rad(n)=\prod_{p\mid n}p$ 是所有不同素因子的乘积。
\item \textbf{平方自由}：分解式中每个指数都 $\le1$，即不被 $4,9,25,\ldots$ 这类平方数整除。例如 $30=2\cdot3\cdot5$ 平方自由，$12=2^2\cdot3$ 不是。
\item \textbf{素数幂}：形如 $Q=p^\alpha$ 的数（$p$ 为素数，$\alpha\ge1$）。本书大量结论先在素数幂下证明，再用中国剩余定理拼回一般模数，这条路线第 3 章起会反复出现。
\end{itemize}

\begin{tip}[“素因子”与“质因数”是同一个东西]
中学课本说“分解质因数”，大学教材与竞赛资料常说“素因子”“素因数”。本书四个词混用，含义相同。类似地，“最大公约数”与“最大公因数”相同。看到新词先怀疑它是不是老词的别名，这是读数学资料省时间的诀窍。
\end{tip}

\begin{tip}[为什么要关心“平方自由”]
因为很多计数公式里会出现一个开关：$\mu(n)$（莫比乌斯函数）。它的定义只有一行——$n=1$ 时为 $1$；$n$ 平方自由且有偶数个素因子时为 $1$、奇数个时为 $-1$；一旦有平方因子就是 $0$。所以“平方自由”不是趣味知识，而是“哪些项会消失”的判断条件。第 7 章与第 8 章大量使用它。
\end{tip}

\section{同余：按余数把整数分堆}
$a\equiv b\pmod m$ 读作“$a$ 与 $b$ 模 $m$ 同余”，定义就是 $m\mid a-b$。它只关心“除以 $m$ 的余数是否相同”，所以两个数同余，等价于它们落在同一堆里。把整数按余数分成 $m$ 堆，这 $m$ 堆合起来称为\textbf{完全剩余系}；每一堆里挑一个代表写出来的数叫\textbf{代表元}，本书默认用 $0,1,\ldots,m-1$（最小非负代表元）。

为什么同余式能加减乘？因为 $(a\pm b)\bmod m$ 只由 $a,b$ 各自的余数决定；$m\mid a-b$ 且 $m\mid c-d$ 时 $m\mid(a+c)-(b+d)$、$m\mid ac-bd$。所以“每步取模”不改变结果。\textbf{除法不行}：除法是同余式里唯一的陷阱，它要求先检查可除性，且往往要把模数一起缩小，见第 2 章的约分定理与易错框。

两个必记词：
\begin{itemize}
\item \textbf{模运算的“负数”}：$-1\bmod 7=6$，不是 $-1$。规范做法是先取余再加模数。C++ 里写成 \texttt{((a\%m)+m)\%m}；本书模板中的 \texttt{norm(a,m)} 干的就是这件事。
\item \textbf{单位（可逆元）}：若存在 $x$ 使 $ax\equiv1\pmod m$，称 $a$ 是模 $m$ 的\textbf{单位}，$x$ 叫 $a$ 的\textbf{逆元}。$m=7$ 时 $1$ 到 $6$ 都是单位；$m=8$ 时只有 $1,3,5,7$ 是单位（偶数不行，因为 $\gcd$ 不为 $1$）。这些单位凑在一起，就是下面 0.8 节说的“单位群”。
\end{itemize}

\begin{tip}[“剩余系”这三个字到底在说什么]
把 $0$ 到 $m-1$ 排成一圈，只保留余数信息：$m=5$ 时的剩余系是 $\{0,1,2,3,4\}$。其中与 $5$ 互质的有 $\{1,2,3,4\}$，这就是“简化剩余系”，它的个数记 $\varphi(5)=4$。$m=8$ 时全部 $8$ 个余数里，与 $8$ 互质的是 $\{1,3,5,7\}$，所以 $\varphi(8)=4$。“乘以 $2$”在 $\{0,1,\ldots,7\}$ 上不是一一对应，因为 $1\cdot2$ 与 $5\cdot2$ 都是 $2$；但“乘以 $3$”在四个单位 $\{1,3,5,7\}$ 上是：$1\to3,3\to1,5\to7,7\to5$，每个余数恰好出现一次——这就是一次“双射”的手算。这个“在互质的那几格里可以一一对应”的事实，就是第 2 章证明欧拉定理时最关键的一步。
\end{tip}

\section{幂、阶乘与组合数}
$a^n$ 表示 $n$ 个 $a$ 相乘，$a^0=1$（包括 $0^0=1$，本书按组合意义约定）。指数律 $a^{m}a^{n}=a^{m+n}$、$(a^m)^n=a^{mn}$ 在整数指数下要额外注意 $a=0$ 与负指数。

$n!=1\cdot2\cdots n$，并约定 $0!=1$。$\binom{n}{k}=\dfrac{n!}{k!(n-k)!}$ 是“从 $n$ 个里选 $k$ 个”的方法数，高中课本里的写法也是 $\mathrm{C}_n^k$；本书用 $\binom{n}{k}$，因为分式形式在“约掉公共因子”时更好用。二项式定理：$(x+y)^n=\sum_{k}\binom{n}{k}x^{k}y^{n-k}$。

\textbf{幂塔}（“高德纳箭号”的最小情形）：$a^{b^{c}}$ 按\textbf{右结合}读，等于 $a^{(b^c)}$，而不是 $(a^b)^c=a^{bc}$。$3^{3^{3}}=3^{27}$ 已经约 $7.6\times10^{12}$，所以处理幂塔的核心矛盾是：指数巨大，而余数只有 $m$ 种可能。第 2 章的“扩展欧拉定理”与“截断塔接口”就是为此准备的。

\begin{tip}[“降幂”是什么意思]
“降幂”不是初中因式分解里的“降次”。在数论里，它指：要算 $a^E\bmod m$，但 $E$ 大到写不下来，于是把 $E$ 换成一个更小、余数不变的指数（例如 $E\bmod\varphi(m)$）。这个替换\textbf{只在特定条件下合法}，条件不满足时必须保留“$E$ 是否已经很大”这一位信息。第 2 章用 $2^4\bmod 8$ 这个反例演示了乱降幂的错误。
\end{tip}

\begin{tip}[$\nu_p(n)$：一个只有“数因子”动作的记号]
$\nu_p(n)$ 读作“$n$ 的 $p$-进赋值”，做法是：拿 $n$ 不断除以 $p$，能除几次就是几。例如 $\nu_2(24)=3$（$24\to12\to6\to3$，停了），$\nu_3(24)=0$，$\nu_5(100)=2$。它满足两条：$\nu_p(xy)=\nu_p(x)+\nu_p(y)$（乘法变加法），以及 $\nu_p(x+y)\ge\min(\nu_p(x),\nu_p(y))$（加法只能“更好”，除非两项赋值相同才可能大幅抵消）。第 5 章的升幂引理（LTE）本质上就是“什么时候恰好等于最小值”的判据。看到“赋值”这个词，先在心里换成“这个素因子出现了几次”。
\end{tip}

\section{数论函数、前缀和与卷积}
\textbf{数论函数}只是“定义在正整数上的函数”$f:\N_{>0}\to\Z$（或 $\mathbb R$）：$\varphi(n)$ 数互质个数、$\tau(n)$ 数约数个数、$\mu(n)$ 是上面那个开关。别把它想得太高深，它其实就是一张表：$f(1),f(2),f(3),\ldots$

\begin{itemize}
\item \textbf{积性}与\textbf{完全积性}：若 $\gcd(a,b)=1$ 时总有 $f(ab)=f(a)f(b)$，叫 $f$ 是\textbf{积性}的；若不要求互质也成立，叫\textbf{完全积性}的。例：$\id_t(n)=n^t$ 完全积性；$\varphi$、$\tau$、$\mu$ 只是积性（拿 $f=\tau$，$a=b=2$ 试：$\tau(4)=3\ne\tau(2)\tau(2)=4$，一眼看出“互质”这条不能丢）。有了积性，只要知道 $f$ 在每个素数幂上的值，整张表就全知道了。
\item \textbf{狄利克雷卷积}：给两个数论函数定义一种新乘法
$(f*g)(n)=\sum_{d\mid n}f(d)g(n/d)$。它不是多项式乘法，也不是逐点相乘。手算一次就懂：$f=\one$（常一函数，$\one(n)\equiv1$）、$g=\id_1(n)=n$，则 $(\one*\id_1)(6)=1\cdot6+1\cdot3+1\cdot2+1\cdot1$ 按 $d=1,2,3,6$ 取 $\id_1(6/d)$，等于 $6+3+2+1=12$，正好是 $6$ 的约数和。本书用 $\sum_{d\mid n}$ 表示“对 $n$ 的每个约数 $d$ 求和”。
\item \textbf{单位函数} $\eps$：$\eps(1)=1$、$\eps(n)=0$（$n>1$）。它扮演“数字 $1$”的角色：$f*\eps=f$。于是“卷积逆”就是满足 $f*f^{-1}=\eps$ 的那个 $f^{-1}$——可以把它完全类比成“倒数”，只是这里的乘法换成了上面的卷积。
\item \textbf{前缀和}：$S_f(n)=\sum_{i=1}^{n}f(i)$，即“表的前 $n$ 项之和”，约定 $S_f(0)=0$。第 9 章“杜教筛”要解决的问题只有一个：$n$ 高达 $10^{10}$，表建不出来，怎么仍算出 $S_f(n)$。
\end{itemize}

\begin{tip}[为什么要费劲定义一种“新乘法”]
因为很多数论恒等式会变成“乘法等式”。经典例子：$\sum_{d\mid n}\varphi(d)=n$，用卷积语言就是 $\one*\varphi=\id_1$。一旦写成这样，两边同时“乘 $\mu$”（$\mu$ 是 $\one$ 的卷积逆）就自动得到 $\varphi=\mu*\id_1$，反演公式不再需要灵光一现。第 7 章与第 8 章的整套机械就是把这个类比推到底。初学时你要掌握的只有三件事：怎么按定义展开、$*$ 满足交换律与结合律、$\one$ 的逆是 $\mu$。
\end{tip}

\section{“群、环、域、同态”：这一页只需听懂，不必学会}
本书不要求你会群论，但正文会用这些词压缩语言。下面是够用的最小版本：

\begin{itemize}
\item \textbf{群}：一个集合 $G$ 加一种运算（写成乘法），满足：任取两元素运算结果还在 $G$ 里（封闭）；$(ab)c=a(bc)$（结合律）；有“什么都不做”的元素 $e$（单位元）；每个元素都有能把结果变回 $e$ 的逆元。例：$\Z$ 与加法、非零有理数与乘法。
\item \textbf{子群}：$G$ 的一部分，自己也是群。例：$\Z$ 里的偶数在加法下是子群。
\item \textbf{阶}：有限群的元素个数叫群的阶；元素 $a$ 的阶是使 $a^r=e$ 的最小正整数 $r$。本书第 4 章的“乘法阶”就是它在模 $m$ 乘法下的具体版本。
\item \textbf{循环群}：整个群由一个元素不断自乘得到（如模 $7$ 的单位 $\{1,2,3,4,5,6\}$，取 $3$：$3,2,6,4,5,1$ 恰好遍历）。那个能遍历全体的元素就叫\textbf{原根}。
\item \textbf{单位群}：$(\Z/m\Z)^{\times}$ 表示“模 $m$ 的可逆余数全体，配上乘法”。它的阶是 $\varphi(m)$。第 2 章说“乘以 $a$ 是单位群上的双射”，意思就是“在这张余数表上乘同一个可逆数，元素只是换了位置，一个不多一个不少”。
\item \textbf{环}：可以做加、减、乘（不必能除）的体系，如 $\Z$、模 $m$ 的余数集合。
\item \textbf{域}：环上还额外保证“非零元素都能除”，如 $\mathbb Q$、$\mathbb R$，以及素数模下的余数 $\mathbb F_p=\Z/p\Z$。这就是为什么“$d$ 次多项式至多 $d$ 个根”在 $\mathbb F_p$ 成立、在 $\Z/8\Z$ 不成立（$x^2=1$ 模 $8$ 有 $4$ 个根）。
\item \textbf{同态/同构}：“两个体系结构相同”的精确说法。同态是保持运算的映射 $f(xy)=f(x)f(y)$；同构还是一一对应。本书只在“指数把乘法变成加法”这类场合用一次，见第 4 章。
\end{itemize}

\begin{tip}[遇到“群”字时的实用策略]
先把它换成“一组余数 + 一种运算”，再问：这些余数是谁（是不是都得跟模数互质）？运算是取模后的普通乘法吗？90\% 的场合这么一换，句子就懂了。真需要群论工具时（第 12、14 章的少数段落），正文会把要用的性质单独陈述成引理，不会拿“由群论知”糊弄你。
\end{tip}

\section{算法行话：复杂度、预处理、在线、对拍、溢出}
\begin{itemize}
\item \textbf{大 $O$}：$T(n)=O(f(n))$ 表示存在常数 $C$ 与 $n_0$，使 $n\ge n_0$ 时 $T(n)\le C f(n)$。它只关心增长快慢，不承诺绝对时间，也不管常数。本书还写 $O,\Omega,\Theta$（下界、同阶）与“期望 $O(\cdot)$”（随机算法在平均意义下的界）。
\item \textbf{字长模型}：竞赛里默认 $64$ 位整数加减乘比较都是 $O(1)$。因此“$n$ 位数”的高精度乘法是 $O(n^2)$，而不是 $O(1)$。本书说某算法“多项式时间”时，变量是位数 $\log n$；$O(n^{1/4})$ 关于位数\textbf{不是}多项式。
\item \textbf{预处理 / 查询}：预处理是开始前先算好一张表，查询是每次问一行。“$O(N)$ 预处理、$O(1)$ 查询”与“$O(1)$ 预处理、$O(\log N)$ 查询”是两种不同取舍，总时间相同也可能一个内存爆一个不爆。本书凡给复杂度都分开报这两项。
\item \textbf{离线 / 在线}：在线是“来一问答一问，不许先看后面”；离线是“所有问题一次读完，可以重排、分块、按模数分组”。同一道题，两种口径的最优做法可能完全不同。
\item \textbf{记忆化}：递归时把算过的结果存进哈希表或数组，下次直接查。它把“指数级重复计算”压成“状态数个”，是第 9 章杜教筛不可缺的一步。
\item \textbf{摊还}：单步可能很贵，但 $m$ 步的总代价除以 $m$ 很小。典型：并查集、单调栈。别把“摊还 $O(1)$”读成“每步 $O(1)$”。
\item \textbf{对拍}：写一个笨但显然正确的程序（暴力枚举），与要提交的程序在同一批随机数据上比答案。它是“找反例”的工程化版本，本书 \texttt{validation/} 与 \texttt{code/selftest.cpp} 干的就是这件事。
\item \textbf{阈值}：某个“从此以后行为不同”的数值。例如扩展欧拉定理里的 $\varphi(m)$：指数不到它不能乱换，超过它才能换。阈值不是“经验值”，它由证明里的不等式给出。
\item \textbf{溢出}：$10^9\cdot10^9=10^{18}$ 还能塞进 \texttt{long long}（上限约 $9.22\times10^{18}$），$(2\cdot10^9)^2$ 就超了。本书模板用 \texttt{\_\_int128} 做中间量、用 \texttt{mul\_mod} 收口；\texttt{unsigned} 溢出是回绕而非未定义行为，别把两者混为一谈。
\item \textbf{取模的正确姿势}：减法后可能为负，要 \texttt{norm}；除法不是“分子分母各自取模”，而是“乘分母的逆元”，且分母必须先确认与模数互质。
\end{itemize}

\section{高中叫法 $\leftrightarrow$ 本书叫法}
\begin{longtable}{>{\raggedright\arraybackslash}p{0.23\textwidth}>{\raggedright\arraybackslash}p{0.23\textwidth}>{\raggedright\arraybackslash}p{0.46\textwidth}}
\toprule
高中课本 / 常见叫法 & 本书叫法 & 差别与注意\\\midrule
因数 & 约数 & 同义；本书的“约数”含 $1$ 与自身，且只在正整数范围讨论。\\
辗转相除法 & 欧几里得算法 / gcd & 同一种算法；本书强调复杂度来自“每两轮规模至少减半”。\\
$(a,b)=1$ & $\gcd(a,b)=1$、互质 & 本书偶尔也用 $(a,b)$ 表 gcd，出现时会说明。\\
分解质因数 & 素因数分解 & 同义；本书还会写“拆素因子”“按素数幂拆开”。\\
“$a$ 除以 $m$ 余 $r$” & $a\bmod m=r$、$a\equiv r\pmod m$ & 前者是运算、后者是关系；负数时本书规定余数落在 $[0,m)$。\\
排列、组合 $\mathrm{C}_n^k$ & $\dbinom{n}{k}$ & 同义；本书用分式形式，便于讨论“分母能否取逆”。\\
二项式定理 & $(1+X)^n$ 的系数 & 同一件事；第 6 章用“模 $p$ 后只剩两项”推出 Lucas 定理。\\
数学归纳法 & 归纳 / 最小反例 & “取最小反例”与归纳等价，本书两种都说。\\
等差数列 & 环上的固定步长 & 在 $0,1,\ldots,n-1$ 上“每次加 $k$”就是绕圈走的等差数列。\\
等比数列 & 几何数列 / 幂序列 & 第 12 章的“指数向量”研究它的周期。\\
充分必要条件 & 当且仅当、$\iff$ & 双向都要证，见 0.2 节的 tip。\\
集合、$\in$、$\subseteq$、$\varnothing$ & 同高中 & 本书默认 $\N$ 含 $0$。\\
对数 $\log$ & $\log$、$\log\log$ & 无特别说明时本书大 $O$ 里的 $\log$ 底数不影响结论；$\log\log$ 增长极慢。\\
绝对值 $|x|$ & 同 & 赋值符号 $\nu_p$ 与绝对值无关，别混淆。\\
函数 & 数论函数（一张表） & 见 0.6 节。\\
“显然” & 偶尔用于省略一层代数 & 你若验算不出来，就当它是作业。\\
\bottomrule
\end{longtable}

\section{证明的三种句式，以及“不跳步”长什么样}
\subsection{双向证明}
要证“$P$ 当且仅当 $Q$”，写两个方向、各起一段，并给每段一个开头词（“必要性.”“充分性.”）。若只证了一边，就把结论降格为“$P$ 推出 $Q$”。

\subsection{反证与最小反例}
要证“所有 $n$ 都满足 $P(n)$”，可设存在反例，取\textbf{最小的}那个 $n_0$，然后由 $n_0$ 造出更小的反例 $n'<n_0$，矛盾。这就是“最小反例法”，它与数学归纳法等价。本书第 1 章证明“欧几里得算法正确”、第 15 章的韦达跳跃都用它，因为它们都需要“$n_0$ 的最小性”这一件武器。

\subsection{归纳}
写清三件事：命题 $P(n)$ 是什么；$P(1)$（或 $P(0)$）成立；“若对所有 $k<n$ 成立则 $P(n)$ 成立”。需要“所有更小的”而不是只有 $n-1$ 时，叫第二归纳法，本书会说“强归纳”。

\subsection{示范：把一段跳步的证明补全}
正文第 1 章的裴蜀定理原文只写了两句“必要性来自……充分性来自……”。下面是它应有的样子——请把它当模板，而不是当额外负担：

\paragraph{示范命题.}设 $a,b$ 为整数，不全为零，$d=\gcd(a,b)>0$，$c$ 为整数。则 $ax+by=c$ 有整数解 $\iff$ $d\mid c$。

\begin{proof}
必要性（有解 $\Rightarrow$ $d\mid c$）。假设存在整数 $x,y$ 使 $ax+by=c$。由 $d$ 是公约数，可写 $a=da'$、$b=db'$，其中 $a',b'$ 是整数。代入得
\[
c=ax+by=da'x+db'y=d(a'x+b'y).
\]
括号里是整数（整数对加、乘封闭），所以 $d$ 整除 $c$。这一步没有用到任何“显然”的性质，只是把 $d$ 提出来。

充分性（$d\mid c$ $\Rightarrow$ 有解）。分两小步。

第一步：先证 $d$ 本身能被写成 $a,b$ 的整数线性组合，即存在 $u,v$ 使 $au+bv=d$。对较小的数归纳：若 $b=0$，则 $d=a$，取 $u=1,v=0$ 即可；若 $b>0$，写 $a=qb+r$、$0\le r<b$。由归纳假设用于更小的对 $(b,r)$（因为 $\gcd(b,r)=\gcd(a,b)=d$，见第 1 章定理），存在 $u_1,v_1$ 使 $bu_1+rv_1=d$。把 $r=a-qb$ 代回去：
\[
d=bu_1+(a-qb)v_1=av_1+b(u_1-qv_1),
\]
于是可取 $u=v_1$、$v=u_1-qv_1$，它们都是整数。归纳完成。

第二步：把右边的 $d$ 放大成 $c$。因为 $d\mid c$，记 $c=kt$（$t$ 为整数）。把第一步的等式两边同乘 $t$：
\[
a(ut)+b(vt)=dt=c.
\]
取 $x=ut$、$y=vt$ 即为所求。注意方向：是“乘 $c/d$”，不是“乘 $d/c$”——原稿在这里写过反，第 1 章的解析会再强调一次。
\end{proof}

\begin{tip}[怎么判断自己“真的读懂了一个证明”]
拿一张白纸，只看定理陈述，自己把证明重建一遍。卡住的地方就是你不懂的地方。第二遍再看书，只读你卡住的那两行。这样读十页，比顺着看五十页有效。
\end{tip}

\section{自查小测}
每题答案不超过两行。做错的题，回到对应小节。

\begin{enumerate}
\item $\floor{-7/2}$ 与 C++ 里 \texttt{-7/2} 的值各是多少？
\item 写出“$\forall x\exists y\,P(x,y)$”的否定。
\item $\tau(72)$、$\mu(72)$、$\nu_2(72)$ 分别是多少？（$72=2^3\cdot3^2$）
\item $-1\bmod 7$ 是多少？$2$ 在模 $8$ 下有逆元吗？在模 $7$ 下呢？
\item $17^{3^{2}}$ 与 $(17^3)^2$ 相等吗？
\item 为什么说“$0!=1$”不任性，而是为了递推 $\binom{n}{0}=1$ 与空积约定一致？
\item 验证“$2$ 在模 $8$ 下没有逆元”，并说出用的是哪一条判别条件。
\item “$f$ 是积性函数”与“$f$ 是完全积性函数”，哪个更强？用 $\tau$ 说明。
\item 一句话说清“预处理”和“查询”的区别，并给一个本书的例子。
\item “取最小反例”与哪种证明方法等价？
\end{enumerate}

\begin{tip}[参考答案]
1) 都是 $-4$ 与 $-3$。2) $\exists x\forall y\,\lnot P(x,y)$。3) $12,0,3$。4) $6$；没有（$\gcd(2,8)=2$）；有，$2^{-1}\equiv4\pmod 7$。5) 不相等，$17^{9}\ne17^{6}$——幂塔右结合。6) 若 $0!=x$，由 $1!=1\cdot0!$ 得 $x=1$。7) 逐个检查 $2x\bmod 8$（$x=0,\ldots,7$）得 $0,2,4,6,0,2,4,6$，从不等于 $1$，故无逆元；依据是“可逆 $\iff$ 与模数互质”，而 $\gcd(2,8)=2$。8) 完全积性更强；$\tau(2\cdot2)=3\ne4$。9) 预处理是开始前建表、查询是每次读表，如第 2 章 $O(m^{2/3})$ 预处理 + $O(1)$ 查询。10) 数学归纳法。
\end{tip}


\chapter{整除、最大公约数与线性不定方程}
\begin{chapterintro}[章节导读]
上一章我们先把本书会反复出现的数学语言统一了；这一章开始真正进入数论。我们先从“一个数能不能整除另一个数”出发，用最大公约数整理这种关系，再进一步把它变成可以求解的线性不定方程。
\end{chapterintro}

\section{从质因数分解理解整除}
对正整数 $n=\prod_p p^{\alpha_p}$，只有有限个 $\alpha_p$ 非零。若 $a=\prod p^{u_p},b=\prod p^{v_p}$，则
\[
 a\mid b\iff\forall p,\ u_p\le v_p,\qquad
 \gcd(a,b)=\prod_p p^{\min(u_p,v_p)},\quad
 \lcm(a,b)=\prod_p p^{\max(u_p,v_p)}.
\]

\begin{tip}[把整数看成\textbf{指数向量}]
唯一分解让你把 $n$ 记成一串指数：$72\leftrightarrow(3,2,0,0,\ldots)$ 对应素数 $2,3,5,7,\ldots$。于是\textbf{$a\mid b$}变成\textbf{每一位都不超过}，\textbf{$\gcd$}变成逐位取小，\textbf{$\lcm$}变成逐位取大。整章的许多结论都能这样\textbf{逐位}证明——看到 $\prod p^{u_p}$ 就先想这个画面。
\end{tip}

因此 $\gcd(a,b)\lcm(a,b)=ab$。实际计算最小公倍数时先除后乘：$a/\gcd(a,b)\cdot b$，仍要检查乘法是否溢出。


\begin{teach}[导入：把“共同部分”分离出来]
在黑板上写 $12=2^2\cdot3$、$18=2\cdot3^2$，请学生分别用“每种素数取较小指数”和“取较大指数”得到 $6,36$。
接着引出最常用换元
\[
 d=\gcd(a,b),\quad a=du,\ b=dv,\quad\gcd(u,v)=1.
\]


换元最后一句不能省略。许多数论题不是缺少定理，而是没有把公共因子从互质部分中分离出来。
\end{teach}

\begin{tip}[为什么：先除后乘]
\textbf{先除后乘}是为了不溢出：$a/\gcd(a,b)\cdot b$ 与 $a\cdot b/\gcd(a,b)$ 数学上相同，但后者的中间值 $ab$ 可能超过 $9.22\times10^{18}$。写成代码时，两个 $10^9$ 级整数相乘就已经贴着上限。
\end{tip}



\section{欧几里得算法及其复杂度}
\begin{theorem}
若 $b>0$，则 $\gcd(a,b)=\gcd(b,a\bmod b)$；终止条件为 $\gcd(a,0)=a$（$a\ge0$）。
\end{theorem}
\begin{proof}
先固定一次带余除法：$a=qb+r$，其中 $q=\floor{a/b}$、$0\le r<b$（$b>0$ 时这样的 $q,r$ 存在且唯一）。

第一步（公约数不会变少）。设整数 $c$ 同时整除 $a$ 与 $b$，写 $a=ca'$、$b=cb'$。由 $r=a-qb$ 得
\[
 r=ca'-q(cb')=c(a'-qb'),
\]
括号中是整数，所以 $c\mid r$。于是 $c$ 也同时整除 $b$ 与 $r$。

第二步（公约数不会变多）。设 $c\mid b$ 且 $c\mid r$，写 $b=cb''$、$r=cr''$。由 $a=qb+r$ 得
\[
 a=q(cb'')+cr''=c(qb''+r''),
\]
故 $c\mid a$，即 $c$ 同时整除 $a$ 与 $b$。

第三步（取最大）。前两步说明\textbf{$a$ 与 $b$ 的公约数集合}和\textbf{$b$ 与 $r$ 的公约数集合}互为子集，因此完全相同。两个集合都非空（$1$ 在内）且有限（公约数不超过 $\min(a,b)$），故最大元素相同，这正是 $\gcd(a,b)=\gcd(b,r)$。

终止条件：$r=0$ 时 $\gcd(a,0)=a$，因为整除 $a$ 与 $0$ 的数就是 $a$ 的约数（$0$ 被任何非零整数整除），其中最大的为 $a$。
\end{proof}
若 $a\ge b>0$，则 $a\bmod b<a/2$：当 $b\le a/2$ 时余数小于 $b$；否则商为 $1$，余数为 $a-b<a/2$。所以每两轮规模至少减半，复杂度为 $O(1+\log\min(a,b))$。更细地，令 $d=\gcd(a,b)$，可写成
\[
 O\!\left(1+\log\min\left(\frac ad,\frac bd\right)\right).
\]

\begin{tip}[口径：\textbf{每两轮规模至少减半}]
\textbf{每两轮规模至少减半}是把两行不等式合起来说的：一轮里 $a$ 变成 $b$（$b$ 可能只比 $a$ 小一点点，比如 $a=b+1$），但下一轮的余数一定小于 $b/2$。所以\textbf{规模}指两个数中的较大者，它每两轮才保证减半。
\end{tip}



\paragraph{手算例.} 对 $(252,198)$ 依次得到
\[
252=198+54,\quad198=3\cdot54+36,\quad54=36+18,\quad36=2\cdot18.
\]
故 gcd 为 $18$。这里不是每一轮两个数都减半，而是隔两轮观察同一位置上的数；此区别是讲解复杂度时的常见含糊点。
\paragraph{追问.} 连续斐波那契数为什么让欧几里得算法递归较深？因为每一步商多为 $1$，规模缩减相对缓慢。可以让学生实际记录 $(89,55)$ 的调用链。


\section{裴蜀定理与扩展欧几里得}
\begin{theorem}[裴蜀定理]
设整数 $a,b$ 不全为零，$d=\gcd(a,b)>0$。方程 $ax+by=c$ 有整数解，当且仅当 $d\mid c$。更一般地，$\sum_i a_ix_i=c$ 有解当且仅当 $\gcd(a_1,\ldots,a_n)\mid c$。
\end{theorem}

\begin{proof}[证明]
先证必要性。假设存在整数 $x,y$ 使 $ax+by=c$。由于 $d=\gcd(a,b)$，所以 $d\mid a$ 且 $d\mid b$，可写成 $a=da'$、$b=db'$。于是
\[
 c=ax+by=d(a'x+b'y),
\]
括号里的数是整数，因此 $d\mid c$。

再证充分性。首先证明更强的结论：存在整数 $u,v$ 使
\[
 au+bv=d.
\]
对欧几里得算法的递归深度归纳。当 $b=0$ 时，$d=\gcd(a,0)=a$，取 $u=1,v=0$ 即可。当 $b>0$ 时写 $a=qb+r$，其中 $0\le r<b$。由 $\gcd(a,b)=\gcd(b,r)=d$，对更小的二元组 $(b,r)$ 应用归纳假设，存在整数 $u_1,v_1$ 使
\[
 bu_1+rv_1=d.
\]
又 $r=a-qb$，代入得到
\[
 d=bu_1+(a-qb)v_1=av_1+b(u_1-qv_1),
\]
因此取 $u=v_1$、$v=u_1-qv_1$ 即得 $au+bv=d$。

最后，由 $d\mid c$，设 $t=c/d\in\mathbb Z$。将 $au+bv=d$ 两边乘以 $t$，得
\[
 a(ut)+b(vt)=dt=c.
\]
所以 $x=ut$、$y=vt$ 即为整数解。

多元情形对 $n$ 归纳。$n=2$ 已证。$n>2$ 时令 $g=\gcd(a_1,\ldots,a_{n-1})$。由归纳假设，可将前 $n-1$ 项表示为 $g$；于是原方程先化为二元方程 $gt+a_nx_n=c$。又 $\gcd(g,a_n)=\gcd(a_1,\ldots,a_n)=d$，故由二元情形可解。将得到的 $t$ 乘回前 $n-1$ 个系数，即得到原方程的整数解。
\end{proof}

\begin{pitfall}[思考：这份证明为什么要这样走？]
已知 $d=\gcd(a,b)$ 并且 $d\mid c$。我们真正要证明的是存在整数 $x,y$，使 $ax+by=c$。这里不要一上来猜 $x,y$；先把目标拆开来看。

\paragraph{为什么先证明 $d$ 能写成 $au+bv$？}
如果能找到整数 $u,v$ 使 $au+bv=d$，那么因为 $d\mid c$，就可以写成 $c=dt$，其中 $t=c/d$ 是整数。把 $au+bv=d$ 整体乘以 $t$，马上得到 $a(ut)+b(vt)=c$。所以充分性的真正关键只有一个：为什么 $d$ 一定能写成 $a,b$ 的整数线性组合？

\paragraph{为什么 $d$ 能写成 $au+bv$？}
这里其实就是把欧几里得算法反着走。先看最简单的情况 $b=0$。这时 $d=\gcd(a,0)=a$，所以直接有 $a\times1+0\times0=a=d$。也就是说，取 $u=1,v=0$ 就可以了，这就是递归的起点。

\paragraph{当 $b>0$ 时怎么办？}
对 $a,b$ 做一次带余除法，写成 $a=qb+r$，其中 $0\le r<b$。例如 $a=17,b=5$ 时，就是 $17=3\times5+2$。欧几里得算法最重要的性质之一是 $\gcd(a,b)=\gcd(b,r)$，因此原来的 $d$ 同时也是 $\gcd(b,r)$。

注意这里真正发生了什么：原来我们面对的是 $a,b$，现在变成了 $b,r$；而且 $r<b$，所以问题规模变小了。这就是为什么可以使用归纳假设。

\paragraph{对更小的 $b,r$，归纳假设告诉我们什么？}
因为 $d=\gcd(b,r)$，根据归纳假设，存在整数 $u_1,v_1$，使 $bu_1+rv_1=d$。

这里特别容易看错。我们现在得到的不是 $au_1+bv_1=d$，因为这一层归纳处理的是 $b,r$，所以得到的自然是 $b,r$ 的线性组合。现在的问题只是：组合中还出现了 $r$，而我们的最终目标需要重新只出现 $a,b$。

\paragraph{怎么把 $r$ 换回 $a,b$？}
刚才有 $a=qb+r$，所以 $r=a-qb$。把它代入 $bu_1+rv_1=d$，得到
\[
 d=bu_1+(a-qb)v_1=av_1+b(u_1-qv_1).
\]
现在右边终于重新变成了 $a,b$ 的整数线性组合。因此令 $u=v_1$、$v=u_1-qv_1$，就有 $au+bv=d$。

这一步就是整段证明最核心的“回去”：欧几里得算法先把问题变小，归纳假设解决更小的问题，然后利用 $r=a-qb$ 把得到的结果重新换回原来的 $a,b$。

\paragraph{到这里，裴蜀定理的充分性其实已经基本结束了。}
我们已经证明了：只要 $d=\gcd(a,b)$，就存在整数 $u,v$ 使 $au+bv=d$。接下来只需利用题目的条件 $d\mid c$。因为 $c=dt$，把前面的等式整体乘以 $t$，就是
\[
 a(ut)+b(vt)=dt=c.
\]
所以令 $x=ut$、$y=vt$，原方程就有整数解。

整份二元证明其实可以压缩成这样一条逻辑链：先用欧几里得算法得到 $d=au+bv$，再利用 $d\mid c$ 把这个等式整体放大 $c/d$ 倍。以后看到 $ax+by=c$，不要一上来猜 $x,y$；先算 $d=\gcd(a,b)$，检查 $d\mid c$，如果成立，就求出 $au+bv=d$，最后乘上 $c/d$。

\paragraph{最后，多元情形又是在做什么？}
对于 $a_1x_1+\cdots+a_nx_n=c$，技巧其实非常直接：把前 $n-1$ 项先压缩成一个数。令 $g=\gcd(a_1,\ldots,a_{n-1})$。根据已经证明的低元情形，存在整数系数使
\[
 a_1x_1+\cdots+a_{n-1}x_{n-1}=g.
\]
于是原来的 $n$ 元方程可以先看成二元方程 $gt+a_nx_n=c$。

而 $\gcd(g,a_n)=\gcd(a_1,\ldots,a_n)=d$，所以这又回到了我们刚才已经解决的二元问题。只要 $d\mid c$，就能得到 $t,x_n$。最后把前面的表示式整体乘以 $t$，得到
\[
 a_1(tx_1)+\cdots+a_{n-1}(tx_{n-1})=gt,
\]
再加上 $a_nx_n$，原来的 $n$ 元方程就被恢复出来了。

所以多元情形没有增加新的核心思想：它只是把前 $n-1$ 项压缩成一个 gcd，先解决一个二元问题，再把答案展开回来。
\end{pitfall}

\begin{tip}[为什么：相邻斐波那契数最慢]
直觉：商越大，减得越多。斐波那契数相邻两项的商永远是 $1$，等于每步只做了一次减法，所以最慢。这也解释了为什么竞赛里说\textbf{$\gcd$ 的最坏输入是相邻斐波那契数}。
\end{tip}

这里只用到一个初等事实：把 $au+bv$ 全体（$u,v$ 取遍整数）排成一列，其中最小的正数就是 $\gcd(a,b)$。

令 $a=qb+r$。若递归已求得 $bu+rv=d$，则
\[
 d=av+b(u-qv),\qquad x=v,\ y=u-qv.
\]
因此 exgcd 只需在 gcd 的递归返回时维护两个系数。


对上例回代：
\[
18=54-36=4\cdot54-198=4\cdot252-5\cdot198.
\]
所以 $252x+198y=36$ 的一个解是 $(8,-10)$。
“求 $ax+by=d$”与“求 $ax+by=c$”应分开：前者是算法接口，后者是检查 $d\mid c$ 后进行缩放。原稿把这两步中的比值写反了。
\begin{lstlisting}
using i128 = __int128_t;
i128 exgcd(i128 a, i128 b, i128 &x, i128 &y) {
    // Contract: a,b >= 0, not both zero.
    if (b == 0) { x = 1; y = 0; return a; }
    i128 u, v;
    i128 d = exgcd(b, a % b, u, v);
    x = v;
    y = u - (a / b) * v;
    return d;
}
\end{lstlisting}
允许负系数时，先对 $|a|,|b|$ 求解，再恢复对应系数的符号。不要对最小负整数直接在原类型里取绝对值，应先转更宽类型。


\section{从一个解得到全部解}
当 $a,b>0$ 且 $d\mid c$，已知特解 $(x_0,y_0)$，则全部整数解为
\[
 x=x_0+\frac bd t,\qquad y=y_0-\frac ad t,\qquad t\in\Z.
\]

\begin{tip}[小心：对 \texttt{INT64\_MIN} 取相反数]
因为 $\texttt{INT64\_MIN}=-2^{63}$ 的绝对值不在 \texttt{long long} 里（$2^{63}>9.22\times10^{18}$），直接取相反数是未定义行为。正确顺序：先转成 \texttt{\_\_int128} 或无符号类型，再取绝对值。
\end{tip}

\begin{proof}
先看清要证的是\textbf{恰好}：既要说这个族里的每个 $t$ 都给解，也要说任何解都在这个族里。

反向（族中皆解）。把 $x=x_0+\frac bd t$、$y=y_0-\frac ad t$ 代入左边：
\[
 a\Big(x_0+\frac bd t\Big)+b\Big(y_0-\frac ad t\Big)=ax_0+by_0+\frac{ab}{d}t-\frac{ab}{d}t=c.
\]
两个含 $t$ 的项恰好抵消，所以任意整数 $t$ 都给出解。

正向（任何解都在族中）。设 $(x,y)$ 也是解，记 $\Delta x=x-x_0$、$\Delta y=y-y_0$。两式相减得 $a\Delta x+b\Delta y=0$，两边除以 $d$：
\[
 \frac ad\,\Delta x=-\frac bd\,\Delta y.
\]
因为 $d=\gcd(a,b)$，$\frac ad$ 与 $\frac bd$ 互质。上式表明 $\frac bd\mid\frac ad\Delta x$，由欧几里得引理（$m\mid nk$ 且 $\gcd(m,n)=1$ 时 $m\mid k$）得 $\frac bd\mid\Delta x$。取整数 $t$ 使 $\Delta x=\frac bd t$，代回 $\frac ad\Delta x=-\frac bd\Delta y$ 并约去非零因子 $\frac bd$（这是整数等式，可以直接约）得 $\Delta y=-\frac ad t$。

两个方向合起来即\textbf{全部解恰为该单参数族}。这也解释了为什么参数只有一个：方程只有一个独立约束。
\end{proof}
若限制 $L_x\le x\le R_x$、$L_y\le y\le R_y$，就将四个不等式转为 $t$ 的上下界，取交集。解的计数问题变成一维整数区间计数。

\begin{problem}{P5656：二元一次不定方程（exgcd）}
给定正整数 $a,b,c$，求 $ax+by=c$ 的正整数解个数以及 $x,y$ 的最小值、最大值。若没有正整数解但有整数解，原题还要求分别报告能在整数解中取得的最小正 $x$、最小正 $y$；这两个值不必来自同一组解。
\end{problem}
\hint 先判定 $d\mid c$，再用通解把正整数约束转为参数区间。

\begin{pitfall}[讲解：先看这题到底在做什么]
这题真正的工作分成两层：先用 exgcd 找到一组整数特解，再处理“正整数”这个额外限制。第一层只关心 $d\mid c$ 是否成立；第二层把所有解写成一个参数 $t$ 的形式，于是正整数条件自然变成 $t$ 的一个区间。讲解时不要把这两层混在一起，否则学生容易以为“求 gcd”本身就在解决正整数约束。
\end{pitfall}

\sol 令 $A=a/d,B=b/d$。正整数要求给出
\[
 t\ge t_L=\ceil{\frac{1-x_0}{B}},\qquad
 t\le t_R=\floor{\frac{y_0-1}{A}}.
\]

\begin{tip}[把\textbf{不等式}变成\textbf{一个参数的区间}]
以 $x=x_0+bt\le R_x$ 为例：移项得 $t\le(R_x-x_0)/b$，取整为 $t\le\floor{(R_x-x_0)/b}$；$x\ge L_x$ 给 $t\ge\ceil{(L_x-x_0)/b}$。四个不等式给 $t$ 一个区间 $[t_{\min},t_{\max}]$，解的个数就是 $\max(0,t_{\max}-t_{\min}+1)$。空区间表示\textbf{无正整数解}，但不表示\textbf{无整数解}——这两件事在 P5656 里是不同的输出分支。
\end{tip}

若 $t_L\le t_R$，则解数为 $t_R-t_L+1$，并且
\[
\begin{array}{ll}
x_{\min}=x_0+Bt_L,&x_{\max}=x_0+Bt_R,\\
y_{\min}=y_0-At_R,&y_{\max}=y_0-At_L.
\end{array}
\]
若区间为空，只说明不能同时为正，并不说明无整数解。单独的最小正 $x$ 是 $x_0\bmod B$ 的正代表元（余数 $0$ 改为 $B$）；$y$ 同理取模 $A$。

\paragraph{完整算例.} $6x+9y=60$ 化为 $2x+3y=20$，取 $(x_0,y_0)=(10,0)$，通解 $x=10+3t,y=-2t$。正整数约束为 $-3\le t\le-1$，三组解是
\[
 (1,6),\ (4,4),\ (7,2).
\]

\begin{tip}[词义：正代表元]
\textbf{正代表元}= 把同余类里的那个类写成 $1,\ldots,B$ 之间的数。余数为 $0$ 时不能留 $0$（题目要正整数），所以换成 $B$。这类\textbf{$0$ 换成模数}的小心，在取模计数题里非常常见。
\end{tip}

故答案依次为解数 $3$，$x_{\min}=1,y_{\min}=2,x_{\max}=7,y_{\max}=6$。

\begin{pitfall}[负数整除不等于向下取整]
C++ 的整数除法向零截断。对 $b>0$，设 $q_0$ 为 C++ 中 \texttt{a/b} 的结果、$r_0$ 为 \texttt{a\%b} 的结果；注意 $r_0$ 不一定是本书数学定义中的非负模余数。可以实现
\[
\operatorname{floorDiv}(a,b)=q_0-\iva{r_0<0},\qquad
\operatorname{ceilDiv}(a,b)=q_0+\iva{r_0>0}.
\]
这样还避免了对最小负整数直接求 $-a$。测试 $a=-5,b=3$ 时，应得到 floor 为 $-2$、ceil 为 $-1$。
\end{pitfall}
\cost 每组 $O(\log\min(a,b))$，额外空间为递归深度。$c/d$ 的缩放、区间端点的乘法应使用足够宽的整数。

\src{https://www.luogu.com.cn/problem/P5656}

\begin{problem}{P8255：数学游戏}
给定 $x,z>0$，求满足 $z=xy\gcd(x,y)$ 的最小正整数 $y$，无解则报告无解。原数据规模：$t\le5\times10^5$，$x\le10^9$，$z<2^{63}$。
\end{problem}
\hint 若 $x\nmid z$ 则无解。令 $w=z/x$，考察 $\gcd(w,x^2)$ 是否为完全平方数。

\begin{pitfall}[讲解：先看这题到底在做什么]
先不要急着猜 $y$。题目中的 gcd 可以拆成一个公共因子和两个互质部分；设 $d=\gcd(x,y)$ 后，原式会自然出现一个平方 $d^2$。因此判定的核心变成：从 $z/x$ 里能不能找到一个平方因子，而它恰好就是 $\gcd(z/x,x^2)$。
\end{pitfall}

\sol 令 $d=\gcd(x,y)$，$x=da,y=db$，其中 $\gcd(a,b)=1$。则
\[
 w=\frac zx=bd^2,\qquad x^2=a^2d^2,
\]
所以
\[
 \gcd(w,x^2)=d^2\gcd(b,a^2)=d^2.
\]
注意这里必须是 $x^2$，不是原稿写的 $a^2$。

反过来，若 $g=\gcd(w,x^2)=d^2$ 是完全平方数，则 $d\mid x$、$d^2\mid w$。令 $a=x/d$、$b=w/d^2$，有 $\gcd(a^2,b)=1$，故 $\gcd(a,b)=1$。构造
\[
 y=\frac wd
\]
即满足全部条件。这还证明了\textbf{解若存在就是唯一的}，“最小”不需要额外搜索。

\paragraph{算例与反例.} $x=12,z=864$ 时 $w=72$，$g=\gcd(72,144)=72$ 不是平方，无解。
$x=12,z=1296$ 时 $w=108$，$g=36$，$d=6$，$y=18$，验证 $12\cdot18\cdot6=1296$。

\cost 每组一次 gcd 与整数平方根。可用浮点平方根给出候选，再用整数乘法向两侧校正；不要直接判断浮点数是否为整数。用 $128$ 位验证 $d^2=g$ 和最终乘积。

\src{https://www.luogu.com.cn/problem/P8255}

\section{环上的固定步长}
\begin{theorem}
编号 $0,\ldots,n-1$ 的环，每步加 $k$ 并取模 $n$。设 $d=\gcd(n,k)$，则形成 $d$ 个循环，每个循环长度为 $n/d$，恰对应一个模 $d$ 的同余类。
\end{theorem}

\begin{tip}[图景：围成一圈走 $k$ 步]
把 $n$ 个编号想象成围成一圈的小朋友，每次往右走 $k$ 步。\textbf{形成 $d$ 个循环}就是\textbf{分成 $d$ 个小圈子，圈子之间走不到}。$\gcd(n,k)$ 决定了圈子数，这是\textbf{互质才能走遍全部}的几何版。
\end{tip}

\begin{proof}
把问题写成带余除法的语言：$n$ 个点编号为 $0,1,\ldots,n-1$，从 $i$ 出发走 $t$ 步后位于 $i+tk\pmod n$。

第一步（周期长度）。回到出发点当且仅当 $n\mid tk$。设 $d=\gcd(n,k)$，写 $n=dn'$、$k=dk'$，则 $\gcd(n',k')=1$，条件成为 $dn'\mid t\,dk'$，即 $n'\mid tk'$；再由欧几里得引理得 $n'\mid t$。所以使\textbf{回到起点}成立的正整数 $t$ 恰为 $n'$ 的倍数，最小的是
\[
 t=n'=\frac{n}{\gcd(n,k)}.
\]
这就是每个循环的长度。

第二步（谁和谁在同一个循环）。每步加 $k$，而 $d\mid k$，故编号模 $d$ 的余数始终不变：同一循环中的点全在同一个剩余类里。反过来，同一个剩余类 $\{j:0\le j<n,\ j\equiv i\pmod d\}$ 恰有 $n/d$ 个元素，而由第一步，从 $i$ 出发的循环也恰有 $n/d$ 个点（周期内不重复，否则会得到更小的周期）。前者包含后者且个数相同，所以两者相等：每个剩余类正好是一个循环。

结论：循环数为 $d=\gcd(n,k)$，每个长度 $n/d$。$k=0$ 时 $d=n$，公式给出 $n$ 个长度为 $1$ 的循环，与\textbf{每点自成一环}一致，不必单独讨论。
\end{proof}

\paragraph{课堂小练习.} $n=12,k=8$ 时 $d=4$，四个环分别为
\[
 (0,8,4),\ (1,9,5),\ (2,10,6),\ (3,11,7).
\]
请学生说明：这里的“子环”是置换的循环，不一定是抽象代数里的子环。后文 CF516E 要用这一分组，但“在同一环里”不意味着“距离起点相同”。



\chapter{同余、逆元与指数降幂}
\begin{chapterintro}[章节导读]
上一章，我们用 gcd 和 exgcd 处理整数之间的线性关系；但很多竞赛问题真正关心的不是整数等式，而是除以某个模数以后留下什么。这一章把这种“只看余数”的语言建立起来，并继续讨论逆元、快速幂以及如何把过大的指数压缩下来。
\end{chapterintro}

\section{同余是整数等式的压缩}
对 $m>0$，$a\bmod m$ 是 $[0,m)$ 中与 $a$ 同余的唯一整数；
\[
 a\equiv b\pmod m\iff m\mid a-b.
\]

\begin{tip}[“最小非负代表元”这一句在管什么]
同一堆余数里有无穷多个整数：$-1$、$6$、$13$ 在模 $7$ 下都算\textbf{同一类}。写\textbf{$a\bmod m$ 是 $[0,m)$ 中唯一的整数}就是为了选定一个写法，此后所有比较都只比这个写法。C++ 的 \% 对负数给出负余数，所以本书模板一律用 \texttt{norm(a,m)}。
\end{tip}

$\{0,1,\ldots,m-1\}$ 把可能的余数排成一张表，称为模 $m$ 的\textbf{完全剩余系}；表中与 $m$ 互质的那些元素组成\textbf{简化剩余系}，它的个数记作 $\varphi(m)$，读\textbf{$m$ 的欧拉函数值}。
加、减、乘可以在每一步后取模，交换律、结合律、分配律均保持成立。例如
\[
 (a+b)c\equiv ac+bc\pmod m.
\]

\begin{tip}[剩余系 = 一张余数表]
完全剩余系 = 每个余数出现一次的表；简化剩余系 = 表里只留\textbf{能与模数互质}的那些行。$\varphi(m)$ 就是留下的行数。之后凡是\textbf{在剩余系上乘一个数}\textbf{重排这张表}，都在这里理解。
\end{tip}


\section{乘法逆元与正确约分}
\begin{theorem}
$m\ge2$ 时，$a$ 的模 $m$ 逆元存在当且仅当 $\gcd(a,m)=1$；若存在，则在模 $m$ 意义下唯一。
\end{theorem}

\begin{tip}[把\textbf{除法}翻译成\textbf{解方程}]
所谓 $a^{-1}$，就是找 $x$ 使 $ax\equiv1\pmod m$，也就是找整数对 $(x,y)$ 使 $ax+my=1$。于是\textbf{逆元存不存在}立刻变成\textbf{这个二元一次方程有没有整数解}，答案是\textbf{有当且仅当 $\gcd(a,m)=1$}（第 1 章裴蜀定理）。模 $m=1$ 时同余式恒成立，本书统一约定结果为 $0$。
\end{tip}

\begin{proof}
先把\textbf{逆元}翻译成方程。按同余的定义，
\[
 ax\equiv1\pmod m\iff m\mid ax-1\iff \exists y\in\Z,\ ax+my=1
\]
（取 $y=-(ax-1)/m$ 即得后一式的 $y$；反之把 $my$ 移到右边也一样）。所以\textbf{$a$ 有逆元}与\textbf{$ax+my=1$ 有整数解}是同一句话。

充分性：若 $\gcd(a,m)=1$，由裴蜀定理这样的 $(x,y)$ 存在，于是 $ax\equiv1\pmod m$，逆元存在。
必要性：若有 $x$ 使 $ax\equiv1\pmod m$，设 $g=\gcd(a,m)$。则 $g\mid a$ 与 $g\mid m$ 给出 $g\mid(ax+my)=1$，故 $g=1$。（用到的唯一事实是：$\gcd$ 整除任何整数线性组合。）
唯一性：设 $ax\equiv ax'\equiv1\pmod m$，两式相减得 $a(x-x')\equiv0\pmod m$，即 $m\mid a(x-x')$。因 $\gcd(a,m)=1$，欧几里得引理给 $m\mid x-x'$，即 $x\equiv x'\pmod m$。

所以\textbf{唯一}只是模 $m$ 意义下的唯一；一旦限定 $0\le x<m$，代表元就真正唯一了。$m=1$ 时同余式恒成立，本书把结果统一约定为 $0$，避免与\textbf{不存在}混淆。
\end{proof}
\begin{theorem}[约分后模数的变化]
\[
 ax\equiv bx\pmod m
 \iff a\equiv b\pmod{m/\gcd(x,m)}.
\]
\end{theorem}
\begin{proof}
从定义出发逐步变形。$ax\equiv bx\pmod m$ 意为 $m\mid x(a-b)$，即存在整数 $k$ 使
\[
 x(a-b)=km. \tag{$\ast$}
\]
令 $d=\gcd(x,m)$，写 $x=dx'$、$m=dm'$，则 $\gcd(x',m')=1$。代入 $(\ast)$ 并约去 $d$（这是普通的整数等式，可以约）得
\[
 x'(a-b)=km',
\]
即 $m'\mid x'(a-b)$。由欧几里得引理进一步得 $m'\mid a-b$，也就是 $a\equiv b\pmod{m'}=a\equiv b\pmod{m/d}$。

反向：若 $m\mid a-b$ 显然成立前一式；具体地，$a-b=m's$ 时 $x(a-b)=dx'm's=m(dx's)$，故 $m\mid x(a-b)$。

两个方向都用完了，注意 $(\ast)$ 中的 $k$ 在约分后仍是同一个整数：不要求 $d\mid k$。这正是\textbf{同除一个公因子时，被约的是模数，不是那些商}的原因。
\end{proof}

\begin{pitfall}[把整数除法直接搬到同余式中会出错]
$2\cdot1\equiv2\cdot4\pmod6$，但 $1\not\equiv4\pmod6$。正确结论是 $1\equiv4\pmod3$。
一般若 $d\mid a,b,m$，则 $a\equiv b\pmod m$ 等价于 $a/d\equiv b/d\pmod{m/d}$。对等式 $a=km+b$ 同除 $d$ 得到的是 $a/d=k(m/d)+b/d$，不需要 $d\mid k$。
\end{pitfall}
\paragraph{算例.} $12x\equiv18\pmod{30}$，先除 gcd $6$ 得 $2x\equiv3\pmod5$，于是 $x\equiv4\pmod5$。若在模 $30$ 下列出全部解，则是 $4,9,14,19,24,29$；新模数为 $5$ 并不意味着在模 $30$ 下只有一个解。


\section{费马小定理与欧拉定理}
\begin{definition}
$\varphi(n)=\sum_{1\le i\le n}\iva{\gcd(i,n)=1}$，其中 $\varphi(1)=1$。当 $p$ 为素数时 $\varphi(p)=p-1$。
\end{definition}
\begin{theorem}[欧拉定理与费马小定理]
若 $m\ge2$ 且 $\gcd(a,m)=1$，则 $a^{\varphi(m)}\equiv1\pmod m$。
特别地，素数 $p\nmid a$ 时 $a^{p-1}\equiv1\pmod p$。
\end{theorem}
\begin{proof}
取定 $R=\{r:0\le r<m,\ \gcd(r,m)=1\}$，它有 $\varphi(m)$ 个元素。分四步。

第一步（乘完还在 $R$ 里）。设 $\gcd(a,m)=1$ 且 $\gcd(r,m)=1$。若有素数 $p\mid\gcd(ar,m)$，则 $p\mid m$ 且 $p\mid ar$；$p\nmid a$（否则 $p\mid\gcd(a,m)$），故 $p\mid r$，与 $\gcd(r,m)=1$ 矛盾。所以 $\gcd(ar,m)=1$。

第二步（乘 $a$ 是一一对应）。若 $ar\equiv ar'\pmod m$，则 $m\mid a(r-r')$，由欧几里得引理得 $m\mid r-r'$；而 $r,r'\in[0,m)$，故 $r=r'$。有限集合上的单射自动是双射（抽屉原理：像的元素数与原集合相同，又不能漏，只能恰好覆盖）。于是
\[
 \{\,ar\bmod m:r\in R\,\}=R
\]
作为集合相等，只是顺序被打乱。

第三步（顺序无关，可以整表相乘）。两个集合相同意味着把它们各自的元素全部乘起来结果相同：
\[
 \prod_{r\in R}(ar)\equiv\prod_{r\in R}r\pmod m.
\]
左边重新整理成
\[
 \prod_{r\in R}(ar)=a^{|R|}\prod_{r\in R}r=a^{\varphi(m)}P,\qquad P:=\prod_{r\in R}r .
\]

第四步（约去 $P$）。$m\mid P\big(a^{\varphi(m)}-1\big)$。因为 $P$ 的每个因子都与 $m$ 互质，$P$ 也与 $m$ 互质（按素因子看：任何 $p\mid m$ 都不整除 $R$ 中元素，故 $p\nmid P$）。用上节定理，取 $x=P$、$a'=a^{\varphi(m)}$、$b'=1$，由 $\gcd(P,m)=1$ 得
\[
 a^{\varphi(m)}\equiv1\pmod m.
\]
费马小定理就是 $m=p$ 为素数、$\varphi(p)=p-1$ 的特例，不需要另一套证明。
\end{proof}

\begin{tip}[词义：双射（这里为何只验一半）]
\textbf{双射}= 一一对应：既不碰撞（单射）也不遗漏（满射）。有限集合上\textbf{不碰撞}自动推出\textbf{不遗漏}（抽屉原理），所以下面的证明只验证了不碰撞这一半。
\end{tip}

\begin{tip}[小心：约去必须互质]
\textbf{约去}是全文最容易被误用的一步：只有当被约掉的数与模数互质时才合法。反例：模 $8$ 下 $2\cdot1\equiv2\cdot5$（都等于 $2$），但 $1\not\equiv5\pmod 8$。
\end{tip}


\begin{teach}[证明中的两个关卡]
先问“$ar\bmod m$ 是否仍在简化剩余系？”再问“是否可能重复？”只有两问都答对，才能说乘法置换了整个集合。最后约掉的不是一个普通整数，而是一个可逆元素。
以 $m=10,a=3$ 演示：$\{1,3,7,9\}$ 乘 $3$ 得到 $\{3,9,1,7\}$。若取 $a=2$，第一步就失败。
\end{teach}


\section{扩展欧拉定理：为什么要保留“大于阈值”的信息}
设 $a\ge1,m\ge2,k\ge0$，则
\[
 a^k\equiv
 \begin{cases}
 a^{k\bmod\varphi(m)},&\gcd(a,m)=1,\\
 a^k,&k<\varphi(m),\\
 a^{k\bmod\varphi(m)+\varphi(m)},&k\ge\varphi(m)
 \end{cases}\pmod m.
\]
后两行也可不区分互质性使用。核心是：只有在原指数已经达到阈值时，才可以采用“取余再加 $\varphi(m)$”。


\begin{proof}[按素数幂拆开的证明]
写 $m=\prod p_i^{e_i}$。对每个素数幂分别验证，再用 CRT 合并。
若 $p_i\nmid a$，则 $\varphi(p_i^{e_i})\mid\varphi(m)$，指数相差 $\varphi(m)$ 的倍数不改变余数。
若 $p_i\mid a$，则 $k\nu_{p_i}(a)\ge e_i$ 时 $a^k\equiv0\pmod{p_i^{e_i}}$。而
\[
\varphi(m)\ge\varphi(p_i^{e_i})=p_i^{e_i-1}(p_i-1)\ge e_i.
\]


当 $k\ge\varphi(m)$ 时，原指数及替换指数都至少为 $e_i$，所以两边均为 $0$。
\end{proof}

\begin{tip}[阈值这个词]
阈值 = \textbf{从此以后行为改变}的分界数。这里分界是 $\varphi(m)$：指数 $k$ 小于它时不能改，大于等于它时可以改成 $k\bmod\varphi(m)+\varphi(m)$。之所以要\textbf{加回 $\varphi(m)$}，是为了保住\textbf{$k$ 已经够大}这一信息——非互质底数时，这个信息决定结果是不是 $0$。
\end{tip}


\paragraph{反例应先于模板.} 取 $a=2,m=8,k=4$。只取 $k\bmod\varphi(8)=0$ 会算成 $1$，正确值是 $0$；加回 $4$ 后正确。反过来，$k=1<4$ 时不能强行改为 $5$：$2^1\equiv2$，$2^5\equiv0\pmod8$。

\paragraph{指数塔的接口.} 对巨大指数 $E$，递归函数不能只返回 $E\bmod q$，还要返回真假值 $E\ge q$。两个指数 $0$ 与 $q$ 具有相同余数，但在非互质底数下作用不同。
可以维护二元组
\[
 (\operatorname{rem},\operatorname{large})=(E\bmod q,\iva{E\ge q}).
\]

\begin{tip}[怎么做：接口口诀\textbf{余数 + 一位真假}]
记住这个接口口诀：\textbf{余数 + 一位真假}。返回 $(E\bmod q,\ \iva{E\ge q})$，上层用 $\operatorname{rem}+q\cdot\operatorname{large}$ 还原一个\textbf{行为等价的指数}。判断 $E\ge q$ 用截断计算 $\min(E,q)$：一旦超过就停下，绝不构造大数。
\end{tip}

计算上层幂时使用 $\operatorname{rem}+q\operatorname{large}$。
判定大小使用截断计算 $\min(E,q)$，而不是通过取模后的数值猜测。

\paragraph{截断幂与截断塔.} 定义 $\operatorname{powcap}(a,e,C)=\min(a^e,C)$，每次乘法超过 $C$ 就截断。对 $a\ge2$ 的塔 $T_0=1,T_h=a^{T_{h-1}}$，若 $C>1$，先求最小 $r$ 使 $a^r\ge C$，再计算 $\min(T_{h-1},r)$。若结果达到 $r$，直接返回 $C$，否则算出小幂。阈值迅速下降，避免构造天文数字。$a=1$ 与 $C=1$ 单独处理。


\section{五类逆元问题}
\begin{problem}{统一练习：如何选择求逆元的方法？}
\begin{enumerate}[label=(\alph*)]
\item 素数模数，单次求 $a^{-1}$，$m\le10^{18}$。
\item 一般模数，单次求 $a^{-1}$，或报告不存在。
\item 素数模数，求 $1,2,\ldots,n$ 的逆元，$n<m$。
\item 一般模数，离线求很多与 $m$ 互质的数的逆元。
\item 固定素数模数 $m\le10^9$，强制在线处理多达 $10^8$ 次询问。
\end{enumerate}
\end{problem}

\begin{tip}[在线/离线]
在线：来一问答一问，不许偷看后面的问题。离线：先把所有问题读完，可以排序、分组、分块。第 (d)(e) 两种做法只在离线/强制在线意义下有区别——题目一旦说\textbf{强制在线}，分块预处理就不许用了。
\end{tip}

\hint 单次用快速幂或 exgcd；连续区间用线性递推；批量互质数用前缀积。固定模数的大量在线询问可以进一步预处理。


\subsection{前四类：应掌握的基础方法}
(a) 由 $a^{m-1}\equiv1$ 得 $a^{-1}\equiv a^{m-2}\pmod m$，快速幂 $O(\log m)$。
(b) 求 $ax+my=1$；gcd 不为 $1$ 则不存在，反之将 $x$ 规范到 $[0,m)$。
(c) 令 $m=qi+r$，对 $2\le i<m$，有 $r\ne0$ 且
\[
 \operatorname{inv}[1]=1,\qquad
 \operatorname{inv}[i]\equiv-q\operatorname{inv}[r]\pmod m.
\]

\begin{tip}[为什么要分五类]
差别只在三个开关上：模数是素数吗（决定能不能用费马小定理）、是单次还是批量（决定要不要前缀积）、是否强制在线（决定能不能分块预处理）。做逆元题时先把这三个开关写下来，方法基本就唯一了。
\end{tip}

\begin{tip}[前提：$m$ 是素数且 $a\not\equiv0$]
这一步用了\textbf{$m$ 是素数且 $a\not\equiv0$}：此时 $a^{m-1}\equiv1$（费马小定理），于是 $a\cdot a^{m-2}\equiv1$。$m$ 是合数时 $a^{\varphi(m)}\equiv1$ 才成立，且要先确认互质——直接套 $m-2$ 会算错。
\end{tip}

$r<i$ 保证可以从小到大递推，时间 $O(n)$。
(d) 令 $P_0=1,P_i=P_{i-1}a_i\bmod m$。求一次 $P_n^{-1}$，逆序处理
\[
 a_i^{-1}=P_{i-1}P_i^{-1},\qquad P_{i-1}^{-1}=a_iP_i^{-1}.
\]

\begin{tip}[线性递推逆元公式是怎么来的]
写 $m=qi+r$、$0<r<i$，则 $qi+r\equiv0\pmod m$，即 $r\equiv-qi$。两边同乘 $i^{-1}r^{-1}$（两者都需可逆）得 $i^{-1}\equiv-q\,r^{-1}\pmod m$。因为 $r<i$，$r^{-1}$ 已经算过，所以能一遍扫过去。注意它要求模数是素数（保证每个 $r<i<m$ 可逆），合数时会断，见下面的易错框。
\end{tip}

时间 $O(n+\log m)$，空间 $O(n)$。所有 $a_i$ 可逆才保证总乘积可逆；并不要求模数是素数。

\paragraph{小例.} 模 $17$ 求 $3,5,7$ 的逆元。前缀积为 $1,3,15,3$；$3^{-1}=6$，逆推得到 $7^{-1}=5$、$5^{-1}=7$、$3^{-1}=6$。

\begin{pitfall}
线性公式在合数模数下不能照抄，即使当前 $i$ 可逆也未必有 $r$ 可逆。例如模 $8$、$i=3$ 时余数 $r=2$ 不可逆，但 $3^{-1}=3$ 确实存在。
\end{pitfall}

\subsection{在线方案一：每步至少减半}
设 $m=ua+v$，$0<v<a$。除了 $a^{-1}\equiv-u v^{-1}$，还有
\[
 (u+1)a\equiv a-v\pmod m,
 \qquad a^{-1}\equiv(u+1)(a-v)^{-1}\pmod m.
\]

\begin{tip}[前缀积批量求逆元的想法]
只花一次\textbf{昂贵的求逆}，其余全是乘法和查表：先算前缀积 $P_i=a_1\cdots a_i$，求 $P_n^{-1}$ 一次；然后 $a_i^{-1}=P_{i-1}\cdot P_i^{-1}$ 从后往前推，因为 $P_i^{-1}=a_i^{-1}P_{i-1}^{-1}$。这就是\textbf{用一次 $O(\log m)$ 换 $n$ 次 $O(1)$}。
\end{tip}

\begin{tip}[这个恒等式从哪来]
$m=ua+v$ 意味着 $ua\equiv-v$，两边再加 $a$ 得 $(u+1)a\equiv a-v$。为什么规模减半：新的递归对象是 $\min(v,a-v)$，它不超过 $a/2$。所谓\textbf{每步至少减半}就是把\textbf{取余数}与\textbf{取补数}中较小者递下去，至多 $\log_2 m$ 步进入预处理表。
\end{tip}

在 $v$ 与 $a-v$ 中选较小者，递归规模至少减半。预处理 $[1,L]$ 的逆元后，单次询问最多进行 $O(1+\log(m/L))$ 次递归。取 $L=2\times10^7,m\le10^9$，规模在至多 $6$ 次减半后落入表内；若 $L\ge m$，只预处理到 $m-1$。
这不是一般意义上的渐近 $O(1)$，而是数据范围内步数很小。

\subsection{在线方案二：\texorpdfstring{$O(m^{2/3})$}{O(m2/3)} 预处理与常数次查询运算}
这是选修内容。令 $B=U=\ceil{m^{1/3}}$、$L=\ceil{m/U}$。将 $a$ 写为 $qB+r$，$0\le r<B$。对每个块 $q$ 寻找
\[
1\le u_q\le U,\qquad c_q=qBu_q-k_qm,\qquad |c_q|\le L.
\]
为什么存在？把 $0,qB,\ldots,UqB$ 的模 $m$ 余数放入 $U$ 个长度不超过 $L$ 的区间，有两者落在同一区间，相减即可。
于是
\[
 a u_q\equiv c_q+ru_q=t,\qquad |t|\le L+BU=O(m^{2/3}),
\]

\begin{tip}[抽屉原理一句话]
$U+1$ 个数放进 $U$ 个盒子，必有两个同盒。这里把 $0,qB,\ldots,UqB$ 这 $U+1$ 个余数放进 $[-L,L]$ 划出的 $U$ 段，取同段两者相减，就造出一个\textbf{小的 $qBu-km$}。看到\textbf{存在性证明里出现 $+1$}，基本都是抽屉原理。
\end{tip}

并且 $t\not\equiv0\pmod m$。预处理这一小范围的正负逆元，即得
\[
 a^{-1}\equiv u_q\,t^{-1}\pmod m.
\]
查询只需除法、查表、乘法。对很小的 $m$ 直接预处理全部非零余数。

\paragraph{如何在正确复杂度内找到全部块的系数？}
枚举 $u=1,\ldots,U$，再枚举 $k=0,\ldots,u+1$。满足条件的 $q$ 组成整数区间
\[
 \ceil{\frac{km-L}{uB}}\le q\le\floor{\frac{km+L}{uB}},
\]

\begin{tip}[口径：这里的\textbf{常数次数}指哪几步]
这一步的查表范围是 $[-L,L]$，所以正负都要有逆元：$(-t)^{-1}\equiv-(t^{-1})$，不必另算。"常数次数"指的是：一次整数除法、两次乘法、两次查表——不是 $O(\log m)$。
\end{tip}

与合法块编号范围相交。每个 $u$ 的区间个数为 $O(u)$，包含的候选总数为 $O(u+L/B)=O(U)$；总预处理 $O(U^2)=O(m^{2/3})$。已找到系数的块不必覆盖。

\paragraph{算例.} 模 $101$，$B=U=5,L=21$。对 $a=73=14\cdot5+3$，选 $u=3$，$c=14\cdot5\cdot3-2\cdot101=8$，于是 $t=8+3\cdot3=17$。查表 $17^{-1}=6$，得到 $73^{-1}=3\cdot6=18$。
\cost 这是字长模型中的 $O(m^{2/3})$ 时间、空间预处理与 $O(1)$ 查询；必须把内存规模纳入选择，而不是一看到“大量询问”就使用最复杂方案。



\chapter{中国剩余定理与同余约束的合并}
\begin{chapterintro}[章节导读]
上一章我们已经能够在单个模数下处理同余、逆元和幂；当问题同时给出多个模数的限制时，新的困难就变成“这些限制能不能合在一起”。这一章利用中国剩余定理，把多个局部条件重新拼成一个整体，并处理模数不互质时需要额外检查的情况。
\end{chapterintro}

\section{互质模数：构造局部为一、其余为零的系数}
考虑 $x\equiv a_i\pmod{m_i}$，其中 $m_i\ge2$ 两两互质。令
\[
 M=\prod_i m_i,\qquad M_i=M/m_i,\qquad t_i\equiv M_i^{-1}\pmod{m_i}.
\]

\begin{tip}[为什么：把内存规模一起算进选择]
$O(m^{2/3})$ 在 $m=10^9$ 时约 $10^6$，表很小；但 $m=10^{18}$ 时约 $10^{12}$，直接内存爆炸。所以\textbf{更高级的做法}可能反而不可用，选方法要同时看时间与内存。
\end{tip}

\begin{theorem}[中国剩余定理]
上述方程组在模 $M$ 意义下有唯一解：
\[
 x\equiv\sum_i a_iM_it_i\pmod M.
\]
\end{theorem}

\begin{proof}
\begin{tip}[思考：CRT 的构造到底在做什么？]
先不要急着记公式。我们真正需要的是：对每一个模数 $m_i$，造出一个“开关”$E_i$，使它模 $m_i$ 等于 $1$，而模其他 $m_j$ 都等于 $0$。这样，把 $E_i$ 乘上目标余数 $a_i$，最后把所有这些项相加，就能一次得到所有同余条件。

以三个模数为例，想让一个数模 $m_1$ 为 $1$、模 $m_2,m_3$ 为 $0$，最自然的起点就是 $M_1=m_2m_3$。因为它本身已经在模 $m_2,m_3$ 下为 $0$；又因为三个模数两两互质，$M_1$ 与 $m_1$ 互质，所以只需要乘上一个系数 $t_1$，把 $M_1$ 在模 $m_1$ 下调成 $1$。这就是 $t_1\equiv M_1^{-1}\pmod{m_1}$。

于是 $M_1t_1$ 就是我们要的第一个“开关”。其他模数同理。最后令 $x=\sum_i a_iM_it_i$，每一项在自己对应的模数下留下 $a_i$，在其他模数下全部消失，所以所有同余条件同时成立。下面的正式推导只是把这个想法逐个模数写清楚。
\end{tip}
存在性与唯一性分开证，并把\textbf{逆元为什么存在}这一步补上。

先补一句：$M_i=\prod_{j\ne i}m_j$ 与 $m_i$ 互质（模数两两互质，$M_i$ 的每个因子都与 $m_i$ 互质），所以 $t_i\equiv M_i^{-1}\pmod{m_i}$ 确实存在，依据是第 2 章的\textbf{可逆 $\iff$ 互质}。

存在性。把 $x=\sum_j a_jM_jt_j$ 代入第 $i$ 个同余式。当 $j\ne i$ 时 $M_j$ 含因子 $m_i$，故 $M_jt_j\equiv0\pmod{m_i}$，这些项全部消失；当 $j=i$ 时 $M_it_i\equiv1\pmod{m_i}$，该项留下 $a_i$。于是 $x\equiv a_i\pmod{m_i}$，对每个 $i$ 都成立。

唯一性。设 $x,y$ 都满足方程组，则每个 $i$ 都有 $m_i\mid x-y$。两两互质时反复使用\textbf{$p\mid n$、$q\mid n$、$\gcd(p,q)=1$ 蕴含 $pq\mid n$}（把 $n=pu$ 代入 $q\mid pu$，由互质得 $q\mid u$，故 $pq\mid n$），得到 $M=\prod_i m_i\mid x-y$，即 $x\equiv y\pmod M$。

把结论说准确：解集是模 $M$ 的一个剩余类，所以\textbf{唯一}只在 $[0,M)$ 内成立，不是\textbf{只有一个整数}。
\end{proof}

\paragraph{板书例.} 解
\[
 x\equiv2\pmod3,\quad x\equiv3\pmod5,\quad x\equiv2\pmod7.
\]
$M=105$，$M_i=35,21,15$，相应逆元为 $2,1,1$，所以
\[
 x\equiv2\cdot35\cdot2+3\cdot21+2\cdot15=233\equiv23\pmod{105}.
\]
讲解时应让学生先验算 $70$ 模三个模数的余数为 $(1,0,0)$，再展示整式。
\begin{pitfall}
$M_i$ 不是 $m_i$。$t_i$ 的逆元模数必须写明为 $m_i$。所谓“唯一解”是唯一的同余类，不是唯一整数。
\end{pitfall}


\section{不互质模数：合并两个等差数列}
现在考虑
\[
 x\equiv a\pmod m,\qquad x\equiv b\pmod n.
\]
写 $x=a+mt$，得到 $mt\equiv b-a\pmod n$。设 $d=\gcd(m,n)$：
\begin{itemize}
\item $d\nmid b-a$ 时无解；
\item 否则解 $(m/d)t\equiv(b-a)/d\pmod{n/d}$，得到 $t_0$；
\item 合并结果为 $x\equiv a+mt_0\pmod{\lcm(m,n)}$。
\end{itemize}
重复合并即可得到扩展 CRT。模数为 $1$ 的约束不限制整数，可跳过。


\paragraph{为什么没有漏解？}
关于 $t$ 的所有解相差 $n/d$，所以对应 $x$ 相差 $mn/d$。反过来，每个满足两式的 $x$ 都有 $t=(x-a)/m$，因此每一个合法 $x$ 都被表示到了。
\paragraph{可解与不可解的对照.}
\[
 x\equiv2\pmod6,\quad x\equiv5\pmod9
 \Longrightarrow 6t\equiv3\pmod9
 \Longrightarrow t\equiv2\pmod3,
\]

\begin{tip}[词义：扩展 CRT 的\textbf{扩展}]
\textbf{扩展}只是指模数不必两两互质。做法：每次只合并两个式子，把结果当成一个新式子再和下一个合并。合并一次可能无解（相容性 $d\mid b-a$），也可能把模数变成 $\lcm$。顺序不影响最终解集，但影响中间数值大小。
\end{tip}

故 $x\equiv14\pmod{18}$。把第二个余数改为 $4$，则 $3\nmid(4-2)$，无解。

\begin{lstlisting}
// Contracts: m,n > 0; all intermediate products fit i128.
i128 norm(i128 a, i128 m) { a %= m; return a < 0 ? a + m : a; }
bool merge_crt(i128 a, i128 m, i128 b, i128 n,
               i128 &r, i128 &mod) {
    a = norm(a,m); b = norm(b,n);
    i128 u,v, d = exgcd(m,n,u,v);
    if ((b-a) % d != 0) return false;
    i128 q = n/d;
    i128 t = norm(((b-a)/d) * u, q);
    mod = m*q;
    r = norm(a+m*t, mod);
    return true;
}
\end{lstlisting}
若输入规模使乘积超出 $128$ 位，此段代码不再适用，应改用大整数或安全模乘。exCRT 并不会自动解决整数溢出。

\paragraph{附带一个下界.} 若最终解为 $x\equiv r\pmod M$，$0\le r<M$，还要求 $x\ge L$，则最小解是
\[
 r+M\max\!\left(0,\ceil{\frac{L-r}{M}}\right).
\]

\begin{tip}[词义：大整数（要换掉的不止乘法）]
\textbf{大整数}= 用数组或 \texttt{boost::multiprecision::cpp\_int} 存任意长度。竞赛里能不用就不用（慢一个量级）；实在躲不掉时，要连带把 gcd、取模、快速幂都换成高精度版本，不能只换乘法。
\end{tip}

\begin{tip}[下界公式是怎么推出来的]
解集是 $r,r+M,r+2M,\ldots$。要挑第一个不小于 $L$ 的，就是求最小 $t\ge0$ 使 $r+tM\ge L$，即 $t\ge(L-r)/M$，故 $t=\max(0,\ceil{(L-r)/M})$。写成代码时注意 $L\le r$ 时不能让 $\ceil{\text{负数}}$ 把 $t$ 压成负数，所以外面套 $\max(0,\cdot)$。
\end{tip}

同余约束解决的是周期，大小约束决定选周期中的哪一项。


\begin{problem}{P4774：屠龙勇士}
依次攻击每条龙。面对生命值 $a_i$，从当前剑的多重集合中选攻击力不高于 $a_i$ 的最大剑；若不存在，则选最小剑。设选出的攻击力为 $c_i$。攻击固定的 $x$ 次后，龙可以不断恢复 $p_i$ 生命值；要求恰好恢复到 $0$。击败后所用剑被消耗，并得到奖励剑。求可通关的最小 $x$。
\end{problem}
\hint 有序多重集合确定每次武器；每条龙给出 $c_ix\equiv a_i\pmod{p_i}$ 和 $c_ix\ge a_i$，不能只保留同余式。

\begin{pitfall}[讲解：先看这题到底在做什么]
先把“选什么剑”和“攻击多少次”拆开。武器选择只由当前生命值和武器集合决定，因此可以先独立模拟；确定每条龙使用的攻击力 $c_i$ 后，剩下的问题才是同时满足同余式和攻击次数下界，再用 exCRT 合并。
\end{pitfall}

\sol 分四步完成。
\paragraph{第一步：武器序列独立于攻击次数.}
选剑只与当前集合、$a_i$ 和奖励剑有关。用 \texttt{multiset} 的 \texttt{upper\_bound(a[i])} 寻找第一个大于 $a_i$ 的元素，若它不是 begin，则退一步；否则选择 begin。删除时只删除这一个迭代器，不能删除所有相等攻击力的剑。

\paragraph{第二步：准确翻译死亡条件.}
攻击后生命值 $a_i-c_ix$，只能增加 $p_i$，不能减少。因此存在整数 $k_i\ge0$ 满足
\[
 a_i-c_ix+k_ip_i=0
\]

\begin{tip}[两个容器术语]
\texttt{upper\_bound(x)} 返回\textbf{第一个严格大于 $x$ 的元素}的位置（迭代器），不是位置编号；所以\textbf{不高于 $a_i$ 的最大值}= 先找第一个大于的，再往回退一步，退不了（已是 begin）就取 begin。\texttt{erase(迭代器)} 只删这一个，\texttt{erase(值)} 会把所有相等的都删掉——本题恰好在有重复剑时出错。
\end{tip}

等价于
\[
 c_ix\equiv a_i\pmod{p_i},\qquad x\ge\ceil{a_i/c_i}.
\]
若没有第二条，可能输出一个攻击伤害还不够的“同余解”。

\paragraph{第三步：把带系数同余化为标准式.}
令 $d_i=\gcd(c_i,p_i)$。若 $d_i\nmid a_i$，立即无解；否则除去 $d_i$ 后求逆元，得到
\[
 x\equiv \frac{a_i}{d_i}\left(\frac{c_i}{d_i}\right)^{-1}
 \pmod{p_i/d_i}.
\]
将这些式子逐个 exCRT 合并，并记录 $L=\max_i\ceil{a_i/c_i}$。
\paragraph{第四步：选择不小于下界的第一项.}
用前面的等差数列公式即可。原题保证所有 $p_i$ 的最小公倍数不超过 $10^{12}$，但中间乘法仍可能溢出 $64$ 位。

\paragraph{把官方说明中的 $59$ 推出来.}
初始剑 $\{1,9,1000\}$，生命值依次为 $3,5,7$，恢复量为 $4,6,10$，前两次奖励剑为 $7,3$。选出的剑是 $1,7,3$，得到
\[
 x\equiv3\pmod4,\quad7x\equiv5\pmod6,\quad3x\equiv7\pmod{10}.
\]
即 $x\equiv3\pmod4,x\equiv5\pmod6,x\equiv9\pmod{10}$，先合并前两式得 $x\equiv11\pmod{12}$，再合并第三式得 $x\equiv59\pmod{60}$。下界为 $3$，故答案 $59$。

\cost 选剑 $O(n\log m)$，合并 $O(n\log V)$；集合大小始终为 $m$。当 $p_i=1$ 时同余式没有约束，但伤害下界仍不能丢。
\begin{teach}[学生容易漏掉的三件事]
选剑条件是“不高于”，不是“严格小于”；恢复只能向上，不是任意取模操作；$a_i\le p_i$ 并非所有测试点都具有的统一条件，应使用一般下界处理。
\end{teach}

\src{https://www.luogu.com.cn/problem/P4774}


\chapter{乘法阶、原根与离散对数}
\begin{chapterintro}[章节导读]
上一章解决的是多个同余约束如何合并；这一章换一个方向，研究一个数不断乘下去会形成怎样的周期。乘法阶告诉我们周期有多长，原根刻画整个单位群，而离散对数则反过来要求我们从结果找回指数。
\end{chapterintro}

\section{周期与乘法阶}
\begin{definition}
若 $m\ge2$ 且 $\gcd(a,m)=1$，满足 $a^r\equiv1\pmod m$ 的最小正整数 $r$ 称为阶，记作 $\ord_m(a)$。
\end{definition}
\begin{theorem}
$\ord_m(a)$ 存在当且仅当 $\gcd(a,m)=1$。设其为 $r$，则
\[
 a^k\equiv1\pmod m\iff r\mid k,\qquad r\mid\varphi(m).
\]
若 $a^x\equiv b\pmod m$ 有解 $x_0$，则其非负整数解为 $x\equiv x_0\pmod r$ 中非负的那些整数。
\end{theorem}
\begin{proof}
逐条展开。

阶的存在性等价于互质。若某个 $r\ge1$ 使 $a^r\equiv1\pmod m$，则 $a\cdot a^{r-1}\equiv1$，说明 $a$ 有逆元 $a^{r-1}$，由上节\textbf{可逆 $\iff$ 互质}得 $\gcd(a,m)=1$。反之若 $\gcd(a,m)=1$，欧拉定理给 $a^{\varphi(m)}\equiv1\pmod m$，所以集合 $\{t\ge1:a^t\equiv1\pmod m\}$ 非空；正整数集的非空子集有最小元（最小数原理），记作 $r=\ord_m(a)$。

整除刻画。设 $a^k\equiv1$。用带余除法写 $k=qr+s$、$0\le s<r$。则
\[
 1\equiv a^k=(a^r)^qa^s\equiv1^q a^s=a^s\pmod m.
\]
若 $s>0$，$s$ 就是比 $r$ 更小、又能使幂为一的正指数，与 $r$ 的最小性矛盾。故 $s=0$，即 $r\mid k$。反向也写出来：若 $r\mid k$，$k=rt$，则 $a^k=(a^r)^t\equiv1$。

$r\mid\varphi(m)$。取 $k=\varphi(m)$（此时 $a^k\equiv1$ 由欧拉定理保证，前提 $\gcd(a,m)=1$ 已由阶的存在给出），套上一条即得。

离散对数的解结构。设 $a^{x_0}\equiv b\pmod m$ 有解 $x_0$。若 $a^x\equiv b$，两边同乘 $a^{x_0}$ 的逆元 $a^{r-1}$（注意 $x-x_0$ 可能为负，此时用 $a^{x-x_0}\equiv a^{x-x_0+r}$ 化归非负指数）得 $a^{x-x_0}\equiv1$，由整除刻画得 $r\mid x-x_0$。反向：$x\equiv x_0\pmod r$ 时 $x=x_0+rt$，$a^x\equiv a^{x_0}(a^r)^t\equiv b$。所以全部非负整数解正是同余类 $x_0\bmod r$ 中的非负元——这也说明\textbf{最小非负解}唯一存在。
\end{proof}

\begin{tip}[为什么：求阶前先查 $\gcd(a,m)=1$]
$a^r\equiv1$ 与\textbf{$a$ 可逆}是一件事的两面：能回到 $1$ 就自动有逆元 $a^{r-1}$；可逆才有欧拉定理给出周期。所以任何\textbf{求阶}的代码，第一行都该是检查 $\gcd(a,m)=1$，检查不过就报\textbf{不存在}，而不是返回一个错误的最小值。
\end{tip}


\paragraph{具体周期.} 模 $7$ 的 $2^x$ 依次为 $1,2,4,1,\ldots$，阶为 $3$；不是每个非零余数都能表示成 $2$ 的幂，例如 $3$ 不行。模 $8$ 的 $2^x$ 是 $1,2,4,0,0,\ldots$，只有最终周期，不存在乘法阶。
\paragraph{求阶.} 已知 $\varphi(m)$ 的质因数分解，令 $r=\varphi(m)$。对每个素因子 $q$，只要 $q\mid r$ 且 $a^{r/q}\equiv1$，就令 $r\gets r/q$，直至不能继续。循环结束时 $r$ 正好是阶：若仍大于真阶，商中必有某个素因子可以再除。
预处理含求 $\varphi(m)$ 及其分解；每次求阶通常用 $O(\log^2m)$ 次字长运算作粗上界。不能把分解的成本藏在一次 $O(\log^2m)$ 查询里。


\section{有限域上的根数与原根}
\begin{theorem}[多项式根数界，亦称此处的拉格朗日定理]
在 $\mathbb F_p$ 上，首项系数非零的 $d$ 次多项式至多有 $d$ 个不同根。
\end{theorem}
\begin{proof}
对次数 $d$ 归纳，并把\textbf{可逆}用在哪一行标出来。

$d=0$：非零常数多项式没有根，至多 $0$ 个根，成立。设 $d\ge1$ 且结论对 $d-1$ 次成立。若 $f$ 无根，结论已成立；否则取一根 $r$，即 $f(r)=0$。

第一步（提取因式）。用带余除法（首项系数非零保证可除）：$f(X)=(X-r)g(X)+c$，$\deg g=d-1$、$c$ 为常数。代入 $X=r$ 得 $0=f(r)=0\cdot g(r)+c$，故 $c=0$。

第二步（其余根落到 $g$ 上）。设 $s\ne r$ 也是 $f$ 的根，则
\[
 0=f(s)=(s-r)g(s).
\]




$s-r\ne0$，而在 $\mathbb F_p$ 中每个非零元素都有逆元；两边乘 $(s-r)^{-1}$ 得 $g(s)=0$。

第三步（收尾）。于是\textbf{$f$ 除 $r$ 以外的根}全是 $g$ 的根，由归纳假设至多 $d-1$ 个，故 $f$ 至多 $1+(d-1)=d$ 个。

注意第二步是唯一用到\textbf{域}的地方。在 $\Z/8\Z$ 中 $2\ne0$ 却不可逆，且 $2\cdot4=0$，同样的推理立刻断掉——这就是易错框里四个根的来历。
\end{proof}

\begin{tip}[图景：$0$ 是黑洞]
$0$ 是\textbf{黑洞}：一旦乘出 $0$ 就再也回不到 $1$。所以模 $8$ 的 $2$ 只有\textbf{最终周期}（尾巴进 $0$），没有乘法阶——阶要求 $a^r\equiv1$。这正是求阶之前必须查互质的原因。
\end{tip}


\begin{tip}[提醒：两个\textbf{拉格朗日定理}不是一回事]
这条定理的名字容易撞车：群论里还有一条\textbf{子群的阶整除群的阶}也叫拉格朗日定理，与这里无关。本书尽量用\textbf{多项式根数界}这种描述性名字，避免靠人名记忆。
\end{tip}


\begin{tip}[怎么做：从 $\varphi(m)$ 往下除出最小周期]
把做法翻译成人话就是：先猜一个周期 $r=\varphi(m)$，再逐个素因子试着往下除，\textbf{除了之后还是一}、且还能继续除，就除掉。最后剩下的一定是最小正周期。前提是你已经知道 $\varphi(m)$ 和它的分解——这两项成本要单独计入预处理。
\end{tip}


\begin{pitfall}
这不是任意合数模数下都成立的性质。例如 $X^2-1\equiv0\pmod8$ 有 $1,3,5,7$ 四个根。“次数为二”只有在相应域或整环条件下才能限制根数。本节的名称也不要与群论中“子群阶整除群阶”的拉格朗日定理混淆。
\end{pitfall}


\begin{definition}
满足 $\ord_m(g)=\varphi(m)$ 的元素 $g$ 称为模 $m$ 的原根。
\end{definition}

\begin{tip}[原根一句话]
原根就是\textbf{一个数不断自乘，能走遍所有与 $m$ 互质的余数}。模 $7$ 下 $3$ 的幂给出 $3,2,6,4,5,1$——六个互质余数一个不落，所以 $3$ 是原根。这样的数存在与否完全由 $m$ 的形状决定（$2,4,p^\alpha,2p^\alpha$）。
\end{tip}

\begin{theorem}
对 $m\ge2$，原根存在当且仅当
\[
 m=2,4,p^\alpha,2p^\alpha\quad(p\text{ 为奇素数},\ \alpha\ge1).
\]
若 $\gcd(g,m)=1$，则 $g$ 是原根当且仅当对每个\textbf{不同素因子} $q\mid\varphi(m)$ 都有 $g^{\varphi(m)/q}\not\equiv1\pmod m$。
（注意枚举的对象是 $\varphi(m)$ 的素因子，不是它的约数；$\gcd(g,m)=1$ 这条前提也不能省，否则逆元不存在，测试无从谈起。）
\end{theorem}
原根存在时，简化剩余系为 $1,g,\ldots,g^{\varphi(m)-1}$，原根个数为 $\varphi(\varphi(m))$；且
\[
 \ord_m(g^t)=\frac{\varphi(m)}{\gcd(t,\varphi(m))}.
\]

\paragraph{判定式为什么只查素因子？}
若真阶 $r<\varphi(m)$，则整数 $\varphi(m)/r>1$ 有某个素因子 $q$，于是 $r\mid\varphi(m)/q$，记 $\varphi(m)/q=rt$，得 $g^{\varphi(m)/q}=(g^r)^t\equiv1$，会被检测到。反向：若存在素因子 $q$ 使 $g^{\varphi(m)/q}\equiv1$，则阶整除 $\varphi(m)/q<\varphi(m)$，$g$ 不是原根。
原根存在性分类的完整证明需要初等数论之外的工具（对 $2^\alpha$ 与奇素数幂分别构造，再由中国剩余定理合并），篇幅超过一节；基础课把它当作已知定理使用，判定式、原根计数与阶公式则应当课堂证明。
\paragraph{模 $7$ 的例子.} $\varphi(7)=6$，只需测指数 $3,2$。$3^3\equiv6$，$3^2\equiv2$，故 $3$ 为原根。与 $6$ 互质的指数为 $1,5$，所以原根是 $3^1=3$ 与 $3^5\equiv5$。
\paragraph{随机找原根的边界.} 若原根存在，已知 $\varphi(t)/t=\Omega(1/\log\log t)$，可得到在 $[1,m)$ 随机抽样找到原根的期望次数为 $O((\log\log m)^2)$。每次还要做若干快速幂，不能把抽样次数当成总时间；小模数单独处理。


\section{BSGS：用平方根空间换时间}
\begin{problem}{离散对数}
求 $a^x\equiv b\pmod m$ 的最小非负整数解，或报告无解。
\end{problem}
当 $\gcd(a,m)=1$ 时，若有解，则存在 $0\le x<\varphi(m)\le m$ 的解。这由欧拉定理得到，不需要误用“合数模数下的费马小定理”。
设 $B=\ceil{\sqrt m}$，写 $x=iB+j$，$0\le j<B$，则
\[
 a^j\equiv b(a^{-B})^i\pmod m.
\]

\begin{tip}[口径：\textbf{抽几次}不等于\textbf{跑多久}]
\textbf{期望抽几次}和\textbf{总共跑多久}是两件事：每次抽样后要做 $\omega(\varphi(m))$ 次快速幂，还要有 $\varphi(m)$ 的分解。小模数直接枚举反而快；把渐近式当成实现建议是常见的坑。
\end{tip}

预处理 baby steps $a^j$，再枚举 giant steps 的右边并查表，时间、空间均为 $O(\sqrt m)$（哈希期望）。


\paragraph{保证最小解的细节.}
哈希表的同一个余数只保留最小 $j$。按照 $i=0,1,\ldots$ 的顺序查询，那么各块的指数区间互不重叠且递增，首次命中的 $iB+j$ 即为最小解。即使枚举区间略超过 $m$，首次命中仍然不会错；无解需完整检查足够覆盖 $[0,m)$ 的块数。
若 $b$ 与 $m$ 不互质而 $a$ 与 $m$ 互质，则立即无解，因为任何 $a^x$ 都是单位。

\paragraph{手算例.} 解 $2^x\equiv5\pmod{13}$。取 $B=4$，baby 表为
\[
1\mapsto0,\quad2\mapsto1,\quad4\mapsto2,\quad8\mapsto3.
\]

\begin{tip}[BSGS 的核心一次代换]
把未知的指数按块拆开：$x=iB+j$，$0\le j<B$。等式 $a^{iB+j}\equiv b$ 两边同乘 $a^{-iB}$，变成\textbf{只有 $j$ 在左、只有 $i$ 在右}：$a^j\equiv b(a^{-B})^i$。左边全部预先算好放进表，右边一个个枚举去查表。这就是\textbf{用空间换时间}的标准形状，$\sqrt m$ 来自 $B\cdot B\approx m$。
\end{tip}

\begin{tip}[为什么：表里只留最小的 $j$]
查表法求\textbf{最小解}时，插入策略和枚举顺序共同决定正确性：表里保留最小的 $j$，外层 $i$ 从小到大，于是第一次命中就是字典序最小的 $(i,j)$，也就是最小 $x$。若表被后来的大 $j$ 覆盖，答案可能偏大。
\end{tip}

$2^4\equiv3$，逆元为 $9$；giant 依次为 $5,6,2$，在 $i=2$ 时命中 $j=1$，得 $x=9$。

\paragraph{多次询问如何分块？}
若\textbf{底数 $a$ 与模数 $m$ 都固定}，可以预处理 $a^{iB}$ 的表，每问枚举 $ba^j$，将 $x=iB-j$ 还原。预处理 $O(m/B)$，$Q$ 次查询总计 $O(QB)$，平衡时 $B\approx\sqrt{m/Q}$；若该值小于 $1$，取 $B=1$。总时间应写成 $O(\sqrt{Qm}+Q)$ 的相应界，并说明取整与内存限制。
如果只有模数固定、底数改变，原表一般不能复用。求最小解时，对同一 $j$ 找到最小可行 giant 指数，再对所有 $j$ 的候选取最小值，不能随意用哈希表覆盖重复值。


\section{exBSGS：剥离非互质部分}
先处理 $m=1$、$x=0$、$a=0$ 等边界。令 $d=\gcd(a,m)>1$。对 $x\ge1$，若 $d\nmid b$ 则无解；否则
\[
 \frac ad a^{x-1}\equiv\frac bd\pmod{m/d}.
\]

\begin{tip}[exBSGS 在剥什么]
$a$ 与 $m$ 不互质时 $a$ 不可逆，$a^{-1}$ 无从谈起。技巧是\textbf{约掉一个公因子再降一层}：$a^x\equiv b\pmod m$ 且 $d=\gcd(a,m)\mid b$ 时，可化为 $\frac ad a^{x-1}\equiv\frac b d\pmod{m/d}$。每剥一层模数变小、指数减一；剥到互质为止，再交给 BSGS。
\end{tip}

\begin{tip}[\textbf{平衡}两个成本的做法]
总时间形如 $T(B)=m/B+QB$（预处理一项、查询一项）。两项一个随 $B$ 减、一个随 $B$ 增，令相等 $m/B=QB$ 得 $B=\sqrt{m/Q}$，此时 $T=O(\sqrt{Qm})$。也可以用求导验证；再检查取整、内存上限、以及 $B<1$ 时要夹回 $1$。
\end{tip}

由于 $\gcd(a/d,m/d)=1$，该系数可逆，可以递归到更小的模数。递归前每层都要检查零指数解。最终底数与新模数互质时使用 BSGS。


更适合实现的形式是维护已剥离次数 $c$ 与系数 $k$，循环不变量为
\[
 k a^{x-c}\equiv b\pmod m,\qquad\gcd(k,m)=1.
\]

\begin{tip}[什么叫\textbf{循环不变量}]
每次循环开头都必须成立的一句话，例如\textbf{剥掉 $c$ 层后，余下的式子是 $ka^{x-c}\equiv b$，且 $k$ 可逆}。写代码前先写不变量，循环结束时的不变量加退出条件就自动是答案；调试时也照这句话打表检查。
\end{tip}

初始 $c=0,k=1$。每次先判断 $b=k$，若成立则 $x=c$ 就是最小解；再令 $d=\gcd(a,m)$。
若 $d=1$，解 $a^{x-c}\equiv bk^{-1}\pmod m$；若 $d>1$ 且 $d\nmid b$ 则无解。否则
\[
 m\gets m/d,\quad b\gets b/d,\quad
 k\gets k(a/d)\bmod m,\quad c\gets c+1.
\]
这里 $a/d$ 使用本轮的 $d$；底数 $a$ 保持不变。$m$ 至少减半，剥离最多 $O(\log m)$ 轮。

\paragraph{算例.} $4^x\equiv8\pmod{14}$。先检查 $x=0$ 不满足。剥离 $d=2$ 后 $m=7,b=4,k=2,c=1$，于是需解 $4^{x-1}\equiv2\pmod7$，最小 $x-1=2$，答案 $3$。验证 $4^3=64\equiv8\pmod{14}$。
\begin{pitfall}
离散对数的零指数是合法候选；$a\equiv0$ 时，$x=0$ 给 $1$，$x\ge1$ 给 $0$。非互质情形的解集可能带有限前缀，不能一开始就写成 $x_0+k\ord_m(a)$。
\end{pitfall}


\begin{problem}{P5605：小 A 与两位神仙}
$m=p^q$，$p$ 为奇素数。多次询问单位 $x,y$，判断是否存在 $a\ge0$ 使 $x^a\equiv y\pmod m$。$m\le10^{18}$，询问数不超过 $2\times10^4$，允许使用快速因数分解工具。
\end{problem}

\begin{tip}[\textbf{快速因数分解工具}指什么]
指 Miller--Rabin 判素 + Pollard--Rho 分解（附录模板里的 \texttt{is\_prime} 与 \texttt{factor}）。它们在 $64$ 位范围内实践上很快，但 Rho 的速度依赖随机函数的\textbf{表现}，没有严格多项式界；而分解只是把 $m$ 拆成素数幂的工具，拆完之后每一步都是确定的整数运算。
\end{tip}

\hint 单位群是循环群，条件恰为 $\ord_m(y)\mid\ord_m(x)$；无需真正求离散对数。

\begin{pitfall}[讲解：先看这题到底在做什么]
题目的关键转化是把单位群里的元素写成某个原根的幂。这样 $x^a=y$ 就从乘法问题变成指数上的线性同余：若 $x=g^u,y=g^v$，就只需研究 $au\equiv v\pmod{\varphi(m)}$。因此真正要检查的是这个线性同余是否可解，而不是去真的求出离散对数。
\end{pitfall}

\sol 取原根 $g$，令 $x=g^u,y=g^v$，则需解
\[
 au\equiv v\pmod{\varphi(m)}.
\]

\begin{tip}[这条提示里的每个词]
\textbf{单位}= 与 $m$ 互质的余数（可逆）；\textbf{单位群是循环群}= 存在一个 $g$，使得每个单位都写成 $g^t$；\textbf{阶}$\ord_m(x)$ = 使 $x^r\equiv1$ 的最小正指数。把 $x=g^s$、$y=g^t$ 代入 $x^a\equiv y$，得 $sa\equiv t$，这类方程有解的条件正是 $\ord_m(y)\mid\ord_m(x)$——所以根本不必真的求出离散对数。
\end{tip}

裴蜀定理给出 $\gcd(u,\varphi(m))\mid v$。等价地，它整除 $\gcd(v,\varphi(m))$；使用阶公式便得到所述整除关系。
也可以从循环群子群唯一性理解：$\langle x\rangle$ 是阶为 $\ord_m(x)$ 的唯一子群。

\paragraph{为什么条件不能推广到任意模数？}
模 $8$，$3$ 与 $5$ 的阶都为 $2$，但 $3$ 的幂只能取 $1,3$，不能得到 $5$。单位群不是循环群，“阶整除”仅必要而不充分。
\cost 分解 $m$、计算并分解 $\varphi(m)$ 后，每问两次求阶。因数分解的 $O(m^{1/4})$ 应理解为后文 Pollard--Rho 的常用启发式估计，不是无条件最坏界。

\src{https://www.luogu.com.cn/problem/P5605}

\begin{problem}{P6736：白泽教育——高阶运算取模}
定义 $a\uparrow^1 b=a^b$，且 $a\uparrow^n0=1$；$n>1,b>0$ 时
\[
 a\uparrow^n b=a\uparrow^{n-1}(a\uparrow^n(b-1)).
\]
给定 $1\le a\le10^9,1\le n\le3,0\le b<m\le10^9$，求最小非负 $x$ 使 $a\uparrow^n x\equiv b\pmod m$。
\end{problem}
\hint $n=1$ 用 exBSGS；$n=2$ 利用欧拉链计算有限高度的塔并证明稳定；$n=3$ 利用高度增长极快而不是构造大整数。先特判 $x=0$、$m=1$、$a=1$。

\begin{pitfall}[讲解：先看这题到底在做什么]
幂塔题最容易犯的错是直接对指数取模。正确顺序是先判断底数、模数和指数是否处于“可以降幂”的条件，再沿欧拉链逐层处理；当底数和模数不互质时，还要保留“指数是否已经足够大”这一位信息。
\end{pitfall}

\sol 给出一条以正确性易于讲清为优先的算法路线。令 $T_0=1,T_h=a^{T_{h-1}}$，它是 $a\uparrow^2h$。例如底数 $2$ 的前几项为
\[
1,\ 2,\ 4,\ 16,\ 65536,\ 2^{65536},\ldots.
\]

\begin{tip}[为什么这三个特判不能省]
$x=0$：任何 $a^0=1$，若 $b\equiv1$ 答案就是 $0$，而按\textbf{指数为正}的循环会漏掉；$m=1$：所有余数都是 $0$，答案恒为 $0$，但 $\varphi(1)$、逆元之类的记号在此都退化；$a=1$：$1$ 的幂永远是 $1$，要么立刻命中要么永远不命中，递归降幂会打转。边界情形不特判，正确的算法也会给出错误答案，这是竞赛里最常见的失分方式之一。
\end{tip}

\paragraph{有限枚举的依据：稳定上界.}
设欧拉链 $m_0=m,m_{i+1}=\varphi(m_i)$，到 $m_L=1$。取
\[
 R=\ceil{\log_2 m},\qquad H=R+L+1.
\]
对 $a\ge2$，$T_h\ge2^h$。当 $h\ge R$，在链上所有模数处，“是否达到阈值”的标记都为真。底层模 $1$ 的余数恒定，再向上逐层归纳：每向上一层多需要一层塔高；因此 $T_h\bmod m$ 在 $h\ge H$ 时稳定。
于是得到一个安全的 $O(\log m)$ 高度上界，不追求最紧的常数。

$n=2$ 时从 $h=0$ 到 $H$ 枚举，使用扩展欧拉递归与截断塔计算 $T_h\bmod m$，第一次命中就是答案；都不命中则永远不命中。\textbf{不能仅凭两次顶层余数相同就提前宣布稳定}，应使用上述证明或完整状态的稳定判据。

\paragraph{$n=3$ 为什么只需极少候选？}
记 $P_x=a\uparrow^3x$，则 $P_0=1,P_{x+1}=T_{P_x}$。当 $a\ge3$，
\[
 P_1=a,\quad P_2=T_a\ge3^{27}>H.
\]
故 $P_3=T_{P_2}$ 已使用稳定的塔高，此后结果相同，只需检查 $x=0,1,2,3$。
当 $a=2$，$P_1=2,P_2=4,P_3=T_4=65536>H$，故还需检查 $x=4$。在本题 $m\le10^9$ 时，上述 $H$ 小于 $100$，比较非常安全。

计算 $P_2$ 或 $P_3$ 的“高度”时，只保留截断到 $H$ 的结果；高度达到 $H$ 即可使用稳定值。原稿中 $3\uparrow^33=3^{3^{27}}$ 的等式不正确：前者是高度 $3^{27}$ 的 $3$ 的幂塔，而右边仅是高度 $4$ 的塔。

\cost 每问预处理欧拉链，缓存重复模数的 $\varphi$，再做有限次快速幂。上述教学实现应按“枚举高度数 $\times$ 欧拉链深度 $\times$ 快速幂成本”计时；不是一句“答案很小”就能省略全部复杂度。大量测试下可复用链、预筛试除质数，并改为自底向上批量计算。$n=1$ 的瓶颈仍可能是平方根级 BSGS。

\src{https://www.luogu.com.cn/problem/P6736}


\chapter{升幂引理：把巨大幂变成指数的加法}
\begin{chapterintro}[章节导读]
上一章我们研究了模意义下的周期和指数；现在继续追问一个更细的问题：两个非常大的幂相减，究竟含有多少个素因子 p？p 进赋值把“乘法中的指数”变成可加的量，升幂引理则把这个观念直接用于幂的差值。
\end{chapterintro}

\section{\texorpdfstring{$p$}{p}-进赋值与使用前的检查}
对非零整数 $x$，$\nu_p(x)=\max\{e\ge0:p^e\mid x\}$。基本性质为
\[
 \nu_p(xy)=\nu_p(x)+\nu_p(y),\qquad
 \nu_p(x+y)\ge\min(\nu_p(x),\nu_p(y)),
\]

\begin{tip}[幂塔与高度]
$a\uparrow\uparrow h$ 表示 $h$ 层的塔 $a^{a^{\cdot^{\cdot^a}}}$，\textbf{高度}就是层数。塔的计算顺序是从最上层往下（右结合），所以\textbf{高度 $5$ 的 $2$ 的塔}是 $2^{2^{2^{2^2}}}=2^{65536}$，而不是 $((2^2)^2)^2$。
\end{tip}

\begin{tip}[箭号记号的两条约定]
$a\uparrow^1 b=a^b$，$a\uparrow^{k}1=a$，$a\uparrow^{k}b=a\uparrow^{k-1}(a\uparrow^{k}b)$ 的塔式展开：$3\uparrow^2 3=3^{3^3}$ 高度为 $3$；再升一级 $\uparrow^3$ 就是\textbf{高度的高度}。凡是这种记号，先代最小的数值（高度 $1,2,3$）确认自己算的和定义一致，再往下看结论。
\end{tip}

若 $\nu_p(x)\ne\nu_p(y)$，后一式取等。
下文统一设 $p$ 为素数，$p\nmid xy$，$n$ 为正整数。涉及差为 $0$ 的情况时可用 $\nu_p(0)=+\infty$，实际算题通常先排除零。


\paragraph{直观解释.} $\nu_p$ 是“能连续除以 $p$ 多少次”。乘法变加法；加法却可能发生抵消，所以不能写 $\nu_p(x+y)=\nu_p(x)+\nu_p(y)$。例如 $\nu_2(6)=1,\nu_2(10)=1$，但 $\nu_2(16)=4$。
\begin{teach}[LTE 的使用检查表]
\begin{enumerate}
\item 要计算的是幂的差还是幂的和？指数是否为正？
\item 素数是否整除 $x-y$，或整除 $x+y$？
\item $p\nmid x$、$p\nmid y$ 是否都满足？若不满足，先提出共同素数幂。
\item $p=2$ 还是奇素数？指数是奇数还是偶数？
\end{enumerate}
只背“加一个 $\nu_p(n)$”容易误用到二进制的偶指数情况。
\end{teach}


\section{与素数互质的指数}
\begin{theorem}
若 $p\mid x-y$ 且 $p\nmid n$，则
\[
 \nu_p(x^n-y^n)=\nu_p(x-y).
\]



若 $p\mid x+y$、$n$ 为奇数且 $p\nmid n$，则
\[
 \nu_p(x^n+y^n)=\nu_p(x+y).
\]
\end{theorem}

\begin{tip}[口径：为什么规定 $\nu_p(0)=+\infty$]
规定 $\nu_p(0)=+\infty$ 是为了让\textbf{$p^e\mid x$ 对任意 $e$ 成立}这句话有地方安放，从而 $\nu_p(xy)=\nu_p(x)+\nu_p(y)$ 在 $y=0$ 时不用另写情形。做题时代入具体数字前先检查 $x=y$，否则会得到\textbf{$\nu_p(0)$ 等于某个有限值}的错误等式。
\end{tip}


\begin{tip}[前提：$p\nmid xy$ 为什么不能少]
\textbf{$p\nmid xy$}是三条结论共同的前提，它等价于 $p\nmid x$ 且 $p\nmid y$。少了它，$x,y$ 都能被 $p$ 整除时左边右边同时虚增，公式就不成立了：例如 $p=3$、$x=3,y=6,n=2$，左边 $\nu_3(9-36)=2$，右边 $\nu_3(-3)+\nu_3(2)=1$。
\end{tip}

\begin{proof}
先写出恒等式
\[
 x^n-y^n=(x-y)S,\qquad S=\sum_{i=0}^{n-1}x^{\,n-1-i}y^{\,i}.
\]
它只是多项式恒等式（把右边展开、逐项相消即得），对任何整数都成立，不依赖整除条件。

由赋值的可加性（第 0 章 0.6 节，或本章开头那条 $\nu_p(xy)=\nu_p(x)+\nu_p(y)$）：
\[
 \nu_p(x^n-y^n)=\nu_p(x-y)+\nu_p(S).
\]
所以只剩一件事要证：$\nu_p(S)=0$，即 $p\nmid S$。

模 $p$ 计算 $S$。条件 $p\mid x-y$ 给 $x\equiv y\pmod p$，于是 $S$ 的每一项都满足 $x^{n-1-i}y^{i}\equiv x^{n-1-i}x^{i}=x^{n-1}\pmod p$，共 $n$ 项，故
\[
 S\equiv n x^{\,n-1}\pmod p.
\]
题设有 $p\nmid n$；又 $p\nmid x$（若 $p\mid x$，由 $p\mid x-y$ 得 $p\mid y$，与 $p\nmid xy$ 矛盾）。两者相乘仍不被 $p$ 整除，故 $p\nmid S$，即 $\nu_p(S)=0$。结论得证。

和的情形：$n$ 为奇数时 $x^n+y^n=x^n-(-y)^n$，且 $p\mid x-(-y)=x+y$。把已证的\textbf{差}情形用于 $-y$ 即得
$\nu_p(x^n+y^n)=\nu_p(x+y)+\nu_p(n)$。$n$ 必须为奇数，否则 $(-y)^n=y^n$，第一步就不成立——这就是分类的原因。
\end{proof}

\section{奇素数的完整形式}
\begin{theorem}[LTE，奇素数]
若 $p$ 为奇素数且 $p\mid x-y$，则
\[
 \nu_p(x^n-y^n)=\nu_p(x-y)+\nu_p(n).
\]
若 $p\mid x+y$ 且 $n$ 为奇数，则
\[
 \nu_p(x^n+y^n)=\nu_p(x+y)+\nu_p(n).
\]
\end{theorem}

\begin{proof}[先把指数乘一次 $p$，再反复提升]
令 $s=\nu_p(x-y)\ge1$，写 $x=y+h$，$\nu_p(h)=s$。二项式展开
\[
 x^p-y^p=py^{p-1}h+\sum_{j=2}^{p-1}\binom pj y^{p-j}h^j+h^p.
\]
首项的赋值恰为 $s+1$。中间各项赋值至少为 $1+2s\ge s+2$；末项赋值为 $ps\ge s+2$，这里用了 $p\ge3$。因此最低赋值只有首项取得，不会抵消，故
\[
\nu_p(x^p-y^p)=s+1.
\]
写 $n=p^e u$，$p\nmid u$。先用上一节处理 $u$，再连续 $e$ 次把指数乘 $p$，每次增加一，得到 $s+e$。幂的和仍以 $y\mapsto-y$ 转为差。
\end{proof}
\paragraph{教学重点.} 原稿的二重求和证明可以成立，但容易让学生看不见真正原因。用“唯一最低赋值项”说明为什么不抵消，之后只需归纳，结构更清楚。


\section{素数二的特殊形式}
\begin{theorem}[LTE，$p=2$]
设 $x,y$ 为奇数。若 $n$ 为奇数，则
\[
 \nu_2(x^n-y^n)=\nu_2(x-y),\qquad
 \nu_2(x^n+y^n)=\nu_2(x+y).
\]


若 $n$ 为偶数，则
\[
 \nu_2(x^n-y^n)=\nu_2(x-y)+\nu_2(x+y)+\nu_2(n)-1.
\]
特别地，若 $4\mid x-y$，则对所有正整数 $n$ 有
\[
 \nu_2(x^n-y^n)=\nu_2(x-y)+\nu_2(n).
\]
\end{theorem}

\begin{tip}[方法：唯一最低赋值项]
\textbf{唯一最低赋值项}是一种很常用的论证格式：把一项写成若干项之和，证明其中某一项所含 $p$ 比其它任何一项都少，那么这个\textbf{少}无法被抵消，整个和的赋值就等于它。用这个格式，二重求和可以省掉。
\end{tip}


\begin{proof}
奇指数已由第一种形式处理。偶数写为 $n=2^e u$，$u$ 为奇数、$e\ge1$。先去掉奇数部分：
\[
\nu_2(x^n-y^n)=\nu_2(x^{2^e}-y^{2^e}).
\]
继续因式分解
\[
 x^{2^e}-y^{2^e}=(x-y)(x+y)\prod_{i=1}^{e-1}(x^{2^i}+y^{2^i}).
\]
奇数的平方模 $8$ 为 $1$，因此每个 $i\ge1$ 对应的括号模 $4$ 为 $2$，恰贡献一个因子 $2$，一共 $e-1$ 个。相加得到公式。
若 $4\mid x-y$，则 $x+y\equiv2\pmod4$，所以 $\nu_2(x+y)=1$，偶数公式简化；奇数时 $\nu_2(n)=0$，也一致。
\end{proof}
\paragraph{另一个边界.} 若 $x,y$ 为奇数且 $n$ 偶，则 $x^n+y^n\equiv2\pmod4$，因此赋值为 $1$。不要把“幂的和的 LTE”用于偶数指数。


\section{分层例题与解析}
\begin{problem}{例 1：原稿中的两个练习}
求 $\nu_7(10^7-3^7)$ 与 $\nu_2(7^{10}-3^{10})$。
\end{problem}
\hint 第一个是奇素数差的形式；第二个必须用 $p=2$ 的偶指数形式。

\begin{pitfall}[讲解：先看这题到底在做什么]
第一题直接套奇素数版本的 LTE，但必须先把条件逐条检查清楚；第二题是 $p=2$ 的特殊情形，不能照搬奇素数公式。这里最值得讲的是“先识别定理，再核对条件”，而不是只看最后的数字。
\end{pitfall}

\sol 因为 $7\mid10-3$ 且 $7\nmid30$，
\[
 \nu_7(10^7-3^7)=\nu_7(7)+\nu_7(7)=2.
\]

\begin{tip}[怎么用一节里的例题]
每题先只读题面，写下\textbf{用哪条定理、条件是否满足}，再看解析的第一句是否与你猜的一致。对不上时不要往下读，先找出是自己漏了条件还是解析换了方法。
\end{tip}

第二题中 $7,3$ 均为奇数，指数 $10$ 为偶数，故
\[
 \nu_2(7^{10}-3^{10})=\nu_2(4)+\nu_2(10)+\nu_2(10)-1=2+1+1-1=3.
\]
注意两个 $\nu_2(10)$ 来自不同对象：一个是底数之和，一个是指数。


\begin{problem}{例 2：先分组再使用 LTE}
求 $\nu_3(2^{2025}+1)$ 与 $\nu_5(3^{100}-1)$。
\end{problem}
\hint 第一题为奇指数幂的和；第二题先把 $3^{100}$ 看成 $81^{25}$。

\begin{pitfall}[讲解：先看这题到底在做什么]
第一题先看成幂的和；第二题则先把指数拆成一个让底数之幂与 $1$ 发生整除关系的部分，再使用 LTE。学生容易记住公式却忘记“先找到合适的 $a^d-1$”，所以这里应强调结构而不是算术。
\end{pitfall}

\sol $2025=3^4\cdot25$，所以
\[
\nu_3(2^{2025}+1)=\nu_3(2+1)+\nu_3(2025)=1+4=5.
\]
$5\nmid3-1$，不能直接对底数 $3,1$ 使用差的 LTE。但 $3^4-1=80$ 可被 $5$ 整除，因此
\[
\nu_5(3^{100}-1)=\nu_5(81^{25}-1)=\nu_5(80)+\nu_5(25)=1+2=3.
\]
此例展示“先用阶找到能被 $p$ 整除的基本差，再用 LTE 提升”。


\begin{problem}{例 3：最小指数与素数幂模数下的阶}
求最小正整数 $n$ 使 $5^4\mid3^n-1$。
\end{problem}
\hint 先由模 $5$ 的阶得到 $4\mid n$，再比较赋值。

\begin{pitfall}[讲解：先看这题到底在做什么]
先解决“$n$ 至少应该是什么倍数”这个粗条件，再处理“还需要多少个 $p$ 因子”这个精条件。前者来自模 $p$ 的阶，后者来自 LTE；两个条件合起来，最小指数就不再需要试枚举。
\end{pitfall}

\sol 模 $5$ 下 $3$ 的阶为 $4$，设 $n=4t$。LTE 给出
\[
 \nu_5(3^{4t}-1)=\nu_5(3^4-1)+\nu_5(t)=1+\nu_5(t).
\]
要至少为 $4$，最小 $t=5^3=125$，故答案 $n=500$。
一般对奇素数 $p$、$d=\ord_p(a)$，若 $s=\nu_p(a^d-1)$，则
\[
 \ord_{p^k}(a)=d\,p^{\max(0,k-s)}.
\]
证明正是上面的两步。它把阶和升幂引理连接起来；$p=2$ 需另行分类。



\chapter{阶乘、素数幂与组合数取模}
\begin{chapterintro}[章节导读]
上一章的赋值方法能告诉我们一个数里含多少个 p；这一章把它应用到阶乘。因为阶乘本身会迅速变得巨大，我们要把“p 的指数”和“去掉 p 后留下的单位部分”拆开，最终得到素数幂模数下组合数的计算方法。
\end{chapterintro}

\section{Wilson 定理及单位群乘积}
\begin{theorem}[Wilson 定理]
整数 $n>1$ 为素数，当且仅当 $(n-1)!\equiv-1\pmod n$。
\end{theorem}
\begin{proof}
分\textbf{素数时成立}与\textbf{合数时不成立}两个方向。

正向（$p$ 为素数）。$p=2$：$(2-1)!=1\equiv-1\pmod2$（此时 $1$ 与 $-1$ 是同一余数，需单独看一眼）。设 $p$ 为奇素数，考虑集合 $U=\{1,2,\ldots,p-1\}$，其中每个元素都可逆。把 $U$ 按\textbf{$x$ 与 $x^{-1}$ 相同吗}分成两类：
若 $x\ne x^{-1}$，把 $x$ 与 $x^{-1}$ 配成一对，每对的乘积 $\equiv1$。这样的配对不会重叠，因为 $(x^{-1})^{-1}=x$。
若 $x=x^{-1}$，即 $x^2\equiv1\pmod p$，则 $p\mid(x-1)(x+1)$。$p$ 是素数，故 $p\mid x-1$ 或 $p\mid x+1$（欧几里得引理），得 $x\equiv1$ 或 $x\equiv-1$。反过来这两者确实自逆。
于是
\[
 (p-1)!\equiv\Big(\prod_{\text{对}}1\Big)\cdot1\cdot(-1)\equiv-1\pmod p .
\]



这里\textbf{素数}用在两处：每个非零元可逆（才能配对），以及 $p\mid AB$ 蕴含 $p\mid A$ 或 $p\mid B$（才能说自逆元只有 $\pm1$）。

反向（$n$ 为合数则 $(n-1)!\not\equiv-1$）。取 $1<d<n$ 且 $d\mid n$。
情形一：$d<n-1$，即 $n>2d$。此时 $d$ 出现在 $1\cdot2\cdots(n-1)$ 的因子中，故 $d\mid(n-1)!$。若同时 $(n-1)!\equiv-1\pmod n$，则 $n\mid(n-1)!+1$，从而 $d\mid(n-1)!+1$；两式相减得 $d\mid1$，矛盾。
情形二：$n=2d$ 且 $d$ 是大于 $2$ 的素数（即 $n$ 是 $2$ 倍奇素数）。此时取 $d$ 本身，$d\ge3$ 与 $d\le n-2$ 使 $d$ 与 $2d$... 更简单的做法是：$n$ 有平方因子时（$n=d^2$ 或 $d^2\mid n$）$d$ 与 $n/d$ 是两个不同且都 $\le n-1$ 的因子，都出现在 $(n-1)!$ 中，故 $d^2\mid(n-1)!$，于是 $n\mid(n-1)!$；若又有 $(n-1)!\equiv-1\pmod n$，则 $n\mid1$，矛盾。剩下 $n=2p$（$p$ 为奇素数）时直接验算：$(n-1)!\equiv0\pmod p$，而 $-1\not\equiv0\pmod p$，故不成立。$n=4$ 时 $(4-1)!=6\equiv2\not\equiv-1\pmod4$。

综合：$(n-1)!\equiv-1\pmod n\iff n$ 为素数。
\end{proof}

\begin{tip}[一节里其实在算什么乘积]
所有\textbf{把一堆余数乘起来}的问题（$(p-1)!$、单位群全体、约数之积）都用同一套：把元素按互为逆元配对，配对项乘积为 $1$，剩下自逆元单独数。所以真正要算的总是\textbf{解 $x^2\equiv1$ 有几个解}。
\end{tip}


\begin{tip}[\textbf{单位群乘积}在算什么]
模 $m$ 下所有可逆余数的乘积。做法永远是同一套：把元素按\textbf{与自己互为逆元}配对，乘积为 $1$，剩下的只有自逆元（解 $x^2\equiv1$）。所以真正要算的是\textbf{自逆元有几个}，这由 $m$ 的形状决定；模素数时恰有 $\pm1$。
\end{tip}

对一般 $m>1$，所有单位的乘积模 $m$ 为 $-1$ 或 $1$。取 $-1$ 的模数类型是 $2,4,p^\alpha,2p^\alpha$（$p$ 奇素数），其他取 $1$；模 $2$ 下 $1=-1$，数值上重合。
特别地，令 $Q=p^\alpha$，则
\[
 C_Q=\prod_{\substack{1\le i\le Q\\p\nmid i}}i\equiv
 \begin{cases}1,&p=2,\alpha\ge3,\\-1,&\text{其他}\end{cases}\pmod Q.
\]


\paragraph{推广形式的理由.}
在任意有限阿贝尔群中，将非自逆元素成对消去，剩下的是满足 $u^2=1$ 的元素。它们构成若干个二阶循环群的直积：若只有一个非单位自逆元素，乘积是该元素；若有至少两个独立的二阶方向，每个坐标取 $1$ 的次数为偶数，乘积是单位元。
模奇素数幂的 $x^2=1$ 只有 $\pm1$：$p$ 不可能同时整除 $x-1,x+1$，故一个因子吸收全部 $p^\alpha$。模 $2^\alpha$（$\alpha\ge3$）则有四个自逆元素。再用 CRT 组合，就得到分类。
\begin{pitfall}
$C_Q$ 是\textbf{删去 $p$ 的倍数后}的单位乘积，而 $(Q!)_p$ 是对每个整数\textbf{除尽自身包含的 $p$ 因子}后再相乘，两者不相同。
例如 $p=2,Q=8$：$C_8=1\cdot3\cdot5\cdot7=105$，而 $(8!)_2=315$。
\end{pitfall}


\section{Legendre 公式：阶乘中的素数指数}
\begin{theorem}
对素数 $p$ 与非负整数 $n$，
\[
 \nu_p(n!)=\sum_{j\ge1}\floor{n/p^j}
 =\frac{n-s_p(n)}{p-1},
\]


其中 $s_p(n)$ 是 $p$ 进制各位数字之和。$\nu_p(0!)=0$。
\end{theorem}

\begin{tip}[一句话版]
$\nu_p(n!)=\sum_{j\ge1}\floor{n/p^j}$：先数 $p$ 的倍数，再数 $p^2$ 的倍数（它们多贡献一个），一路数到超过 $n$ 为止。记不住公式时，用\textbf{$10!$ 里有几个 $2$}现场推一遍只要十秒。
\end{tip}

\begin{proof}
先把\textbf{按倍数数}写成等式，再化成数位和。

$n!$ 中恰好含有多少个因子 $p$？逐个看 $1,2,\ldots,n$ 各贡献 $\nu_p$，所以
\[
 \nu_p(n!)=\sum_{m=1}^{n}\nu_p(m).
\]
把每个 $m$ 的贡献按\textbf{它是 $p$ 的几次倍数的倍数}拆开：$\nu_p(m)=\sum_{j\ge1}\iva{p^j\mid m}$（$m$ 被 $p^j$ 整除就有第 $j$ 个 $p$）。代进去并交换求和：
\[
 \nu_p(n!)=\sum_{j\ge1}\ \sum_{m=1}^{n}\iva{p^j\mid m}=\sum_{j\ge1}\floor{\frac{n}{p^j}},
\]
因为 $1$ 到 $n$ 中 $p^j$ 的倍数恰有 $\floor{n/p^j}$ 个（即 $p^j,2p^j,\ldots$）。$j>\log_p n$ 时项为 $0$，和是有限的。

现在写 $n=\sum_i a_ip^i$、$0\le a_i<p$（这就是 $p$ 进制表示）。由 $\floor{n/p^j}=\sum_{i\ge j}a_ip^{i-j}$ 得
\[
 \sum_{j\ge1}\floor{\frac{n}{p^j}}
 =\sum_{j\ge1}\sum_{i\ge j}a_ip^{i-j}
 =\sum_{i}a_i\sum_{j=1}^{i}p^{\,i-j}
 =\sum_i a_i\big(1+p+\cdots+p^{i-1}\big)
 =\sum_i a_i\frac{p^i-1}{p-1}.
\]
而 $n=\sum_i a_ip^i$、数位和 $s_p(n)=\sum_i a_i$，于是
\[
 \nu_p(n!)=\frac{n-s_p(n)}{p-1}.
\]
最后这个式子就是 Kummer 定理要用到的形态：分子\textbf{$n$ 减数位和}正是这些位上的借位总量。
\end{proof}

\paragraph{算例.} $\nu_2(10!)=5+2+1=8$；二进制 $10=(1010)_2$，数位和为 $2$，也得到 $10-2=8$。$\nu_5(100!)=20+4=24$，所以 $100!$ 的十进制末尾零数为 $\min(\nu_2,\nu_5)=24$。
实现时反复令 $n\gets\floor{n/p}$ 并累加，而不是不断计算 $p^j$；后者可能溢出。


\section{Kummer 定理：进位就是损失的数位和}
\begin{theorem}
对 $0\le k\le n$，
\[
\nu_p\binom nk=\frac{s_p(k)+s_p(n-k)-s_p(n)}{p-1}.
\]


它等于在 $p$ 进制下把 $k$ 与 $n-k$ 相加时的进位次数，也等于从 $n$ 减 $k$ 的借位次数（逐位、包含连续借位）。
\end{theorem}

\begin{tip}[它的用法是\textbf{判零}而不是\textbf{算值}]
实战里几乎不用它算出 $\nu_p$ 的具体值，只用它的一句话推论：$p\nmid\binom nk\iff$ 竖式 $k+(n-k)$ 在 $p$ 进制下不进位。Lucas 定理就是这个推论的另一面。
\end{tip}


\begin{proof}
Legendre 公式（上一条）说 $\nu_p(m!)=\dfrac{m-s_p(m)}{p-1}$，所以
\[
 \nu_p\binom{n}{k}=\nu_p(n!)-\nu_p(k!)-\nu_p\big((n-k)!\big)
 =\frac{s_p(k)+s_p(n-k)-s_p(n)}{p-1}.
\]
只剩一句话要说清：为什么 $s_p(k)+s_p(n-k)-s_p(n)=(p-1)\times$（进位次数）。

看竖式加法 $k+(n-k)=n$。设从低到高共有 $c$ 次进位。若不进位，每一位的数码相加就等于该位结果；每发生一次\textbf{第 $i$ 位向第 $i+1$ 位进 $1$}，第 $i$ 位上留下的和就少 $p$，而第 $i+1$ 位多 $1$，因此两个加数的数位之和比结果的数位之和多 $p-1$。多个进位互不抵消：可以把进位链按发生顺序逐个消除，每消除一次总数位和减少 $p-1$，而进位次数也减少一次。归纳于 $c$ 即得
\[
 s_p(k)+s_p(n-k)=s_p(n)+c(p-1).
\]
代回上式即 $\nu_p\binom nk=c$。

这也给出一个实用推论：$\binom nk$ 不被 $p$ 整除，当且仅当竖式加法一次进位都没有，即每一位 $k_i\le n_i$——这正是下一条 Lucas 定理的内容。
\end{proof}
\paragraph{手算例.} $\binom{10}{4}=210$。二进制 $4=(0100)_2$，$6=(0110)_2$，相加得到 $(1010)_2$，在 $2^2$ 位发生一次进位，因此 $\nu_2(210)=1$。三进制 $4=(11)_3$、$6=(20)_3$ 相加为 $(101)_3$，也有一次进位，所以 $\nu_3(210)=1$。
\begin{teach}[连接到 Lucas]
无进位等价于每一位 $k_i\le n_i$，也等价于组合数不被 $p$ 整除。Lucas 计算余数，Kummer 计算被 $p$ 整除的深度，两者描述的是同一个数字结构的不同侧面。
\end{teach}


\section{Lucas 定理：逐位计算组合数}
约定 $\binom nk=0$（$k<0$ 或 $k>n$）。
\begin{theorem}[Lucas]
$p$ 为素数，$n=\sum_i n_ip^i,k=\sum_i k_ip^i$，则
\[
 \binom nk\equiv\prod_i\binom{n_i}{k_i}\pmod p.
\]


等价的递归形式为
\[
 \binom nk\equiv\binom{\floor{n/p}}{\floor{k/p}}
 \binom{n\bmod p}{k\bmod p}\pmod p.
\]
\end{theorem}

\begin{tip}[ Lucas 只在素数模数下成立]
证明依赖 $(1+X)^p\equiv1+X^p\pmod p$，这一步要求 $p$ 为素数（保证 $\binom pi$ 被 $p$ 整除）。模 $4$ 时它直接失效：$\binom42=6\equiv2\pmod4$，按\textbf{逐位相乘}却会得 $\binom{10}{10}$ 型的 $1$。合数模数要走\textbf{素数幂 + CRT}那条路。
\end{tip}


\begin{proof}[生成函数系数法]
二项式展开中 $\binom pi$（$0<i<p$）均被 $p$ 整除，所以
\[
(1+X)^p\equiv1+X^p\pmod p.
\]
不断使用此式，得到
\[
 (1+X)^n\equiv\prod_i(1+X^{p^i})^{n_i}.
\]
要取到 $X^k$，各因子应分别选 $X^{k_ip^i}$。因为每层指数小于 $p$，这个选择由 $p$ 进制表示唯一确定；系数相乘即为结论。若某位 $k_i>n_i$，相应系数为零。
\end{proof}

\begin{tip}[\textbf{生成函数}在这里其实就是多项式]
$\sum_k\binom nk X^k=(1+X)^n$：把\textbf{选法条数}当成一个多项式的系数。于是\textbf{模 $p$ 的 $\binom nk$}变成\textbf{多项式在 $\pmod p$ 下的某个系数}。用多项式的好处是可以相乘、可以只看某一项，而不用讨论组合意义。
\end{tip}

\paragraph{完整算例.} 模 $3$ 求 $\binom{10}{3}$。$10=(101)_3$、$3=(010)_3$，中间位出现 $1>0$，故余数为 $0$。模 $7$ 求 $\binom{10}{3}$，$10=(13)_7,3=(03)_7$，得到 $\binom10\binom33=1$，与 $120\bmod7=1$ 一致。
\paragraph{二进制结论.} $\binom nk$ 为奇数当且仅当 $k$ 的所有二进制 $1$ 位都包含于 $n$，可测试 $(k\mathbin{\&}n)=k$。
\cost 当 $p$ 足够小，可预处理 $0,\ldots,p-1$ 的阶乘与逆阶乘，单问 $O(\log_p n)$。若 $p$ 本身非常大，$O(p)$ 预处理不现实，需利用小 $k$ 等附加条件；Lucas 本身不是所有素数模数下的“常数时间组合数”。


\section{去掉素因子后的阶乘递推}
定义 $U_p(n)=(n!)_p$，即
\[
 n!=p^{\nu_p(n!)}U_p(n),\qquad\gcd(U_p(n),p)=1.
\]

\begin{tip}[$(n!)_p$ 与 $U_p(n)$ 的区别]
两个都很像、但不同：$(n!)_p$ 是把 $n!$ 里的 $p$ 因子全部除尽；而\textbf{去因子阶乘}$U_p(n)$ 是对 $1$ 到 $n$ 每个数先除尽 $p$ 因子，再相乘。数值上 $U_p(n)\equiv(n!)_p$，但递归推导时必须用后者的定义，才能把 $p$ 的倍数那部分单独提出来。
\end{tip}

固定 $Q=p^\alpha$，预处理
\[
 F_Q(t)=\prod_{\substack{1\le i\le t\\p\nmid i}}i\pmod Q,\quad0\le t\le Q,\qquad F_Q(0)=1.
\]
则
\[
 U_p(n)\equiv C_Q^{\floor{n/Q}}F_Q(n\bmod Q)
 U_p(\floor{n/p})\pmod Q.
\]
特别地，$Q=p$ 时
\[
 U_p(n)\equiv(-1)^{\floor{n/p}}(n\bmod p)!
 U_p(\floor{n/p})\pmod p.
\]

\begin{proof}
目标是把\textbf{去因子阶乘}$U_p(n)$（把 $1,\ldots,n$ 每个数中所有 $p$ 因子除尽后再相乘，模 $Q=p^\alpha$）写成递归。把 $1,\ldots,n$ 分成两部分。

第一部分：非 $p$ 倍数。模 $Q$ 的完全剩余系中，非 $p$ 倍数的集合是 $Q$ 周期重复的（因为\textbf{是否 $p$ 的倍数}与\textbf{同余模 $Q$}相容：$p\mid(x+tQ)\iff p\mid x$）。所以每连续 $Q$ 个整块的乘积都相同，等于 $F_Q:=\prod_{1\le j\le Q,\,p\nmid j}j\bmod Q$；一共 $\floor{n/Q}$ 块，剩下不足一块的尾部直接用预处理表查。

第二部分：$p$ 的倍数，即 $p,2p,3p,\ldots,\floor{n/p}p$。把其中 $pj$ 的一个 $p$ 除掉，得到 $j$；若 $j$ 还含因子 $p$，那部分与\textbf{$p^2$ 的倍数}在第一层已被处理，递归到 $\floor{n/p}$ 正好继续除尽。故这部分的贡献恰是 $U_p(\floor{n/p})$。

合并两部分：
\[
 U_p(n)\equiv F_Q^{\floor{n/Q}}\cdot(\text{尾部 }U_p(n\bmod Q))\cdot U_p(\floor{n/p})\pmod Q .
\]
深度为 $\log_p n$，每层 $O(1)$ 查表。最后要检查的是：$F_Q$ 只有 $\pm1$（Wilson 型论证：单位集合在取逆下不变，只有自逆元 $\pm1$ 无法配对；当 $Q$ 为 $2,4,p^\alpha,2p^\alpha$ 之外的模数时可能有更多自逆元，此时应直接按预处理表算，而不是套用 $\pm1$）。
\end{proof}
\paragraph{验算.} $p=2,Q=8,n=10$。$C_8=1$，$F_8(2)=1$；递归的 $n$ 为 $10,5,2,1$，得到
\[
 U_2(10)\equiv F_8(2)F_8(5)F_8(2)F_8(1)
 \equiv1\cdot7\cdot1\cdot1=7\pmod8.
\]
实际 $10!/2^8=14175$，余数也为 $7$。
\cost 预处理 $O(Q)$ 时间、空间。由于 $C_Q\in\{1,-1\}$，其幂只看奇偶；一次 $U_p(n)$ 仅 $O(\log_p(n+1))$ 次模乘。若每层仍用普通快速幂计算整块贡献，代码会额外产生一个对数因子，复杂度表必须与实现一致。


\section{扩展 Lucas：先计算素数幂，再 CRT}
对 $M=\prod_i p_i^{\alpha_i}$，分别计算模 $Q_i=p_i^{\alpha_i}$ 的答案。令
\[
 e=\nu_p(n!)-\nu_p(k!)-\nu_p((n-k)!),
\]

\begin{tip}[三步流程]
$Q=p^\alpha$：第一步用 Legendre 算指数 $e=\nu_p$，若 $e\ge\alpha$ 答案就是 $0$；第二步用\textbf{去因子阶乘}递归算单位部分，分母用逆元（单位一定可逆）；第三步对各个 $Q$ 的结果做 exCRT。每一步都要单独检查模数，不能合并成\textbf{一次求阶乘逆元}。
\end{tip}

则
\[
\binom nk\equiv
\begin{cases}
0,&e\ge\alpha,\\
p^eU_p(n)U_p(k)^{-1}U_p(n-k)^{-1},&e<\alpha
\end{cases}\pmod{p^\alpha}.
\]
分母的两个单位部分一定可逆，使用 exgcd 求逆即可。最后按互质素数幂进行 CRT。


\paragraph{从头算 $\binom{10}{4}\bmod12$.}
模 $4$：$e=8-3-4=1$，$U_2(10)\equiv3,U_2(4)\equiv3,U_2(6)\equiv1$，所以
\[
 \binom{10}{4}\equiv2\cdot3\cdot3^{-1}\cdot1^{-1}\equiv2\pmod4.
\]
模 $3$：$e=4-1-2=1\ge1$，余数为 $0$。解 $x\equiv2\pmod4,x\equiv0\pmod3$ 得 $x\equiv6\pmod{12}$，与 $210\bmod12=6$ 一致。

\paragraph{实现流程.}
\begin{enumerate}
\item 若 $M=1$，返回 $0$；若 $k<0$ 或 $k>n$，按约定返回 $0$。
\item 分解 $M$，对每个素数幂建立 $F_Q$。
\item 用 Legendre 求 $e$，先做零值判定，避免无用计算。
\item 求三个 $U$，只对单位求逆，再乘 $p^e$。
\item CRT 合并；禁止在模 $M$ 下先求 $k!$ 的逆元。
\end{enumerate}
\cost 不计模数分解，预处理 $O(\sum_iQ_i)$。使用单位块符号优化后，一次询问的主要成本为
\[
O\!\left(\sum_i\bigl(\log_{p_i}(n+1)+\log Q_i\bigr)\right)
\]
以及 CRT 合并成本。原稿将查询复杂度只写成模数的函数，漏掉了 $n$ 的递归层数。


\section{原稿例题的完整建模}
\begin{problem}{HDU 2973：YAPTCHA}
对 $n\le10^6$，多次求
\[
\sum_{k=1}^n\floor{\frac{(3k+6)!+1}{3k+7}-\floor{\frac{(3k+6)!}{3k+7}}}.
\]
\end{problem}

\begin{tip}[怎么做：把嵌套取整化成判定式]
看到\textbf{嵌套取整}+\textbf{素数判定}的组合（例如把 $\floor{\text{某式}}$ 当成示性函数），第一反应应该是：能否化成 Wilson 型判定，或者 $\gcd$ 型计数。这类题很少真的需要枚举。
\end{tip}

\hint 单项恰为 $[3k+7\text{ 为素数}]$，筛法加前缀和。

\begin{pitfall}[讲解：先看这题到底在做什么]
先观察题目的巨大阶乘表达式，不要真的去算阶乘。把 $(N-1)!$ 除以 $N$ 的余数写出来后，括号里的取整量恰好变成 Wilson 定理的判定条件，于是原题就被翻译成“统计某个区间里的素数”。
\end{pitfall}

\sol 设 $N=3k+7$，$(N-1)!=qN+r$，$0\le r<N$。括号内等于 $(r+1)/N$，取整为 $1$ 当且仅当 $r=N-1$。由 Wilson 的充要条件，这恰等价于 $N$ 为素数。
只需筛到 $3n_{\max}+7$，令 $A_k=A_{k-1}+[3k+7\in\Pp]$。例如 $n=5$ 时考察 $10,13,16,19,22$，答案为 $2$。
\cost 埃氏筛 $O(N\log\log N)$ 或线性筛 $O(N)$，每问 $O(1)$。本题绝不计算阶乘；超大表达式常常是在提示一个判定定理。

\src{https://acm.hdu.edu.cn/showproblem.php?pid=2973}

\begin{problem}{P2183：礼物}
$n$ 件互不相同的礼物送给 $m$ 个有编号的人，第 $i$ 人收到 $a_i$ 件，允许剩余。求方案数模 $M$；若不可能则按原题输出 \texttt{Impossible}。
$n\le10^9,m\le5$，模数的每个素数幂因子不超过 $10^5$。
\end{problem}
\hint $\sum_i a_i>n$ 时无解；否则是一个多项式系数，可写为一串组合数的乘积。

\begin{pitfall}[讲解：先看这题到底在做什么]
把“依次给每个人分配 $a_i$ 件”看成一连串组合选择：第 $i$ 个人从还剩下的物品中选出 $a_i$ 件。这样每一步独立形成一个组合数，乘起来后自然得到阶乘商。关键是先定义前缀总量 $s_i$，把每一步还剩多少说清楚。
\end{pitfall}

\sol 令 $s_0=0,s_i=\sum_{j=1}^i a_j$，则
\[
\mathrm{Ans}=\prod_{i=1}^m\binom{n-s_{i-1}}{a_i}
=\frac{n!}{(n-s_m)!\prod_i a_i!}.
\]

\begin{tip}[提醒：$O(N)$ 不等于省内存]
$O(N)$ 不等于省内存。$N=10^8$ 时，多个 \texttt{int} 数组就是 $400$\,MB 级别。只判素数用位压缩的埃氏筛；要 $\varphi,\mu,lp$ 再上线性筛，并先算内存再写代码。
\end{tip}

逐人选礼物：第一人从 $n$ 件中选，第二人从剩余中选，依此类推。每个完整分配恰计一次，不需要再除 $m!$，因为人有编号。
可以调用 $m$ 次扩展 Lucas，或者直接在每个素数幂上对多项式系数统计单位部分与指数。

\paragraph{算例.} $n=4,m=2,(a_1,a_2)=(1,2)$，答案 $\binom41\binom32=12$。模 $6$ 后为 $0$，但\textbf{这是有方案而余数为零}，不能输出 Impossible。只有数量约束失败才输出 Impossible。

\src{https://www.luogu.com.cn/problem/P2183}

\begin{problem}{P7488：黑色毛衣}
选 $m$ 个整数边长的等腰三角形，腰长 $a\in[1,n]$、底长 $b\in[1,2a-1]$，要求所有腰长互不相同、所有底长也互不相同，忽略三角形的排列顺序。求组数模一般模数 $M$。$m\le n\le10^{18},M\le10^5$。
\end{problem}
\hint 答案是 $\binom nm\,n^{\underline m}=\binom nm^2m!$，先证明计数，再处理不可逆阶乘。

\begin{pitfall}[讲解：先看这题到底在做什么]
先忽略“模数可能不能除阶乘”的实现问题，只解决计数。把腰长看成行、底长看成列，就是一个阶梯棋盘上的不重复选点问题；最后一行是否被使用会给出递推式，再把递推化成组合数。
\end{pitfall}

\sol 把腰长 $a$ 看成棋盘的第 $a$ 行，该行可选列为 $1,\ldots,2a-1$。选三角形等价于在行、列都不能重复的阶梯棋盘上放 $m$ 个车。
记方案数为 $R(n,m)$。不使用最后一行的方案数为 $R(n-1,m)$。使用最后一行时，前面已经放了 $m-1$ 个车，占用 $m-1$ 个不同列；这些列全部在最后一行的范围内，所以还有
\[
 (2n-1)-(m-1)=2n-m
\]
个位置。因此
\[
 R(n,m)=R(n-1,m)+(2n-m)R(n-1,m-1),\quad R(n,0)=1.
\]
将 $\binom nm^2m!$ 代入即可验证递推（先提出公共因子，或使用 Pascal 恒等式），初值也相同，所以由归纳法成立。

\paragraph{小规模核对.} $n=3,m=1$ 时为 $1+3+5=9$；$n=3,m=2$ 时，所选两行分别为 $(1,2),(1,3),(2,3)$，方案数为 $2,4,12$，总计 $18=\binom32^2\cdot2!$。

\paragraph{合数模数下如何算？}
使用
\[
 R(n,m)=\frac{(n!)^2}{m!((n-m)!)^2}.
\]
对每个 $p^\alpha\mid M$，总指数为
\[
 e=2\nu_p(n!)-\nu_p(m!)-2\nu_p((n-m)!),
\]
单位部分相同地相乘求逆，再 CRT。不能先算 $m!\bmod M$，再据此决定整个分式是否为零。
\begin{pitfall}
题意不是要求全部 $2m$ 个“腰长数值与底长数值”两两不同。一个三角形可以是等边三角形，某只的底长也可以等于另一只的腰长。若误删 $b=a$，计数公式就变了。
\end{pitfall}

\src{https://www.luogu.com.cn/problem/P7488}


\chapter{数论函数、狄利克雷卷积与筛法}
\begin{chapterintro}[章节导读]
前面几章主要研究单个整数的结构；这一章开始把视野扩大到定义在所有正整数上的函数。先区分积性与完全积性，再用狄利克雷卷积把约数求和重新组织起来，最后把这些公式转化成可以在线性或近线性时间运行的筛法。
\end{chapterintro}

\section{先分清积性与完全积性}
\begin{definition}
数论函数是定义在正整数上的函数，通常取值于复数或某个模数环。
若 $f(1)=1$，且 $\gcd(a,b)=1$ 时 $f(ab)=f(a)f(b)$，称为积性函数。
若不要求互质也恒有 $f(ab)=f(a)f(b)$，称为完全积性函数。
\end{definition}
积性函数由所有素数幂上的取值唯一确定：$f(\prod p_i^{e_i})=\prod f(p_i^{e_i})$。

\begin{longtable}{p{0.17\textwidth}p{0.37\textwidth}p{0.34\textwidth}}
\toprule
函数 & 定义 & $p^k$ 处的值（$k\ge1$）\\\midrule
$\eps$ & $\eps(n)=[n=1]$ & $0$\\
$\one$ & $\one(n)=1$ & $1$\\
$\mu$ & 有平方因子时为 $0$；否则为 $(-1)^{\omega(n)}$ & $-1$（$k=1$），$0$（$k\ge2$）\\
$\tau$ & 正约数个数 & $k+1$\\
$\sigma_t$ & 所有正约数的 $t$ 次幂之和 & $1+p^t+\cdots+p^{kt}$\\
$\varphi$ & 与 $n$ 互质的正整数个数 & $(p-1)p^{k-1}$\\
$\id_t$ & $\id_t(n)=n^t$ & $p^{kt}$\\
\bottomrule
\end{longtable}

\paragraph{最小反例.} $\varphi$ 积性但不完全积性：$\varphi(4)=2\ne\varphi(2)^2=1$。$\mu(4)=0\ne\mu(2)^2=1$。将积性函数当成完全积性，是后文洲阁筛中最危险的错误之一。
\paragraph{为什么 $\mu$ 长这样？}
对 $n=\prod_{i=1}^r p_i^{e_i}>1$，只有平方自由约数贡献非零项，因此
\[
 \sum_{d\mid n}\mu(d)=\sum_{S\subseteq\{1,\ldots,r\}}(-1)^{|S|}=(1-1)^r=0.
\]

\begin{tip}[词义：积性 / 完全积性]
一句话记法：积性 = \textbf{互质才能拆}，完全积性 = \textbf{随便拆}。检查函数属于哪一种，只要代 $a=b=2$ 试一次；$\tau(4)=3\ne4$ 就说明 $\tau$ 只是积性。后面洲阁筛里把 $x^t$（完全积性）与积性函数混用会直接算错。
\end{tip}

\begin{tip}[词义：把 $\mu(d)$ 当开关]
\textbf{只有……贡献非零项}的意思是：把 $\mu(d)$ 当成开关，非平方自由的 $d$ 一律 $\mu(d)=0$，求和时自动跳过。所以 $\sum_{d\mid n}\mu(d)$ 只取决于 $n$ 有几个不同素因子：$2^r$ 个平方自由约数，一半带 $+$、一半带 $-$，除非 $n=1$。
\end{tip}

\begin{tip}[\textbf{由素数幂确定}有多省]
例：$72=2^3\cdot3^2$，只要知道 $\varphi(8)=4$、$\varphi(9)=6$，就有 $\varphi(72)=24$，不必枚举 $72$ 个数。做题时的动作是：先把 $n$ 分解，再逐素数幂查公式，最后相乘；写成代码就是\textbf{$\log n$ 的因子数}次查表。
\end{tip}

$n=1$ 时只有约数 $1$，结果为 $1$。它本质上是对不同素因子的容斥。


\section{欧拉函数的公式与恒等式}
容斥或 CRT 可得
\[
 \varphi(n)=n\prod_{p\mid n}\left(1-\frac1p\right).
\]
因此 $\varphi$ 是积性函数，且对任意正整数 $a,b$、$d=\gcd(a,b)$，
\[
 \varphi(ab)\varphi(d)=\varphi(a)\varphi(b)d.
\]

\begin{proof}[第二个恒等式]
按素因子比较。当 $p$ 只整除 $a,b$ 中一个时，两边相同；当它同时整除两者时，设赋值为 $u,v\ge1$，两边关于 $p$ 的因子分别是
\[
 (p-1)^2p^{u+v+\min(u,v)-2},
\]
仍相同。把各个素数的贡献乘起来即可。
\end{proof}
对 $n=12$，公式为 $12(1-1/2)(1-1/3)=4$，对应 $1,5,7,11$。实现时用 $ans=ans/p\cdot(p-1)$，避免浮点误差。


\section{狄利克雷卷积与逆元}
\begin{definition}
\[
 (f*g)(n)=\sum_{d\mid n}f(d)g(n/d)=\sum_{ab=n}f(a)g(b).
\]
\end{definition}
卷积满足交换律、结合律，并对逐点加法分配；其单位元是 $\eps$，不是 $\one$。
若 $g(1)\ne0$（在模数环中需为可逆元素），则 $g*h=f$ 有唯一解。递推式为
\[
 h(1)=\frac{f(1)}{g(1)},\qquad
 h(n)=\frac{f(n)-\sum_{\substack{d\mid n\\d>1}}g(d)h(n/d)}{g(1)}.
\]

\begin{tip}[自己验证一次只要两行]
交换律：把 $(f*g)(n)=\sum_{d\mid n}f(d)g(n/d)$ 里的 $d$ 换成 $n/d$（这是约数集合上的对合，一一自对应），求和就变成 $\sum_{d\mid n}f(n/d)g(d)=(g*f)(n)$。结合律：两边都等于 $\sum_{abc=n}f(a)g(b)h(c)$。凡是\textbf{显然满足}的说法，都能这样两行落地。
\end{tip}

两个积性函数的卷积仍积性；积性函数的卷积逆及积性函数之间的卷积商仍积性。


\begin{proof}[为什么卷积保持积性？]
若 $\gcd(a,b)=1$，每个 $ab$ 的约数唯一写成 $uv$，$u\mid a,v\mid b$。于是
\[
 (f*g)(ab)=\sum_{u\mid a}\sum_{v\mid b}f(u)f(v)g(a/u)g(b/v)
 =(f*g)(a)(f*g)(b).
\]
对卷积商，先在每个素数幂上递推求值，再由积性延拓得到 $h'$，可验证 $g*h'=f$；唯一性保证 $h'=h$。
\end{proof}
\paragraph{局部生成函数.}
固定一个素数 $p$，令 $F_p(z)=\sum_{k\ge0}f(p^k)z^k$。那么卷积对应形式幂级数乘法
\[
 H_p(z)=F_p(z)G_p(z).
\]
这里说的是形式幂级数，不要求它真的收敛。$\one$ 的局部级数为 $(1-z)^{-1}$，$\mu$ 为 $1-z$，所以 $\one*\mu=\eps$。
\begin{pitfall}
写 $\varphi=\id*\mu$ 是卷积恒等式，不是 $\varphi(n)=n\mu(n)$。
同样，“两个积性函数相加”未必积性，例如 $\one+\one$ 在 $1$ 处为 $2$，已不满足归一化定义。
\end{pitfall}


\section{线性筛：唯一的最小素因子分解}
对合数 $x$，令 $p$ 为其最小素因子、$i=x/p$，则 $p\le\operatorname{minp}(i)$。
线性筛枚举 $i$，再递增枚举素数 $p$，标记 $ip$；一旦 $p\mid i$ 即标记后退出内层循环。每个合数恰被处理一次。
计算 $\varphi,\mu$ 时使用
\[
\begin{array}{c|cc}
 &p\mid i&p\nmid i\\\hline
\varphi(ip)&p\varphi(i)&(p-1)\varphi(i)\\
\mu(ip)&0&-\mu(i)
\end{array}
\]

\begin{tip}[为什么：每个合数只被最小素因子筛一次]
$O(N)$ 的关键是\textbf{每个合数只被它的最小素因子筛掉一次}：循环遇到 $i\bmod p==0$ 就 break，保证不会用更大的素因子去生成同一个数。$\varphi$、$\mu$ 能同步线性求出，也是靠这条唯一分解。
\end{tip}

\begin{tip}[\textbf{形式}两个字的含义]
只看系数、不问 $z$ 取什么值：$\sum a_kz^k$ 就是一张系数表，乘法是\textbf{系数的卷积}。所以 $(1-z)^{-1}=1+z+z^2+\cdots$ 在这里是\textbf{定义}而不是\textbf{收敛计算}——你要验证的是两边系数相乘后 $z^k$ 的系数为 $1$，而不是级数是否收敛。
\end{tip}

并设置 $\varphi(1)=\mu(1)=1$。


\paragraph{为何要在整除时退出？}
若 $p\mid i$，那么 $p$ 已是 $i$ 的最小素因子；继续枚举更大的 $q$，$iq$ 的最小素因子仍是 $p$，它会在唯一的分解 $p\cdot(iq/p)$ 中被处理。继续做只会重复，并破坏线性复杂度的计数依据。
\begin{lstlisting}
vector<int> primes, lp(N+1), phi(N+1), mu(N+1);
phi[1]=mu[1]=1;
for (int i=2;i<=N;++i) {
    if (!lp[i]) {
        lp[i]=i; primes.push_back(i);
        phi[i]=i-1; mu[i]=-1;
    }
    for (int p:primes) {
        if (p>N/i) break;
        int v=i*p; lp[v]=p;
        if (i%p==0) {
            phi[v]=phi[i]*p; mu[v]=0;
            break;
        }
        phi[v]=phi[i]*(p-1); mu[v]=-mu[i];
    }
}
\end{lstlisting}
\paragraph{扩展到一般积性函数.} 当 $p\nmid i$ 时 $f(ip)=f(i)f(p)$；当 $p\mid i$ 时应把 $i=p^e u$ 分离出来，维护最小素因子的指数和对应幂，算 $f(p^{e+1})f(u)$。不能假定任意积性函数都存在“除以 $f(p^e)$”的安全转移：该值可能为零或不可逆。
\cost 时间 $O(N)$，常规数组空间 $O(N)$。$N=10^8$ 时四个 int 数组已约 $1.6$ GB，应按需要只保留必要数据；若只筛素数，位图埃氏筛可能更合适。


\section{整除分块：相同商的最大区间}
对 $1\le l\le n$，令 $q=\floor{n/l}$，则商保持 $q$ 的区间右端点为
\[
 r=\floor{n/q}.
\]

\begin{tip}[$r=\floor{n/q}$ 是怎么想到的]
固定 $l$、$q=\floor{n/l}$，要找最大的 $r$ 使 $\floor{n/x}=q$ 对整段成立。条件等价于 $q\le n/x<q+1$，左半边给出 $x\le n/q$，所以 $r=\floor{n/q}$（右半边给出 $x\ge\floor{n/(q+1)}+1$，由 $l$ 已经满足）。代码里 $l=r+1$ 跳到下一块，块数 $O(\sqrt n)$ 来自\textbf{小 $l$ 至多 $\sqrt n$ 个，大的商也不超过 $\sqrt n$}。
\end{tip}

这是因为 $\floor{n/x}\ge q\iff x\le\floor{n/q}$。所有不同商只有 $O(\sqrt n)$ 个：$x\le\sqrt n$ 的位置不超过 $\sqrt n$，其余位置的商也不超过 $\sqrt n$。

\paragraph{例子.} $n=10$ 的商块为
\[
[1,1]\mapsto10,\ [2,2]\mapsto5,\ [3,3]\mapsto3,\ [4,5]\mapsto2,\ [6,10]\mapsto1.
\]
若同时含 $\floor{n/i},\floor{m/i}$，共同右端点应取
\[
 r=\min\left(\floor{n/\floor{n/l}},\floor{m/\floor{m/l}}\right),
\]
并仅枚举到 $\min(n,m)$。不能用其中一个商的区间跨过另一个商的变化点。
\begin{lstlisting}
for (long long l=1,r;l<=n;l=r+1) {
    long long q=n/l;
    r=n/q;
    // sum over [l,r] = (prefix[r]-prefix[l-1]) * F(q)
}
\end{lstlisting}


\section{狄利克雷前缀和、后缀和与差分}
“前缀”在此按整除偏序，不是普通区间前缀。
\[
 Z_f(n)=\sum_{d\mid n}f(d)=(f*\one)(n),\qquad
 U_f(n)=\sum_{\substack{k\le N\\n\mid k}}f(k).
\]

\begin{tip}[这三个词在说什么]
把 $f\mapsto g$，$g(n)=\sum_{d\mid n}f(d)$（每个位置吸收它的约数）叫狄利克雷\textbf{前缀和}；$g(n)=\sum_{n\mid d}f(d)$（每个位置吸收它的倍数）叫\textbf{后缀和}；它们的逆运算统称\textbf{差分}。实现就是枚举 $d$ 或枚举倍数，复杂度 $O(N\log N)$，本质是调和级数。
\end{tip}

利用素数逐层传播，可在 $O(N\log\log N)$ 完成这些变换；反向操作给出差分。


\begin{lstlisting}
// a[1..N] initially stores f. primes contains all primes <=N.
// Divisor zeta: a[n] <- sum_{d|n} f[d]
for (int p:primes)
    for (int i=1;i<=N/p;++i) a[i*p]+=a[i];

// Its inverse (run on the transformed array):
for (auto it=primes.rbegin();it!=primes.rend();++it)
    for (int i=N/(*it);i>=1;--i) a[i*(*it)]-=a[i];

// Multiple zeta: a[n] <- sum_{n|k, k<=N} f[k]
for (int p:primes)
    for (int i=N/p;i>=1;--i) a[i]+=a[i*p];
\end{lstlisting}
\paragraph{为何前缀递增、后缀递减？}
前缀时希望 $p^2$ 能收到已经汇总到 $p$ 的贡献，故从小到大；后缀时希望 $1$ 能收到已经汇总到 $p$ 的贡献，故从大到小。
固定 $p$，整数按“不含 $p$ 的部分”分成链 $u,up,up^2,\ldots$，上述循环恰在每条链上做普通前缀和或后缀和。不同素数维度依次处理，正好枚举每个约数一次。

\paragraph{推广：任意函数卷一个积性函数.}
已知积性 $g$ 的素数幂点值，逐个素数在上述幂链上做一维卷积即可。长度为 $\ell+1$ 的链即使朴素 $O(\ell^2)$，固定 $p$ 的非平凡链总成本仍为 $O(N/p)$；因此总计 $O(N\log\log N)$。这里默认素数幂点值可常数时间读取，且只枚举起点 $u\le N/p,p\nmid u$，不能每个素数都扫完整个 $[1,N]$。



\chapter{莫比乌斯反演与 gcd 型求和}
\begin{chapterintro}[章节导读]
上一章我们学会了用卷积组织“对所有约数求和”的关系；有了这套语言，就可以反过来把被混在一起的贡献拆开。莫比乌斯反演正是这种“消掉约数层层贡献”的工具，它会进一步把许多看似复杂的 gcd 求和变成可以交换顺序、直接计算的形式。
\end{chapterintro}

\section{两种基础“拆贡献”恒等式}
\[
\sum_{d\mid n}\mu(d)=[n=1],\qquad
\sum_{d\mid n}\varphi(d)=n.
\]

\begin{tip}[词义：拆贡献]
\textbf{拆贡献}= 不直接按 $i$ 算，而是按\textbf{谁对答案造成多少}重新分组。两种写法：按约数拆 $\sum_{d\mid n}$、按倍数拆 $\sum_{d\mid k}$。选哪一种，取决于化简后哪个和式容易算——通常先各自试两行，再决定。
\end{tip}

卷积写法为 $\mu*\one=\eps$、$\varphi*\one=\id$，从而 $\varphi=\id*\mu$。
\begin{theorem}[莫比乌斯反演]
\[
 F(n)=\sum_{d\mid n}f(d)
 \iff f(n)=\sum_{d\mid n}\mu(d)F(n/d).
\]
若函数在有限范围外为零，则还有倍数版本
\[
 F(n)=\sum_{k\ge1}f(kn)
 \iff f(n)=\sum_{k\ge1}\mu(k)F(kn).
\]
\end{theorem}

\begin{proof}[因数版本]
两个方向都写全，且把\textbf{哪些和式合法}说清楚。

因数版本：若 $g(n)=\sum_{d\mid n}f(d)$，则
\[
 \sum_{d\mid n}\mu(d)g(n/d)
 =\sum_{d\mid n}\mu(d)\sum_{e\mid n/d}f(e).
\]
右边是有限和，可以自由换序。$(d,e)$ 跑遍\textbf{$d\mid n$ 且 $e\mid n/d$}，等价于\textbf{$e\mid n$ 且 $d\mid n/e$}（因为 $de\mid n$）。于是
\[
 \sum_{e\mid n}f(e)\sum_{d\mid n/e}\mu(d).
\]
内层和是经典判别式：$\sum_{d\mid m}\mu(d)=\iva{m=1}$（按 $\mu$ 的定义，$m=\prod p^{u_p}$ 时内层为 $\prod_p(1+(-1)+\cdots+(-1)^{u_p})$，只要某个 $u_p\ge1$ 就有因子 $1-1=0$；$m=1$ 时为 $1$）。因此只剩 $n/e=1$，即 $e=n$ 的项，得到 $f(n)$。

倍数版本：若 $g(n)=\sum_{n\mid d,\ d\le N}f(d)$，则
\[
 \sum_{k\ge1}g(kn)\mu(k)=\sum_{k\ge1}\mu(k)\sum_{\ell\ge1}f(k\ell n)
 =\sum_{m\le N,\,n\mid m}f(m)\!\!\sum_{k\mid m/n}\!\!\mu(k)=f(n),
\]
最后一步同样是把 $k\ell n$ 合成一个变量 $m$，并用 $\sum_{d\mid m/n}\mu=\iva{m/n=1}$。

两个版本的区别只在于\textbf{$f$ 是挂在约数上还是挂在倍数上}，用哪一版取决于你先把原式整理成哪一种。
\end{proof}
\paragraph{欧拉恒等式的组合证明.}
把 $1,\ldots,n$ 按 $\gcd(i,n)=g$ 分类。写 $i=gj$，该类恰有 $\varphi(n/g)$ 个数；各类总数为 $n$，换元 $d=n/g$ 得到公式。
\begin{teach}[怎样想到反演？]
看到 $[\gcd(i,j)=1]$，替换成 $\sum_{d\mid i,d\mid j}\mu(d)$；看到 $\gcd(i,j)$，替换成 $\sum_{d\mid i,d\mid j}\varphi(d)$；看到任意 $C(\gcd(i,j))$，先求 $D=C*\mu$，再写 $C=\one*D$。反演的目的不是“制造更多求和号”，而是让 $i,j$ 分离。
\end{teach}


\section{先处理单变量与互质计数}
\[
 \sum_{i=1}^n\gcd(i,n)
 =\sum_{d\mid n}\varphi(d)\floor{n/d}
 =\sum_{d\mid n}d\varphi(n/d).
\]

\begin{tip}[\textbf{拆贡献}的三步套路]
第一步把要求的东西写成 $\sum_{i}f(\gcd(\cdots))$ 之类；第二步用 $f(n)=\sum_{d\mid n}c(d)$ 把 $f$ 换掉，于是 $f(\gcd)$ 变成\textbf{$d$ 整除各变量}的 $\iva{}$ 条件；第三步交换求和顺序，内层只剩\textbf{倍数个数}。整套动作没有灵感要求，练十道题就会背下来。
\end{tip}

以及
\[
 \#\{(i,j):1\le i\le n,1\le j\le m,\gcd(i,j)=1\}
 =\sum_{d=1}^{\min(n,m)}\mu(d)\floor{n/d}\floor{m/d}.
\]

\paragraph{第一式的每一步.}
\[
\sum_i\gcd(i,n)
=\sum_i\sum_{d\mid i,d\mid n}\varphi(d)
=\sum_{d\mid n}\varphi(d)\sum_i[d\mid i].
\]
原稿在插入约数求和后仍保留 $\gcd(i,n)$，会把同一个 gcd 重复累加约数个数次；这里必须替换为 $\varphi(d)$。
例如 $n=6$，直接求和 $1+2+3+2+1+6=15$；约数形式给出 $6\cdot1+3\cdot1+2\cdot2+1\cdot2=15$。
\paragraph{互质计数验算.} $n=m=3$ 时 $9-1-1=7$；九个有序数对中只有 $(2,2),(3,3)$ 不互质。注意计的是有序数对，不能再额外乘二。


\begin{problem}{P2398：GCD SUM}
求 $\sum_{i=1}^n\sum_{j=1}^n\gcd(i,j)$。
\end{problem}
\hint 答案为 $\sum_{d=1}^n\varphi(d)\floor{n/d}^2$，可筛 $\varphi$ 并利用前缀和做整除分块。

\begin{pitfall}[讲解：先看这题到底在做什么]
看到双重求和里的 gcd，先把 gcd 用欧拉展开成“所有公共约数的权重之和”。交换求和后，两个坐标只剩下“有多少个数同时被 $d$ 整除”，这就是 $\floor{n/d}$，于是二维问题变成一维整除分块。
\end{pitfall}

\sol 把 gcd 的欧拉展开代入，交换三层求和：
\[
\sum_{i,j}\sum_{d\mid i,d\mid j}\varphi(d)
=\sum_d\varphi(d)\left(\sum_i[d\mid i]\right)
\left(\sum_j[d\mid j]\right).
\]
两个括号都为 $\floor{n/d}$，得到公式。
$n=3$ 时矩阵为
\[
\begin{pmatrix}1&1&1\\1&2&1\\1&1&3\end{pmatrix},
\]
总和 $12$，公式也给 $9+1+2=12$。
\cost 线性筛预处理 $O(N)$ 后，每问分块 $O(\sqrt n)$。若只有一个不大的 $n$，直接累加 $n$ 项也是 $O(n)$，不用为了形式上“高级”而特意分块。前缀和、商的平方和最终答案都要用足够宽的整数。

\src{https://www.luogu.com.cn/problem/P2398}

\begin{problem}{一般加权 gcd 求和}
给定三个长度为 $n$ 的数组 $A,B,C$，求
\[
 \sum_{i=1}^n\sum_{j=1}^n A(i)B(j)C(\gcd(i,j)),\qquad n\le2\times10^6.
\]
\end{problem}

\begin{tip}[怎么做：先做莫比乌斯差分，再数倍数]
凡是答案形如 $\sum f(\gcd(a_i,b_j))$，机械做法都一样：把 $f$ 写成 $f=\sum_{d\mid n}c(d)$（即对 $f$ 做莫比乌斯差分），于是 $\gcd$ 变成\textbf{$d$ 同时整除两个数}的计数，后面全是倍数计数与容斥。
\end{tip}

\hint 求 $D=C*\mu$，以及 $A,B$ 的倍数后缀和，答案变成一维点积。

\begin{pitfall}[讲解：先看这题到底在做什么]
这里把上一题的核心机制抽象出来：只要一个函数能够写成约数卷积形式 $C=\one*D$，那么 $C(\gcd(i,j))$ 就可以按公共因子 $d$ 拆贡献。之后分别预处理两个数组中“被 $d$ 整除”的元素总和，答案就成为一维点积。
\end{pitfall}

\sol 由 $C=\one*D$，
\[
\begin{aligned}
\mathrm{Ans}
&=\sum_{i,j}A(i)B(j)\sum_{d\mid i,d\mid j}D(d)\\
&=\sum_{d=1}^nD(d)
 \left(\sum_{d\mid i}A(i)\right)
 \left(\sum_{d\mid j}B(j)\right).
\end{aligned}
\]
$A,B,C$ 都不要求积性。上一章的素数维度变换适用于任意数组：对 $C$ 做因数前缀的逆变换，得到 $D$；对 $A,B$ 各做倍数后缀，最后逐点相乘求和。
\cost 时间 $O(n\log\log n)$，空间 $O(n)$。若用朴素枚举所有倍数，也能得到 $O(n\log n)$ 的更易写版本，应根据时限选择。
\paragraph{追问.} 取 $A=B=1,C(n)=n$ 会恢复哪题？恢复 GCD SUM，因为 $D=\varphi$。取 $C(n)=[n=1]$ 则 $D=\mu$，恢复互质数对计数。


\section{多询问与不规则预处理表}
\begin{problem}{P4240：毒瘤之神的考验}
多次求
\[
 \sum_{i=1}^n\sum_{j=1}^m\varphi(ij)\pmod{998244353},
\]


其中 $T\le10^4$，$n,m\le N=10^5$。
\end{problem}

\begin{tip}[词义：不规则的预处理表]
\textbf{不规则}指：表不能建成整齐的 $N\times N$，只能按 $\floor{n/k}$ 的取值分段存储。做法是把\textbf{大参数}（$>K$）的情况列成一张稀疏表，查询时先定位区间端点，再取前缀差。
\end{tip}

\hint 先用 $\varphi(ij)=\varphi(i)\varphi(j)\gcd(i,j)/\varphi(\gcd(i,j))$ 转成加权 gcd；再预处理倍数方向的部分和。不能每次询问都做 $O(N)$ 全扫。

\begin{pitfall}[讲解：先看这题到底在做什么]
先处理函数本身：把 $k/\varphi(k)$ 写成卷积形式，从而制造出适合公共因子拆分的系数。完成这一步后，题目就退化成一般加权 gcd 求和；实现上真正需要优化的是倍数方向的预处理，而不是重新设计一套计数公式。
\end{pitfall}

\sol 把函数 $A(k)=k/\varphi(k)$ 拆为 $A=\one*f$，则
\[
 f(k)=\sum_{d\mid k}\frac{d\mu(k/d)}{\varphi(d)}.
\]
套用一般加权 gcd 公式，得到
\[
\boxed{\mathrm{Ans}(n,m)=\sum_{k=1}^{\min(n,m)}f(k)
G\left(k,\floor{n/k}\right)G\left(k,\floor{m/k}\right)},
\]
其中 $G(k,t)=\sum_{i=1}^t\varphi(ki)$。

\paragraph{第一层优化：化简原稿中未继续化简的 $f$.}
$A$ 积性，因此 $f=A*\mu$ 积性；在素数幂处
\[
 f(p)=\frac p{p-1}-1=\frac1{p-1},\qquad
 f(p^e)=0\ (e\ge2).
\]
所以
\[
 f(k)=\frac{\mu(k)^2}{\varphi(k)}.
\]
在本题模数下所有 $\varphi(k)\le N<998244353$ 都可逆。这个形式可以线性筛或直接用 $\mu,\varphi$ 表得到。

\paragraph{第二层优化：只存实际存在的 $G$ 项.}
对于固定 $k$，只需 $0\le t\le\floor{N/k}$，递推
\[
 G(k,0)=0,\quad G(k,t)=G(k,t-1)+\varphi(kt).
\]
总项数 $\sum_k\floor{N/k}=O(N\log N)$。不要开 $N\times N$ 矩阵；使用按 $k$ 分行的变长数组或一段连续内存。

\paragraph{第三层优化：小 $k$ 直接算，大 $k$ 按商分块.}
取阈值参数 $1\le B\le N$，令 $K=\floor{N/B}$。$k\le K$ 时每问直接枚举；$k>K$ 时两个商均不超过 $B$。
对 $1\le u,v\le B$ 预处理
\[
 H_{u,v}(r)=\sum_{K<k\le r}f(k)G(k,u)G(k,v),
 \quad K\le r\le\floor{N/\max(u,v)}.
\]
对没有合法 $r$ 的 $(u,v)$ 不开表。若询问的共同商区间为 $[l,r]$、$u=\floor{n/l},v=\floor{m/l}$，该块贡献就是
\[
 H_{u,v}(r)-H_{u,v}(l-1).
\]
共同右端点为 $\min(n/u,m/v)$，所以索引一定处于该行合法范围。

\paragraph{为什么这张三维表开得下？}
若对所有 $(u,v,r)$ 开等长数组，需要 $O(NB^2)$，但实际 $r$ 的上界依赖 $\max(u,v)$。只开有效范围，项数不超过
\[
\sum_{u=1}^B\sum_{v=1}^B\floor{N/\max(u,v)}
\le N\sum_{w=1}^B\frac{2w-1}{w}=O(NB).
\]
等价地，对每个 $k>K$ 枚举 $u,v\le N/k$，成本为
$\sum_{k>K}O((N/k)^2)=O(N^2/K)=O(NB)$。
$H_{u,v}=H_{v,u}$ 还可以再减半存储。

\cost 总预处理时间、空间为 $O(N\log N+NB)$；每次询问 $O(N/B+\sqrt N)$。例如 $N=10^5,B\approx100$ 时，直接枚举约 $1000$ 项，其余使用商块；内存按真实表项计算后选择阈值。这里的界对应\textbf{变长有效表}实现，不能套在未经压缩的三维数组上。

\paragraph{小例验算.} $n=2,m=3$，直接求和为
$\varphi(1)+\varphi(2)+\varphi(3)+\varphi(2)+\varphi(4)+\varphi(6)=9$。
公式中 $k=1$ 贡献 $2\cdot4=8$，$k=2$ 贡献 $1\cdot1\cdot1=1$，合计 $9$。
\begin{teach}[这道题的三个教学收获]
把复杂二维函数转为 gcd 权重；在素数幂处继续化简卷积系数；开多维表前先写出每一维的合法范围。原稿“前面后面分一下”不足以支撑实现，应把三个函数的参数含义与存储上界完整写出。
\end{teach}

\src{https://www.luogu.com.cn/problem/P4240}


\chapter{积性函数求和：杜教筛与 Powerful Number}
\begin{chapterintro}[章节导读]
上一章的反演解决了如何拆贡献，但如果 n 本身已经大到不能逐项枚举，新的瓶颈就变成“这些函数值怎么快速求前缀和”。这一章从商集和倍增估计出发，再进入杜教筛与 Powerful Number，用结构化的递归和稀疏支撑代替完整枚举。
\end{chapterintro}

\section{先看规模，再看倍增估计}
目标是求 $S_f(n)=\sum_{i\le n}f(i)$。当 $n$ 很大时，还经常要求所有 $v=\floor{n/k}$ 的 $S_f(v)$，本讲义称其为商集上的前缀和表（原稿称“块筛”）。不同 $v$ 只有 $O(\sqrt n)$ 个，但计算它们仍需要算法。

对非负项，若区间 $[x,2x)$ 内各项同阶，可按 $[2^j,2^{j+1})$ 倍增分块估计。例如
\[
 \sum_{i\le n}i^{-\alpha}=
 \begin{cases}
 \Theta(n^{1-\alpha}),&0\le\alpha<1,\\
 \Theta(\log(n+1)),&\alpha=1,\\
 \Theta(1),&\alpha>1.
 \end{cases}
\]

\begin{tip}[\textbf{倍增估计}是什么意思]
把求和项按\textbf{商相同}分成块：在 $[l,r]$ 内 $\floor{n/i}$ 都一样，于是整块可用\textbf{块长 $\times$ 值}估计。写成 $O$ 或 $\Theta$，不要写成逐步等号——这一步只保证同阶，不保证精确。
\end{tip}


\paragraph{证明应写成估计而不是等式.}
$0\le\alpha<1$ 时，第 $j$ 块有 $\Theta(2^j)$ 项，每项为 $\Theta(2^{-\alpha j})$，故总和为 $\Theta(\sum_j2^{(1-\alpha)j})=\Theta(n^{1-\alpha})$。
把一整块的所有函数值替换为端点，只是常数倍估计，不能像原稿一样连写精确等号。
\paragraph{算法选择问题.} 知道 $f(p^k)$ 并不自动意味着能亚线性求和。还需要至少一种结构：容易求和的卷积关系、与另一函数在素数处相同、或者素数处的值能由可筛的完全积性函数组合。


\section{用卷积把求和变成递归}
若 $f=g*h$，则
\[
S_f(n)=\sum_{ab\le n}g(a)h(b)
=\sum_{a=1}^n g(a)S_h(\floor{n/a}).
\]

\begin{tip}[选 $g$ 的原则]
要的是 $f*g=h$，其中 $h$ 的前缀和比 $S_f$ 好算，且 $g$ 的前缀和也好算（通常挑 $g=\one$、$g=\mu$、$g=\id_t$ 之一，或它们的组合）。选定后写出 $S_f(n)=\sum_{d}h(d)-\sum_{l\ge2}(\cdots)S_f(\floor{n/l})$，剩下的就是复杂度分析。
\end{tip}

如果能常数时间读取 $S_g,S_h$ 在商集所需位置的值，整除分块可在 $O(\sqrt n)$ 求出上式。
反过来，若 $f*g=h$ 且 $g(1)\ne0$，则
\[
\boxed{S_f(n)=\frac{S_h(n)-\sum_{j=2}^n g(j)S_f(\floor{n/j})}{g(1)}}.
\]
在模数环中需要 $g(1)$ 可逆；积性函数的 $g(1)=1$，分母可省。


\paragraph{杜教筛的关键不是“套递归”.}
它把未知的 $S_f(n)$ 从卷积前缀里单独抽出来，剩下的递归参数 $\floor{n/j}$ 都小于 $n$，所以可以自顶向下记忆化。每个商块 $[l,r]$ 的贡献为
\[
 (S_g(r)-S_g(l-1))\,S_f(\floor{n/l}).
\]

\begin{tip}[记忆化为什么够用]
递归的自变量只有 $\floor{n/j}$ 这一族值，而不同的商最多 $2\sqrt n$ 个（小于 $\sqrt n$ 的商有 $\sqrt n$ 个，大于的对应 $j\le\sqrt n$ 也只有 $\sqrt n$ 个）。所以\textbf{状态数}不是猜的，是可以数出来的——这就是杜教筛复杂度分析的核心一行。
\end{tip}

其中 $S_g$、$S_h$ 必须容易计算；随便找一个卷积恒等式并不一定有用。方法本身不要求 $f$ 积性，积性只常用于构造关系和快速预处理。

\paragraph{为何状态仍在原商集中？}
\[
\floor{\frac{\floor{n/a}}b}=\floor{\frac n{ab}}.
\]
对子问题的子问题继续整除，仍落在原商集中，因此记忆化不会生成所有整数 $1,\ldots,n$。

\subsection{无预处理与带预处理的复杂度}
若每个规模 $v$ 的子问题花费 $O(\sqrt v)$，则总成本可估计为
\[
 \sum_{v\le\sqrt n}O(\sqrt v)
 +\sum_{i\le\sqrt n}O(\sqrt{n/i})=O(n^{3/4}).
\]
这不是 $O(\sqrt n)$：状态数为 $O(\sqrt n)$ 不等于总工作量。
若线性筛预处理到 $B\ge\sqrt n$，大状态至多为 $n/i$（$i\le n/B$），所以
\[
 O\left(B+\sum_{i\le n/B}\sqrt{n/i}\right)
 =O(B+n/\sqrt B).
\]

\begin{tip}[提醒：状态数 $\neq$ 工作量]
状态数只告诉你\textbf{有多少个格子要填}，不告诉你\textbf{填一格要多久}。这里一格要 $O(\sqrt v)$，于是总时间是 $\sum_{v}O(\sqrt v)=O(n^{3/4})$。凡是递归带记忆化的复杂度，都按这个格式老老实实求和。
\end{tip}

取 $B\asymp n^{2/3}$ 得 $O(n^{2/3})$。空间是 $O(B+n/B)$，其中后者为记忆化表的量级。对 $v\le B$ 直接返回已存的 $S_f(v)$，绝不是返回 $\sqrt v$。


\begin{problem}{P4213：杜教筛}
给定 $n\le10^9$，求 $S_\mu(n)$ 和 $S_\varphi(n)$。
\end{problem}
利用 $\mu*\one=\eps,\varphi*\one=\id$，得到
\[
 S_\mu(n)=1-\sum_{j=2}^n S_\mu(\floor{n/j}),
\]

\begin{tip}[$\asymp$ 与\textbf{平衡}]
$A\asymp B$ 表示两者同阶（差常数倍，忽略对数）。做法同上：总代价 $B+n/\sqrt B$，让两项同阶解出 $B=n^{2/3}$，再代回得 $O(n^{2/3})$。这里还要额外报空间 $O(B+n/B)$，因为记忆化表是真会占内存的。
\end{tip}

\[
 S_\varphi(n)=\frac{n(n+1)}2-\sum_{j=2}^nS_\varphi(\floor{n/j}).
\]

\begin{pitfall}[讲解：先看这题到底在做什么]
杜教筛的核心不是“记一个公式”，而是让原本需要扫到 $n$ 的卷积关系重新排列。选择一个可以快速求前缀和的 $g$，使 $f*g$ 的前缀值容易计算，然后把大的 $S_f(n)$ 表达成若干个更小的 $S_f(\floor{n/k})$，再用记忆化递归。
\end{pitfall}

\sol 两式中 $g(j)=1$，块系数直接为块长。先线性筛出 $\mu,\varphi$ 和普通前缀，再分别记忆化大状态即可。
\begin{lstlisting}
// S is either the Mertens prefix or the phi prefix.
// H(n) is 1 for Mertens; n*(n+1)/2 for phi.
i128 solve(long long n) {
    if (n<=B) return prefix[n];
    auto it=memo.find(n);
    if (it!=memo.end()) return it->second;
    i128 ans=H(n);
    for (long long l=2,r;l<=n;l=r+1) {
        long long q=n/l;
        r=n/q;
        ans-=(i128)(r-l+1)*solve(q);
    }
    return memo[n]=ans;
}
\end{lstlisting}
$S_\mu$ 可能为零或负数，所以不能用 \texttt{if (memo[n])} 判断是否已计算；也不能把零当作“未访问”。

\paragraph{手算 $n=10$.}
$\varphi(1),\ldots,\varphi(10)$ 为 $1,1,2,2,4,2,6,4,6,4$，前缀为 $32$。递归式给出
\[
55-S_\varphi(5)-S_\varphi(3)-2S_\varphi(2)-5S_\varphi(1)
=55-10-4-4-5=32.
\]
$\mu(1),\ldots,\mu(10)=1,-1,-1,0,-1,1,-1,0,0,1$，前缀为 $-1$。
\cost 多询问可按最大 $n$ 预处理并复用表，但不能简单将单次复杂度当成所有任意询问的总复杂度；记忆化表的状态并集仍需分析。

\src{https://www.luogu.com.cn/problem/P4213}

\section{Powerful Number：稀疏的高次素因子结构}
\begin{definition}
正整数 $x$ 是 Powerful Number（PN），若每个整除 $x$ 的素数 $p$ 都满足 $p^2\mid x$。$1$ 也属于 PN。
\end{definition}
每个 PN 可唯一写为 $x=a^2b^3$，其中 $b$ 为平方自由数：对每个指数，偶数写成 $2u$，奇数（至少为 $3$）写成 $2u+3$。
因此
\[
 \#\{x\le n:x\in\mathrm{PN}\}
 \le\sum_{b\le n^{1/3}}\sqrt{n/b^3}
 =O(\sqrt n).
\]

\begin{tip}[词义：Powerful Number]
Powerful Number = 每个素因子指数都 $\ge2$ 的数，等价于\textbf{形如 $a^2b^3$}。$\le n$ 的个数约 $\zeta(3/2)$ 级、约 $O(\sqrt n)$：因为平方部分至多 $\sqrt n$ 种，剩下每个 $a$ 的立方部分至多 $\sqrt n$ 种，相乘即 $O(\sqrt n)$。\textbf{支撑稀疏}就是靠这个数出来的。
\end{tip}

若不要求 $b$ 平方自由，则表示不唯一，但仍可以用于上界计数。


\paragraph{哪些数属于 PN？}
$1,4,8,9,16,25,27,32,36$ 属于 PN；$12=2^2\cdot3$、$18=2\cdot3^2$ 不属于。定义要求\textbf{每一个}出现的素因子指数都至少为二，不是“存在平方因子”。
\paragraph{枚举而不重复.} 递归状态记录当前乘积与最后使用的素数下标，后续只加更大的素数，每个新素数指数至少为二。用 $p\le\sqrt{n/\text{cur}}$ 判断是否能扩展，不先计算可能溢出的 $\text{cur}\cdot p^2$。每个 PN 的质因数分解唯一，因此恰访问一次；递归同时携带函数值。


\section{把差异留给 PN 支撑}
\begin{theorem}[PN 求和转化]
设 $f,g$ 积性，且对所有素数 $p$ 有 $f(p)=g(p)$。令 $f=g*h$，则 $h(p)=0$，从而 $h$ 的非零支撑只可能落在 PN 上。于是
\[
 S_f(n)=\sum_{\substack{d\le n\\d\in\mathrm{PN}}}h(d)S_g(\floor{n/d}).
\]
\end{theorem}

\begin{tip}[怎么做：差值为何只落在 PN 上]
套路：构造 $g$ 使 $\hat g(p)=\hat f(p)$（素数处一致），于是 $h=f*g^{-1}$ 在每个素数处为 $0$，由积性，$h$ 在含单次素因子的数上也全为 $0$。差值只能落在\textbf{所有素因子至少二次}的数上——那里的数很少，枚举得起。
\end{tip}


\begin{proof}
先把\textbf{支撑}讲明白：说某个函数 $h$ 的支撑是集合 $A$，意思就是\textbf{$h(d)$ 只在 $d\in A$ 时可能非零}。这里要证的是 $h$ 只在\textbf{完全平方数}上非零。

第一步（素数处）。$f=g*h$ 在 $n=p$ 处展开只有两项（$p$ 的约数是 $1$ 与 $p$）：
\[
 f(p)=g(p)h(1)+g(1)h(p)=g(p)+h(p),
\]
用了 $g,h$ 积性所必需的 $g(1)=h(1)=1$（非零积性函数在 $1$ 处取值 $1$：$g(1)=g(1\cdot1)=g(1)^2$，且 $g$ 不恒为零）。所以 $h(p)=f(p)-g(p)$ 被两边在素数处的差唯一决定——这正是\textbf{把差异全部交给 $h$}的做法。

第二步（含单次素因子的数不贡献）。设 $d$ 有某个素因子 $p$ 的指数恰为 $1$，写 $d=p\,e$、$p\nmid e$。积性给出
\[
 h(d)=h(p)h(e)=\big(f(p)-g(p)\big)h(e).
\]
若我们构造 $g$ 使得 $g(p)=f(p)$（例如让 $\hat g(p)=\hat f(p)$ 在生成函数层面一致），则 $h(p)=0$，从而 $h(d)=0$。也就是说：$h$ 在一个\textbf{某个素因子只出现一次}的数上必为零，于是非零只可能出现在每个素因子指数都 $\ge2$ 的数上，即完全平方数（更严格地说，是 Powerful Number 集合）。

第三步（换序得到公式）。$f=g*h$ 即 $f(n)=\sum_{d\mid n}g(n/d)h(d)$，把求和限制在支撑上：
\[
 f(n)=\sum_{\substack{d\mid n\\ d\ \text{为平方数}}} g(n/d)h(d).
\]
要求 $f(n)$ 就只需要枚举 $n$ 的平方数约数 $d$（个数很少，见下一节的计数），并且 $h(d)$ 已经由素数处的差递推出来。
\end{proof}
固定素数 $p$，局部系数由
\[
 h(p^e)=f(p^e)-\sum_{j=1}^e g(p^j)h(p^{e-j}),\qquad h(1)=1
\]
递推。这里不能把 $h(p^e)$ 写成 $f(p^e)-g(p^e)$，因为还有交叉项。
\cost 若 $S_g$ 的商集前缀已备好、局部系数已预处理且可常数时间读取，则 PN 枚举阶段为 $O(\sqrt n)$；总时间还必须加上 $g$ 的前缀预处理与局部系数构造。不能不加前提就宣布任意这样的函数都有 $O(\sqrt n)$ 总算法。

\begin{problem}{入门例：平方自由数的个数}
求 $\sum_{i\le n}\mu(i)^2$。
\end{problem}
\begin{pitfall}[讲解：先看这题到底在做什么]
先把“平方自由”翻译成“没有任何平方因子”。莫比乌斯函数正好提供了一个按平方因子做容斥的系数，因此最后只需要枚举 $d^2\le n$。这个例子非常适合先讲 PN 的支撑为什么会变稀疏。
\end{pitfall}

\sol 取 $g=\one$。局部级数 $(1+z)/(1-z)^{-1}=1-z^2$，所以 $h(d^2)=\mu(d)$，其他位置为零，得到
\[
 \sum_{i\le n}\mu(i)^2=\sum_{d\le\sqrt n}\mu(d)\floor{n/d^2}.
\]
筛到 $\sqrt n$ 后直接求和。$n=10$ 时 $10-2-1=7$，排除的是 $4,8,9$。


\begin{problem}{LOJ 6053：异或定义的积性函数}
积性函数满足 $f(p^k)=p\mathbin{\oplus}k$，$\oplus$ 表示按位异或。求 $S_f(n)$，$n\le10^{10}$。
\end{problem}
\hint 奇素数处 $f(p)=p-1=\varphi(p)$；素数 $2$ 是唯一例外。拆成 $f=\varphi*h$ 后枚举奇数 PN 与二的幂。

\begin{pitfall}[讲解：先看这题到底在做什么]
先看奇素数：因为奇数的最低位是 $1$，所以 $p\oplus1=p-1=\varphi(p)$，于是卷积商在奇素数处消失。唯一特殊的是 $2$，所以 PN 枚举只需要在“奇数部分 + 二的幂”这个结构上展开。
\end{pitfall}

\sol 当 $p>2$，最低二进制位为一，所以 $p\oplus1=p-1$。因此 $h(p)=0$（奇素数）。在 $p=2$ 处 $f(2)=3$，$\varphi(2)=1$，故 $h(2)=2$，不能把 $2$ 也当成普通 PN 因子。
$h$ 有值的位置形如 $2^e u$，其中 $u$ 为奇数 PN。它们的总数量仍不超过
\[
 \sum_{e\ge0}O\!\left(\sqrt{n/2^e}\right)=O(\sqrt n).
\]
用杜教筛求 $S_\varphi$，再按上述支撑枚举。作为局部计算示例，
\[
 h(3)=0,\qquad h(3^2)=f(3^2)-\varphi(3^2)-\varphi(3)h(3)=1-6=-5.
\]
卷积商可能出现负数，模运算中应及时规范。素数幂 $p^k$ 的点值是 $p\oplus k$，不是 $(p^k)\oplus k$。

\src{https://loj.ac/p/6053}


\chapter{洲阁筛与素数贡献的分层统计}
\begin{chapterintro}[章节导读]
上一章我们处理的是一般积性函数的前缀和；进一步地，如果最重要的贡献集中在素数及其幂上，就可以把“素数部分”和“合数部分”分层处理。这一章沿着这个方向建立洲阁筛的状态，并特别说明哪些完全积性条件是这种拆分能够成立的基础。
\end{chapterintro}

\section{先把适用条件说清楚}
本章讨论原稿“洲阁筛”所用的两阶段思路：先在商集上筛出素数位置的函数和，再按最小素因子补入合数。它与常见 Min\_25 类筛法的若干实现共享同一核心结构；不同资料的命名、状态和具体复杂度可能不同，不能光看名字就混用模板。

适用于以下条件：
\begin{enumerate}
\item $f$ 积性，素数幂 $f(p^e)$ 可高效计算；
\item $f(p)$ 可表示成常数个\textbf{完全积性函数}在素数处取值的线性组合；
\item 这些完全积性函数的普通前缀和可快速计算。
\end{enumerate}
最常见的是 $f(p)$ 为关于 $p$ 的低次多项式，基函数为 $1,p,p^2,\ldots$。

\begin{problem}{P5325：Min\_25 筛模板所对应的积性函数}
$f(1)=1$，$f(p^e)=p^e(p^e-1)$，求 $\sum_{i\le n}f(i)$，$n\le10^{10}$，答案按题目要求取模。
\end{problem}
\hint 素数处 $f(p)=p^2-p$，先分别求素数的一次幂和与平方和。不要误认为所有 $n$ 上都有 $f(n)=n(n-1)$。

\section{第一阶段：把合数从简单前缀和中筛掉}
令 $p_1,p_2,\ldots,p_s$ 为 $\sqrt n$ 以内的素数。对 $t=0,1,2$ 等所需次数，令
\[
 G_t(v,j)=\sum_{2\le x\le v}
 [x\text{ 为素数或 }\operatorname{minp}(x)>p_j]x^t.
\]

\begin{tip}[前提：洲阁筛的要求与它的输出]
洲阁筛（Min\_25）要求 $f$ 是积性函数、$f(p)$ 是关于 $p$ 的多项式（且能在素数幂上递推）。它的输出是\textbf{$\sum_{i\le n}f(i)$}，不是\textbf{$\sum_{p\le n}f(p)$}；后者只是它的中间量 $G$。
\end{tip}

\begin{tip}[两阶段各自在做什么]
第一阶段：从\textbf{所有整数的 $\sum p^t$}出发，一点点扣掉合数贡献，得到 $G_t(v,j)=$\textbf{最小素因子 $>p_j$ 的数（含素数与 $1$）}的 $p^t$ 和。第二阶段：按最小素因子枚举，把合数一块一块补回来，此时才用到 $f$ 的真实值。第一阶段只依赖完全积性，所以能提前算。
\end{tip}

初始 $j=0$ 时没有筛掉任何 $x\ge2$，所以
\[
G_0(v,0)=v-1,\quad G_1(v,0)=\frac{v(v+1)}2-1,
\]
\[
G_2(v,0)=\frac{v(v+1)(2v+1)}6-1.
\]
记 $P_t(j)=\sum_{i=1}^jp_i^t$。当 $p_j^2\le v$ 时，
\[
 G_t(v,j)=G_t(v,j-1)-p_j^t
 \left(G_t(\floor{v/p_j},j-1)-P_t(j-1)\right).
\]
若 $p_j^2>v$，不需要转移。最后 $G_t(v,s)$ 只剩素数贡献。


\begin{proof}[为什么恰好扣掉最小素因子为 $p_j$ 的合数？]
这样的合数唯一写成 $x=p_jy$，其中 $y\ge p_j$ 且 $y$ 没有小于 $p_j$ 的素因子。$G_t(\floor{v/p_j},j-1)$ 已保留这些 $y$，但还含小于 $p_j$ 的素数，因此减去 $P_t(j-1)$。函数 $x^t$ 完全积性，故 $(p_jy)^t=p_j^ty^t$，可提出系数。该步允许 $p_j\mid y$，所以普通积性不够用。
\end{proof}
\paragraph{例子.} $v=10,t=1$，初始 $G_1=54$。用 $2$ 筛时扣掉 $2(2+3+4+5)=28$，剩 $26$，对应 $2,3,5,7,9$；用 $3$ 筛时再扣 $9$，剩素数和 $17$。

\subsection{只保留商集与滚动数组}
需要的 $v$ 只有 $\mathcal V=\{\floor{n/k}:k\ge1\}$，因为每次下标都变为 $\floor{v/p}$，仍在原商集中。可对小值用编号 $v$、大值用编号 $n/v$ 映射，或先离散化；不要对 $10^{10}$ 开数组。
外层按素数递增，内层 $v$ 必须\textbf{递减}，这样源状态 $\floor{v/p}<v$ 仍是处理本素数之前的值。若递增原地更新，会读到已经扣过同一素数的状态。
\begin{lstlisting}
for (int j=0;j<(int)primes.size();++j) {
    long long p=primes[j];
    for (long long v:values_descending) {
        if (p>v/p) break;
        G[id(v)]-=weight(p)*(G[id(v/p)]-prime_prefix[j]);
        // prime_prefix[j] contains primes strictly before p.
    }
}
\end{lstlisting}
初始和式除以 $2,6$ 时，如果工作模数下分母不可逆，应先在整数域精确约分，再取模；本题常见素数模数下可使用逆元，但仍要写清条件。


\section{第二阶段：用素数幂重建一般积性函数}
记 $P_f(x)=\sum_{p\le x}f(p)$，它已由第一阶段的线性组合得到。定义
\[
 H(v,j)=\sum_{2\le x\le v}
 [x\text{ 为素数或 }\operatorname{minp}(x)>p_j]f(x).
\]
起点为 $H(v,s)=P_f(v)$，按 $j=s,s-1,\ldots,1$ 反向加入最小素因子为 $p_j$ 的合数。对 $p_j^2\le v$，
\[
\begin{aligned}
 H(v,j-1)=H(v,j)+\sum_{\substack{e\ge1\\p_j^e\le v}} f(p_j^e)
 \biggl(&[e\ge2]+H(\floor{v/p_j^e},j)\\
 &-P_f(\min(p_j,\floor{v/p_j^e}))\biggr).
\end{aligned}
\]
最终答案为 $1+H(n,0)$，其中 $1=f(1)$。


\paragraph{把公式拆成三类对象.}
每个待补合数唯一写作 $p_j^e u$，$p_j\nmid u$：
\begin{itemize}
\item $u=1$ 时必须 $e\ge2$，否则是已经存在的素数 $p_j$；这给出 $[e\ge2]$。
\item $u>1$ 时要求 $\operatorname{minp}(u)>p_j$；$H$ 中仍包含所有素数，所以要减去不超过 $p_j$ 的素数贡献。
\item $p_j^e$ 与 $u$ 互质，故普通积性足以提出 $f(p_j^e)$。这里和第一阶段使用完全积性的理由不同。
\end{itemize}
大于 $v$ 的素数幂不能进入求和，否则常数项 $[e\ge2]$ 会产生不存在的数。

原地滚动时，素数要按递减处理，而 $v$ 仍按\textbf{递减}处理，以保证读到的是本轮之前的 $H(\floor{v/p^e},j)$。两次筛的素数顺序相反，商值顺序却都相同。

\paragraph{一个覆盖所有类型的核对.}
原题 $n=10$ 时，
\[
(f(1),\ldots,f(10))=(1,2,6,12,20,12,42,56,72,40),
\]
总和为 $263$。例如 $f(6)=f(2)f(3)=12$，而不是 $6\cdot5=30$；$f(4)=12$，而不是 $f(2)^2=4$。前者提醒“公式只给素数幂”，后者提醒“不是完全积性”。

\subsection{复杂度说明与备课时的取舍}
每个商值 $v$ 的素数转移数大致为 $\pi(\sqrt v)$。把所有商值相加，可先得到保守的 $O(n^{3/4})$ 数量级；再利用素数计数上界并精细求和，可得常见的
\[
 O(n^{3/4}/\log n)
\]
量级（默认固定次数的基函数、素数幂点值可常数时间读取、滚动表常数时间访问）。第二阶段的额外幂次枚举可按
\[
 \sum_{p\le\sqrt v}\floor{\log_p v}
 =2\pi(\sqrt v)+\sum_{e\ge3}\pi(v^{1/e})
\]
计数，不会改变主导量级。小 $n$ 的对数表达应理解为渐近说明。
空间为常数张 $O(\sqrt n)$ 的商集表和素数前缀表。若改成递归搜索、使用平衡树映射、每次转移重新求 $f(p^e)$，实际复杂度需要重新计入这些成本。

\begin{teach}[不要在第一次课就背完整转移]
先用 $t=0$ 实现素数个数筛；再用 $t=1$ 算素数和；最后才加入一般积性函数的第二阶段。要求学生分别解释“减去较小素数前缀”和“补一个纯素数幂”两个容易漏掉的项。
\end{teach}

\src{https://www.luogu.com.cn/problem/P5325}


\chapter{二维积性函数：消去主要贡献以得到稀疏支撑}
\begin{chapterintro}[章节导读]
前面的筛法大多只处理一个整数；这一章把对象提升到两个坐标。二维卷积看起来更复杂，但核心仍然是“先匹配主要贡献，再让剩余部分变得稀疏”，因此我们会先消掉坐标轴，再处理对角线，最终把庞大的状态空间压缩成可枚举的支撑。
\end{chapterintro}

\section{问题、定义与二维卷积}
\begin{problem}{UOJ 885：红场阅兵}
已知积性函数 $S$，满足 $S(1)=1$，对素数 $p$ 与 $k\ge1$，
\[
 S(p^k)=w_0+w_1p^k+w_2p^{2k}.
\]


求 $\sum_{i=1}^n\sum_{j=1}^nS(ij)\bmod998244353$，$n\le10^9$。
\end{problem}

\begin{tip}[词义：二维积性]
\textbf{二维}= 自变量是数对 $(x,y)$，且固定一个变量后对另一个是积性的。处理方式与一维一样：先把\textbf{对角线/坐标轴}上的一维积性部分消掉，剩下的非零项只出现在稀疏的坐标对上。
\end{tip}

\begin{definition}
二维函数 $f$ 满足 $f(1,1)=1$，且 $\gcd(ab,cd)=1$ 时有
\[
 f(ac,bd)=f(a,b)f(c,d),
\]
则称 $f$ 为二维积性函数。等价地，若 $x=\prod p^{\alpha_p},y=\prod p^{\beta_p}$，则
\[
 f(x,y)=\prod_p f(p^{\alpha_p},p^{\beta_p}).
\]
二维卷积定义为
\[
 (f*g)(x,y)=\sum_{u\mid x}\sum_{v\mid y}f(u,v)g(x/u,y/v).
\]
\end{definition}
常数项为可逆元素时可以逐项求卷积逆。二维积性函数的卷积、卷积商仍为二维积性函数；固定素数后对应二元形式幂级数乘除。


\begin{teach}[从一维过渡时必须补上的条件]
两个坐标不是分别要求 $\gcd(a,c)=\gcd(b,d)=1$，而是要求第一对坐标所含的所有素因子与第二对坐标所含的所有素因子完全分离，即 $\gcd(ab,cd)=1$。否则可能把同一个素数拆成两个独立局部因子，积性分解就不成立。
此外，素数幂公式只用于 $k\ge1$。$S(p^0)=S(1)=1$，不是 $w_0+w_1+w_2$。
\end{teach}


\section{算法一：先消去坐标轴}
令 $f(x,y)=S(xy)$、$g(x,y)=S(x)S(y)$，定义 $f=g*h$。$g$ 的二维前缀易算：
\[
 S_g(A,B)=S_S(A)S_S(B),\qquad S_S(t)=\sum_{i\le t}S(i).
\]
由于 $f(p^k,1)=g(p^k,1)$ 以及另一坐标轴上的对应等式，
\[
 h(p^k,1)=h(1,p^k)=0\quad(k\ge1).
\]
因此只有 $x,y$ 含相同素因子集合时 $h(x,y)$ 才可能非零，答案为
\[
 \sum_{x\le n}\sum_{y\le n}h(x,y)
 S_S(\floor{n/x})S_S(\floor{n/y}).
\]
直接按共同素因子递归枚举这些位置，不必遍历 $n^2$ 个点。


\subsection{局部系数怎样实际求？}
固定素数 $p$，令 $s_0=1,s_k=S(p^k)$（$k\ge1$），$h_{a,b}=h(p^a,p^b)$，则
\[
 h_{0,0}=1,\qquad
 h_{a,b}=s_{a+b}-\sum_{\substack{0\le i\le a,\,0\le j\le b\\i+j>0}}
 s_i s_j h_{a-i,b-j}.
\]
按 $a+b$ 递增计算。例：
\[
 h_{1,0}=h_{0,1}=0,\qquad h_{1,1}=S(p^2)-S(p)^2.
\]
局部二维多项式除法的成本应单独预处理或缓存，不要在每个递归节点从头做一次。

\subsection{为何只有 \texorpdfstring{$O(n)$}{O(n)} 个位置？一个更直接的证明}
将 $x,y\le n$ 放宽为 $xy\le n^2$，不是原稿误写的 $xy\le n$。
令 $a(t)$ 为满足 $xy=t$、$x,y$ 具有相同素因子集合的有序数对个数。它是积性的，且
\[
 a(p^k)=\begin{cases}0,&k=1,\\k-1,&k\ge2.\end{cases}
\]
$k-1$ 来自正指数拆分 $k=1+(k-1)=\cdots=(k-1)+1$。

令 $T_2$ 为完全平方数的示性函数，写 $a=T_2*b$。局部生成函数相乘可得
\[
 b(p)=b(p^2)=0,\qquad b(p^k)=2\ (k\ge3).
\]
于是
\[
\sum_{t\le X}a(t)=\sum_{d\le X}b(d)\floor{\sqrt{X/d}}
\le\sqrt X\sum_{d\ge1}\frac{b(d)}{\sqrt d}.
\]
右边的级数可分解成
\[
 \prod_p\left(1+2\sum_{k\ge3}p^{-k/2}\right).
\]
因为每个局部因子为 $1+O(p^{-3/2})$，且 $\sum_p p^{-3/2}$ 收敛，整个乘积有限。因此 $\sum_{t\le X}a(t)=O(\sqrt X)$。代入 $X=n^2$，便得到 $O(n)$ 的支撑上界。
这个证明不需要“随便写出一个高次卷积支配函数后省略检查”。

\paragraph{复杂度边界.} 算法一的枚举阶段为 $O(n)$，另加 $S_S$ 的商集前缀与局部系数的预处理。适合中等规模或作为正确性基准，$n=10^9$ 时需要进一步优化。


\section{算法二：再消去局部对角线}
定义一维积性函数 $D(t)=h(t,t)$，以及只在对角线上有值的二维函数
\[
 g_3(x,y)=[x=y]D(x).
\]
令 $h=g_3*h'$。逐素数在对角线上作一维除法可得
\[
 h'(p^a,p^b)=0\quad\text{当 }a=0<b,\ b=0<a,\text{或 }a=b>0.
\]
所以对 $h'(x,y)\ne0$ 的位置，每个素数要么在两边都不出现，要么两边指数都为正且不相等。
答案转为
\[
 \sum_{x,y\le n}h'(x,y)\,L(\floor{n/x},\floor{n/y}),
\]
其中
\[
 L(A,B)=\sum_{d\le\min(A,B)}D(d)
 S_S(\floor{A/d})S_S(\floor{B/d}).
\]
对两个商共同分块，在已知 $S_D$ 前缀时，单次计算需要
\[
 O\bigl(\min(A,B,\sqrt A+\sqrt B)\bigr)
\]
次操作。


\paragraph{为什么前缀预处理仍可做？}
\[
 D(p)=S(p^2)-S(p)^2
 =w_0+w_1p^2+w_2p^4-(w_0+w_1p+w_2p^2)^2,
\]
是关于 $p$ 的四次多项式。因此 $S,D$ 的素数值均可拆成常数个幂函数，能用上一章方法求商集前缀。$D(p^k)=h(p^k,p^k)$ 由局部表计算。
由于 $A=\floor{n/x}$，$\floor{A/d}=\floor{n/(xd)}$，所有需要的 $S_S$ 参数仍在原商集中；共同块的端点来自 $A$ 或 $B$ 的整除商，$S_D$ 也只需原商集所对应的值。

\subsection{稀疏支撑数为 \texorpdfstring{$O((AB)^{1/3})$}{O((AB)1/3)}}
令 $v(x,y)$ 表示上述“每个出现的素数指数均正且不同”的支撑，且 $v(1,1)=1$。
定义积性 $v_1$：除 $(1,1)$ 的局部常数项外，只在局部 $(p,p^2),(p^2,p)$ 处为 $1$；定义 $v_2$ 在其他合法非零局部位置为 $1$，两处被删除的位置为 $0$，并令 $v_2(1,1)=1$。在非负系数意义下有 $v\le v_1*v_2$。

$v_1$ 的位置可表示成 $(a^2b,ab^2)$。实际还对 $a,b$ 有平方自由及互质限制，去掉这些限制只会增加数量，所以
\[
 S_{v_1}(A,B)\le\sum_a\min\left(\frac A{a^2},\sqrt{\frac Ba}\right).
\]
两项在 $a_0=A^{2/3}/B^{1/3}$ 附近相等；在此处分段求和，得到 $O((AB)^{1/3})$。若分点落在合法范围外，使用 $O(A)$ 或 $O(B)$ 的更粗界也满足结论。

$v_2$ 的非零局部项满足两坐标指数之和至少为 $4$，故
\[
 \sum_{x,y\ge1}\frac{v_2(x,y)}{(xy)^{1/3}}
 =\prod_p\left(1+O(p^{-4/3})\right)<\infty.
\]
展开卷积前缀并使用上界，得到
\[
 S_v(A,B)\le\sum_{x,y}v_2(x,y)
 O\!\left((AB/(xy))^{1/3}\right)=O((AB)^{1/3}).
\]
这说明剥掉最低总次数为二的对角项后，主要的规模指数从 $1/2$ 变成了 $1/3$。

\subsection{将单次商块成本求和}
使用不等式
\[
 \min(A,B,\sqrt A+\sqrt B)
 \le\min(A,\sqrt B)+\min(B,\sqrt A).
\]
两部分对称，只分析第一部分。把最小值写成层数计数，有
\[
\begin{aligned}
 \sum_{x,y\le n}v(x,y)\min\left(\frac nx,\sqrt{\frac ny}\right)
 &\ll\sum_{t\le\sqrt n}S_v(\floor{n/t},\floor{n/t^2})\\
 &\ll\sum_{t\le\sqrt n}\left(\frac{n^2}{t^3}\right)^{1/3}
 =O(n^{2/3}\log n).
\end{aligned}
\]
因此\textbf{除前缀表及局部系数预处理以外}，算法二耗时 $O(n^{2/3}\log n)$。如果前缀用洲阁筛计算，总复杂度还要加对应的预处理项，不能一笔抹去。


\section{推广：第一个非零局部系数决定主导规模}
设非负整数序列 $a_k$ 为常数次多项式增长，$a_0=1$，令 $r$ 是第一个满足 $a_r>0$ 的正下标。积性函数 $F$ 若满足 $F(p^k)=a_k$，则常见的上界为
\[
 S_F(n)=O\left(n^{1/r}(\log(n+1))^{a_r-1}\right),
\]

\begin{tip}[怎么做：先代一个小例子]
\textbf{$r$ 是第一个 $a_r>0$ 的下标}的意思是：局部生成函数 $a_0+a_1z+\cdots$ 前面连续有 $r$ 个零。这些零把每一项能出现的素因子次数抬高，于是和式的主导规模从 $n$ 变成 $n^{1/r}$。先代 $a_0=1,a_1=0,a_2=1$（$r=2$）算一个具体例子，再读一般式子。
\end{tip}

其中 $r,a_r$ 视为常数。这里给出的是在上述归一化、非负整数系数与多项式增长条件下的结论，不能无条件推广到任意系数列。

\paragraph{证明思路.} 定义 $T_r(n)=[n\text{ 是完全 }r\text{ 次方}]$。从局部级数中提出 $(1-z^r)^{-a_r}$，相当于 $F=T_r^{*a_r}*H$。剩余局部级数从 $z^{r+1}$ 才开始，因此
\[
 \sum_{d\ge1}|H(d)|d^{-1/r}<\infty.
\]
而 $T_r^{*a_r}$ 的前缀和等于 $a_r$ 重约数函数在 $n^{1/r}$ 的前缀和，为 $O(n^{1/r}\log^{a_r-1}(n+1))$；卷积上界便给出结论。
例如 $\tau(n)^2$ 的局部值为 $(k+1)^2$，$r=1,a_1=4$，得到 $O(n\log^3 n)$。若 $a_r$ 非整数或取负值，需要另行使用分析数论工具，不能直接用“做 $a_r$ 次卷积”的证明。


\section{算法三：去掉枚举阶段的对数（研究拓展）}
原稿最后提到可继续优化。本讲义保留这一方向，但把所需的新结论单独列出来，不将它冒充前两种算法的自然常数优化。
如果能在合适预处理后，把 $L(A,B)$ 的一次询问做到 $O(\min(A,B)^\rho)$，其中固定 $\rho<2/3$，那么稀疏枚举阶段就能降到 $O(n^{2/3})$。

\paragraph{为什么是阈值 $2/3$？}
令 $w(t)=t^\rho-(t-1)^\rho=O(t^{\rho-1})$。总询问成本不超过常数倍
\[
\sum_{t\ge1}w(t)S_v(\floor{n/t},\floor{n/t})
\ll n^{2/3}\sum_{t\ge1}t^{\rho-1-2/3}.
\]
只有 $\rho<2/3$ 时该级数收敛；例如取 $\rho=3/5$ 即可。

\paragraph{新的难点是什么？}
需要批量维护所有商值 $A$ 对应的加权前缀
$\sum_{d\le t}D(d)S_S(\floor{A/d})$，而不是每个 $(A,B)$ 都独立重做分块。
对大 $A$ 可单独预处理商块；对小 $A$ 需要利用相邻商值的 $\floor{A/i}$ 数组变化总量，维护小下标项及按商聚合项的两套前缀。这涉及额外的数据结构与摊还证明。
本节给出目标和总复杂度归约，完整高性能实现及变化量分析参见下列延伸资料；课堂上可将证明 $\sum w(t)S_v(n/t,n/t)$ 作为专题作业。采用哪种前缀预处理也影响\textbf{完整算法}的时间界。

\src{https://uoj.ac/problem/885}
\src{https://www.cnblogs.com/zkyJuruo/p/18204106}


\chapter{综合专题：循环、指数向量、图与期望}
\begin{chapterintro}[章节导读]
前面几章已经把同余、阶、函数与筛法中的工具分别建立起来；现在把这些工具放回真正复杂的综合题里。这里不再只有一个公式，而是会同时出现循环结构、指数向量、图上的状态传播以及概率计数，需要把前面学到的结构识别能力组合起来。
\end{chapterintro}

\section{CF516E：快乐传播不是只保留最小起点}
\begin{problem}{Drazil and His Happy Friends}
$n$ 个男生、$m$ 个女生分别编号从 $0$ 开始。第 $t$ 天男生 $t\bmod n$ 与女生 $t\bmod m$ 见面；只要有一人快乐，另一人也会变快乐，快乐不可逆。给定初始快乐的男生、女生集合，求所有人都快乐的最早日期，或报告不可能。
$n,m\le10^9$，初始快乐人数总计不超过 $2\times10^5$，保证至少有一个人初始不快乐。
\end{problem}
\hint 用 $d=\gcd(n,m)$ 分成独立传播分量；在每一侧转成步长环上的\textbf{带不同初始时间的多源最短路}，利用环坐标排序与前后缀最小值求最晚到达时间。

\subsection{传播时间的精确形式}
将两侧初始快乐的编号合并为多重集合 $Z$。能发生有效传播的时刻恰是
\[
 t=z+vn+wm,\qquad z\in Z,\quad v,w\ge0.
\]
必要性沿传播链归纳：男生再次见面时间增加 $n$ 的倍数，女生增加 $m$ 的倍数。
充分性：若 $z$ 来自男生，在 $z+vn$ 时该男生使对应女生快乐，再在 $z+vn+wm$ 时由该女生传播；若 $z$ 来自女生，交换顺序即可。

对某个初始不快乐的男生，$vn$ 不改变其编号且只会推迟时间，因此最早时间可写为
\[
 F(i)=\min_{\substack{z\in Z,w\ge0\\z+wm\equiv i\pmod n}}(z+wm).
\]
女生同理交换 $n,m$。初始快乐的人无需等到第一次见面，应从“最晚变快乐”统计中排除。

\subsection{修正原稿的关键反例}
取 $n=5,m=3$，初始快乐男生为 $\{4\}$，女生为 $\{0\}$。同一个环中最小源为 $z=0$，它单独向男生 $2$ 传播需要时间 $12$；而源 $z=4$ 在时间 $7$ 就能到达男生 $2$。实际全员快乐的日期为 $7$。
因此“同一子环上所有初始不快乐的点都由最小 $z$ 最先到达”是错误的。小编号源可能离目标更远，必须保留多源竞争。

\subsection{把步长变成加一}
对男生侧，令 $d=\gcd(n,m),L=n/d$。每个余数 $r\pmod d$ 是独立环。$L>1$ 时取
\[
 u=(m/d)^{-1}\pmod L.
\]
编号为 $i\equiv r\pmod d$ 的点映射到环坐标
\[
 t_i=((i-r)/d)u\bmod L.
\]
每走一步 $+m$，新坐标恰加一。$L=1$ 时只有坐标零，无需求逆。
源 $z$ 的坐标用 $z\bmod n$ 或直接用上述同余式计算，初始代价仍是\textbf{整数 $z$}，而不是 $z\bmod n$。
若某个余数类没有源，两侧对应的人都永远无法快乐。无需开 $d$ 大小数组：哈希分组后，检查出现的余数类个数是否为 $d$。

\subsection{环上的带权多源扫线}
一个环上把不同源坐标排序为 $t_1<\cdots<t_s$，同坐标只保留最小初始代价 $z_i$。设
\[
 c_i=z_i-mt_i.
\]
对于目标坐标 $t$，最早时间为
\[
 F(t)=mt+\min\left(
 \min_{t_i\le t}c_i,\quad
 mL+\min_{t_i>t}c_i\right).
\]
第一项表示从左侧源不绕环到达，第二项表示从右侧源先绕一整环。
预处理 $c_i$ 的前缀最小值和后缀最小值，每个源之间的区间都能 $O(1)$ 确定常数项；区间内斜率为正数 $m$，最晚时间在其最右端。

按区间 $[0,t_1-1]$、$[t_i,t_{i+1}-1]$、$[t_s,L-1]$ 检查右端点，并排除初始快乐男生。由于每个初始快乐男生本身就是源，区间内部没有需要额外跳过的初始快乐点；只有长度为一、右端点就是该初始快乐源的区间需要跳过。
女生侧对称处理，取两侧初始不快乐者时间的最大值。

\cost 令初始快乐总人数为 $K$。使用排序，确定性排序部分为 $O(K\log K)$，gcd 与逆元为 $O(\log\max(n,m))$，额外空间 $O(K)$。分组可用哈希，或者按余数排序获得不依赖哈希的实现。不沿用原稿未经成立支配结论支撑的 $O(K)$ 宣称。
\begin{teach}[备课验证]
先让学生手工模拟 $n=5,m=3$ 的反例，再画出步长坐标；最后才给前缀、后缀公式。随附的小规模对拍脚本将该公式与逐日模拟比较，覆盖 $1\le n,m\le12$ 的随机初始集合。
\end{teach}

\src{https://codeforces.com/problemset/problem/516/E}

\section{CF571E：把等比数列交集转到指数空间}
\begin{problem}{Geometric Progressions}
给定 $n$ 个等比数列 $a_i,a_ib_i,a_ib_i^2,\ldots$，求最小公共正整数，答案模 $10^9+7$；无公共项则报告无解。$n\le100,a_i,b_i\le10^9$。
原稿还提出 $n\le10^5,a_i,b_i\le10^{18}$ 的 Bonus，作为研究拓展保留。
\end{problem}
\hint 分解质因数后，一个等比数列变为非负整数参数的一条指数向量射线。合并两条射线时可能无交、一个点或一条新射线；不同质因子的指数必须共享同一个参数。

\subsection{为什么不能对每个素数独立 CRT？}
设素数全集为 $p_1,\ldots,p_s$，把 $a,b$ 分别写成指数向量 $A,B$。数列为
\[
 A+kB,\qquad k\in\Z_{\ge0}.
\]
所有坐标里的 $k$ 是同一个数。若分别为每个素数求一个可行指数，却没有检查它们对应相同的 $k$，会生成并不存在于数列中的整数。

\subsection{两条射线的完整分类}
合并 $A+kB$ 与 $C+\ell D$，需要解所有坐标上的
\[
 B_jk-D_j\ell=C_j-A_j,\qquad k,\ell\ge0.
\]
\begin{enumerate}
\item \textbf{零方向.} $B=0$ 或 $D=0$ 表示原数列公比为一，是单点。直接检查这个点是否能用另一条射线的同一个非负整数参数表示。
\item \textbf{方向不成比例.} 找到两个坐标使对应 $2\times2$ 系数矩阵行列式非零。用精确整数消元求出唯一有理数解 $(k,\ell)$；检查分子能否被分母整除、参数是否非负，再代入\textbf{全部}坐标。全通过则交集是唯一一个点，否则为空。
\item \textbf{方向成比例且均非零.} 选一个 $B_j,D_j>0$ 的坐标，解
$B_jk-D_j\ell=C_j-A_j$。令 $d=\gcd(B_j,D_j)$。若无整数解则为空；否则所有解形如
\[
 k=k_0+(D_j/d)t,\qquad\ell=\ell_0+(B_j/d)t.
\]
先检查特解是否满足其余坐标，再取
\[
 t_{\min}=\max\left(\ceil{-k_0/(D_j/d)},\ceil{-\ell_0/(B_j/d)}\right).
\]
最小可行参数给出新起点 $E=A+k_{\min}B$，新方向为 $F=(D_j/d)B$，交集仍是 $E+tF$（$t\ge0$）。
\end{enumerate}
成比例情形中，改变 $t$ 不会破坏其余坐标等式，因为方向差为零。新方向各分量非负，因此新射线起点对应最小公共整数。

\paragraph{例子.} $2\cdot4^k$ 与 $8\cdot8^\ell$ 转成 $1+2k=3+3\ell$，最小 $(k,\ell)=(1,0)$，公共数列为 $8,512,32768,\ldots$，公比 $64$。
$2\cdot6^k$ 与 $12\cdot9^\ell$ 则在素数 $2,3$ 上分别给出 $1+k=2,k=1+2\ell$，唯一 $k=1,\ell=0$，公共项只有 $12$。

\cost 反复合并直到处理所有数列。指数与射线步长可能增长很大，应使用大整数或严格证明范围；最后对素因子做快速幂相乘。原题各素因子小于 $10^9+7$，指数可模 $10^9+6$ 降幂；Bonus 中若素因子等于输出模数，必须先判断其指数是否为正，不能直接当成可逆底数。
原范围可先筛到 $\sqrt{10^9}$ 试除分解。Bonus 不能简单把原算法换成 Pollard--Rho 就宣称通过：稀疏指数向量的合并次数、总维数和高精度运算都需要进一步优化。

\src{https://codeforces.com/problemset/problem/571/E}

\section{AGC031F：把巨大的状态图缩成每点六个状态}
\begin{problem}{Walk on Graph}
无向连通图有 $n$ 点、$m$ 边，模数 $M$ 为奇数。路径按经过次序的边权为 $w_0,\ldots,w_{k-1}$，权值为 $\sum_i2^iw_i$。多次询问是否存在从 $s$ 到 $t$、权值模 $M$ 为 $r$ 的游走，允许重复边。
$n,m,Q\le5\times10^4,M\le10^6$。
\end{problem}

\begin{tip}[怎么做：压成\textbf{生成树 + 回边}]
遇到这类\textbf{图上求权值模 $M$}的题，通常先把图压成\textbf{一棵生成树 + 若干回边}。树上两点之间的权值差由树唯一决定，回边给出可自由组合的生成元。本题还要额外处理\textbf{乘 $2$}不是可逆映射这件事——这就是六状态与奇偶性的来源。
\end{tip}

\hint 反转路径，将状态转移变成 $(u,x)\to(v,2x+w)$；证明状态可逆和同余类平移，再将乘二的指数只保留奇偶性，配合并查集。

\subsection{第一步：反向读路径消去长度变量}
从 $t$ 出发，倒序处理原路径的边，每走一条边令 $x\gets2x+w$，从 $x=0$ 得到的最终值正好是原题权值。因此询问等价于 $(t,0)$ 能否到达 $(s,r)$。
沿同一条边连续走 $h$ 次，数值变为
\[
 2^hx+(2^h-1)w.
\]

\begin{tip}[为什么要\textbf{反转路径}]
正着走：从 $s$ 出发到达 $t$ 的权值难合并。反着走：把每条边的贡献写成\textbf{乘 $2$ 再加权值}的仿射变换，$s$、$t$ 各自到公共点的变换可以复合、比较，于是可达性等价于两个变换像的交集。这就是把\textbf{路径问题}变成\textbf{状态问题}。
\end{tip}

$M$ 为奇数，取偶数 $h=2\varphi(M)$，既回到原顶点又回到原数值。所以一次状态转移可通过继续走同一条边反向撤销，状态可达关系可作为无向连通关系处理。$M=1$ 直接回答 YES。

\subsection{第二步：推出同顶点的平移等价}
顶点 $u$ 邻接两条权为 $a,b$ 的边，来回走一次分别产生
\[
 x\mapsto4x+3a,\qquad x\mapsto4x+3b.
\]
先撤销前一个变换再执行后一个，得到平移 $x\mapsto x+3(b-a)$。沿其他路径运输这个平移，只会多乘 $2$ 的可逆幂；生成的加法子群不变。
连通图中任意两边的权差可由相邻边权差相加得到，故在所有顶点都可进行由全部 $3(w_i-w_1)$ 生成的平移。
令
\[
 g=\gcd(M,w_2-w_1,\ldots,w_m-w_1),\qquad M'=\gcd(M,3g).
\]
由裴蜀定理，同一个顶点上模 $M'$ 相同的状态彼此可达，所以可以无损地把模数缩为 $M'$。\textbf{这不表示同顶点状态可达当且仅当模 $M'$ 相同}，之后仍有乘四产生的等价。
因为 $g\mid M$，$M'$ 只能是 $g$ 或 $3g$。

\subsection{第三步：统一边权余数并只保留六类形式}
令 $z=w_1\bmod g$，所有边权写为 $w_i=z+c_ig$。用新坐标 $X=x+z$，转移变成
\[
 X\mapsto2X+c_ig\pmod{M'}.
\]
来回走同一边给 $X\mapsto4X+3c_ig\equiv4X$，所以同一点的 $X$ 与 $4X$ 等价。
从初始值 $z$ 出发，任意形式可写成
\[
 2^ez+kg\pmod{M'}.
\]
$k$ 只需模 $3$；乘四不改变 $kg$ 的模 $M'$ 值，因此指数只需保留奇偶性。为每个顶点建立状态 $(u,\epsilon,k)$，$\epsilon\in\{0,1\},k\in\{0,1,2\}$，一共 $6n$ 个并查集点。
每条边 $u\leftrightarrow v$、系数 $c$ 建立两种方向的合并
\[
 (u,\epsilon,k)\sim(v,\epsilon\oplus1,(2k+c)\bmod3).
\]
即使 $M'=g$ 导致某些状态数值重复，保留这六种形式并在询问时枚举也不影响正确性。

\subsection{第四步：把询问数值匹配到形式}
预处理布尔表
\[
 \operatorname{vis}[\epsilon][a]=[\exists e\ge0,\ e\bmod2=\epsilon,\ 2^ez\equiv a\pmod{M'}].
\]
迭代状态 $(e\bmod2,2^ez\bmod M')$ 直到首次重复即可，至多 $2M'$ 个状态。
询问 $s,t,r$ 时，枚举六个 $(\epsilon,k)$。若
\[
\operatorname{find}(t,0,0)=\operatorname{find}(s,\epsilon,k)
\]
且 $\operatorname{vis}[\epsilon][(r+z-kg)\bmod M']$ 为真，就回答 YES，否则 NO。
负余数必须规范化；原模数为 $M$ 的 $r$ 在这里自然按 $M'$ 处理。

\paragraph{最小例子.} 只有两个顶点、一条边权为 $1$，$M=7$。从一侧走到另一侧必须走奇数条边，余数为 $2^h-1$（$h$ 奇）。可得到 $1,0,3$，不能得到所有余数。这个例子能检查“只看连通图就全可达”的错误直觉。
\cost 预处理 $O(M+(n+m)\alpha(n))$，每问常数次并查集查询；空间 $O(M+n+m)$。这里 $\alpha$ 是反阿克曼函数，不是素数幂指数。
\begin{teach}[为何这道题适合作为同余代数的结课题]
可逆性用于把有向动态改为无向等价；裴蜀用于生成平移子群；取模用于压缩数值；乘四的周期用于压缩状态。四个证明缺一不可，仅写“开六倍点并查集”无法解释算法。
\end{teach}

\src{https://atcoder.jp/contests/agc031/tasks/agc031_f}
\src{https://www.luogu.com.cn/article/hz11gwzl}

\section{WC2020 猜数游戏：把期望拆成分量贡献}
\begin{problem}{P6730：猜数游戏}
给定互不相同的非零余数 $a_1,\ldots,a_n$，模数 $M$ 为奇素数幂。随机等概率选择一个非空子集。询问一个被选中的 $a_i$ 时，可得知子集中所有能够写成 $a_i^k\bmod M$（$k$ 为正整数）的数。题目对每个固定子集考察其最少所需询问次数，求该最小值的期望乘 $(2^n-1)$ 后模 $998244353$ 的结果。
$n\le5000,M\le10^8$。
\end{problem}
\hint 建立“一个数的正整数幂能到另一个数”的传递关系；将强连通分量视为一组，统计它被选中且所有严格祖先都未选中的概率。单位部分用阶，非单位部分直接枚举短幂链。

\subsection{先分清题目中的最优含义}
这里求的是每个子集的\textbf{最小覆盖询问数}的平均值，不是在未知答案下设计一个自适应策略、使平均猜测次数最小。原题样例中子集只含 $1$ 时最少次数为一，正体现了这一口径。把它误建成在线决策问题会得到不同答案。

建边 $i\to j$ 当 $a_i^k\equiv a_j\pmod M$ 对某个 $k\ge1$ 成立。幂再取幂仍为幂，所以该关系传递。缩点后是一个传递 DAG。
对一个固定选中子集，只有“被选中且没有任何被选中的严格祖先”的分量必须各询问一次；其他选中分量都由某个上方的选中分量覆盖。同一强连通分量中只需选一个被选中的点询问。

\subsection{期望线性性给出封闭计数式}
对分量 $C$，设其大小为 $c_C$，能够到达它但不属于它的原图点数为 $u_C$。它贡献一次询问，当且仅当自己至少选一个点、全部严格祖先都不选，其余任意选。
满足该事件的子集数是
\[
 (2^{c_C}-1)2^{n-c_C-u_C}.
\]

\begin{tip}[口径：两个\textbf{平均询问数}不是一回事]
\textbf{对所有子集，求最小询问数的平均}与\textbf{设计一个策略，使平均询问数最小}是两个不同的量。前者是固定对象的平均，可以直接换序计数；后者要证明策略最优。读题时先把这句话写完，再决定用哪种概率语言。
\end{tip}

\begin{tip}[为什么不需要\textbf{独立}]
期望线性性：$\mathbb E[\sum_iX_i]=\sum_i\mathbb E[X_i]$，对任意随机变量都成立，不要求独立。这也正是它能化简\textbf{互相牵连的覆盖关系}的原因：把每个强连通分量当成一个指示变量，只需算\textbf{它被计入}的概率，不必管分量之间是否相关。
\end{tip}

这些子集一定非空，因此把所有分量贡献相加就是题目要求的已乘分母结果：
\[
 \boxed{\sum_C(2^{c_C}-1)2^{n-c_C-u_C}\pmod{998244353}}.
\]
这里不用再求 $(2^n-1)^{-1}$。

\subsection{不用真的建出平方规模的图}
单位的幂仍是单位；非单位在素数幂模数下的正整数幂始终为非单位，故两部分之间没有边。
\begin{itemize}
\item 单位部分：$M$ 有原根，所以 $i\to j$ 等价于 $\ord_M(a_j)\mid\ord_M(a_i)$。阶相同的元素构成一个分量，按阶计数后，$u_C$ 等于所有出现的\textbf{严格倍数阶}的元素总数。直接两两检查阶的整除关系也只有 $O(n^2)$。
\item 非单位部分：写 $M=p^\alpha$，$\nu_p(a_i)=v>0$。$a_i^k$ 的赋值在变为零之前为 $kv$，严格递增，因此不同非单位不可能相互可达，每个点自成分量。从每个 $a_i$ 枚举 $k=1,\ldots,\ceil{\alpha/v}-1$ 的非零幂，用值到下标的哈希表查找；命中其他点就给其祖先数加一。$k=1$ 的自己不计入严格祖先。
\end{itemize}

\paragraph{样例一完整核对.}
$M=7$，数为 $1,3,4,6$，阶分别为 $1,6,3,2$。对应祖先数为 $3,0,1,1$，分量都为单点，贡献为 $1,8,4,4$，总计 $17$。
\paragraph{样例二核对.}
$M=9$，数为 $1,\ldots,8$。单位的阶 $6,3,2,1$ 的组大小为 $2,2,1,1$，贡献分别为 $192,48,32,4$；非单位 $3,6$ 各贡献 $128$，总计 $532$。
\cost 分解并求阶后，关系计数 $O(n^2+n\alpha)$，额外空间可做到 $O(n+\alpha)$，无需存全部边。若显式建图再缩点，虽然概念清楚，也应把 $O(n^2)$ 存储写出来。

\src{https://www.luogu.com.cn/problem/P6730}


\chapter{素性判定与因数分解}
\begin{chapterintro}[章节导读]
上一章的综合题告诉我们，很多数论问题最后都绕不开“这个数的结构到底是什么”。这一章把问题直接拉回素数与因数分解：先从试除和费马测试开始，再进入 Miller--Rabin 与 Pollard--Rho，建立从“判断是不是素数”到“真正拆出因子”的完整工具链。
\end{chapterintro}

\section{试除、费马测试与 Carmichael 数}
试除只需检查到 $\sqrt n$：若 $n$ 合数，至少有一个非平凡因子不超过 $\sqrt n$。对单个 $64$ 位整数，试除通常太慢；对一批不太大的数，预先筛素数仍然有效。
费马小定理给出必要条件：素数 $n$、$\gcd(a,n)=1$ 时 $a^{n-1}\equiv1\pmod n$。但其逆命题不成立。
Carmichael 数是对所有\textbf{与 $n$ 互质}的 $a$ 都通过费马测试的合数。这里不能把条件改成 $n\nmid a$。


\paragraph{例：$561$.} $561=3\cdot11\cdot17$，$2,10,16$ 都整除 $560$。对所有与 $561$ 互质的 $a$，费马小定理分别给出 $a^{560}\equiv1$ 模 $3,11,17$，CRT 得到模 $561$ 也为 $1$。所以随机重复费马测试无法可靠识别这一类合数。
\begin{teach}[从必要条件寻找更强必要条件]
对奇素数，平方等于一的元素只有 $\pm1$。不仅要检查最终 $a^{n-1}$，还要检查通过连续平方到达一的过程里是否出现了非法的“非平凡平方根”。这正是 Miller--Rabin 的核心。
\end{teach}

\begin{tip}[词义：非平凡平方根]
$x^2\equiv1\pmod n$ 的解里，$x\not\equiv\pm1$ 的那个叫非平凡平方根。它的价值在于：$n\mid(x-1)(x+1)$ 但 $n$ 不整除任何一个因子，于是 $\gcd(x-1,n)$ 必是真因子——一句话把\textbf{求平方根}变成\textbf{求 gcd}。
\end{tip}



\section{Miller--Rabin 的强伪素数测试}
对奇数 $n>2$，写
\[
 n-1=2^s d,\qquad d\text{ 为奇数}.
\]
选取底数 $a$，先算 $x=a^d\bmod n$。本轮通过，当且仅当
\[
 a^d\equiv1\pmod n
 \quad\text{或}\quad
 \exists j\in\{0,\ldots,s-1\},\ a^{d2^j}\equiv-1\pmod n.
\]
不通过就能确定 $n$ 为合数；通过只说明该底数未发现合数证据。


\begin{proof}[为何素数一定通过？]
素数情形最终 $a^{d2^s}=a^{n-1}=1$。若起点已为一，第一种情况通过；否则看首次由非一平方得到一的地方，前项是 $X^2=1$ 的根，只能为 $-1$。因此上述两种情况必有一个成立。
\end{proof}
\paragraph{完整示例：底数二检测 $561$.}
$560=2^4\cdot35$，连续平方链为
\[
 2^{35}\equiv263,\quad263^2\equiv166,\quad166^2\equiv67,\quad67^2\equiv1\pmod{561}.
\]
没有出现 $-1\equiv560$，起点也不是一，所以发现合数。$67$ 是非平凡平方根，$\gcd(67-1,561)=33$，还能顺便得到因子。
\cost 一轮只做一次快速幂和至多 $s-1$ 次平方，为 $O(\log n)$ 次模乘，不要为每个指数重新快速幂而变成 $O(\log^2n)$。


\section{随机误差与确定性底数}
经典强伪素数计数定理给出：对固定奇合数 $n$，均匀随机底数的一轮误判概率不超过 $1/4$。独立测试 $k$ 轮，误判概率不超过 $4^{-k}$。这是对每个固定输入、每轮重新独立抽样的概率结论。
在 $0\le n<2^{64}$ 范围内，可用固定底数
\[
 2,325,9375,28178,450775,9780504,1795265022
\]

\begin{tip}[口径：随机的是底数还是模数]
\textbf{对固定 $n$、随机底数}和\textbf{对固定底数、随机 $n$}是两个不同的陈述，不能混用。固定底数（如 $2$）时存在反例，所以本书模板用一组固定的小素数当底数，在 $2^{64}$ 范围内是\textbf{确定}正确的判定，而不是概率判定。
\end{tip}

\begin{tip}[读一条测试链的方法]
先算 $a^d$，再一路平方到 $a^{d2^s}=a^{n-1}$。判定规则是：起点是 $1$（通过），或首次到达 $1$ 的前一项是 $-1$（通过）；否则 $n$ 是合数，而且那个前一项就是非平凡平方根，顺手 $\gcd$ 出因子。$561$ 是最小的反例（Carmichael 数），所以费马测试单独用不安全。
\end{tip}

得到确定性判定。每个底数先对 $n$ 取模，余数为零时跳过该轮；该有限范围保证不能外推到任意大整数。


\subsection{课堂可证明的较弱上界：\texorpdfstring{$1/2$}{1/2}}
为了理解误差来源，下面证明一个后文也会用到的引理；$1/4$ 的更紧界作为经典结论使用。
\begin{lemma}[已知群指数倍数时的平方根分裂]
设奇数 $n$ 至少含两个不同素因子，$T=2^s d$ 为 $\varphi(n)$ 的倍数，$d$ 为奇数。在模 $n$ 单位中均匀选择 $a$，序列 $a^d,a^{2d},\ldots,a^{2^sd}=1$ 中，最后一个非一元素存在且不为 $-1$ 的概率至少为 $1/2$。
\end{lemma}
\begin{proof}
先把这一段用到的每个词落到具体计算上，读者可按行核对。

设定。写 $n=\prod_i q_i^{e_i}$（$q_i$ 为奇素数，因为固定底数情形只处理奇合数）。模 $q_i^{e_i}$ 的单位群是循环群，阶为
\[
 \varphi(q_i^{e_i})=q_i^{e_i-1}(q_i-1)=:2^{s_i}t_i,\qquad t_i\ \text{奇},\ s_i\ge1 .
\]
$d$ 由 $n-1=2^s d$（$d$ 奇）给出，$T$ 是 $\varphi(n)$ 的倍数，而 $\varphi(n)=\prod_i2^{s_i}t_i$，故每个奇部分 $t_i$ 都整除 $d$。

单个素数幂上的失败概率。设 $a_i$ 在单位群中均匀分布，群为循环群 $C_{2^{s_i}t_i}\cong C_{2^{s_i}}\times C_{t_i}$。$a_i^d$ 的奇部分被 $d$（它是 $t_i$ 的倍数）消掉，所以 $a_i^d$ 落在 $2$-部分 $C_{2^{s_i}}$ 上且均匀。定义 $J_i$ 为\textbf{从 $a_i^d$ 开始连续平方，第一次得到 $1$ 所需的次数}（$0$ 表示它已经是 $1$）。在 $C_{2^{s_i}}$ 中，$2^{s_i-j}$ 阶元素恰有 $\varphi(2^{s_i-j})=2^{s_i-j-1}$ 个（$j\ge1$），故
\[
 \Pr(J_i=0)=2^{-s_i},\qquad \Pr(J_i=j)=2^{-1}\cdot2^{-s_i+j-1}\cdot2=2^{\,j-1-s_i}\ (1\le j\le s_i).
\]
两式都 $\le1/2$，因为 $s_i\ge1$。这正是\textbf{一个局部测不出问题的概率至多一半}的来源。

多个素数幂同时失败。中国剩余定理说\textbf{各素数幂上的余数可以独立选取}，故 $J_1,\ldots,J_r$ 相互独立。测试通过要求：要么某个 $a_i^d=1$（全局 $a^d=1$），要么所有局部序列在同一个 $j$ 处首次到达 $-1$ 后再平方变 $1$——即诸 $J_i$ 必须全部相等。固定 $J_1=j$ 后，其余 $J_i$ 各自以不超过 $1/2$ 的概率落到这个值，因此\textbf{全相等}的概率至多 $(1/2)^{r-1}$。当 $n$ 是素数幂（$r=1$）时，用\textbf{$X^2\equiv1$ 在奇素数幂模数下只有 $\pm1$ 两个根}直接给出同一个界。

合起来：一轮通过的概率 $\le1/2$，而经典计数定理给出更强的 $1/4$；本书按 $1/4$ 叙述，因为它对任意奇合数成立且证明可查（下一段的 $1/2$ 是无条件版本，二者不矛盾，只是强弱不同）。
\end{proof}
若 $n=q^e$ 为奇素数的真幂，费马测试通过的单位比例仅 $1/q^{e-1}\le1/3$。
若 $n$ 至少有两个不同素因子，取 $T=(n-1)\varphi(n)$。通过 Miller--Rabin 的底数，其各局部二幂阶同步；把奇数部分再乘 $\varphi(n)$ 的奇数部分不会改变这些二幂阶，所以也属于上面引理中“不分裂”的集合。于是得到不超过 $1/2$ 的较弱误判上界。
原稿的长证明中模数 $q_i$ 与 $q_i^{e_i}$、二进制指数与奇数部分多次混淆，这里统一用局部循环群的分布表示。

\begin{pitfall}[随机结论与安全工程]
随机底数不能固定为随意挑出的几个小素数后仍声称对任意大整数有 $4^{-k}$ 的误差界。$64$ 位的固定底数集是另一种有适用范围的结论。竞赛用随机数发生器也不是密码学随机源；本讲义模板不直接用作安全协议实现。
\end{pitfall}

\src{https://cp-algorithms.com/algebra/primality_tests.html}

\section{Pollard--Rho：在未知模素因子下找碰撞}
若 $n$ 合数，寻找 $1<d<n,d\mid n$，再递归分解 $d,n/d$。寻找非平凡因子前先做素性判定，并处理 $1$、偶数与小素因子。
设 $p$ 是 $n$ 的最小素因子。若找到 $x\not\equiv y\pmod n$ 但 $x\equiv y\pmod p$，则 $\gcd(|x-y|,n)$ 就是非平凡因子。
采用伪随机迭代
\[
 x_{i+1}=x_i^2+c\pmod n.
\]
模 $p$ 的投影也遵循同样的迭代。用 Floyd 快慢指针或 Brent 判环，在不知道 $p$ 的情况下通过 gcd 检查是否发生局部碰撞。


\subsection{生日悖论如何给出规模估计}
若抽取近似均匀的模 $p$ 状态，$r$ 个状态没有碰撞的概率近似
\[
 \prod_{j=0}^{r-1}(1-j/p)\approx\exp(-r(r-1)/(2p)).
\]
所以 $r$ 达到 $\sqrt p$ 量级便有常数碰撞概率。但若把全部数对都检查，需要 $O(r^2)=O(p)$ 次 gcd；原稿把这一朴素两两比较的成本写小了。
Floyd 每轮更新慢指针一次、快指针两次，检查两者差的 gcd。只要投影迭代进入循环，两指针在一个环长内相遇，因此用线性于序列长度的检查数代替了全部数对。

\paragraph{为什么有时会得到 $n$？}
如果两个状态模 $n$ 也相同，gcd 就是 $n$，并没有得到真因子。应换常数 $c$、初始值后重试，不能把 $n$ 原样递归下去，否则无限递归。

\subsection{批量 gcd 与故障回退}
累积一批差值的乘积
\[
 P=\prod_i|x_i-y_i|\pmod n,
\]
再求一次 $\gcd(P,n)$。结果为 $1$ 说明整批没有发现因子；在 $(1,n)$ 之间就成功；等于 $n$ 则回退到该批起点逐项求 gcd，或重新启动。
如果两项分别含不同素因子，乘起来也可能让 gcd 变成 $n$，所以“乘积 gcd 为 $n$”不等于这一批没有可用因子。
设批长为 $B$，常见估计为
\[
 O\left(\sqrt p+\frac{\sqrt p}{B}\log n+B\log n\right).
\]
取 $B\asymp\log n$ 后应先写成 $O(\sqrt p+\log^2n)$，恢复成本不能不加说明就消失。实际实现常用固定批长减少常数。

\subsection{应怎样表述复杂度}
在把多项式迭代视作足够随机的启发式假设下，常见的寻找因子规模为 $\widetilde O(\sqrt p)$，因 $p\le\sqrt n$，通常用约 $n^{1/4}$ 的经验／启发式量级描述。
这不是已证明对每个输入和每个随机参数成立的最坏界。每个递归子问题变小也不单独证明总成本由第一轮支配，需要结合因子的大小与各层成本分析。

\begin{teach}[给学生的实验任务]
比较 $n$ 为小素数乘大素数、两近似相等素数的乘积、素数幂时的运行情况。固定并记录随机种子用于调试；统计模乘、gcd、重启次数，而不是仅报告某次运行耗时。判素与因数分解验证应独立：输出因子必须都通过判素，乘回必须等于输入。
\end{teach}


\section{位数模型与大整数题的边界}
输入一个整数 $n$ 需要 $\Theta(\log n)$ 位。通常意义下的多项式时间是关于位数的多项式，而不是关于数值 $n$ 的多项式。
对一般整数分解，目前没有已知的经典多项式时间算法；量子计算模型下存在 Shor 算法，这与竞赛环境中的经典算法不同。素性判定则已有确定性多项式算法，不能把“判素难”和“分解难”混为一谈。

\paragraph{大整数实现提醒.} $64$ 位模乘可用无符号 $128$ 位中间值；$150$ 位输入不能套用该实现。模乘、快速幂、开整数根、gcd 都要换成高精度版本，并把它们的位复杂度计算进去。下一章有额外信息的分解题可获得随机多项式算法，关键在于已知一个与真实因子群阶相关的倍数，而不是 Pollard--Rho 的直接加速。



\chapter{特殊信息如何把分解变成随机多项式问题}
\begin{chapterintro}[章节导读]
上一章我们讨论的是在信息有限时怎样主动寻找因子；如果题目额外告诉我们 φ(n) 等信息，问题的结构就会发生变化。利用这些额外信息，可以先恢复与因子有关的代数关系，再把分解转化成随机选择参数、寻找非平凡平方根的过程。
\end{chapterintro}

\section{QOJ 4920：已知 \texorpdfstring{$n$}{n} 与 \texorpdfstring{$\varphi(n)$}{phi(n)}}
\begin{problem}{给定整数与欧拉函数值，分解质因数}
已知 $N$ 和 $\varphi(N)$，$N\le2^{150}$。不要求使用固定宽度整数实现，设计关于输入位数的随机多项式时间分解算法。
\end{problem}

\begin{tip}[这一节为什么\textbf{特殊信息}是关键]
$N\le2^{150}$ 不可能用 $O(\sqrt N)$ 枚举。但已知 $\varphi(N)$ 时，$T=(N-1)\varphi(N)$ 是单位群阶的倍数，于是\textbf{求平方根再 $\gcd$}的每一步都在 $\operatorname{poly}(\log N)$ 时间内完成。缺了这个倍数，就退回 Pollard--Rho 的启发式，不能宣称多项式时间。
\end{tip}

\hint 已知 $T=\varphi(N)$，因此对每个递归因子 $n\mid N$ 的单位 $a$ 都有 $a^T\equiv1\pmod n$。通过连续平方找非平凡平方根，再用 gcd 分裂。


\subsection{先理解一个二素数例子}
若已知 $N=pq$、$p,q$ 为不同素数，则 $\varphi(N)=N-p-q+1$，所以
\[
 p+q=N-\varphi(N)+1,\qquad pq=N.
\]
直接解一元二次方程即可。例如 $N=77,\varphi(N)=60$，和为 $18$、积为 $77$，两根为 $7,11$。一般 $N$ 可能含很多素因子，不能再把它当成二次方程，但“已知群阶的倍数”依然可以使用。

\subsection{保持原始的群阶倍数}
在所有递归中固定 $T=\varphi(N)$。若 $n\mid N$，则 $\varphi(n)\mid\varphi(N)$：按素数幂公式比较，删除素因子或降低指数后，欧拉函数都是原值的约数。因此 $T$ 仍适用于每个子问题。
先处理 $n=1$、素数、二的幂因子；若 $n=b^e$ 是完全幂，则递归分解 $b$ 并将结果指数乘 $e$。整数开根可枚举 $e\le\log_2n$，用二分或整数根算法验证，不使用浮点近似作最终判定。
剩余 $n$ 为奇合数且至少有两个不同素因子。

\subsection{一次随机尝试}
\begin{enumerate}
\item 随机 $1\le a<n$，若 $1<\gcd(a,n)<n$，已经得到因子。
\item 否则 $a$ 为单位，写 $T=2^sd$、$d$ 奇，计算 $a^d,a^{2d},\ldots,a^{2^sd}=1$。
\item 若在链中发现 $V\not\equiv\pm1$ 且 $V^2\equiv1\pmod n$，取 $g=\gcd(V-1,n)$。
\end{enumerate}
由于 $n\mid(V-1)(V+1)$ 而 $n$ 不整除两因子中的任何一个，$g$ 是非平凡因子。上一章的平方根分裂引理保证，在单位上一次尝试成功概率至少为 $1/2$。

\paragraph{数值示范.} $n=77,T=60=4\cdot15$，取 $a=2$，$2^{15}\equiv43\pmod{77}$，$43^2\equiv1$，且 $43\ne1,76$。于是 $\gcd(42,77)=7$，得到另一个因子 $11$。
\cost 每次尝试的模幂、平方、gcd 都是关于 $\log N$ 的多项式操作，每个有效分裂严格减小子问题，递归树规模为 $O(\log N)$。可使用确定性多项式判素保证理论正确性；实践中常使用高精度 Miller--Rabin，并对整体误差预算单独分配。截断重试时，用联合界控制全树失败概率。

\src{https://qoj.ac/problem/4920}

\section{QOJ 1860：由 \texorpdfstring{$M=N\varphi(N)$}{M=Nphi(N)} 还原 \texorpdfstring{$N$}{N}}
\begin{problem}{构造满足 $M=N\varphi(N)$ 的整数}
给定正整数 $M\le10^{36}$，保证存在正整数 $N$ 使 $M=N\varphi(N)$，求这个 $N$。有多组数据。
\end{problem}
\hint 先证明 $N$ 唯一，再选择性地寻找 $N$ 所含素因子。未知 $N$ 时仍知道 $M$ 是 $\varphi(N)$ 的倍数，但它不一定是递归整数 $m$ 的欧拉函数的倍数。


\subsection{第一步：解确实唯一}
若 $N=\prod q_i^{e_i}$，则
\[
 M=\prod_i(q_i-1)q_i^{2e_i-1}.
\]

\begin{tip}[为什么要先证唯一性]
构造类题目（\textbf{求满足条件的 $N$}）如果解不唯一，输出格式就必须规定取哪一个；题面只让输出一个数，等于在提示你\textbf{解唯一}。先证唯一（这里靠最大素因子的指数 $2e-1$ 被唯一确定），算法就能按\textbf{从大到小逐个恢复}写，不必回溯搜索。
\end{tip}

设 $q$ 为 $N$ 的最大素因子。因为 $q_i-1<q_i\le q$，它们不会引入比 $q$ 更大的素因子，也不会额外贡献 $q$；所以 $q$ 同时是 $M$ 的最大素因子，并且
\[
 \nu_q(M)=2e_q-1.
\]
$q$ 的指数唯一确定。除去 $(q-1)q^{2e_q-1}$ 后，又得到剩余部分的同型问题，依次从大到小恢复即可。$M=1$ 时 $N=1$。

\paragraph{示例.} $M=48$，最大素因子 $3$ 的指数为一，故 $N$ 中有一个 $3$；除去 $3(3-1)=6$ 后剩 $8$。$2$ 的指数为三，所以 $N$ 中有两个 $2$；最终 $N=3\cdot4=12$。

\subsection{第二步：不必完全分解原数}
$M$ 中可能含仅来自 $q_i-1$、却不整除 $N$ 的素因子。例如 $N=15$ 时 $M=120$，其中 $2\nmid15$。
算法只需要收集到 $N$ 的所有不同素因子，允许额外收集少量不需要的素数。递归处理当前整数 $m\mid M$，始终固定指数 $T=M$。

用于证明的量 $d_0=\gcd(m,N)$ 是未知的，\textbf{算法不能计算它}。不过 $\varphi(d_0)\mid\varphi(N)\mid T$，因此对所有与 $m$ 互质的 $a$，
\[
 a^T\equiv1\pmod{d_0}.
\]
这就是能够安全剪枝的原因。

\subsection{一次递归怎样处理？}
\begin{enumerate}
\item $m=1$ 返回；$m$ 为素数就记录；提出因子 $2$ 并记录它。若 $m=b^e$ 为完全幂，只需递归 $b$ 以收集不同素因子。
\item 对剩余奇合数随机 $a\in[1,m)$。若 gcd 已为非平凡因子，直接分裂并递归。
\item 否则计算 $D=\gcd(a^T-1,m)$。因为 $d_0\mid D$：
\begin{itemize}
\item $D=1$：说明 $\gcd(m,N)=1$，整个 $m$ 中没有需要的素因子，可以安全丢弃；
\item $1<D<m$：得到非平凡分裂，递归 $D,m/D$；
\item $D=m$：说明本次 $a^T\equiv1\pmod m$，进入连续平方链寻找非平凡平方根。
\end{itemize}
\item 在第三种情况中，写 $T=2^sd$、$d$ 奇，从 $a^d$ 不断平方。若得到 $V^2\equiv1$、$V\not\equiv\pm1$，使用 $\gcd(V-1,m)$ 分裂；若未得到，重新随机。
\end{enumerate}
\begin{pitfall}[原稿中一个使算法直接失效的符号错误]
当 $V\equiv-1\pmod m$ 时，$\gcd(V+1,m)=m$，并非非平凡因子。必须要求 $V$ 是既非 $1$ 也非 $-1$ 的平方根，才用 $V\pm1$ 分裂。
\end{pitfall}

\subsection{第三步：为什么仍有常数成功概率？}
这里不能直接说 $\varphi(m)\mid T$，因为它未必成立。正确做法是对通过 $a^T=1$ 的底数集合进行条件分析。
模 $m$ 的单位群中，满足 $a^T=1$ 的元素构成子群。按每个奇素数幂分量来看，这个子群的大小为 $\gcd(T,\varphi(q^e))$。当 $M>1$ 时 $T=M$ 为偶数，因此每个局部子群的二幂部分至少有两个元素。
条件于 $a^T=1$ 后，CRT 仍使不同局部分量独立均匀；平方链首次变为一的层数分布与上一章引理相同，每项概率至多 $1/2$。至少有两个不同素因子时，找到非平凡平方根的条件成功率至少 $1/2$。
若不满足 $a^T=1$，前面的 $D<m$ 已经使算法安全剪枝或分裂。因此每次尝试无法推进的概率至多为 $1/2$，不需要错误地假设 $T$ 是整个 $m$ 的群阶倍数。

\subsection{最后按降序恢复答案}
将收集到的不同素数从大到小排序，令剩余值 $R=M$、答案 $A=1$。遇到 $q\nmid R$ 就跳过；若 $q\mid R$，则当前最大的尚未处理相关素数一定是剩余 $N$ 的最大素因子。计算 $e=\nu_q(R)$，应为奇数，令 $k=(e+1)/2$，更新
\[
 A\gets A q^k,\qquad R\gets\frac{R}{(q-1)q^e}.
\]
直到 $R=1$。最终用收集到的质因数分解计算 $\varphi(A)$，验证 $A\varphi(A)=M$。
剪枝不会遗漏 $N$ 的素因子；额外收集的素数要么已随更大的 $q-1$ 被消去，要么最终不整除 $R$，因此不会造成错误恢复。

\cost 这是关于 $\log M$ 的随机多项式路线，依赖高精度运算、整数完全幂检测与可靠判素。理论上可以采用确定性多项式判素并无限重试形成期望多项式算法；实践中若用概率判素和有限重试，应声明错误预算并验证输出。它不是直接对 $10^{36}$ 运行 $64$ 位 Pollard--Rho。

\src{https://qoj.ac/problem/1860}


\chapter{韦达跳跃与构造不定方程的解}
\begin{chapterintro}[章节导读]
前面的分解算法依赖“把一个整体拆成因子”；这一章换成另一类离散结构：从一个已经存在的整数解出发，如何构造出另一个更小的解。韦达跳跃把二次方程的两个根联系起来，再结合正性与最小性建立递降，从而把许多看似只能试答案的构造题变成系统的证明。
\end{chapterintro}

\section{把一个变量看成二次方程的根}
若 $at^2+bt+c=0$ 的两个根为 $t_1,t_2$，则
\[
 t_1+t_2=-b/a,\qquad t_1t_2=c/a.
\]

\begin{tip}[韦达跳跃到底\textbf{跳}什么]
把方程看成关于某个变量（比如 $z$）的一元二次方程：$z^2-k(x+y)z+(x^2+y^2-1)=0$。已知一个根 $z$，用韦达定理 $z+z'=k(x+y)$ 直接写出另一个根 $z'=k(x+y)-z$——它就是解，且通常更小（降）或更大（升）。构造题里我们反向用它\textbf{从小解生成无穷多个不同解}。
\end{tip}

当方程各系数为整数、首项系数为一且已知一根为整数时，另一根也为整数。这种“固定其余变量、替换一根”的操作称为韦达跳跃。
若新根比旧根小且仍满足约束，可用于无穷递降；反过来可从基本解生成更多解。

\begin{problem}{P1400：Easy Equation}
给定 $2\le k\le1000$、$1\le n\le1000$，输出 $n$ 组正整数解
\[
 x^2+y^2+z^2=k(xy+yz+zx)+1.
\]
全部 $3n$ 个输出整数必须互不相同，每个数最多 $100$ 位，即小于 $10^{100}$。
\end{problem}
\hint 固定两变量，第三变量可替换为 $k$ 倍的另外两数之和减去自身；从基本解建立递降树，反向广度优先生成并检查全局重复值。


\subsection{推导跳跃公式而不是记变换}
固定 $x,y$，将方程改写为
\[
 z^2-k(x+y)z+x^2+y^2-kxy-1=0.
\]
如果 $z$ 为一根，另一根为
\[
 z'=k(x+y)-z.
\]
因此 $(x,y,z')$ 仍满足原方程。该操作是对合：再替换一次就回到 $z$。但另一根可能为零或负数，不能无条件说它仍是正整数解。

\subsection{严格递降与停止条件}
设 $1\le x\le y\le z$，定义上面的二次多项式 $F(t)$。有
\[
 F(y)=(2-k)y^2+x^2-2kxy-1<0.
\]
首项系数为正，故两根分别在 $y$ 两侧；$z$ 是较大根，另一个根满足 $z'<y<z$。
若 $z'>0$，则排序后得到和更小的正整数解；反复操作不可能无限下降。若 $z'\le0$，在正整数解空间中停止。这定义了以边界解为根的递降森林。

\begin{pitfall}
“不停得到更小整数”只有同时保证仍处于研究对象的集合里，才能组成无限递降矛盾。整数性、正性和严格变小是三个独立检查。对本题我们把越过正数边界的点视为基本解，而不是强行继续下降。
\end{pitfall}

\subsection{从基本解向上生成：两个孩子}
从 $1\le x<y<z$ 的解出发，替换两个较小的变量，排序后的两个孩子为
\[
 (y,z,k(y+z)-x),\qquad (x,z,k(x+z)-y).
\]
由于 $k\ge2$，两个新产生的数都严格大于 $z$，所以仍是严格递增的正整数三元组。孩子再替换最大数就回到父亲；父亲由此唯一，所以这棵树不会出现同一三元组从不同父亲重复生成的情况。
然而\textbf{三元组不重复不意味着其中的整数不重复}：父子自然共享两个数，因此输出时还要维护一个全局已使用整数集合。

\subsection{一个适用于所有合法 \texorpdfstring{$k$}{k} 的种子}
非负整数解 $(0,0,1)$ 通过两次跳跃可得到
\[
 (0,1,k)\longrightarrow(1,k,k(k+1)).
\]

\begin{tip}[为什么需要\textbf{统一种子}]
BFS 从一个固定的最小三元组出发，能保证不会漏掉某个 $k$；同时要用\textbf{全局去重}（同一数不能出现在两组解里）来筛掉重复。少了任一条，程序要么少解要么重解，而对拍很难发现——因为它在小数据上经常是对的。
\end{tip}

最后一组是严格递增的正整数解，可作为统一的搜索起点。这样无需先枚举几千万对基本解。

\begin{center}
\begin{tikzpicture}[every node/.style={draw=ink!50,rounded corners,fill=soft,font=\small,inner sep=5pt},>=Stealth]
\node (a) at (0,0) {$(0,0,1)$};
\node (b) at (0,-1.1) {$(0,1,3)$};
\node (c) at (0,-2.2) {$(1,3,12)$};
\node (d) at (-2.7,-3.5) {$(3,12,44)$};
\node (e) at (2.7,-3.5) {$(1,12,36)$};
\draw[->] (a)--(b);\draw[->] (b)--(c);\draw[->] (c)--(d);\draw[->] (c)--(e);
\end{tikzpicture}
\end{center}
图为 $k=3$ 的局部生成关系，由 TikZ 内嵌绘制，无需原稿缺失的 \texttt{tree\_k3.png}。

\subsection{BFS 与全局去重的可实现策略}
\begin{enumerate}
\item 队列初始放入 $(1,k,k(k+1))$，已使用集合为空。
\item 取出一个三元组。若其三个值均未出现在已使用集合，就输出它并登记三个值。
\item \textbf{无论该三元组是否被输出}，都将两个孩子入队；否则会错误截断可用分支。
\item 获得所需组数后停止。
\end{enumerate}
采用大整数，例如 Python 整数或 C++ 的 \texttt{boost::multiprecision::cpp\_int}。计算判定等式时需要平方，不能因为坐标本身似乎较小就假定中间值也适合 $128$ 位。

\subsection{正确性证明与规模证据分开陈述}
前面的代数推导严格保证：每个生成三元组满足方程、三个数为正且不同、父子关系无重复；全局集合保证输出的 $3n$ 个值两两不同。
仍需说明在位数限制内能否找到足够多组，不能用“树是无限的”代替定量保证。

本次整理对本题\textbf{全部 $999$ 个合法参数 $k=2,\ldots,1000$}运行统一种子的 BFS，均生成了 $1000$ 组，并逐组用大整数验算方程与全局去重条件。有限检查结果为：最多访问 $7132$ 个节点，最深输出层为 $12$，输出整数最大为 $43$ 位。任意更小的 $n$ 直接取相应前缀，因此覆盖原题全部参数组合。
验证脚本及逐 $k$ 记录附于资料包。这是可复现的\textbf{有限范围计算验证}，不是关于任意 $k,n$ 的渐近定理；教师应明确区别。可让提高层学生尝试补充适用于无界参数的定量树分析。

\paragraph{一个独立的位数上界核对.}
每次孩子的最大数小于原最大数的 $2k$ 倍，因此深度 $d$ 时最大数小于
\[
 k(k+1)(2k)^d.
\]
结合有限验证得到的 $d\le12,k\le1000$，该式也远小于 $10^{100}$。实际出现的 $43$ 位数已经可能超过有符号 $128$ 位的范围，因此本方案不能沿用原稿“\texttt{\_\_int128} 一定足够”的笼统说法。

\src{https://www.luogu.com.cn/problem/P1400}



% ============================================================================
%  第三部分  组合数学
% ============================================================================
\BeginPart{3}{6}{III}{第三部分}{组合数学}
% !TeX program = xelatex
% ============================================================================
%  第三部分  组合数学
%  第 1 章  计数原理与排列组合
%  第 2 章  二项式定理与组合恒等式
%  第 3 章  容斥原理
%  第 4 章  递推关系与经典计数数列
%  第 5 章  抽屉原理与极端原理
%  第 6 章  格路计数与反射原理
%  第 7 章  组合构造与综合例题
% ============================================================================

\chapter{计数原理与排列组合}
\begin{chapterintro}[章节导读]
组合数学的第一课不是公式，而是“把要数的东西说清楚”。同一道题，把对象理解成“有序”还是“无序”、把接收者理解成“有身份”还是“无身份”，答案会差一个阶乘。本章先把加法原理、乘法原理、排列、组合与插板法这条主线走通，并用“数的是什么”作为每一步的自查口号。
\end{chapterintro}

\begin{tip}[组合数学的总口诀]
\textbf{先说清数的是什么，再选工具。}具体做法：拿出一张纸，写下“一个被计数的对象由哪些独立的选择决定”“这些选择有没有顺序”“有没有重复计数”。这三问答完，方法基本就唯一了。本书第二部分数论里反复强调的“把整数看成指数向量”，本质也是同一件事。
\end{tip}

\section{加法原理与乘法原理}
\begin{theorem}[加法原理]
完成一件事有 $n$ 类互斥的办法，第 $i$ 类有 $a_i$ 种，则总方法数为
\[
S=a_1+a_2+\cdots+a_n.
\]
\end{theorem}
\begin{theorem}[乘法原理]
完成一件事分 $n$ 个步骤，第 $i$ 步有 $a_i$ 种方法，则总方法数为
\[
S=a_1a_2\cdots a_n.
\]
\end{theorem}
两条原理本身平淡无奇，难点全在\textbf{分类是否互斥、分步是否独立}。

\begin{pitfall}[分类不互斥]
把“$1$ 到 $20$ 中能被 $2$ 或被 $3$ 整除的数的个数”算成 $10+6=16$ 就错了：$6,12,18$ 被数了两次，正确答案是 $10+6-3=13$。这正是第三部分第 3 章容斥原理的出发点。\textbf{每次分类后问一句：有没有对象同时落进两类？}
\end{pitfall}

\section{排列与组合}
\begin{definition}[排列数]
从 $n$ 个不同元素中取出 $m$ 个（$m\le n$）排成一列，排列数为
\[
A_n^m=n(n-1)(n-2)\cdots(n-m+1)=\frac{n!}{(n-m)!}.
\]
特别地，$m=n$ 时叫全排列，$A_n^n=n!$。
\end{definition}
\begin{definition}[组合数]
从 $n$ 个不同元素中取出 $m$ 个组成一个集合（不计顺序），组合数为
\[
\binom nm=\frac{A_n^m}{m!}=\frac{n!}{m!(n-m)!}.
\]
读作“$n$ 选 $m$”。规定 $m>n$ 时 $\binom nm=0$。
\end{definition}

\begin{proof}[组合数公式的推导]
$n$ 个人选 $m$ 个排队有 $A_n^m$ 种；若不计顺序，则同一个人数集合被它的 $m!$ 种排列重复计算了 $m!$ 次，故 $\binom nm=A_n^m/m!$。
\end{proof}

\begin{tip}[排列与组合的唯一区别]
\textbf{顺序算不算不同。}排队、排名、放有编号的位置——顺序算不同，用排列；选人、选子集、分组——顺序不算，用组合。做题时先把这句话写在旁边，能避免一半的错误。
\end{tip}

\section{组合数的基本性质}
\begin{theorem}[组合数的五条基本性质]
\begin{align}
\binom nm&=\binom n{n-m}&&\text{（对称性）},\\
\binom nk&=\frac nk\binom{n-1}{k-1}&&\text{（递推，}k\ge1\text{）},\\
\binom nm&=\binom{n-1}{m}+\binom{n-1}{m-1}&&\text{（杨辉三角）},\\
\sum_{k=0}^{n}\binom nk&=2^n&&\text{（二项式定理，}a=b=1\text{）},\\
\sum_{k=0}^{n}(-1)^k\binom nk&=\iva{n=0}&&\text{（二项式定理，}a=1,b=-1\text{）}.
\end{align}
\end{theorem}

\begin{proof}[杨辉恒等式的组合证明]
要证 $\dbinom{n}{m}=\dbinom{n-1}{m}+\dbinom{n-1}{m-1}$，只需分别数“从 $n$ 个人里选 $m$ 个”的方法数。把某一个人（记为甲）单独拿出来：
\begin{itemize}
\item 不含甲的选法：从剩下 $n-1$ 个人里选 $m$ 个，共 $\dbinom{n-1}{m}$ 种；
\item 含甲的选法：甲已选定，再从剩下 $n-1$ 个人里选 $m-1$ 个，共 $\dbinom{n-1}{m-1}$ 种。
\end{itemize}
两类互斥且穷尽，相加即得。这个证明是\textbf{组合证明}的样板：不化简代数式，而是对同一个量给出两种不同的数法。
\end{proof}

\begin{teach}[代数证明与组合证明都要会]
杨辉恒等式也可以纯代数地证：
\[
\binom{n-1}{m}+\binom{n-1}{m-1}
=\frac{(n-1)!}{m!(n-1-m)!}+\frac{(n-1)!}{(m-1)!(n-m)!}
=\frac{n!}{m!(n-m)!}=\binom nm.
\]
两种证明都要求学生写一遍。代数证明训练恒等变形，组合证明训练“把同一个量数两遍”——后者是第三部分第 3 章容斥原理与第 5 章抽屉原理的思想源头。
\end{teach}

\section{插板法}
\begin{theorem}[插板法]
把 $n$ 个\textbf{相同}的元素分成 $k$ 组：
\begin{itemize}
\item 每组至少一个：$\dbinom{n-1}{k-1}$ 种；
\item 每组允许为空：$\dbinom{n+k-1}{k-1}=\dbinom{n+k-1}{n}$ 种；
\item 第 $i$ 组至少 $a_i$ 个（$\sum a_i\le n$）：$\dbinom{n-\sum a_i+k-1}{k-1}$ 种。
\end{itemize}
\end{theorem}
\begin{proof}
$n$ 个相同元素排成一列，中间有 $n-1$ 个空隙。每组至少一个，就是在这 $n-1$ 个空隙里插入 $k-1$ 块板子，每种插法对应一种分组，共 $\dbinom{n-1}{k-1}$ 种。

允许为空时不能直接插（多块板子可插入同一空隙），改用“借元素”：先借 $k$ 个元素，在 $n+k$ 个元素形成的 $n+k-1$ 个空隙里插 $k-1$ 块板子，保证每组至少一个，插完再把借来的 $k$ 个元素从各组中拿走。元素相同，所以这是一个双射，答案仍为 $\dbinom{n+k-1}{k-1}$。

不同下界时令 $x_i'=x_i-a_i$，化为非负整数解，套上一种情形。
\end{proof}

\begin{pitfall}[插板法的前提是“元素相同”]
如果 $n$ 个元素\textbf{互不相同}，分组数是 $k^n$（每个元素独立选一组）或 $\dbinom{n+k-1}{k-1}$ 乘以组内排列，与插板法完全不同。读题时先看“$n$ 个相同的球”还是“$n$ 个不同的人”。
\end{pitfall}

\section{圆排列与多重集}
\subsection{圆排列}
$n$ 个人围成一圈，排列数为
\[
Q_n=\frac{A_n^n}{n}=(n-1)!,
\]
因为一个圆圈从 $n$ 个不同位置断开会得到 $n$ 个不同的队列。部分圆排列（选 $r$ 个围成一圈）为
\[
Q_n^r=\frac{A_n^r}{r}=\frac{n!}{r\,(n-r)!}.
\]
若再要求“翻转后相同”（如项链），还要除以 $2$。

\subsection{多重集的排列数}
设多重集有 $n_1$ 个 $a_1$、$n_2$ 个 $a_2$、……、$n_k$ 个 $a_k$，总数 $n=\sum n_i$。全排列数为
\[
\binom{n}{n_1,n_2,\ldots,n_k}=\frac{n!}{n_1!\,n_2!\cdots n_k!},
\]
这叫\textbf{多项式系数}。道理很简单：先当作 $n$ 个不同元素全排列，再把相同元素之间的 $n_i!$ 种排列视为同一种。

\subsection{多重集的组合数}
从上述多重集中选 $r$ 个元素（$r<n_i$ 对一切 $i$ 成立，即每种都取不完），方案数就是 $x_1+\cdots+x_k=r$ 的非负整数解数：
\[
\binom{r+k-1}{k-1}.
\]
若每种有上限 $x_i\le n_i$，则要用第 3 章的容斥原理。

\begin{tip}[三个“多重”不要混]
\textbf{多重集的排列数}（$\frac{n!}{\prod n_i!}$，计顺序）、\textbf{多重集的组合数}（$\binom{r+k-1}{k-1}$，不计顺序）、\textbf{多项式系数}（$\binom{n}{n_1,\ldots,n_k}$，就是多重集排列数的另一种记号）。三者中前两个是最容易混的。
\end{tip}

\section{分层例题与解析}
\begin{problem}{例 1：乘法原理}
一个密码由 $3$ 个大写字母后跟 $3$ 个数字组成。若字母不许重复、数字可以重复，有多少个密码？
\end{problem}
\hint 分两步：字母部分与数字部分，各自用乘法原理，再相乘。

\sol 字母部分 $26\cdot25\cdot24$，数字部分 $10^3$，共
\[
26\cdot25\cdot24\cdot10^3=15\,600\,000.
\]

\begin{problem}{例 2：组合与排列的界线}
从 $10$ 名选手中选 $3$ 人组成代表队，再指定一名队长，有多少种方案？若改为“选 $3$ 人并排成一列拍照”呢？
\end{problem}
\hint 第一问先组合再选队长；第二问直接排列。

\sol 第一问 $\dbinom{10}{3}\cdot3=120\cdot3=360$。第二问 $A_{10}^3=720$。两问不同：第一问中另外两人没有身份区别，第二问中三人位置不同。\textbf{先问“顺序算不算”}，答案自然分开。

\begin{problem}{例 3：插板法}
方程 $x_1+x_2+x_3+x_4=20$ 有多少组正整数解？多少组非负整数解？
\end{problem}
\hint 正整数解每组至少 $1$；非负解借 $4$ 个元素。

\sol 正整数解 $\dbinom{20-1}{4-1}=\dbinom{19}{3}=969$。非负整数解 $\dbinom{20+4-1}{4-1}=\dbinom{23}{3}=1771$。

\begin{problem}{例 4：不相邻的选择}
从 $1,2,\ldots,10$ 中选 $3$ 个数，任意两个不相邻，有多少种选法？
\end{problem}
\hint 设选出的数为 $a<b<c$，令 $b'=b-1$、$c'=c-2$ 去掉间隔。

\sol 设 $a<b<c$ 且 $b-a\ge2$、$c-b\ge2$。令 $a'=a$、$b'=b-1$、$c'=c-2$，则 $1\le a'<b'<c'\le8$，这是一个从 $8$ 个数里选 $3$ 的组合，共 $\dbinom83=56$ 种。一般地，从 $n$ 个数里选 $k$ 个互不相邻的数有 $\dbinom{n-k+1}{k}$ 种。

\begin{problem}{统一练习：分配问题}
把 $12$ 本相同的书分给 $4$ 个同学，每人至少 $1$ 本，有多少种分法？若这 $12$ 本书互不相同呢？
\end{problem}
\hint 第一问用插板法；第二问每本书独立选人，再用容斥减去“有人没书”的情形。

\sol 第一问：$x_1+\cdots+x_4=12$ 的正整数解数为
\[
\binom{12-1}{4-1}=\binom{11}{3}=165.
\]

第二问：每本书有 $4$ 种选择，共 $4^{12}$ 种；减去“某个人一本书都没分到”的情形。设 $A_i$ 为“第 $i$ 个人没书”，则
\[
\left|\overline{A_1}\cap\overline{A_2}\cap\overline{A_3}\cap\overline{A_4}\right|
=\sum_{j=0}^{4}(-1)^j\binom4j(4-j)^{12}.
\]
其中 $j=4$ 一项为 $0^{12}=0$（$12$ 本书不可能谁都不给），于是
\[
4^{12}-4\cdot3^{12}+6\cdot2^{12}-4\cdot1^{12}
=16\,777\,216-2\,125\,764+24\,576-4=14\,676\,024.
\]

\begin{teach}[从“至少 $1$ 本”到“至少 $2$ 本”]
若把条件改成“每人至少 $2$ 本”：相同书的情形令 $x_i'=x_i-2$，化为 $\sum x_i'=4$，答案 $\dbinom73=35$；但\textbf{不同}书的情形再也不能用一次容斥解决，需要用生成函数或按“先给每人 $2$ 本再任意分剩下 $4$ 本”的组合分解。后者留到第五部分系统处理。
\end{teach}

\chapter{二项式定理与组合恒等式}
\begin{chapterintro}[章节导读]
$(1+x)^n$ 是组合数学里最重要的一个多项式：它的第 $k$ 项系数就是 $\binom nk$。把“选法条数”看成“多项式系数”以后，乘法变成了卷积、求导变成了加权求和、代入特殊值变成了一整类恒等式。本章建立这套语言，它是第五部分生成函数的直接前身。
\end{chapterintro}

\section{二项式定理}
\begin{theorem}[二项式定理]
对任意非负整数 $n$ 与实数 $x,y$，
\[
(x+y)^n=\sum_{k=0}^{n}\binom nk x^{n-k}y^{k}.
\]
\end{theorem}
\begin{proof}[组合证明]
$(x+y)^n$ 是 $n$ 个 $(x+y)$ 相乘。展开后的每一项要从每个因子里选 $x$ 或 $y$。若恰有 $k$ 个因子选了 $y$（其余 $n-k$ 个选 $x$），得到 $x^{n-k}y^k$；选哪 $k$ 个因子有 $\dbinom nk$ 种，故系数为 $\dbinom nk$。
\end{proof}

\begin{theorem}[多项式定理]
设 $n$ 为非负整数，则
\[
(x_1+x_2+\cdots+x_t)^n
=\sum_{\substack{n_1,\ldots,n_t\ge0\\ n_1+\cdots+n_t=n}}
\binom{n}{n_1,n_2,\ldots,n_t}x_1^{n_1}x_2^{n_2}\cdots x_t^{n_t},
\]
其中 $\dbinom{n}{n_1,\ldots,n_t}=\dfrac{n!}{n_1!\cdots n_t!}$ 是多项式系数。并且
\[
\sum_{\substack{n_1+\cdots+n_t=n}}\binom{n}{n_1,\ldots,n_t}=t^n.
\]
\end{theorem}
\begin{proof}[组合证明]
展开后的每一项对应一个“函数”$f:\{1,\ldots,n\}\to\{1,\ldots,t\}$，$f(i)$ 记录第 $i$ 个因子里选了哪个变量。若 $n_j=|f^{-1}(j)|$，则该项为 $x_1^{n_1}\cdots x_t^{n_t}$，而有同样计数向量的函数有 $\dbinom{n}{n_1,\ldots,n_t}$ 个（先排 $n$ 个位置，再按组除）。函数总数为 $t^n$，即得求和式。
\end{proof}

\begin{tip}[生成函数视角]
把 $\sum_k\binom nk x^k$ 记作 $(1+x)^n$ 的\textbf{系数列}。这个视角的用处：两个多项式相乘时，系数做卷积
\[
\left(\sum_k a_kx^k\right)\left(\sum_k b_kx^k\right)=\sum_k\left(\sum_{i+j=k}a_ib_j\right)x^k.
\]
于是“$(1+x)^m(1+x)^n=(1+x)^{m+n}$”翻译过来就是范德蒙恒等式。第五部分会把这一套完整展开。
\end{tip}

\section{常用组合恒等式}
\begin{theorem}[范德蒙恒等式]
对非负整数 $m,n,k$，
\[
\sum_{i=0}^{k}\binom ni\binom m{k-i}=\binom{m+n}{k}.
\]
特别地，取 $m=n=k$ 得
\[
\sum_{i=0}^{n}\binom ni^2=\binom{2n}{n}.
\]
\end{theorem}
\begin{proof}[生成函数证明]
$(1+x)^m(1+x)^n=(1+x)^{m+n}$。左边 $x^k$ 的系数是 $\sum_{i+j=k}\binom mi\binom nj=\sum_{i=0}^{k}\binom ni\binom m{k-i}$，右边是 $\dbinom{m+n}{k}$。取 $m=n=k$ 并注意 $\binom ni=\binom n{n-i}$，第二个等式是直接推论。
\end{proof}

\begin{theorem}[朱世杰恒等式（冰球棍）]
\[
\sum_{l=0}^{n}\binom lk=\binom{n+1}{k+1}\qquad(k\ge0).
\]
\end{theorem}
\begin{proof}[组合证明]
右边数的是“从 $\{a_1,\ldots,a_{n+1}\}$ 中选 $k+1$ 个元素”的子集数。按\textbf{最大元素}分类：若最大元素是 $a_{l+1}$，则其余 $k$ 个元素要从 $\{a_1,\ldots,a_l\}$ 中选，有 $\dbinom lk$ 种。对 $l=k,\ldots,n$ 求和即得左边。
\end{proof}

\begin{theorem}[带权和的两条恒等式]
\[
\sum_{i=0}^{n}i\binom ni=n2^{n-1},\qquad
\sum_{i=0}^{n}i^2\binom ni=n(n+1)2^{n-2}.
\]
\end{theorem}
\begin{proof}[求导证明]
由二项式定理 $\sum_i\binom nix^i=(1+x)^n$。两边对 $x$ 求导：
\[
\sum_i i\binom nix^{i-1}=n(1+x)^{n-1}.
\]
令 $x=1$ 得第一条。再求导一次并令 $x=1$：
\[
\sum_i i(i-1)\binom ni=n(n-1)2^{n-2},
\]
而 $\sum i^2\binom ni=\sum i(i-1)\binom ni+\sum i\binom ni=n(n-1)2^{n-2}+n2^{n-1}=n(n+1)2^{n-2}$。
\end{proof}

\begin{teach}[“按最大元素分类”是一个通用套路]
朱世杰恒等式的证明按最大元素分类，这个思路可以推广到几乎所有“上标求和”的恒等式。让学生把 $\sum_{l}\binom lk$ 讲成“数 $k+1$ 元子集，看谁最大”，比记住公式有用得多。
\end{teach}

\section{二项式反演}
\begin{theorem}[二项式反演]
设 $f_n$ 表示“恰好使用 $n$ 个不同元素形成某结构的方案数”，$g_n$ 表示“从 $n$ 个不同元素中任选若干个（可以 $0$ 个）形成该结构的总方案数”，则
\[
g_n=\sum_{i=0}^{n}\binom nif_i,
\qquad
f_n=\sum_{i=0}^{n}(-1)^{n-i}\binom nig_i.
\]
\end{theorem}
\begin{proof}[代入验证]
把 $g_i=\sum_{j}\binom ijf_j$ 代入右边：
\[
\sum_{i=0}^{n}(-1)^{n-i}\binom ni\sum_{j=0}^{i}\binom ijf_j
=\sum_{j=0}^{n}f_j\sum_{i=j}^{n}(-1)^{n-i}\binom ni\binom ij.
\]
用恒等式 $\dbinom ni\dbinom ij=\dbinom nj\dbinom{n-j}{i-j}$，内层和变为
\[
\binom nj\sum_{i=j}^{n}(-1)^{n-i}\binom{n-j}{i-j}
=\binom nj\sum_{t=0}^{n-j}(-1)^{n-j-t}\binom{n-j}{t}.
\]
由 $\sum_t(-1)^t\binom{N}{t}=\iva{N=0}$，内层在 $n=j$ 时为 $1$、在 $n>j$ 时为 $0$。故整式和为 $f_n$。
\end{proof}

\begin{tip}[二项式反演与莫比乌斯反演的关系]
第二部分第 8 章的\textbf{莫比乌斯反演}处理的是“按约数”的偏序，二项式反演处理的是“按子集”的偏序。两者形式相同、场景不同：看到“$g_n=\sum\binom nif_i$”用二项式反演，看到“$g(n)=\sum_{d\mid n}f(d)$”用莫比乌斯反演。第五部分会用生成函数把两者统一处理。
\end{tip}

\section{分层例题与解析}
\begin{problem}{例 1：用二项式定理求系数}
求 $(2x-1)^6$ 中 $x^3$ 的系数。
\end{problem}
\hint 直接套二项式定理，$x=2x$、$y=-1$。

\sol 通项为 $\dbinom6k(2x)^{6-k}(-1)^k$。要 $x^3$ 需 $6-k=3$，即 $k=3$，系数为
\[
\binom63\cdot2^3\cdot(-1)^3=20\cdot8\cdot(-1)=-160.
\]

\begin{problem}{例 2：范德蒙恒等式}
计算 $\sum_{i=0}^{5}\binom{5}{i}\binom{5}{5-i}$。
\end{problem}
\hint 这就是 $\sum_i\binom5i\binom5{5-i}$，令 $m=n=k=5$。

\sol 由范德蒙恒等式的特殊形式，$\sum_i\binom5i\binom5{5-i}=\sum_i\binom5i^2=\binom{10}{5}=252$。

\begin{problem}{例 3：组合数什么时候为零}
计算 $\dbinom30+\dbinom31+\dbinom32+\cdots+\dbinom{3}{10}$。
\end{problem}
\hint 先看每一项是否为 $0$；$\binom nk=0$ 当 $k>n$。

\sol $\binom nk=0$（$k>n$），故只有前四项非零：
\[
\binom30+\binom31+\binom32+\binom33=1+3+3+1=8.
\]
这与 $(1+1)^3=2^3=8$ 一致。这个例子提醒：\textbf{求和上下界要先与组合数的有效范围对齐}，否则容易把空项当成非零项反复计算。

\begin{problem}{例 4：二项式反演}
$n$ 个元素的集合有多少个真子集（不含自身）？用二项式反演验证。
\end{problem}
\hint 令 $f_n$ 为“恰好用 $n$ 个元素”的方案数（只有 $1$ 种：全选），$g_n$ 为总方案数。

\sol $f_n=1$ 对一切 $n$，故
\[
g_n=\sum_{i=0}^{n}\binom ni\cdot1=2^n,
\]
正是子集总数。反过来
\[
f_n=\sum_{i=0}^{n}(-1)^{n-i}\binom ni2^i=(2-1)^n=1,
\]
自洽。这个练习用来说明：二项式反演的两条公式互为逆运算，验证时先代入一条再代回另一条。

\chapter{容斥原理}
\begin{chapterintro}[章节导读]
“至少满足一个条件”的对象有多少个？直接数会重复，容斥原理给出系统的加减法。它是组合数学里最强大的单一工具：从错位排列到欧拉函数，从格路计数到概率里的并集，全都落在它身上。第二部分第 8 章的莫比乌斯反演，本质上就是容斥原理的约数版本。
\end{chapterintro}

\section{容斥原理的形式}
\begin{theorem}[容斥原理]
设 $A_1,\ldots,A_n$ 是有限集，则
\[
\left|\bigcup_{i=1}^{n}A_i\right|
=\sum_{\varnothing\ne S\subseteq[n]}(-1)^{|S|+1}\left|\bigcap_{i\in S}A_i\right|.
\]
等价地，用补集写：
\[
\left|\overline{A_1}\cap\cdots\cap\overline{A_n}\right|
=\sum_{S\subseteq[n]}(-1)^{|S|}\left|\bigcap_{i\in S}A_i\right|.
\]
\end{theorem}
\begin{proof}[按元素被数的次数验证]
固定一个元素 $x$，设它恰好属于 $A_1,\ldots,A_n$ 中的 $r$ 个（$r\ge1$）。左边只把它数 $1$ 次。右边：含 $x$ 的交集 $\bigcap_{i\in S}A_i$ 恰对应 $S\subseteq T$（$T$ 是那 $r$ 个指标组成的集合），共 $2^r$ 项，其贡献为
\[
\sum_{S\subseteq T}(-1)^{|S|+1}=-\sum_{j=0}^{r}\binom rj(-1)^j=-\iva{r=0}=1\qquad(r\ge1).
\]
恰好 $1$ 次。若 $r=0$（$x$ 不属于任何 $A_i$），左边为 $0$，右边所有交集都不含 $x$，贡献也是 $0$。两边对每个元素的计数相同，故相等。
\end{proof}

\begin{teach}[这个证明值得反复讲]
容斥原理的证明不化简任何代数式，只做一件事：\textbf{看一个固定元素被数了几次}。这个手法在第二部分第 5 章证明 LTE 的“唯一最低赋值项”、第三部分第 5 章的抽屉原理里都会再现。建议让学生在黑板上把 $n=3$ 的情形逐个数一遍，再回到一般证明。
\end{teach}

\section{错位排列}
\begin{definition}[错位排列]
$n$ 个元素的一个排列，若每个元素都不在原来的位置上，称为\textbf{错位排列}。错位排列数记作 $D_n$。
\end{definition}
\begin{theorem}[错位排列数]
\[
D_n=n!\sum_{k=0}^{n}\frac{(-1)^k}{k!}.
\]
特别地，$D_0=1$、$D_1=0$、$D_2=1$、$D_3=2$、$D_4=9$、$D_5=44$。且
\[
\lim_{n\to\infty}\frac{D_n}{n!}=\frac1{\mathrm e}.
\]
\end{theorem}
\begin{proof}
设 $A_i$ 为“第 $i$ 个元素留在原位”的排列集合。则 $|A_i|=(n-1)!$，更一般地 $|A_{i_1}\cap\cdots\cap A_{i_k}|=(n-k)!$。由容斥，
\[
D_n=n!-\left|\bigcup_i A_i\right|
=n!-\sum_{k=1}^{n}(-1)^{k+1}\binom nk(n-k)!
=n!\sum_{k=0}^{n}\frac{(-1)^k}{k!}.
\]
极限由 $\mathrm e^{-1}=\sum_k(-1)^k/k!$ 得到。
\end{proof}

\begin{tip}[错位排列的递推]
$D_n$ 也满足递推 $D_n=(n-1)(D_{n-1}+D_{n-2})$（$n\ge2$）。组合解释：第 $1$ 个元素放到位置 $j$（$n-1$ 种选择），再分“第 $j$ 个元素回到位置 $1$”与“不回去”两种，分别对应 $D_{n-2}$ 与 $D_{n-1}$。两种推导都要会。
\end{tip}

\section{容斥的三种典型应用}
\subsection{有上限的多重集组合}
从 $n_1$ 个 $a_1$、……、$n_k$ 个 $a_k$ 中选 $r$ 个，要求 $x_i\le n_i$。全集是非负整数解 $\dbinom{r+k-1}{k-1}$；“违反第 $i$ 个上限”即 $x_i\ge n_i+1$，令 $x_i'=x_i-n_i-1$ 后化为非负整数解，可逐项计算。

\subsection{欧拉函数的容斥表达式}
第二部分第 2 章会证明
\[
\varphi(n)=n\prod_{p\mid n}\left(1-\frac1p\right),
\]
其推导正是容斥：$\varphi(n)$ 数是 $1,\ldots,n$ 中与 $n$ 互质的整数个数，即“不被任何素因子整除”的个数。

\subsection{格路与障碍物}
第三部分第 6 章的反射原理，是容斥原理在“路径被挡住”场景下的精确形式：把坏路径双射到另一类好数的路径，从而减去。

\section{分层例题与解析}
\begin{problem}{例 1：基本容斥}
某班 $50$ 人，参加数学小组的 $30$ 人，参加编程小组的 $25$ 人，两个都参加的 $15$ 人。两个都没参加的有多少人？
\end{problem}
\hint 用 $|A\cup B|=|A|+|B|-|A\cap B|$。

\sol $|A\cup B|=30+25-15=40$，故都没参加的有 $50-40=10$ 人。

\begin{problem}{例 2：错位排列}
$5$ 个人各写一张纸条放入帽中，再随机取出一张（不是自己的）。有多少种取法？
\end{problem}
\hint 这就是 $D_5$。

\sol $D_5=5!\left(1-1+\dfrac12-\dfrac16+\dfrac1{24}-\dfrac1{120}\right)=120\cdot\dfrac{44}{120}=44$。

\begin{problem}{例 3：容斥计数}
求 $1$ 到 $1000$ 中能被 $6$、$10$ 或 $15$ 至少一个整除的整数个数。
\end{problem}
\hint 用 $|A\cup B\cup C|=|A|+|B|+|C|-|A\cap B|-|A\cap C|-|B\cap C|+|A\cap B\cap C|$；每一处交集都先算最小公倍数。

\sol 记 $A,B,C$ 分别为被 $6,10,15$ 整除的数集。
\[
|A|=\floor{\frac{1000}{6}}=166,\quad |B|=\floor{\frac{1000}{10}}=100,\quad |C|=\floor{\frac{1000}{15}}=66.
\]
交集用最小公倍数：
\[
\lcm(6,10)=30,\quad \lcm(6,15)=30,\quad \lcm(10,15)=30,\quad \lcm(6,10,15)=30,
\]
故三个两两交集各为 $\floor{1000/30}=33$，三者交集也是 $33$。于是
\[
|A\cup B\cup C|=166+100+66-33-33-33+33=266.
\]

\begin{pitfall}[交集用的是最小公倍数]
“能被 $a$ 且能被 $b$ 整除”等价于“能被 $\lcm(a,b)$ 整除”。本题三对最小公倍数恰好都是 $30$，但若换成 $6$ 与 $10$、$14$ 与 $21$ 之类的组合，$\lcm$ 各不相同，\textbf{每一处交集都要单独算}。
\end{pitfall}

\begin{problem}{例 4：容斥求欧拉函数}
用容斥原理求 $\varphi(30)$，并与公式 $n\prod_{p\mid n}(1-\frac1p)$ 对照。
\end{problem}
\hint $30=2\cdot3\cdot5$；在 $1,\ldots,30$ 中数不被 $2,3,5$ 整除的数。

\sol 在 $1,\ldots,30$ 中，被 $2$ 整除的 $15$ 个、被 $3$ 整除的 $10$ 个、被 $5$ 整除的 $6$ 个；被 $2$ 与 $3$ 整除的 $5$ 个、被 $2$ 与 $5$ 整除的 $3$ 个、被 $3$ 与 $5$ 整除的 $2$ 个；被 $2,3,5$ 同时整除的 $1$ 个。于是
\[
\varphi(30)=30-(15+10+6)+(5+3+2)-1=30-31+10-1=8.
\]
公式给出
\[
30\left(1-\frac12\right)\left(1-\frac13\right)\left(1-\frac15\right)
=30\cdot\frac12\cdot\frac23\cdot\frac45=8,
\]
两者一致。

\begin{teach}[容斥的算术要逐项写]
本题只有 $7$ 个非零项，但每一项的分子分母都容易抄错。讲评时要求学生把 $30/2=15$、$30/3=10$、$30/5=6$、$30/6=5$、$30/10=3$、$30/15=2$、$30/30=1$ 七个数值先列成一列，再代入公式。\textbf{容斥的失分几乎都在算术，不在思路。}
\end{teach}

\chapter{递推关系与经典计数数列}
\begin{chapterintro}[章节导读]
很多组合问题的答案不能直接套公式，但满足一个递推关系：大问题的解由小问题的解拼成。本章处理三条最经典的递推——斐波那契、卡特兰、斯特林——并说明“递推”与“闭式”的关系：能写出递推往往就够了，闭式只是为了更快地算。
\end{chapterintro}

\section{怎么建立递推关系}
建立递推的标准动作是\textbf{按最后一个动作分类}：
\begin{enumerate}
\item 指出“一个合法对象”的最后一步有几种可能；
\item 对每种可能，说明去掉最后一步以后得到的是一个（或几个）同类的小对象；
\item 把各类计数相加或相乘，得到递推式；
\item 写出边界条件。
\end{enumerate}

\begin{pitfall}[重复计数与漏计]
递推式最常见的错误是同一个对象被两条路径数到（重复）或某类对象没有路径能到达（漏计）。写完递推式后，用 $n=1,2,3$ 的小值手工枚举核对，是成本最低的验证。
\end{pitfall}

\section{斐波那契数列}
\begin{definition}[斐波那契数列]
$F_1=F_2=1$，$F_{n+2}=F_{n+1}+F_n$。前几项：
\[
1,1,2,3,5,8,13,21,34,55,89,\ldots
\]
\end{definition}
\begin{theorem}[斐波那契的组合解释]
$F_{n+1}$ 等于用 $1\times1$ 与 $1\times2$ 的方块铺满 $1\times n$ 的长条的方法数。
\end{theorem}
\begin{proof}
设铺法数为 $f_n$，$f_0=1$（空条有一种铺法）。按最左边一块分类：若最左是 $1\times1$，剩下 $1\times(n-1)$ 有 $f_{n-1}$ 种；若最左是 $1\times2$，剩下有 $f_{n-2}$ 种。故 $f_n=f_{n-1}+f_{n-2}$，$f_0=1$、$f_1=1$，即 $f_n=F_{n+1}$。
\end{proof}

\begin{tip}[斐波那契的增长速度]
由特征方程 $x^2=x+1$ 得 $F_n=\dfrac{1}{\sqrt5}\left[\left(\dfrac{1+\sqrt5}{2}\right)^n-\left(\dfrac{1-\sqrt5}{2}\right)^n\right]$。于是 $F_n=\Theta(\varphi^n)$，其中 $\varphi=\dfrac{1+\sqrt5}{2}\approx1.618$ 是黄金比。由此可得实用估计：$F_{45}\approx1.1\times10^9$、$F_{90}\approx2.9\times10^{18}$，已经贴近 \texttt{long long} 上限。
\end{tip}

\section{卡特兰数}
\begin{definition}[卡特兰数]
\[
C_n=\frac1{n+1}\binom{2n}{n}=\binom{2n}{n}-\binom{2n}{n+1}.
\]
前几项：$C_0=1,C_1=1,C_2=2,C_3=5,C_4=14,C_5=42,C_6=132$。
\end{definition}
\begin{theorem}[卡特兰数的递推]
\[
C_0=1,\qquad C_{n+1}=\sum_{i=0}^{n}C_iC_{n-i}\quad(n\ge0).
\]
\end{theorem}
\begin{proof}[以括号序列为例]
$C_n$ 等于 $n$ 对括号的正确匹配方案数。设 $f_n$ 为该数，$f_0=1$。按\textbf{第一个左括号与谁匹配}分类：若它与第 $2i+2$ 个字符处的右括号匹配（$0\le i\le n-1$），则这对括号内部有 $i$ 对括号（$f_i$ 种），外部有 $n-1-i$ 对（$f_{n-1-i}$ 种），共 $f_if_{n-1-i}$ 种。求和得
\[
f_n=\sum_{i=0}^{n-1}f_if_{n-1-i},
\]
即 $f_{n+1}=\sum_{i=0}^{n}f_if_{n-i}$。再用生成函数或归纳可证 $f_n=C_n$。
\end{proof}

\begin{tip}[卡特兰数的家族]
以下问题的答案都是 $C_n$：$n$ 对括号的正确匹配；$n$ 个结点能形成的不同二叉树的个数；从 $(0,0)$ 到 $(n,n)$ 不越过对角线的单调格路数；$n$ 个元素依次入栈、全部出栈的合法顺序数；“$n+2$ 边形用不相交对角线分成三角形”的方法数。第三部分第 6 章会详细处理格路版本。
\end{tip}

\section{斯特林数}
\subsection{第二类斯特林数}
\begin{definition}[第二类斯特林数]
$S(n,k)$（也记作 $\begin{Bmatrix}n\\k\end{Bmatrix}$）表示把 $n$ 个\textbf{不同}的元素分成 $k$ 个\textbf{非空无标号}子集的方法数。
\end{definition}
\begin{theorem}[第二类斯特林数的递推]
\[
S(n,k)=k\,S(n-1,k)+S(n-1,k-1),
\]
边界 $S(0,0)=1$、$S(n,0)=0$（$n\ge1$）、$S(0,k)=0$（$k\ge1$）。
\end{theorem}
\begin{proof}[按第 $n$ 个元素分类]
把 $n$ 个元素分成 $k$ 个非空无标号子集。看第 $n$ 个元素：
\begin{itemize}
\item 它单独成组：其余 $n-1$ 个元素分成 $k-1$ 组，共 $S(n-1,k-1)$ 种；
\item 它加入已有的某一组：先把前 $n-1$ 个元素分成 $k$ 组（$S(n-1,k)$ 种），再把第 $n$ 个元素放入 $k$ 个组中的一个（$k$ 种选择），共 $kS(n-1,k)$ 种。
\end{itemize}
两类互斥，相加即得。
\end{proof}

\begin{tip}[“无标号”与“有标号”只差一个 $k!$]
若 $k$ 个盒子\textbf{有标号}，方案数是 $k!\,S(n,k)$。这是斯特林数应用中最常见的陷阱：读题时先看“盒子有没有编号”。
\end{tip}

\subsection{第一类斯特林数}
\begin{definition}[第一类斯特林数]
$s(n,k)$（带符号）与 $c(n)=\begin{bmatrix}n\\k\end{bmatrix}$（不带符号）表示把 $n$ 个不同元素排成 $k$ 个不相交的\textbf{圆排列}的方法数。
\end{definition}
\begin{theorem}[第一类斯特林数的递推]
\[
\begin{bmatrix}n\\k\end{bmatrix}=(n-1)\begin{bmatrix}n-1\\k\end{bmatrix}+\begin{bmatrix}n-1\\k-1\end{bmatrix},
\]
边界 $\begin{bmatrix}0\\0\end{bmatrix}=1$。它们还是下降阶乘的系数：
\[
x(x-1)(x-2)\cdots(x-n+1)=\sum_{k=0}^{n}s(n,k)x^k.
\]
\end{theorem}
\begin{proof}[递推的组合解释]
看第 $n$ 个元素：若它单独成一个圆排列，其余 $n-1$ 个排成 $k-1$ 个圆，有 $\begin{bmatrix}n-1\\k-1\end{bmatrix}$ 种；若它插入已有的某个圆排列，该圆有 $m$ 个元素时有 $m$ 个插入位置，对所有圆求和共 $n-1$ 个位置，故有 $(n-1)\begin{bmatrix}n-1\\k\end{bmatrix}$ 种。
\end{proof}

\section{分层例题与解析}
\begin{problem}{例 1：铺砖}
用 $1\times1$ 与 $1\times2$ 的方块铺满 $1\times6$ 的长条，有多少种方法？
\end{problem}
\hint 就是 $F_7$。

\sol $F_7=13$。可以直接列：$f_0=1,f_1=1,f_2=2,f_3=3,f_4=5,f_5=8,f_6=13$。

\begin{problem}{例 2：卡特兰}
$3$ 对括号有多少种正确匹配？$4$ 个结点能形成多少种不同的二叉树？
\end{problem}
\hint 都是 $C_3$ 与 $C_4$。

\sol $C_3=\dfrac16\dbinom63=\dfrac{20}{6}=5$；$C_4=\dfrac15\dbinom84=\dfrac{70}{5}=14$。

枚举验证 $3$ 对括号的 $5$ 种匹配：
\[
\texttt{((()))},\quad \texttt{(()())},\quad \texttt{(())()},\quad \texttt{()(())},\quad \texttt{()()()}.
\]
只有这 $5$ 种，不多不少。

\begin{teach}[枚举必须做到“不重不漏”]
让学生自己列出 $3$ 对括号的全部匹配，几乎一定会出现重复或遗漏。正确的做法是按第 4 章的分类：第一个左括号与谁配对。若它与第 $2i+2$ 个字符处的右括号配对（$i=0,1,2$），则内部有 $i$ 对、外部有 $2-i$ 对，于是
\[
f_2=f_0f_1+f_1f_0+f_2f_0=3,\qquad
f_3=f_0f_2+f_1f_1+f_2f_0=2+1+2=5.
\]
\textbf{按分类枚举才不会重、不会漏。}上面列出的 $5$ 种正是按“最外层括号内含 $0$ 对、$1$ 对、$2$ 对”分三类的代表。
\end{teach}

\begin{problem}{例 3：第二类斯特林数}
把 $5$ 个人分成 $3$ 个非空小组（组无编号），有多少种分法？
\end{problem}
\hint 就是 $S(5,3)$；用递推算。

\sol 列表计算：
\[
S(1,1)=1;\quad S(2,1)=1,\ S(2,2)=1;
\]
\[
S(3,1)=1,\ S(3,2)=2\cdot1+1=3,\ S(3,3)=1;
\]
\[
S(4,1)=1,\ S(4,2)=2\cdot3+1=7,\ S(4,3)=3\cdot1+3=6,\ S(4,4)=1;
\]
\[
S(5,1)=1,\ S(5,2)=2\cdot7+1=15,\ S(5,3)=3\cdot6+7=25.
\]
答案 $25$。若小组有编号，则 $3!\cdot25=150$。

\begin{problem}{例 4：递推的建立}
一条楼梯有 $n$ 级，每次可以上一级或两级，有多少种上楼方法？若每次可以上一、二或三级呢？
\end{problem}
\hint 按最后一步分类。

\sol 设 $f_n$ 为 $n$ 级楼梯的上法数，$f_0=1$。每次一或两级：$f_n=f_{n-1}+f_{n-2}$，即 $f_n=F_{n+1}$。每次一、二或三级：$f_n=f_{n-1}+f_{n-2}+f_{n-3}$，$f_0=1,f_1=1,f_2=2$，于是 $f_3=4,f_4=7,f_5=13$，这就是\textbf{ Tribonacci 数列}。

\begin{problem}{统一练习：卡特兰的又一副面孔}
$n$ 个元素依次入栈，求所有可能的出栈序列数。用 $n=3$ 验证答案为 $C_3$。
\end{problem}
\hint 把入栈记作左括号、出栈记作右括号；出栈数不能超过入栈数。

\sol 入栈序列记作 $1$，出栈记作 $-1$。任意前缀中 $1$ 的个数不少于 $-1$ 的个数，这正是括号匹配的条件，故答案为 $C_n$。$n=3$ 时有 $5$ 种：$123,132,213,231,321$。注意 $312$ 不可能（$3$ 第一个出栈要求 $1,2$ 已在栈中，那么出栈顺序只能是 $321$ 或 $23\ldots$，读者可自行核对）。

\begin{teach}[递推章的教学要点]
\begin{enumerate}
\item 每道题都要求学生先写“按最后一个动作分类”的那句话，再写递推式。
\item 递推式写完后必须用 $n=1,2,3$ 手工核对；对不上就重写分类。
\item 闭式公式（比奈公式、$\frac1{n+1}\binom{2n}{n}$）只作为“算得快”的工具，不要求会推导。
\end{enumerate}
\end{teach}

\chapter{抽屉原理与极端原理}
\begin{chapterintro}[章节导读]
前四章都在“数”，本章开始“证存在”。抽屉原理（鸽巢原理）说：东西比格子多，就有一格不止一个。它是组合数学里最朴素也最强大的工具，而它的两个常见变体——平均值原理与极端原理——分别处理“平均”和“最值”型的存在性论证。
\end{chapterintro}

\section{抽屉原理}
\begin{theorem}[抽屉原理]
把 $n+1$ 个物体放进 $n$ 个抽屉，至少有一个抽屉里有不少于 $2$ 个物体。
\end{theorem}
\begin{proof}[反证]
若每个抽屉至多 $1$ 个物体，则总数至多 $n$，与 $n+1$ 矛盾。
\end{proof}

\begin{theorem}[抽屉原理的加强形式]
把 $mn+1$ 个物体放进 $n$ 个抽屉，至少有一个抽屉里有不少于 $m+1$ 个物体。等价地：若 $n$ 个整数的平均值大于 $a$，则至少有一个整数大于 $a$。
\end{theorem}

\begin{teach}[抽屉原理的三步讲法]
\begin{enumerate}
\item \textbf{造抽屉}：把对象按某个性质分类，说明类别数少于对象数。这是最难的一步——题目从不告诉你抽屉是什么。
\item \textbf{数物体}：说明物体数比抽屉数多（或用加强形式给出 $m$）。
\item \textbf{得结论}：把“某类里有两个物体”翻译回原问题的语言。
\end{enumerate}
第一步没有通法，只能靠题目积累。常见的造抽屉方式有：按余数分类（模 $n$ 的剩余类）、按数值区间分块、按图的度数分层、按坐标奇偶性分类。
\end{teach}

\section{平均值原理}
\begin{theorem}[平均值原理]
若 $n$ 个实数的平均值为 $\bar a$，则至少有一个数不小于 $\bar a$，也至少有一个数不大于 $\bar a$。
\end{theorem}
\begin{proof}[反证]
若所有数都小于 $\bar a$，则平均数小于 $\bar a$，矛盾。大于的情形同理。
\end{proof}
平均值原理是抽屉原理的“加权版”：把 $n$ 个数看成物体、把“是否达到平均”看成两个抽屉。它在图论里最常用：度数平均值 $\dfrac{2m}{n}$ 保证存在度数不小于平均值的顶点。

\section{极端原理}
\begin{theorem}[极端原理]
在一组有限的对象中，取“最大”“最小”“最长”“最短”的那个，利用它的极端性导出矛盾或构造。这与第一部分第 1 章的“最小反例法”是同一件事。
\end{theorem}

\begin{tip}[什么时候用极端原理]
当题目问“证明存在一个……”，而直接构造很困难时，取一个极端对象（例如面积最小的反例、距离最远的一对点、字典序最小的解），它的极端性往往正好提供你需要的那条不等式。第二部分第 15 章的韦达跳跃，就是取“最小的解”然后造出更小的解来完成递降。
\end{tip}

\section{染色与覆盖}
染色是抽屉原理的几何版本：把对象染成少数几种颜色，利用“同色对象必有某种关系”得出结论。经典模式：
\begin{itemize}
\item \textbf{二染色}：把棋盘染成黑白两色，数两种颜色的格子数差。
\item \textbf{模 $k$ 染色}：把整数按模 $k$ 分类，用于处理周期性与整除。
\item \textbf{覆盖论证}：用若干个小图形覆盖大图形，比较面积，得出“至少有两个小图形重叠”。
\end{itemize}

\begin{pitfall}[“至少”与“至多”别写反]
“$n+1$ 个物体放进 $n$ 个抽屉”得到的是“\textbf{至少有一个}抽屉不少于 $2$ 个”，不是“所有抽屉都不少于 $2$ 个”，也不是“至多有一个”。写结论时把量词写全。
\end{pitfall}

\section{分层例题与解析}
\begin{problem}{例 1：基本抽屉}
证明：任意 $13$ 个人中，至少有两个人的生日在同一月份。
\end{problem}
\hint $12$ 个月是抽屉。

\sol $13$ 个人放进 $12$ 个月份，由抽屉原理至少有一个月份有两人以上。注意\textbf{抽屉是月份、物体是人}，不要说反。

\begin{problem}{例 2：加强形式}
证明：任意 $n+1$ 个正整数中，至少有两个数的差能被 $n$ 整除。
\end{problem}
\hint 按模 $n$ 的余数造抽屉。

\sol 两个数的差能被 $n$ 整除 $\iff$ 它们模 $n$ 同余。余数只有 $n$ 种，$n+1$ 个数必有同余的两个。

\begin{problem}{例 3：平均值原理}
证明：任意一群人，认识人数最多的人与最少的人，其认识人数之差不超过总人数减一。（把“认识”视为对称关系。）
\end{problem}
\hint 认识人数是整数，范围在 $0$ 与 $n-1$ 之间。

\sol 设共 $n$ 个人。每个人的认识人数是 $0$ 到 $n-1$ 之间的整数，但 $0$ 与 $n-1$ 不能同时出现（若有人认识所有人，则没有人谁也不认识）。故实际取值至多 $n-1$ 种，极差至多 $n-2$。这个例子用来说明：\textbf{先检查极端值能否同时取到}，是平均值与抽屉问题的常见陷阱。

\begin{problem}{例 4：极端原理}
设 $a_1,\ldots,a_n$ 是正整数，证明其中存在若干个（可以一个）的和能被 $n$ 整除。
\end{problem}
\hint 考虑前缀和 $S_k=a_1+\cdots+a_k$；若有两个前缀和同余则成功，否则 $n$ 个前缀和占据全部 $n$ 个余数类。

\sol 考虑前缀和 $S_0=0,S_1,\ldots,S_n$ 共 $n+1$ 个数。若其中两个模 $n$ 同余，设 $0\le i<j\le n$，则 $S_j-S_i=a_{i+1}+\cdots+a_j$ 被 $n$ 整除。若两两不同余，则 $S_1,\ldots,S_n$ 恰好取遍 $0,\ldots,n-1$ 每个余数一次，特别地某个 $S_k\equiv0\pmod n$，即前 $k$ 项和被 $n$ 整除。两种情况都得结论。

\begin{problem}{例 5：染色}
在 $8\times8$ 的棋盘上任意染黑 $13$ 个格子，证明必有两黑格同行或同列。
\end{problem}
\hint 用两次抽屉，或用行列计数。

\sol 若没有两黑格同行，则每行至多一个黑格，$8$ 行至多 $8$ 个黑格，矛盾。故必有同行。同列的论证完全一样。更强的结论：$13$ 个黑格必有某一行或某一列含不少于 $2$ 个黑格——这正是抽屉原理的直接应用。

\begin{problem}{统一练习：Ramsey 数 $R(3,3)=6$}
证明：任意 $6$ 个人中，必有 $3$ 人两两认识，或 $3$ 人两两不认识。
\end{problem}
\hint 固定一个人 $A$，看 $A$ 与其余 $5$ 人的关系；用抽屉原理得 $A$ 至少认识或不认识 $3$ 人。

\sol 固定 $A$。$A$ 与其余 $5$ 人每人要么认识要么不认识，由抽屉原理，$A$ 至少对其中 $3$ 人有相同的关系。设 $A$ 认识 $B,C,D$（不认识的情形对称）。若 $B,C,D$ 中有两人互相认识，这两人与 $A$ 构成三人两两认识；否则 $B,C,D$ 三人两两不认识。两种情况都成立。

这个证明是 Ramsey 理论的入门例子。它展示了抽屉原理的深层用法：\textbf{先固定一个对象，把“关系”变成“抽屉”}。

\begin{teach}[本章的离堂检查]
让学生独立完成“证明任意 $n+2$ 个正整数中必有两个数的差能被 $n$ 或能被 $n+1$ 整除”这类变体（提示：造 $n+1$ 个抽屉而不是 $n$ 个，或者用加强形式）。做不出就回到例 2 重讲“按余数造抽屉”。
\end{teach}

\chapter{格路计数与反射原理}
\begin{chapterintro}[章节导读]
格路是组合数学里最具体的对象之一：从 $(0,0)$ 出发，每步向右或向上走一格。它的计数看似简单，却能把卡特兰数、选票问题、反射原理串联成一个整体，并直接对应到第二部分数论里的路径分解思想。本章把这条线走完。
\end{chapterintro}

\section{单调格路的基本计数}
\begin{definition}[单调格路]
从 $(0,0)$ 到 $(m,n)$、每步只能向右 $(+1,0)$ 或向上 $(0,+1)$ 的路径叫\textbf{单调格路}。
\end{definition}
\begin{theorem}[单调格路计数]
从 $(0,0)$ 到 $(m,n)$ 的单调格路数为
\[
\binom{m+n}{m}.
\]
\end{theorem}
\begin{proof}
一条路径由 $m+n$ 步组成，其中恰有 $m$ 步向右。选哪 $m$ 步向右有 $\dbinom{m+n}{m}$ 种，其余步自然向上，与路径一一对应。
\end{proof}

\begin{tip}[格路是“序列”的几何化身]
把右步记作 $0$、上步记作 $1$，一条格路就是一个长度为 $m+n$、含 $m$ 个 $0$ 的 $01$ 序列。所以格路计数与“选位置”是同一件事。这个对应在第五部分处理多项式系数时会再次出现。
\end{tip}

\section{不越过对角线的格路与卡特兰数}
\begin{theorem}[不越过对角线的格路]
从 $(0,0)$ 到 $(n,n)$、每步向右或向上、且\textbf{不越过直线 $y=x$}（允许接触）的单调格路数为卡特兰数
\[
C_n=\frac1{n+1}\binom{2n}{n}.
\]
\end{theorem}
\begin{proof}[反射原理]
先数全部路径 $\dbinom{2n}{n}$，再减掉越线的坏路径。一条路径越线，意味着它第一次到达直线 $y=x+1$（即第一次上步比右步多 $1$）。把这一步之后的路径\textbf{关于直线 $y=x+1$ 作镜像}：右步变上步、上步变右步。镜像后路径从 $(0,0)$ 到 $(n-1,n+1)$。

反过来，任意一条从 $(0,0)$ 到 $(n-1,n+1)$ 的路径必然穿过 $y=x+1$，取第一次穿过的位置之后作同样的镜像，就得到一条从 $(0,0)$ 到 $(n,n)$ 的坏路径。于是坏路径与“到 $(n-1,n+1)$ 的路径”一一对应，共 $\dbinom{2n}{n-1}$ 条。故
\[
C_n=\binom{2n}{n}-\binom{2n}{n-1}
=\binom{2n}{n}-\frac n{n+1}\binom{2n}{n}
=\frac1{n+1}\binom{2n}{n}.
\]
\end{proof}

\begin{teach}[反射原理的两步核对]
\begin{enumerate}
\item \textbf{双射}：坏路径与“到 $(n-1,n+1)$ 的路径”之间的映射必须是双射（互为逆）。让学生对 $n=2$ 枚举验证。
\item \textbf{边界}：“不越过”允许接触对角线；“不接触”（严格在下方）是另一个问题，答案为 $C_{n-1}$。两种题面的区别务必讲清。
\end{enumerate}
\end{teach}

\section{选票问题}
\begin{theorem}[选票问题]
一次选举中候选人甲得 $p$ 票、乙得 $q$ 票，$p>q$。若计票过程中甲的票数\textbf{始终严格多于}乙，则计票顺序数为
\[
\frac{p-q}{p+q}\binom{p+q}{p}.
\]
\end{theorem}
\begin{proof}[反射原理]
把每一票记作一步：甲票记右步、乙票记上步。计票顺序就是一条从 $(0,0)$ 到 $(p,q)$ 的单调格路，“甲始终严格多于乙”即路径始终不触碰直线 $y=x$（只能在其下方）。

坏路径是那些触碰 $y=x$ 的。取第一次触碰 $y=x$ 的位置，把之前的路径关于 $y=x$ 镜像，得到一条从 $(0,0)$ 到 $(q,p)$ 的路径；这是一个双射。故好路径数为
\[
\binom{p+q}{p}-\binom{p+q}{q}
=\binom{p+q}{p}-\frac{q}{p}\binom{p+q}{p}
=\frac{p-q}{p+q}\binom{p+q}{p}.
\]
\end{proof}

\begin{tip}[选票问题是卡特兰的一般化]
取 $p=q+1=n+1$、$q=n$，得 $\dfrac{1}{2n+1}\dbinom{2n+1}{n}=C_n$，正是第 4 章的卡特兰数。这个特例说明“不越过”与“始终严格多于”两种边界条件对应不同的题面，用同一个反射原理处理。
\end{tip}

\section{带障碍的格路}
若格路上有若干禁入点，可以用\textbf{容斥原理}：总路径数减去经过禁入点的路径数。若禁入点可以按偏序排列（例如从左下到右上），经过它们的路径要交错的并集，用容斥展开。

\begin{problem}{统一练习：一个禁入点}
从 $(0,0)$ 到 $(4,4)$ 的单调格路中，不经过 $(2,2)$ 的有多少条？
\end{problem}
\hint 总数减去经过 $(2,2)$ 的；经过 $(2,2)$ 的路径 = 到 $(2,2)$ 的路径数 $\times$ 从 $(2,2)$ 到 $(4,4)$ 的路径数。

\sol 总数 $\dbinom84=70$。经过 $(2,2)$ 的路径数为 $\dbinom42\cdot\dbinom42=6\cdot6=36$。故答案为 $70-36=34$。

\section{分层例题与解析}
\begin{problem}{例 1：基本格路计数}
从 $(0,0)$ 到 $(5,3)$ 有多少条单调格路？
\end{problem}
\hint 共 $8$ 步，选 $5$ 步向右。

\sol $\dbinom85=56$。

\begin{problem}{例 2：反射原理}
求从 $(0,0)$ 到 $(4,4)$ 不越过对角线 $y=x$ 的单调格路数。
\end{problem}
\hint 就是 $C_4$。

\sol $C_4=\dfrac15\dbinom84=\dfrac{70}{5}=14$。也可以用反射原理直接算：$\dbinom84-\dbinom8{3}=70-56=14$。

\begin{problem}{例 3：选票问题}
一次选举甲得 $5$ 票、乙得 $3$ 票，求计票过程中甲始终严格多于乙的顺序数。
\end{problem}
\hint 代入公式。

\sol $\dfrac{5-3}{5+3}\dbinom85=\dfrac28\cdot56=14$。可以枚举验证：$5$ 个 $A$、$3$ 个 $B$ 的序列中，每个前缀里 $A$ 的个数都大于 $B$ 的个数，共 $14$ 个。

\begin{problem}{例 4：容斥处理两个禁入点}
从 $(0,0)$ 到 $(5,5)$ 的单调格路，不经过 $(2,2)$ 与 $(3,3)$ 的有多少条？
\end{problem}
\hint 用容斥；经过 $(2,2)$ 的路径数 = 到 $(2,2)$ 的路径数 $\times$ 从 $(2,2)$ 到 $(5,5)$ 的路径数。

\sol 总数 $\dbinom{10}5=252$。记 $A$ 为经过 $(2,2)$ 的路径集，$B$ 为经过 $(3,3)$ 的路径集。
\[
|A|=\binom42\binom{6}{3}=6\cdot20=120,\qquad
|B|=\binom63\binom42=20\cdot6=120.
\]
交集 $A\cap B$：路径同时经过 $(2,2)$ 与 $(3,3)$，由单调性必先经过 $(2,2)$ 再经过 $(3,3)$，路径数为
\[
\binom42\binom21\binom42=6\cdot2\cdot6=72.
\]
故
\[
|A\cup B|=120+120-72=168,\qquad 252-168=84.
\]
答案为 $84$。

\begin{pitfall}[单调性给出禁入点的天然顺序]
在单调格路上，若 $(a,b)$ 与 $(c,d)$ 满足 $a\le c$、$b\le d$，则任何经过这两点的路径必先经过 $(a,b)$。所以容斥里的交集不需要再分顺序——这是格路容斥比一般容斥简单的地方。但\textbf{若两点不可比较}（例如 $(2,3)$ 与 $(3,2)$），一条路径至多经过其中一个，交集为空。
\end{pitfall}

\chapter{组合构造与综合例题}
\begin{chapterintro}[章节导读]
计数回答“有多少”，构造回答“能不能造一个”。竞赛里的构造题往往不需要精巧的想法，只需要把条件逐条满足；真正难的是“证明你造的确实合法”。本章给出构造题的通用流程，并用不变量、图论计数与综合例题收束第三部分。
\end{chapterintro}

\section{构造题的标准流程}
\begin{enumerate}
\item \textbf{写下全部条件}：把题面转写成带量词的命题，明确“任意”与“存在”。
\item \textbf{找必要条件}：先推出一些“必须成立”的简单条件（奇偶性、计数和、度数），它们既是构造的指引，也是“为什么某些情形无解”的答案。
\item \textbf{按必要条件构造}：多数题在满足必要条件后就能构造；常用套路有：贪心、归纳（从 $n-1$ 到 $n$）、对称构造、分解成独立小块。
\item \textbf{验证合法性}：逐条核对条件，并检查边界（$n=1$、空集、相等的情形）。
\end{enumerate}

\begin{pitfall}[“看起来合法”不算证明]
构造题最常见的失分是“给了构造但没验证”。评分时要求：每一组输出都要能说明它满足每一条约束。第四部分博弈论第 2 章“下模仿棋”是一个范例——构造只是“镜像”，但必须论证“镜像永远合法”。
\end{pitfall}

\section{不变量}
\begin{definition}[不变量]
在操作过程中保持不变的量（或保持某种单调性的量）叫\textbf{不变量}。若能找到一个不变量在初始状态与目标状态取值不同，则目标不可达。
\end{definition}
常见不变量：奇偶性、模 $k$ 余数、总和、异或和、图的度数奇偶性。第二部分第 4 章“环上的固定步长”用的就是“编号模 $d$ 的余数不变”。

\begin{problem}{例：倒水问题}
有容量 $3$ 升与 $5$ 升的两个空壶，能量出 $4$ 升水吗？
\end{problem}
\hint 每次操作使某壶装满、倒空或倒入另一壶至满；不变量是“总水量”与“各壶水量都是整数”。

\sol 可以：$(0,0)\to(0,5)\to(3,2)\to(0,2)\to(2,0)\to(2,5)\to(3,4)$，此时 $5$ 升壶中剩 $4$ 升。一般结论：容量 $a,b$（互质）能量出 $1$ 到 $\max(a,b)$ 之间任意整数升——这正是裴蜀定理的构造性版本，第二部分第 1 章会给严格证明。

\section{图上的计数}
\begin{theorem}[握手引理]
任意无向图中，所有顶点的度数之和等于边数的两倍：
\[
\sum_v\deg(v)=2|E|.
\]
因此度数为奇数的顶点个数必为偶数。
\end{theorem}
\begin{proof}
每条边恰好为两个端点的度数各贡献 $1$，求和时每条边被数两次。
\end{proof}

\begin{theorem}[树的边数]
$n$ 个顶点的树恰有 $n-1$ 条边。
\end{theorem}
\begin{proof}[归纳]
$n=1$ 时 $0$ 条边。$n\ge2$ 时树必有叶子（度数为 $1$ 的顶点，否则沿边走会形成环或无限路径，与有限且无环矛盾）。删去一片叶子及其边，剩下一棵 $n-1$ 个顶点的树，由归纳假设有 $n-2$ 条边，加回删除的边得 $n-1$ 条。
\end{proof}

\begin{tip}[欧拉公式]
连通平面图满足 $V-E+F=2$（$F$ 含外部面）。它是平面图计数的基本工具，也能推出“平面简单图边数不超过 $3V-6$”。这部分内容在算法竞赛中出现频率不高，列在此处作为图论计数的入口。
\end{tip}

\section{综合例题}
\begin{problem}{综合 1：二项式系数求和}
计算 $\sum_{k=0}^{n}k\binom nk$ 与 $\sum_{k=0}^{n}\dfrac1{k+1}\binom nk$。
\end{problem}
\hint 第一问用求导；第二问积分，或注意到 $\dfrac1{k+1}\binom nk=\dfrac1{n+1}\dbinom{n+1}{k+1}$。

\sol 第一问 $n2^{n-1}$。第二问：
\[
\sum_{k=0}^{n}\frac1{k+1}\binom nk
=\frac1{n+1}\sum_{k=0}^{n}\binom{n+1}{k+1}
=\frac1{n+1}\left(2^{n+1}-1\right).
\]
最后一步用了 $\sum_{j}\dbinom{n+1}{j}=2^{n+1}$ 并去掉 $j=0$ 项。

\begin{problem}{综合 2：容斥 + 递推}
用 $1\times2$ 与 $2\times1$ 的骨牌完全覆盖 $2\times n$ 的棋盘，有多少种方法？
\end{problem}
\hint 按最左边一列怎么被覆盖分类。

\sol 设 $f_n$ 为覆盖 $2\times n$ 的方法数。最左一列要么被一个竖放的 $2\times1$ 覆盖（剩下 $2\times(n-1)$，$f_{n-1}$ 种），要么被两个横放的 $1\times2$ 覆盖（剩下 $2\times(n-2)$，$f_{n-2}$ 种）。故
\[
f_n=f_{n-1}+f_{n-2},\qquad f_0=1,\ f_1=1,
\]
即 $f_n=F_{n+1}$。

\begin{problem}{综合 3：抽屉原理与奇偶拆分}
证明：任意 $n+1$ 个不超过 $2n$ 的正整数中，必有两个数，其中一个整除另一个。
\end{problem}
\hint 把每个数写成“奇数部分 $\times$ $2$ 的幂”；$1$ 到 $2n$ 中只有 $n$ 个不同的奇数部分。

\sol 任取一个正整数 $a\le2n$，它可唯一写成
\[
a=2^{t}\cdot u,\qquad u\ \text{为奇数}.
\]
在 $1,\ldots,2n$ 中，奇数只有 $1,3,5,\ldots,2n-1$，共 $n$ 个，所以可能的 $u$ 只有 $n$ 种。现有 $n+1$ 个数，由抽屉原理，必有两个数 $a=2^{t_1}u$、$b=2^{t_2}u$ 具有相同的奇数部分 $u$。设 $t_1\le t_2$，则 $a\mid b$。得证。

这个证明是“\textbf{先找正确的分类，再套抽屉}”的典范：分类标准不是数的大小，而是\textbf{剥离 $2$ 的幂以后剩下的奇数核心}。

\begin{teach}[综合 3 讲评要点]
让学生先自己造抽屉。常见的错误尝试有：按“$\le n$ 与 $>n$”分两类（不奏效，因为两类内部没有整除关系）、按“是否偶数”分（同样不奏效）。等到他们发现这两条路都走不通，再给出奇数核心的分类。\textbf{抽屉造不出来，往往是因为分类标准选错了维度。}
\end{teach}

\begin{problem}{综合 4：计数与验证并重}
证明：任意 $n\ge2$ 个人中，若每两个人要么认识要么不认识，则“认识人数为奇数”的人数是偶数。
\end{problem}
\hint 握手引理。

\sol 由握手引理，$\sum_v\deg(v)=2|E|$ 是偶数。设度数为奇数的顶点集合为 $O$、度数为偶数的为 $E_v$，则
\[
\sum_{v\in O}\deg(v)=2|E|-\sum_{v\in E_v}\deg(v)
\]
是偶数减偶数，仍为偶数。而奇数个奇数相加为奇数、偶数个奇数相加为偶数，故 $|O|$ 必为偶数。

\begin{teach}[第三部分的结课检查]
让学生独立完成三道小题：
\begin{enumerate}
\item 证明 $\dbinom{2n}{n}$ 是偶数（$n\ge1$）——提示用 $\dbinom{2n}{n}=\dbinom{2n-1}{n-1}+\dbinom{2n-1}{n}$；
\item 求 $n$ 对括号的正确匹配数，并用反射原理从格路角度再证一次；
\item 判断命题“任意 $n+2$ 个不超过 $2n$ 的正整数中，必有两个数的差能被 $n$ 整除”的真假，并证明或给出反例。
\end{enumerate}
三题全对即可进入第四部分。第三题需要谨慎：先在 $n=2$、$n=3$ 上试探，确认命题真假再下笔。
\end{teach}


% ============================================================================
%  第四部分  博弈论与概率论
% ============================================================================
\BeginPart{5}{4}{IV}{第四部分}{博弈论与概率论}
% !TeX program = xelatex
% 本文件由 assemble.py 自 OI-Docs 原稿抽取生成，请勿手改（会被覆盖）。
% 源文件：Game Theory/game_theory_teacher.tex
% 源行号：143--1641（已删 1638 的 \appendix）
% 说明：已删去独立成册所需的导言区与页码重置；章节正文与交叉引用原样保留。


\chapter{博弈问题简介}
\section{一些约定}
在进入具体内容之前，先说明整本书的立场：我们研究的是\textbf{最优博弈}，即假定双方都足够聪明、每一步都采取最优策略、绝不犯错，然后问：从某个初始状态出发，胜负能否被预先确定？这里的“必胜”是严格的数学概念——存在一种策略，使得无论对手如何应对，按此策略行动的一方最终一定获胜。把希望寄托在对手失误上的“策略”不在讨论范围内。

为什么“最优策略”是良定义的呢？关键在于游戏的\textbf{确定性}与\textbf{有限性}：局面向局面的转移完全由玩家的决策决定，且游戏必然在有限步内结束。于是可以把每个合法局面看成一个节点、把一步合法操作看成一条边，得到一张有向图，胜负就能从“终局”倒着递推回初始局面。全篇几乎所有结论，本质上都是在这张状态图上做文章。出于同样的原因，以下约定非常方便：
\begin{itemize}
\item “$A$”“$B$”分别表示初始状态下的先后手，它们一旦定下就不变。
\item “先手”“后手”分别表示\textbf{当前状态}下的先后手：谁下一步行动，谁就是当前先手。推导中常出现“某方把局面做成 P 态”的说法，意思就是把当前先手让给对方，让对方面对一个先手必败的局面。叙述具体问题时通常用固定的 $A$、$B$，而推导胜负时通常用随局面变化的“先手”“后手”。
\item 对于公平组合游戏，“P 态”“N 态”分别表示后手必胜态和先手必胜态（Previous player wins / Next player wins）。P、N 的判定见 1.3 节的递推刻画与第 3 章的 PN 态 DP。
\end{itemize}

\begin{teach}[导入：从“必胜”这个词开始]
黑板上写“$3$ 堆各 $1$ 颗，轮流取，取到最后一颗者胜”，请学生口述必胜策略；再写“各 $2$ 颗”。学生会自然说出“不管对面怎么走我都……”，顺势指出这句话就是最优策略的数学定义雏形。强调两点排除：不研究“对手会犯错”，不研究“胜率”；先把这两件事排除，能省掉大量无效讨论。
\end{teach}
\begin{pitfall}[两套称呼不要混用]
“$A$／$B$”是初始的、固定的；“先手／后手”是当前的、随局面翻转的。说“$A$ 把局面做成 P 态”，意思是 $A$ 走完后，轮到行动的 $B$（此时 $B$ 是当前先手）面对必败局面。推导时切换称呼是最常见的笔误来源。
\end{pitfall}

\section{组合游戏}
OI 中研究的博弈问题大部分是\textbf{组合游戏}。这类游戏之所以“好研究”，是因为它的五个特征恰好保证了上一节所说的状态图方法可行、“最优策略”良定义且理论上可以求出。组合游戏的特征：
\begin{itemize}
\item \textbf{完全信息}：任何玩家在所有时候都知道游戏的全部状态。\\
反例：斗地主（玩家看不到其他玩家的手牌）。完全信息使得决策只取决于局面本身，而不取决于“我猜对方有什么”；一旦掺入隐藏信息，策略就与信念、试探和误导挂钩，局面不再能用一张简单的状态图刻画，博弈的难度会陡增（例题 [CF98E] Help Shrek and Donkey 就是一个在少量隐藏信息下仍可求解的例子）。
\item \textbf{无随机性}：游戏进程完全由玩家的决策决定。\\
反例：大富翁（骰子的点数是随机的）。没有随机因素，才谈得上“必胜策略”；一旦允许随机，通常只能讨论“胜率最大的策略”，这属于第 7 章纳什均衡一节的混合策略范畴。
\item \textbf{轮流行动}：玩家严格地轮流进行操作，同一时刻至多一名玩家做决策。这与“所有人同时决策”的静态博弈（见 1.4 节）有本质区别：前者有“状态”和“先后手”，后者没有。
\item \textbf{非合作（零和）}：玩家之间不能合作，通常是两个玩家，一方所得即另一方所失。收益之和为常数使得“我要赢”与“让对手输”完全等价，博弈成为纯粹的对立；一旦出现三人以上或非零和的局面，玩家之间还可能结盟或互相坑害，OI 中几乎不出现。
\item \textbf{确定性的结果}：游戏一定会结束，并根据最终状态产生一个明确的输赢或结果（比如本局得分）。“一定会结束”不是白说的：若允许无限进行下去，就要额外引入“平局”这一第三种结果，判定会复杂得多（例题 [省选联考 2023] 过河卒就是典型例子）。
\end{itemize}

另外要留意“结果”的两种常见形式：一是\textbf{无法行动者负}（全篇绝大多数游戏），二是\textbf{无法行动者胜}（反 Nim 游戏等）。两者只差一个字，结论却可能截然不同；拿到题先看清规则，再决定套用哪套理论。

\section{公平组合游戏}
在组合游戏中，还有一类性质更强的游戏，称为\textbf{公平组合游戏}。

\begin{definition}[公平组合游戏]
一个组合游戏是公平的，若对于同一个局面，无论轮到哪一个玩家移动，其决策空间都是相同的。
\end{definition}

通俗地说，规则不能“认人”：同一局面下甲能走的每一步，乙也一定能走。拿石子、走格子这类游戏天然是公平的；反过来，国际象棋、围棋这类“各自的棋子各自走”的游戏不是公平的——同一个局面下，轮到白方与轮到黑方时能操作的棋子不同。这类游戏双方的地位不对等，需要超现实数这样更精细的工具来刻画（见第 6 章）。

正因为公平，同一局面下两名玩家处于完全对称的位置，于是每个状态只可能属于两类。

\begin{definition}[P 态与 N 态]\label{def:pn}
\textbf{P 态}（后手必胜态）：当前先手无论怎么走，当前后手都有应对策略，使先手最终必败；\\
\textbf{N 态}（先手必胜态）：当前先手存在一步合法走法，使走完后的局面成为 P 态，从而立于不败之地。
\end{definition}

\begin{theorem}[递推刻画]\label{thm:recursion}
在“游戏一定在有限步内结束”的前提下，每个状态恰属于 P、N 之一，且有：没有合法走法的局面是 P 态（当前先手直接判负）；一般地，局面 $u$ 是 N 态当且仅当存在后继 $v\in\Np(u)$ 是 P 态；$u$ 是 P 态当且仅当它的所有后继都是 N 态。
\end{theorem}
\begin{proof}
从终局反向归纳：无后继的状态按定义是 P 态。对一般状态 $u$，若存在 P 态后继，先手走到它即立于不败之地，故 $u$ 是 N 态；若所有后继都是 N 态，则先手无论怎么走都把必胜局面交给对方，故 $u$ 是 P 态。由于每一步转移都严格消耗“剩余步数”，归纳是良基的，每个状态恰好获得一个标签。
\end{proof}

这个递推刻画是第 3 章 PN 态 DP 的理论基础，也是全篇最核心的判定工具，后面几乎所有结论都可以看成它的“手算加速”。

前面已经提到，对于公平组合游戏，状态只分为 P 态和 N 态两类。那么对于非公平组合游戏呢？可以发现状态分为四类：先手必胜、后手必胜、玩家甲必胜、玩家乙必胜——因为“对谁有利”取决于轮到谁走，还取决于“棋子的归属”，两者独立组合出四种结果。

对于感兴趣的选手，超现实数是研究非公平组合游戏的一个有力工具。接下来要讨论的 Sprague--Grundy 理论正是它在“公平游戏”这个特殊情况下的投影；完整的图景要等第 6 章再展开。

\begin{teach}[P/N 递推刻画的课堂三问]
第一问：终局（无路可走）标 P 还是 N？（标 P：轮到的人直接输。）第二问：一个点是 N 态，需要“所有后继是 P”还是“存在后继是 P”？（存在；量词混淆是最高频错误。）第三问：为什么“有限步内结束”不可缺少？（否则无限链上的归纳不闭合，可能出现既非 P 也非 N 的状态。）课堂追问：若规则允许平局，这套二值刻画还够用吗？
\end{teach}
\section{静态博弈}
OI 中研究的博弈问题除了组合游戏外，还有一类是\textbf{静态博弈}问题。
\begin{itemize}
\item \textbf{静态}：所有玩家同时作出决策，最后结算。显然静态博弈问题是无所谓状态的——没有“轮到谁”的先后，也就没有状态图可言。
\end{itemize}
最典型的例子是石头剪刀布：双方同时出拳，同时出完立即结算，不存在“看到对方出拳再应对”。OI 中一些“双方各构造一个答案然后互相对抗”的题目也属于静态博弈。因为缺少先后手结构，静态博弈不能套用 P／N 态，需要的是第 7 章“纳什均衡”那一套工具。

每个静态博弈问题都能表示成一个超长方体。具体地，考虑一个 $n$ 名玩家的静态博弈，第 $i$ 名玩家有 $k_i$ 种决策，则这个静态博弈可以表示成一个 $k_1\times k_2\times\cdots\times k_n$ 的超长方体，超长方体的每个单位立方体里写着一个 $n$ 维向量，表示每名玩家的收益。以最常见的两人博弈为例，它就是一张 $k_1\times k_2$ 的\textbf{收益矩阵}：行对应玩家 $A$ 的决策，列对应玩家 $B$ 的决策，格子 $(i,j)$ 里写两人各自拿到的收益。下文讨论纳什均衡时说的“收益矩阵”就是指这张表。

在静态博弈中，还有一类性质更强的博弈，称为\textbf{匿名静态博弈}。
\begin{definition}[匿名静态博弈]
每名玩家的决策空间相同，且对于任意一种可能的博弈结果，决策相同的玩家的收益相同。
\end{definition}
也就是说，玩家之间除了“自己选了哪个决策”之外没有任何身份区别：收益只取决于“我的决策”和“其余人各自选的决策”，其余人之间可以任意交换身份而不影响我的收益。例如“每人报一个 $1\sim10$ 的整数，报出最众数的人赢”就是匿名的；“走私者与检察官”这类角色分工明确的问题则不是匿名的。

\chapter{充要条件：模仿棋、策略窃取与设计定式}
本节介绍判定“初始状态是 P 态还是 N 态”最朴素的手段：直接构造必胜（或必败）的论证。要证明一个状态是 N 态，只需构造先手的第一步乃至整局策略；要证明它是 P 态，则要说明无论先手怎么走，后手都有办法应对到底。常见套路有三板斧，分别对应下面三个小节：
\begin{itemize}
\item \textbf{下模仿棋}：利用局面的对称性，后手（或先手）“照抄”对手的操作，始终与对手保持某种对称；
\item \textbf{策略窃取}：非构造性地证明“后手不可能有必胜策略”，从而先手必胜；
\item \textbf{设计定式}：预先设计一组对应关系（如匹配、配对），把对手的任意一步都映射成自己的合法回应，并把整个博弈化归到已经会做的问题上（通常是 Nim 游戏）。
\end{itemize}
很多题目表面千差万别，最终的解法思路都会落回这三句话上，下面逐个展开。

\section{下模仿棋}
\begin{problem}{经典问题：保龄球瓶}
地上放着一排 $n$ 个保龄球瓶，一击可以击倒一个瓶子或相邻两个瓶子，不能行动者负。判定初始状态。
\end{problem}

\sol \textbf{初始状态是 N 态。}当 $n$ 是偶数时，$A$ 先击倒正中间的两个瓶子；当 $n$ 是奇数时，$A$ 先击倒正中间的一个瓶子。以后 $A$ 不断操作 $B$ 操作的位置的对称位置即可。

\paragraph{为什么这能保证 $A$ 必胜.} 关键在于第一步之后场面的对称性：剩下的瓶子被分成左右两段，两段长度相等（$n$ 为偶数时各 $n/2$ 个，$n$ 为奇数时各 $(n-1)/2$ 个），关于正中间完全对称。于是 $A$ 采取镜像策略：$B$ 在某一段击倒了某个瓶子或相邻两个瓶子，$A$ 就到另一段的对称位置击倒对称的一组。注意这里的操作“击倒一个或相邻两个瓶子”在复制到对称位置后仍是合法操作，且由于两侧的消耗永远同步，只要 $B$ 能击倒的瓶子还存在，其对称位置的那组瓶子也一定还立着——所以 $A$ 的每一步回应都合法。

\paragraph{严格收尾.} 每步至少击倒一个瓶子，游戏必在有限步内结束；又因为 $A$ 的行动是对 $B$ 的合法行动的回应，$B$ 能走一步，$A$ 就永远能跟着走一步，所以先走投无路的一定是 $B$。用博弈论的话说，$A$ 的“能动次数”不少于 $B$，而失败者是轮到自己却不能动的人，故 $A$ 必胜，初始状态是 N 态。$n=1,2$ 的极端情形也被上述方案覆盖：第一步直接击倒全部瓶子，游戏立即结束。

从本题可以提炼出下模仿棋的适用条件：\textbf{操作对象必须能两两配对，且“对称回应”永远合法}。保龄球问题能镜像，正是因为击倒正中间后，左右两侧的剩余部分互不干扰；若操作会跨越中线（例如可以一次击倒任意一段连续瓶子），镜像策略就会失效。

\section{策略窃取}
\begin{problem}{经典问题：擦因数}
黑板上写着 $1\ldots n$ 的每个数，两人博弈，每人要选择黑板上的一个数，擦掉其所有因数，不能行动者负。判定初始状态。
\end{problem}

\sol \textbf{初始状态是 N 态。}$n$ 很大时我们根本不知道必胜策略的第一步该走什么，甚至无从下手验证。策略窃取的思想是\textbf{非构造性}地证明先手必胜：假设后手存在必胜策略，再说明先手可以“偷”来这个策略为己所用，从而导出矛盾。

先给出通用模板：如果先手能找到一步“走了等于没走”的棋（不改变局面的实质、只消耗自己一个回合），而“后手必胜”的假设成立，那么先手可以先走这步不亏的棋，然后完整照搬假设中的后手必胜策略——这样一来先手必胜，与“后手必胜”矛盾，因此后手根本不可能有必胜策略，初始状态必为 N 态。换言之，“存在一步不亏的棋”就足以证明 N 态，尽管我们仍然不知道真正的必胜策略长什么样。

本题中“选 $1$”就是这步不亏的棋：$1$ 的因数只有 $1$ 自己，选 $1$ 只会把 $1$ 从黑板上擦掉，对集合 $\{2,3,\ldots,n\}$ 没有任何影响。于是分两种情况：
\begin{itemize}
\item 如果 $\{2,3,\ldots,n\}$ 是 P 态，则 $A$ 选择 $1$ 即可——$A$ 相当于白走一步后，把一个 P 态交给 $B$，$B$ 成为当前先手必败。
\item 如果 $\{2,3,\ldots,n\}$ 是 N 态，设该状态下先手要取胜可以选 $x$，则 $A$ 第一手就选 $x$。因为“走 $x$ 能赢”只与走完后的局面有关，与谁走无关，所以 $A$ 抢先走掉 $x$ 后，$B$ 面对的局面恰是“先手已走完必胜第一步的 $\{2,3,\ldots,n\}$”，即一个 P 态，$B$ 必败。这里不必担心 $x$ 的因数包含 $1$：即使 $A$ 选 $x$ 也擦掉了 $1$，局面依然对应“$\{2,3,\ldots,n\}$ 中先手走 $x$ 之后”的同一 P 态，论证不受影响。
\end{itemize}
两种情况都推出 $A$ 必胜，故初始状态是 N 态。注意整个证明自始至终没有给出具体的必胜策略，这正是策略窃取的典型特征：\textbf{只能证“存在”，不能给“方案”}。因此它在 OI 中适合用来推导“判定胜负”类的结论；一旦题目要求输出必胜策略（如例题 [CF1147F] ZigZag Game），通常还需要结合其他构造手段。

\section{设计定式}
\begin{problem}{经典问题：台阶游戏（阶梯 Nim）}
有 $n$ 级台阶，第 $i$ 级台阶上有 $a_i$ 个石子。每次可以选择一级台阶（设为第 $i$ 级台阶），把第 $i$ 级台阶上至少一个石子移到第 $i-1$ 级台阶上（当 $i=1$ 时，移到地面上），不能行动者负。判定初始状态。
\end{problem}

\sol 这个游戏看起来比 Nim 游戏多出“向下搬运”的结构，其实两者只隔一层窗户纸：\textbf{只有奇数级台阶上的石子才真正参与博弈，偶数级台阶只是临时堆放区（垃圾桶）}。下面用“不变量”的语言把两个视角统一起来，设 $S$ 为奇数级台阶石子数的异或和。

\paragraph{后手视角.} 后手 $B$ 的策略是：每走一步都把局面恢复成 $S=0$。按 $A$ 的行动分类：
\begin{itemize}
\item $A$ 动了偶数级台阶 $i$：石子被搬到奇数级 $i-1$，$S$ 被破坏。$B$ 不慌不忙地把同一批石子再往下搬一级（从 $i-1$ 搬到 $i-2$；若 $i=2$ 则直接搬下地面），于是奇数级台阶上的石子分布原样恢复，$S$ 回到 $0$。这一来一回等于 $A$ 往垃圾桶里扔了一趟石子，对真正的局面毫无影响。
\item $A$ 动了奇数级台阶 $i$：石子被搬到偶数级 $i-1$（$i=1$ 时直接落地），这等价于从 Nim 的某一堆取走若干石子，$S$ 变成某个非零值。$B$ 把奇数级台阶看作 Nim 游戏的堆，用 Nim 的最优策略走一步（把一个奇数级台阶上的若干石子搬下去），把 $S$ 重新归零。
\end{itemize}
两类回应都合法且把 $S=0$ 交还给 $A$。归纳可知：只要初始 $S=0$，$B$ 就永远不会无路可走；游戏必然结束（石子总数不断减少），而 $A$ 每次行动后 $B$ 都能行动，因此最后无法行动的一定是 $A$。

\paragraph{先手视角.} 反之，如果 $S\ne0$，先手 $A$ 第一手就按 Nim 游戏的最优策略，把某个奇数级台阶上的石子搬到合适的数量使 $S=0$，然后转用上面“后手”的维持策略。

\begin{theorem}[阶梯 Nim 判定]\label{thm:stairnim}
台阶游戏中，先手必胜当且仅当奇数级台阶石子数的异或和不为 $0$。一句话总结：偶数级台阶是垃圾桶，只看奇数级台阶，胜负判定与 Nim 游戏完全相同。
\end{theorem}

\begin{pitfall}[只有奇数级参与异或]
初学者最常见的错误是把所有台阶（包括偶数级）都参与异或，务必记住只有奇数级有效。例题 [SDOI2019] 移动金币和 Gra 最后都会化归到这个结论上。
\end{pitfall}

\begin{problem}{经典问题：开箱子}
有 $n$ 箱石子，第 $i$ 箱石子有 $a_i$ 个。初始时每个箱子都没开启。有两种操作：
\begin{itemize}
\item 打开至少一个箱子。
\item 选择一个已经打开的箱子，取走至少一个石子。
\end{itemize}
不能行动者负。判定初始状态。
\end{problem}

\sol \paragraph{后手视角.} 先建立两个基本观察。第一，所有已打开的箱子整体构成一个 Nim 游戏：已打开的箱子随时可以取石，未打开的箱子只是“待解锁的储备”，取石规则与 Nim 完全一致。第二，后手 $B$ 给自己设定的不变量是：自己每步行动结束后，所有已打开箱子的石子数的异或和为 $0$。

在“不存在异或和为 $0$ 的非空子集”的前提下，$A$ 无论何时打开一批新箱子，这一整批的异或和都不为 $0$（否则这批箱子本身就是一个异或和为 $0$ 的非空子集，矛盾）。于是 $B$ 的回应总是可行的，按 $A$ 的行动分两类：
\begin{itemize}
\item $A$ 打开了一批新箱子：设新批的异或和为 $x\ne0$，则开箱后全体已打开箱子的异或和变成 $x$（原来为 $0$）。$B$ 把已打开箱子当作 Nim 局面，用 Nim 最优策略取石，把异或和归回 $0$。
\item $A$ 从某个已打开的箱子取石：这一步只改变一堆，异或和必变为非零（理由见 Nim 游戏的证明 2，即定理~\ref{thm:bouton} 的证明），$B$ 同样一步取石归零。
\end{itemize}
注意 $B$ 的归零操作只动一个已打开箱子，必然合法；且上面第一条用到条件“任何非空子集异或和不为 $0$”——这正是本题 P 态判定的核心。用归纳可以确认：$B$ 的每一步都有回应，并把 $S=0$ 的局面交还给 $A$。还要排除一种意外：$A$ 会不会一步把局面清空从而取胜？$A$ 一步只动一个已打开的箱子，若 $B$ 归零后的局面只剩一个非空已开箱子，其异或和就等于该箱石子数、不可能为 $0$，所以这种情况不会出现；$A$ 永远无法一步清场，而石子总数有限，游戏最终在 $B$ 的某次回应后结束（全部箱子打开且取空），此时轮到 $A$ 无路可走。

原文把 $A$ 历次开箱称为“第一批”“第二批”，意思正是：无论 $A$ 开多少批，每一批之后 $B$ 都能把不变量补回来。

\paragraph{先手视角.} 反之，如果存在 $a$ 的异或和为 $0$ 的非空子集，则 $A$ 开局直接打开一个极大的这样的子集（不能再并入任何未开箱子而保持异或和为 $0$），把已开箱子的异或和做成 $0$。此后 $B$ 只要打开新的一批 $U$，由极大性知 $U$ 的异或和必不为 $0$（否则 $S^*\cup U$ 是更大的异或和为 $0$ 的子集），$A$ 就按 Nim 策略归零；$B$ 若取石，$A$ 同样归零。角色与后手视角完全互换，于是 $A$ 必胜。

\begin{theorem}[开箱子判定]\label{thm:boxnim}
初始状态是 P 态，当且仅当 $a$ 的任何非空子集的异或和都不为 $0$（即 $\{a_i\}$ 在 $\Z_2$ 上线性无关）。
\end{theorem}

本题与台阶游戏的思路一脉相承：都是先手的“开启动作”会打破不变量，而后手（或赢得主动权的一方）负责把不变量修回来。区别只在于这里的“开启”不是搬运而是解锁，因此判定的条件从“奇数级异或和”变成“是否存在异或和为 $0$ 的子集”。

\section{分层例题与解析}

以下例题与本章的三板斧对应，两题的解析最终都落在“设计定式”上；解析默认读者已读过本章相应小节。

\begin{problem}{[CF1147F] ZigZag Game}
这是一道交互题。

给定一个包含 $2n$ 个节点的完全二分图，每一侧各有 $n$ 个节点。节点 $1$ 到 $n$ 位于二分图的一侧，节点 $n+1$ 到 $2n$ 位于另一侧。你还会得到一个 $n\times n$ 的矩阵 $a$，描述了边的权值。$a_{i,j}$ 表示节点 $i$ 与节点 $j+n$ 之间的边的权值。每条边的权值都不相同。

Alice 和 Bob 在这个图上玩一个游戏。首先，Alice 选择自己要扮演“递增”还是“递减”，Bob 获得剩下的选择。然后，Alice 可以把一个棋子放在图上的任意一个节点。接着，Bob 沿着该节点的任意一条相邻边移动棋子。之后，双方轮流进行如下操作，Alice 先手。

当前玩家必须将棋子从当前顶点移动到某个相邻且未访问过的顶点。设 $w$ 为上一次经过的边的权值。如果玩家扮演的是“递增”，则所经过的边的权值必须严格大于 $w$；如果扮演的是“递减”，则必须严格小于 $w$。无法行动的玩家判负。

给定 $n$ 以及图中各边的权值。你可以选择扮演 Alice 或 Bob，并将与评测机对战。你必须赢得所有对局，答案才会被判为正确。

数据范围：$n\le50$。
\end{problem}

\hint 先想清楚该扮演谁：评测机的先手选项更多，稳妥的是扮演后手；再把“永远有路可走”落实成一组可以预处理出来的完美匹配，匹配的性质与稳定婚姻完全一致。

\sol 先想清楚该扮演谁。评测机（先手 $A$）先挑“递增”或“递减”并把棋子放在某个节点上；它能掌控的因素比我们多，因此更稳妥的是扮演后手 $B$，让 $A$ 先犯错。扮演后手意味着我们的策略要无论 $A$ 怎么开局都能赢，而状态图有 $2n$ 个点、对局不可能无限长——双方每次都必须移到未访问点——所以只要证明“$B$ 永远不会无路可走”，就自动推出“先无路可走的一定是 $A$”。那么 $B$ 怎么保证永远有路？考虑设计定式：预处理一组满足一些性质的完美匹配，使得 $B$ 一直走匹配边合法。

不妨设先手 $A$ 选择扮演“递增”，且 $A$ 选择的起点是左部点（其余情形的化归套路如下，细节从略：$A$ 扮演“递减”时，把所有权值取反即回到“递增”的讨论，匹配按取反后的偏好重新构造即可；$A$ 的起点在右部点时，利用 $B$ 第一步不受单调约束的自由，把局面镜面对称地化归到下面的框架）。设想一局的标准走向：$A$ 的起点是左部点 $u$，$B$ 第一步走匹配边 $u\to v$（$v$ 是右部点）；之后 $A$ 只能从右部走到左部（边 $v\to w$），$B$ 再从左部走匹配边 $w\to x$ 回右部；如此往复，$B$ 的步总是“左$\to$右的匹配边”，$A$ 的步总是“右$\to$左的自由边”。因为 $B$ 扮演“递减”、$A$ 扮演“递增”，一局中相邻两步的权值必须满足：$A$ 的边 $v\to w$ 的权值 $>$ 上一步 $u\to v$ 的权值；$B$ 的边 $w\to x$ 的权值 $<$ 上一步 $v\to w$ 的权值。于是匹配 $M$ 需要保证：只要 $A$ 合法地走了 $v\to w$，就一定有 $a_{w,x}<a_{v,w}$。

匹配要满足怎么样的性质？观察到只要满足：对每条从左部点出发、到达右部点的长度为 $3$ 的路径 $u\to v\to w\to x$，$a_{u,v}<a_{w,v}$ 和 $a_{w,v}<a_{w,x}$ 不同时成立，其中 $(u,v)$、$(w,x)$ 是匹配边，$v\to w$ 是 $A$ 的自由边；$a_{w,v}=a_{v,w}$ 只是记号换了个下标。

这条性质很像稳定婚姻匹配具有的性质：不存在两个人 $u,v$ 满足：对 $u$ 来说 $v$ 比其当前配偶更优，且对 $v$ 来说 $u$ 比其当前配偶更优。对照一下：若 $a_{u,v}<a_{w,v}$ 且 $a_{w,v}<a_{w,x}$ 同时成立，则右部点 $v$（现任配偶是 $u$）觉得左部点 $w$ 比 $u$ 好——因为右部点偏爱权值大的边，而 $a_{w,v}>a_{u,v}$；左部点 $w$（现任配偶是 $x$）也觉得 $v$ 比 $x$ 好——因为左部点偏爱权值小的边，而 $a_{w,v}<a_{w,x}$。两人互相觉得对方更好，$(v,w)$ 就是一个破坏匹配的私奔对；反之亦然。所以“稳定婚姻匹配”恰好就是我们要找的匹配。考虑设计一个优先级，使得任意一组稳定婚姻匹配都满足所需性质。不难发现让左部点优先匹配权值小的边、右部点优先匹配权值大的边即可。

稳定婚姻匹配一定存在，下面给出一个构造性证明。不难想到以下贪心算法：考虑依次加入每个右部点 $u$。对于 $u$，按照从 $u$ 最喜欢到最不喜欢的顺序依次考虑左部点 $v$。如果 $v$ 没有匹配，就让 $v$ 匹配 $u$，停止；如果 $v$ 相比其当前匹配点 $match_v$ 更喜欢 $u$，就交换 $u$ 与 $match_v$；否则考虑下一个 $v$。算法的语义很直白：每个新来的右部点按自己的喜好顺序“追求”左部点，被追求的左部点若觉得对方更好就“劈腿”，被甩的右部点继续往下追——这就是教科书里的延迟接受算法（Gale--Shapley），必然在有限步内终止（每次被甩的右部点都会转而考虑更差的选择，总次数有限），且终止时每个左部点至多一个配偶、每个右部点要么配对要么追完所有左部点被剩。需要说明完美性：由于两侧重数相等，且每个右部点一旦配对就不会再变单身（只会被甩给下一个左部点……准确说是继续追求直到配对），可以归纳证明算法结束时两侧全部配对。省略的细节读者可以自行补齐。

为什么贪心算法求出的匹配是稳定的？观察到对每个 $v$，在它匹配后，每次更换匹配点都是因为有了更喜欢的。所以，对每个 $u$，如果它打不过 $v$ 当前的匹配点，就不可能打过 $v$ 未来的匹配点——$v$ 只会越换越好，被 $u$ 拒绝过的左部点（或者更准确地说，输给 $v$ 当时配偶的左部点）以后更不可能被接受。所以，$u$ 的当前匹配对它来说就是最优的：任何未配对成功的 $u,v$ 若互相想私奔，那么在 $v$ 最后一次考虑时 $u$ 一定排在 $v$ 的现任之后，矛盾。也就是说贪心算法求出的匹配就是稳定的。

\paragraph{收尾.} 回到博弈：$B$ 拿到稳定婚姻匹配 $M$ 后，策略为“永远沿匹配边从当前左部点走到右部点”。上面已经论证：$A$ 的每一步合法行动都会把局面置于使 $B$ 的匹配边恰好合法的位置，而匹配的一一对应保证了 $B$ 不会走到已访问的右部点（$A$ 永远落在左部点，$B$ 的落点两两不同），因此 $B$ 永不卡死；图只有 $2n$ 个点，$A$ 迟早会找不到新的左部点而告负。

\cost 稳定性与贪心构造都是 $O(n^2)$ 的（$n\le50$ 绰绰有余），实现时把每条边的权值看成“左部点对右部点的排序”即可。

\src{https://codeforces.com/problemset/problem/1147/F}

\begin{problem}{[Luogu P8851] T2nz.}
这是一道交互题。

小 X 陷入了一个奇怪的梦。在梦境里，她在和小 Q 下一种奇怪的棋。

这是一个 $2^{2n}\times 2^{n}$ 的棋盘，小 X 执黑先行，小 Q 执白后行。

每次操作，需要在当前未满的第一行内，任意选择一格下棋。一格内只能有一个棋子。

下满之后，共有 $2^{2n}$ 行棋子，小 X 的得分为本质不同的行数。

小 X 想最大化她的得分，但小 Q 想最小化小 X 的得分。

你的任务是，扮演小 X 或小 Q，最大化或最小化得分。

若你是小 X，在满足最大化得分 $ans$ 的同时，你也要最大化前 $ans$ 行中本质不同的行数。

数据范围：$n\le7$。
\end{problem}

\hint 每行恰好有 $n$ 个黑子与 $n$ 个白子，一行的样式由黑白对的选择决定；后手用完美匹配压上限，先手用“每步排除一半候选行”的抽屉原理抬下限。

\sol \paragraph{题面翻译.} 棋盘共 $2^{2n}$ 行、每行 $2^{n}$ 列；行是从上到下逐行填满的（每次只能落在当前最上面未满的行里）。一行填满需要 $2^{n}$ 步，小 X（黑）与小 Q（白）轮流落子，因此每一行恰好有 $n$ 个黑子和 $n$ 个白子，一行由 $2^{n}$ 个格子的颜色完全决定。X 的得分是本质不同的行数，X 想最大化、Q 想最小化。题目的第二段话（“最大化前 $ans$ 行中本质不同的行数”）是因为 $ans$ 达到上界时 $2^{n}$ 种行都可能出现，为了判定输出是否最优而加上的。

答案预告：均衡下得分恰为 $2^{n}$，下面分别给出 Q 的“不超过”论证与 X 的“至少”论证。

\paragraph{后手视角（上界）.} 考虑设计定式。对于 $2^{n}$ 个格子，预先设计一个完美匹配（把每行的格子两两配成 $2^{n-1}$ 对）。$A$（小 X）每选择一个点，$B$（小 Q）就选择那个点的匹配点。为什么这样能把行数压到 $2^{n}$？关键是匹配点被白子占掉后，X 就没法再落同一对里的另一个格子：他每落一个黑子，Q 立刻把该黑子的配偶格染白，于是每一对格子至多被 X 染黑一次。而每行恰好有 $n$ 个黑子、恰好有 $2^{n-1}$ 对格子，所以每行黑子必然每对恰好一个；更妙的是白子的位置也由黑子唯一决定（白子恰在匹配位置），于是整行的样式完全由“$2^{n-1}$ 对中每一对选左端还是右端”决定，总共 $2^{n}$ 种。这样能够保证 $ans\le2^{n}$。

\paragraph{先手视角（下界）.} 猜测答案就等于 $2^{n}$，下面给出一个构造性证明。X 的策略目标：让前 $2^{n}$ 行两两不同。由上面的上界，前 $2^{n}$ 行两两不同就意味着一共只有 $2^{n}$ 种行、每种恰好出现一次，得分直接拉满且自动满足题目的二级目标。设当前考虑到了第 $i$ 行（$i\le2^{n}$），我们论证 X 可以让第 $i$ 行与前面所有行都不相同。考虑让每次决策保证第 $i$ 行与至少一半的行不相同，这样 $n$ 个回合后就能保证第 $i$ 行在之前没有出现过。

把“前面可能与第 $i$ 行相同的行”称为候选行。X 落第 $j$ 个黑子（$1\le j\le n$）之前，第 $i$ 行已填 $2(j-1)$ 格，还空着 $2(n-j+1)$ 列；候选行在前 $2(j-1)$ 列上与第 $i$ 行完全一致，所以候选行每行剩下的 $n-j+1$ 个白子全都落在这些未填列里。设当前还有 $i'$ 行可能与第 $i$ 行相同。那么，这 $i'$ 行在第 $i$ 行未填的列中，一共有 $i'\times(n-j+1)$ 枚白子。由抽屉原理，其中白子最多的列至少有
\[
 \ceil{\frac{i'\times(n-j+1)}{2(n-j+1)}}=\ceil{\frac{i'}{2}}
\]
个白子。A 在这一列落一个黑子，就能排除至少一半的行——凡是这列是白子的候选行，第 $i$ 行在此列是黑子，二者不再可能相同，因此被排除的候选行至少有 $\ceil{i'/2}$ 个，剩余的候选行至多 $\floor{i'/2}$ 个。候选行数量每落一个黑子至少减半：$n$ 个黑子落完后，剩余候选至多 $\floor{(i-1)/2^{n}}=0$（因为 $i\le2^{n}$），第 $i$ 行是全新的行。对每个 $i\le2^{n}$ 都这样做，前 $2^{n}$ 行两两不同，$ans\ge2^{n}$。

\paragraph{总结.} 上下界相撞得到 $ans=2^{n}$。扮演小 X 时照抄抽屉原理构造；扮演小 Q 时照抄完美匹配回应。实现都是 $O(n2^{n})$ 级别的（$n\le7$，行数 $2^{2n}\le2^{14}$），完全足够。本题是“设计定式”思想的综合演练：后手靠预先设计的匹配限制先手，先手靠“每一步排除至少一半”的势能论证突破限制，两段论证恰好咬合成同一个数 $2^{n}$。

\src{https://www.luogu.com.cn/problem/P8851}

\begin{teach}[这道题的三个教学收获]
上界往往来自“给对手的反应预先建档”（完美匹配）；下界往往来自“每步把坏消息砍一半”（抽屉原理）；两级论证在 $2^{n}$ 处咬合，才是“均衡”二字的具体形象。学生常犯的错是把上界论证当成下界复用，批改时优先检查论证方向。
\end{teach}
\chapter{PN 态 DP}
把每个状态标成 P 态或 N 态的动态规划，称为 \textbf{PN 态 DP}。它是公平组合游戏题最通用的工具：只要状态和转移可以完整枚举，胜负判定就能机械地算出来；前面几节那些漂亮的定式，可以理解为 PN 态 DP 在特殊结构下的“手算加速版”。

\section{从递推刻画到动态规划}
回顾定理~\ref{thm:recursion} 给出的递推刻画。设 $f(u)=1$ 表示状态 $u$ 是 N 态，则：
\begin{itemize}
\item 若 $u$ 没有合法后继，当前玩家直接判负（无法行动者负的规则下），$f(u)=0$；
\item 一般地，$f(u)=1$ 当且仅当存在后继 $v$ 使 $f(v)=0$（先手走到 P 态即必胜）；
\item 等价地，$f(u)=0$ 当且仅当所有后继 $v$ 都有 $f(v)=1$。
\end{itemize}
若状态图是 DAG，按拓扑序（或者对每个状态记忆化搜索）从后往前递推即可，复杂度为 $O(\text{状态数}+\text{转移数})$。常见的题目里每个状态的出度只有 $O(1)$，状态数在 $10^6\sim10^7$ 级别时直接 DP 就能通过。

\begin{center}
\begin{tikzpicture}[every node/.style={draw=ink!50,rounded corners,fill=soft,font=\small,inner sep=5pt},>=Stealth]
\node (d) at (0,0) {$d$（P）};
\node (a) at (3,1.1) {$a$（N）};
\node (b) at (3,-1.1) {$b$（N）};
\node (c) at (6,0) {$c$（P）};
\draw[->] (a)--(d);
\draw[->] (b)--(d);
\draw[->] (c)--(a);
\draw[->] (c)--(b);
\end{tikzpicture}\\[2pt]
{\small 状态图上的逆向归纳：终局 $d$ 无后继，为 P 态；$a,b$ 有 P 态后继，为 N 态；$c$ 的全部后继都是 N 态，故为 P 态。}
\end{center}

因此这类题目的难点几乎都不在 DP 本身，而在两个方面：
\begin{itemize}
\item \textbf{压缩状态}：必须抓住博弈的本质，把显然等价的局面合并。例如例题“长条格子”把连续几万步的操作压缩成 $O(n)$ 个状态，“FollowingNim 改”把整堆的博弈预处理好再参与 Nim 判定，都是状态设计的典范。
\item \textbf{处理环}：状态图一旦有环，上述递推就不封闭——环上的状态既推不出 P 也推不出 N，需要引入第三种结果“平局”并用图论方法求解，见例题 [省选联考 2023] 过河卒。
\end{itemize}

另外，很多题目在 DP 之前需要先用“策略窃取”式的论证删掉对手的部分选项，把看似庞大的状态空间压到可枚举的规模。PN 态 DP 的扩展性很强，本章末的分层例题（过河卒、长条格子、FollowingNim 改等）会反复展示这些套路。

\section{分层例题与解析}

以下例题的共同骨架是状态图：先想清楚状态与转移，再决定 DP、拓扑还是队列。

\begin{problem}{[省选联考 2023] 过河卒}
题意摘要：$n\times m$ 棋盘上有三枚棋子，双方轮流按规则移动其中一枚，率先无路可走者负；游戏可能无限进行，需要判定每个初始局面的胜负或平局。完整题面见原题（洛谷 P9169）。
\end{problem}

\hint 状态是三枚棋子的位置，图不是 DAG；先在反图上拓扑排序定“胜／负／未定”，再用调整法说明未定点全是平局，最后按步数从小到大用队列求出分出胜负的步数。

\sol \paragraph{第一步：把博弈翻译成状态图.} 观察到状态是三枚棋子的位置，只有 $O(n^3m^3)$ 种，且每个点的出度是 $O(1)$ 的，所以可以暴力建图。但是有一件比较麻烦的事情：图并不是 DAG，所以不能套路地按拓扑序的逆序 DP（实际上根本不能 DP，只能尝试使用一些图论算法）。状态图里有环意味着可能出现第三种结果：平局——双方若都选择不冒险，游戏可以无限进行。因此本题的答案本质上是给每个状态打三分类标签：胜负平，附带“分出胜负需要几步”。

\paragraph{先考虑判定胜负.} 先尝试确定所有能确定状态的点的状态。对于点 $u$，如果存在 $v\in\Np(u)$ 满足 $v$ 是 P 态，则 $u$ 是 N 态；如果对每个 $v\in\Np(u)$ 都满足 $v$ 是 N 态，则 $u$ 是 P 态。这个过程可以通过在反图上拓扑排序高效地实现：把“状态尚未确定”的点放进队列，每次取出一个点，看它的所有后继是否都已经确定——若存在 P 态后继则标 N，若后继全是 N 态则标 P；标完后扫描它的前驱，更新“前驱的已确定后继数”，一旦某个前驱的后继全部确定且都不是 P 态，它就可以被标成 P 态入队。这本质上就是 P/N 递推刻画（定理~\ref{thm:recursion}）的“拓扑排序实现”，区别只在容许环的存在。

对每个状态尚未确定的点 $u$，一定满足：
\begin{itemize}
\item 对每个 $v\in\Np(u)$，$v$ 要么是 N 态，要么状态尚未确定。
\item 存在 $v\in\Np(u)$，$v$ 的状态尚未确定。
\end{itemize}
第一句话来自“如果所有后继都能确定，$u$ 自己就能确定”的逆否；第二句话因为若所有后继都是 N 态，$u$ 早就被标成 P 态了。于是从任何一个未确定点出发，总能沿着“未确定的后继”一直走下去。

考虑调整法，从 $u$ 出发走到 N 态不优，所以只需考虑状态尚未确定的点的导出子图 $G'$。为什么走到 N 态不优？因为把局面交给一个 N 态等于让对方有必胜策略，理性玩家不会这么选；双方都只会在 $G'$ 内部移动。因为 $G'$ 满足每个点的出度都大于等于 $1$，所以无论玩家如何行动，游戏都会无限进行下去，所以每个点都是平局。这里“每个点出度 $\ge1$”是环上博弈出平局的关键：双方都没有被迫离开 $G'$ 的出口。

\paragraph{再求分出胜负的步数.} 对每个 P 态或 N 态的点 $u$，求分出胜负的步数 $f_u$：
\begin{itemize}
\item 对于 P 态点：$f_u=\max(f_v)+1$。P 态点的先手（必败方）会尽量拖延败局到来的时间，所以从所有后继中挑 $f$ 最大者走；后手反正必胜，在每一个后继（N 态）上都会尽快结束。
\item 对于 N 态点：$f_u=\min(f_v)+1$。N 态点的先手（必胜方）会尽快取胜，从所有后继中挑 $f$ 最小者走。
\end{itemize}
考虑按照 $f$ 从小到大的顺序求出每个点的 $f$。则一个 P 态点的 $f$ 只有在出点的 $f$ 都确定之后才能确定（依赖 $\max$），而一个 N 态点的 $f$ 只要存在一个出点的 $f$ 确定就可以确定了（依赖 $\min$，取已确定的最小值）。为了保证按照 $f$ 从小到大的顺序求出每个点的 $f$，可以用队列维护 $f$ 可以确定但还未确定的所有点：P 态点等所有出点的 $f$ 出炉后入队，N 态点等第一个出点的 $f$ 出炉即入队；由于 $f$ 小者先被算出，$\min$ 与 $\max$ 都能在入队时正确结算。

\cost 建图与两轮拓扑排序都是 $O(\text{状态数}+\text{转移数})=O(n^3m^3)$（状态数约 $1.6\times10^{6}$ 级别，可过）。本题是“带环状态图的 PN 判定”的教科书例题，框架值得背下来：先拓扑定胜负平，再队列算步数。

\src{https://www.luogu.com.cn/problem/P9169}

\begin{problem}{长条格子}
有一个 $1\sim n$ 的长条格子，第 $i$ 个格子有正整数 $a_i$。假设一个 bot 初始在 $s$，$A$ 和 $B$ 轮流操作，$A$ 先手，每次可以进行其中一个行动（假设 bot 当前位置为 $p$）：
\begin{itemize}
\item 将 bot 移动到 $[p+1,\min(p+k,t)]$ 中的任意一个格子；
\item 给 $a_p$ 减 $1$（只有 $a_p>0$ 时才能进行此行动）。
\end{itemize}
其中 $k$ 是一个给定的参数，$t$ 是 bot 的终点。无法行动的人会输掉游戏。

给出了 $q$ 次操作，每次操作是以下两种之一：
\begin{itemize}
\item \texttt{1 l r d} 表示对于所有正整数 $i\in[l,r]$，将 $a_i$ 变为 $a_i+d$。
\item \texttt{2 s t} 表示当 bot 初始点在 $s$，终点在 $t$ 时，谁有必胜策略。
\end{itemize}
数据范围：$1\le n,k,q\le2\times10^{5}$，$1\le l\le r\le n$，$1\le s\le t\le n$，$0\le a_i\le10^{4}$，$0\le d\le300$。
\end{problem}

\hint 先用镜像论证说明“减 $1$”操作不参与状态、只贡献奇偶条件；再观察 PN 态 DP 中 false 点构成的稀疏链，从右往左跳跃定位，用分块把单次查询压到 $O(\sqrt n)$。

\sol \paragraph{读懂终态.} 先搞清楚什么时候会输：若 bot 在 $p=t$（无法前进）且 $a_p=0$（无法减 $1$），当前玩家就无路可走。因此“给 $a_p$ 减 $1$”这个操作的意义是：它不移动 bot，只是把“谁被迫移动 bot”的决定权推来推去。

\paragraph{考虑策略窃取：消去减 1 操作.} 考虑先手视角，如果可以将 bot 向后移动到达一个 P 态，则初始状态是 N 态——这一步是平凡的：直接把 bot 挪进一个先手必败的局面交给对方。否则（所有可达的下一步都是 N 态），局面就完全由 $a_p$ 的奇偶性支配：如果 $2\nmid a_p$，则先手可以逼迫后手将 bot 向后移动，所以初始状态是 N 态；如果 $2\mid a_p$，则后手可以反过来逼迫先手移动 bot，所以初始状态是 P 态。这里的“逼迫”用的是镜像论证：两人在同一个格子上轮流减 $1$，$a_p$ 为奇数时先手会执行最后一次减 $1$、把“移动 bot”的包袱甩给后手；为偶数时则相反。综上，初始状态是 N 态，当且仅当可以将 bot 向后移动到达一个 P 态，或者 $2\nmid a_i$。于是我们消除了给 $a_p$ 减 $1$ 这个操作带来的影响——它不再参与状态，只贡献一个奇偶条件。

\paragraph{考虑 PN 态 DP.} 现在状态数是 $O(n)$：设 $f_i$ 表示 $p=i$ 时先手是否必胜（其余参数 $a,k,t$ 固定）。转移：
\[
 f_i=(2\nmid a_i)\lor\Bigl(\bigvee_{j=1}^{k}\lnot f_{i+j}\Bigr).
\]
读法：站在 $i$，先手赢有三种来源——$a_i$ 是奇数（上面镜像论证的“甩包袱”招），或者能把 bot 直接挪到某个 $f=false$ 的格子（P 态）$i+j$。$t$ 之后没有格子，转移里 $f_{>t}$ 一律视为 $true$ 即可。

\paragraph{换一种表示法（这是本题真正的考点）.} 上面 $O(n)$ 个状态的递推本身不难，难点在于 $a$ 会被区间加 $d$ 修改、还要回答 $O(q)$ 次任意 $(s,t)$ 的查询，每次重新 DP 是 $O(n)$ 的，不能接受。观察 $f$ 的结构：$f_i=false$ 当且仅当 $a_i$ 为偶数且 $i+1,\ldots,i+k$ 全是 $true$（N 态）。从右往左看，false 点之间的间距至少是 $k+1$：若 $x<y$ 都是 false 且 $y-x\le k$，则 $x$ 的后继窗口里有 false 点 $y$，$f_x$ 应为 $true$，矛盾。于是 false 点构成一条稀疏的链，可以从右往左“跳”着找。考虑从右往左枚举每个 $f$ 为 false 的点，设当前已找到的最右 false 点是 $p$，下一个 false 点必是位置 $\le p-k-1$ 的第一个偶数位置：先看偶数的条件——false 点必须 $a$ 为偶数；再看右窗——任何 $x\le p-k-1$ 的窗口 $(x,x+k]$ 都在 $p$ 的左边，而 $(p-k,p-1]$ 之间即便有偶数也不会是 false 点（它们的窗口里躺着 $p$），更左的偶数又不可能比“第一个偶数”更靠右，所以“从 $p-k-1$ 向左第一个偶数”的右窗是干净的，它确实就是下一个 false 点。初始时 bot 在终点 $p=t$：若 $a_t$ 为偶数，$t$ 自己就是最右的 false 点（$t$ 没有后继）；若 $a_t$ 为奇数，$t$ 不是 false 点，此时最右的 false 点是 $t$ 左侧离 $t$ 最近的偶数位置（它的右窗至多含 $t$，而 $t$ 不是 false 点，右窗干净），应把它作为第一个 $p$。每次 $p$ 变成 $a_{p-k-1},a_{p-k-2},\ldots$ 中第一个偶数的下标。这样从最右的 false 点出发，可以逐个向左访问整条 false 链，直到遇到第一个 $\le s$ 的链上点。判定只有两种可能：如果 $p=s$ 则后手必胜——因为 $s$ 自己就是 false 点，即 $f_s=false$；如果 $p<s$（或链耗尽）则先手必胜——因为 $s$ 不是 false 点，它或者 $a_s$ 为奇数，或者右侧 $k$ 步以内已经躺着一个 false 点（这正是 $s$ 没被选进链的原因），两种情况都给出 $f_s=true$。

\cost 这个过程可以通过分块优化：把“从 $p-k-1$ 向左第一个偶数”预处理成一步跳跃（或并查集式跳表），每步跳跃至少前进 $k+1$，再用分块把连续跳跃压缩，使每次查询降至 $O(\sqrt n)$。时间复杂度 $O(q\sqrt n)$，配合区间加的维护（块内打标记即可处理 $d\le300$ 的修改）即可通过 $2\times10^{5}$ 的数据。

\begin{problem}{FollowingNim 改}
有 $n$ 堆石子，第 $i$ 堆大小为 $a_i$，每次可以选一堆石子取走至少一个，不能行动者负。

有一个额外限制：如果上一名玩家取走的石子数量 $x\in S$，则下一名玩家必须选择与上一名玩家同一堆石子。

数据范围：$a_i\le10^{4}$。
\end{problem}

\hint 新元素是“锁定”：取到 $S$ 中的数量就把双方锁死在同一堆，直到有人取到 $S$ 之外的数解锁。预处理“锁内子局”的胜负表，冷堆压缩成排名后异或判定。

\sol \paragraph{先看清规则的新元素.} 普通 Nim 之外这里多了一条“跟随”约束：上一手取走的数量 $x\in S$ 时，下一手被迫留在同一堆。这带来一个全新的子结构——\textbf{锁定}：一旦有人在一堆上取了属于 $S$ 的数量，游戏就被锁死在这堆上，双方轮流从这堆取 $S$ 中的数量，直到一方无法行动（堆空，或堆里剩下的石子不足 $S$ 中任何可取的数量）而告负；只有某人取走一个不属于 $S$ 的数量，锁定才会解除、游戏回到“自由 Nim”。

\paragraph{考虑策略窃取.} 如果存在 $i,j$ 满足从第 $i$ 堆石子取走 $j$ 个后到达一个 P 态（即使 $j\in S$，也不对 B 的下一步操作施加限制），则初始状态是 N 态——这里 P 态按完整规则（包含锁定机制）理解：先手只需要一步走进 P 态，之后无论锁定与否都是后手吃亏，故先手必胜。否则（一步到不了 P 态），分析 A 的第一步必须是什么。若 A 取 $j\notin S$：这一手不产生锁定，局面回到普通的自由博弈；而自由局面下任何一步都到不了 P 态（这正是“否则”的含义），意味着 A 把某个 N 态交给了 B，B 有必胜应对，所以 A 的第一步不可能取 $j\notin S$。因此 A 第一次取走的石子数必须属于 $S$，第一步必然锁定一堆。设 A 取了第 $i$ 堆，则接下来 B 也要取第 $i$ 堆。如果 B 接下来取走的石子数不属于 $S$，则会到达一个 N 态，A 胜——因为解锁后的自由局面是 N 态、轮到 A 先手；所以 B 接下来取走的石子数必须属于 $S$，同理 A 接下来也要取第 $i$ 堆，取走的石子数量也必须属于 $S$，依次类推，直到某一方无法行动。也就是说：一旦开局锁定，理智的双方都不会主动解锁（解锁等于把 N 态送给对方），游戏在锁内走到尽头。

\paragraph{考虑 DP 预处理.} 设 $f_i$ 表示两人轮流从一个大小为 $i$ 的堆中取石子，每次取走的石子数量必须属于 $S$，是否先手必胜——这是“锁内子局”的胜负表：$f_0$ 恒为 $false$（无石可取），一般地 $f_i=\bigvee_{x\in S,x\le i}\lnot f_{i-x}$。锁内子局的“先手”指第一个在锁内取石的人（即发起锁定的玩家）。

\paragraph{回到原局.} 若存在 $i$ 满足 $f_{a_i}=true$，则初始状态是 N 态：A 开局就在第 $i$ 堆取一个属于 $S$ 的数量，发起锁定并充当锁内先手；由于 $f_{a_i}$ 为 $true$，B 若老实留在锁内必败。B 唯一的脱身方式是取一个不属于 $S$ 的数量主动解锁——注意 B 的解锁手仍必须落在这同一堆上，于是 A 的锁定手与 B 的解锁手可以“合并”成 A 的一步（一次取走两者之和），也就是说 B 解锁后交出的自由局面其实都是初始局面的一步后继；而“否则”的前提保证不存在一步可到的 P 态，所以这些局面全是 N 态、轮到 A 必胜。B 无论留在锁内还是解锁都输，故 A 必胜。（若想给出完全严谨的表述，可把“某堆正处于锁定状态”也纳入 P/N 归纳一并处理，此处从略。）

否则（所有堆都是“冷堆”，$f_{a_i}$ 全为 $false$），A 不能指望靠开局锁定取胜：开局锁定后 A 是锁内先手，而 $f_{a_i}=false$ 正说明这个锁内先手必败（B 只要留在锁内，A 就输；B 解锁反而可能给 A 机会，所以 B 乐于留锁）。此时初始状态是 N 态，当且仅当存在 $i,j$ 满足从第 $i$ 堆石子取走 $j$ 个后到达一个 P 态。显然，这要求 $f_{a_i-j}=false$——取完后第 $i$ 堆仍是冷堆或空堆：若取到热堆，则 B 可以反过来锁定它占便宜，A 无利可图。反过来，只要这个要求满足，后续博弈过程就变成了 Nim 游戏：谁先碰热堆谁吃亏，双方都只会在冷堆之间移动，而“把一个冷堆改成另一个更小的冷堆”与 Nim 取石一一对应。所以，预处理 $g_i$ 表示 $i$ 是第几个 $f$ 为 $false$ 的数（把冷堆“压缩”成它在冷堆序列中的排名，排名只和冷堆的相对大小有关），则初始状态是 N 态当且仅当 $\bigoplus g_{a_i}\ne0$。\footnote{原稿此处写作“$\bigoplus g_{a_i}=0$”，按 Nim 游戏结论应为 $\ne0$，整理时已订正。}

\paragraph{为什么压缩成排名 $g$ 而不是直接用 $a_i$？} 因为 Nim 的一步要求“能把一个堆改成任意更小的值”；对冷堆而言，可改到的值只有那些 $f=false$ 的数（改到热堆等于送死），而“从排名 $g_i$ 改到任意更小的排名”恰好一一对应 Nim 里“从大小为 $g_i$ 的堆取到更小”，于是异或判定照搬。

\cost $f,g$ 的预处理是 $O(\max a_i)$ 级别，判定每组堆 $O(n)$，完全满足 $a_i\le10^{4}$ 的限制。

\begin{problem}{[BJOI2018] 双人猜数游戏}
本题为提交答案题。

Alice 和 Bob 是一对非常聪明的人，他们可以算出各种各样游戏的最优策略。现在有个综艺节目《最强大佬》请他们来玩一个游戏。主持人写了三个正整数 $s,m,n$，然后一起告诉 Alice 和 Bob $s\le m\le n$ 以及 $s$ 是多少。（即，$s$ 是接下来要猜的 $m,n$ 的下限。）

之后主持人单独告诉 Alice $m\times n$ 是多少，单独告诉 Bob $m+n$ 是多少。

当然，如果一个人同时知道 $m\times n$ 以及 $m+n$ 的话就能很容易地算出 $m$ 和 $n$ 分别是多少，但现在 Alice 和 Bob 只分别知道其中一个，而且他们只能回答主持人的问题，不能交流。主持人从他们中的一人开始，轮流询问回答者“知不知道 $m$ 和 $n$ 分别是多少”（回答者只能回答知道/不知道）。

为了节目效果，以及显示出 Alice 和 Bob 的绝顶聪明，主持人希望 Alice 和 Bob 一共说了 $t$ 次“不知道”以后两个人都知道 $m$ 和 $n$ 是多少了。现在主持人找到你，希望让帮他构造一组符合条件的 $m$ 和 $n$。

数据范围：$1\le s\le200$，$2\le t\le15$，保证有解。
\end{problem}

\hint 这不是传统意义的“博弈”，而是公共知识推理：每句“我不知道”都在剔除候选对。按时间顺序分层 DP，答案就是“恰好第 $t+1$ 轮被唯一确定”的对。

\sol \paragraph{这是什么类型的题？} 本题不是传统意义的“博弈”，而是公共知识推理：Alice 每说一次“我不知道”，Bob 就获得一条信息——“谜底不在 Alice 的排除集里”；Bob 随后说“我不知道”又反过来帮 Alice 缩小范围。如此迭代，双方在轮流把候选 $(m,n)$ 剔除出去。题目要我们构造的，正是“恰好剔除 $t$ 轮后候选只剩一对”的谜底。注意“两人都绝顶聪明”在这里的含义：他们能立刻完成任意复杂的推理，因此每一句“我不知道”都等价于“在我掌握的信息下，候选不止一个”。

\paragraph{考虑按时间顺序 DP.} 设 $f_{t,n,m}$ 表示如果谜底是 $(n,m)$，是否可能连续说 $t$ 次“我不知道”——即前 $t$ 轮的回答与谜底 $(n,m)$ 相容。边界 $f_{0,n,m}=true$（没说任何话时一切皆有可能）。考虑转移。不妨设第 $t$ 次说“我不知道”的人是 A（若本轮轮到 B，把“积”换成“和”即可，推理完全对称）。A 掌握的信息是积 $P=n\cdot m$。在第 $t$ 轮之前，A 已经通过前 $t-1$ 句话（以及自己每轮的“我不知道”）知道：谜底必须是某个满足 $f_{t-1,u,v}=true$ 的对 $(u,v)$。于是 A 心里列出与自己看到的积相同、且尚未被排除的所有候选：$\{(u,v):uv=P,\ f_{t-1,u,v}=true\}$。如果这个集合只有一个元素，A 就能确定谜底，本轮应该说“我知道”；如果集合里至少还有两对，A 只能继续“我不知道”。翻译回 DP：那么，对于固定的 $n\times m$，如果只有一对 $(n,m)$ 满足 $f_{t-1,n,m}=true$（精确地说，积相同且 $f_{t-1}=true$ 的候选只有这一对），则如果谜底是 $(n,m)$，则第 $t$ 次 A 应该说“我知道”才对，所以 $f_{t,n,m}=false$。在其他情况下（候选至少两对，A 确实说不出所以然），$f_{t,n,m}=f_{t-1,n,m}$——相容性原样延续。实现时按 $t$ 分层滚动：对每个积 $P$ 统计该层候选个数即可，每层 $O(B^2)$。

\paragraph{表不能无穷大.} 但是我们不可能维护一个无穷大的表。合理猜测存在一个阈值 $B$，满足当 $n$ 或 $m$ 大于 $B$ 时，AB 双方在 $t+1$ 轮内不可能知道正确的 $(n,m)$。直觉上，大数对的分解与拆分方式太多，几轮“不知道”根本不足以把它们区分开；$s\le200$、$t\le15$ 的小范围也说明可行解不会依赖太大的数。实测取 $B=500$ 即可通过。本题出成提交答案题可能是因为出题人无法证明 $B$ 的上界——这提醒我们：做提交答案题时，“猜一个足够大的界 + 暴力验证”本身就是正解的一部分。

求出 $f$ 表后，扫描所有对 $(n,m)$ 找“前 $t$ 次都说不知道、第 $t+1$ 次两人都知道”的目标（即 $f_{t,n,m}=true$ 且第 $t+1$ 轮被唯一确定），输出一组即可。

\begin{teach}[为何这道题适合作为状态图题的结课例]
过河卒同时考到状态设计（三棋子位置）、环的处理（平局三分类）与两轮队列（先定胜负平、再算步数）。结课时可让学生只写“三步走”的伪代码框架而不写实现，能写对框架即视为达标；反图、入队时机等实现细节留作上机。
\end{teach}
\chapter{Nim 游戏及其变体}
Sprague--Grundy 理论（习惯上简称 SG 理论）的作用，其实是把 Nim 游戏的结论推广到更多的游戏上。因此我们先研究 Nim 游戏，再研究 Sprague--Grundy 理论。本章的四份“题意”彼此同构：同为取石子，差别只在“每次动几堆”与“无法行动者负还是胜”两处，正好用来观察结论如何随规则变化。

\section{Nim 游戏}
\begin{problem}{Nim 游戏}
有 $n$ 堆石子，第 $i$ 堆有 $a_i$ 个，每次可以选择一堆，取走至少一个石子，不能行动者负。判定初始状态。
\end{problem}

\begin{theorem}[Bouton 定理]\label{thm:bouton}
设 $S=a_1\oplus a_2\oplus\cdots\oplus a_n$。Nim 游戏中先手必胜，当且仅当 $S\ne0$。
\end{theorem}

这个结论是整个 Sprague--Grundy 理论的基石。为什么答案偏偏是异或而不是别的？一个直观的解释是：一次操作恰好改变一堆，而“把一堆从 $a_i$ 减到 $a_i'$”在二进制下等价于翻转若干个二进制位；异或恰好刻画了“每个二进制位上、值为 $1$ 的堆数的奇偶性”如何被一步扰动。至于它为什么能成为判定条件，看下面的两段式证明。

证明分两个部分：N 态的后继中存在至少一个 P 态；P 态的后继全是 N 态。这种“两段式”写法在本书会反复出现（反 Nim、Nim-k、SG 定理同款），本质是定理~\ref{thm:recursion} 递推刻画的直接应用：第一段说 $S\ne0$ 是 N 态，第二段说 $S=0$ 是 P 态。

\begin{proof}[证明 1（$S\ne0$ 的后继中有 P 态）]
设 $k=\topbit(S)$。因为 $S$ 的第 $k$ 位为 $1$，所以存在 $i$ 使得 $(a_i)_k=1$。那么 $a_i\oplus S<a_i$，把第 $i$ 堆石子取到只剩下 $a_i\oplus S$ 个即可。
\end{proof}

\begin{teach}[证明中的两个关卡]
先问“为什么必存在某堆第 $k$ 位为 $1$”（$S$ 的最高位 $1$ 来自奇数个这样的堆），再问“为什么 $a_i\oplus S<a_i$ 是合法取石”（最高位变小即可，低位任意）。两关都过，学生才算拿到 Bouton 定理。可当场演示 $\{3,4,5\}$：$S=2$，最高位是第 $1$ 位，只有 $3$（$011$）该位为 $1$，取成 $3\oplus S=1$，得 $\{1,4,5\}$，验证异或和为 $0$。
\end{teach}
\paragraph{白话翻译.} 异或和 $S$ 的最高位 $k$ 为 $1$，说明二进制第 $k$ 位为 $1$ 的堆有奇数个，于是必存在某堆 $a_i$ 的第 $k$ 位也是 $1$。对这堆做“异或替换”$a_i\leftarrow a_i\oplus S$：由于 $S$ 的最高位恰为 $k$，而 $a_i$ 的第 $k$ 位为 $1$，替换后第 $k$ 位变成 $0$ 且更高位不变，故 $a_i\oplus S<a_i$，确实是合法的取石操作；取完后新的异或和为 $S\oplus a_i\oplus(a_i\oplus S)=0$，一步到达 P 态。若只有一堆石子，则 $S=a_1$，$a_1\oplus S=0$，这一手就是把整堆取完。

\begin{proof}[证明 2（P 态的后继全是 N 态）]
操作一次后 $S$ 一定会发生变化。
\end{proof}
\paragraph{补全.} 证明 2 只有一句话，值得写全：一次操作只把某一堆从 $a_i$ 变成 $a_i'<a_i$，其余堆不动，于是新的异或和
\[
 S'=(a_1\oplus\cdots\oplus a_{i-1})\oplus a_i'\oplus(a_{i+1}\oplus\cdots\oplus a_n)=S\oplus a_i\oplus a_i'.
\]
因为 $a_i'\ne a_i$，所以 $a_i\oplus a_i'\ne0$，进而 $S'\ne S$。特别地，从 $S=0$（P 态）出发，任何一步都会把 $S$ 变成非零值，即 P 态的后继全是 N 态。

\paragraph{归纳收尾.} 每步石子总数严格减少，游戏必然结束；先手面对 $S\ne0$ 时按证明 1 走一步，把 $S=0$ 的局面交给后手；后手无论怎么走，$S$ 都会变回非零（证明 2），先手再归零，如此反复。由于石子数有限，最终先手把局面做成全零（$S=0$ 的终局）时后手无路可走。故 $S\ne0\iff$ 初始为 N 态。

\section{反 Nim 游戏}
\begin{problem}{反 Nim 游戏}
有 $n$ 堆石子，第 $i$ 堆有 $a_i$ 个，每次可以选择一堆，取走至少一个石子，\textbf{不能行动者胜}。判定初始状态。
\end{problem}

直觉上，游戏判定胜负的规则被取反，会导致问题本质的变化，所以反 Nim 游戏的结论会与 Nim 游戏截然不同，但仔细分析后可以发现事实并非如此。规则反转真正改变的是终局附近的行为：正常情况下取走最后一颗石子的人获胜，反规则下他却输了；但只要场上还留有一堆超过 $1$ 个的石子，先手就能把“要不要取最后一颗”的主动权握在自己手里。

\begin{teach}[反规则题的教学顺序]
先让学生猜“规则反转后结论如何变化”，多数会猜“全部反转”；再用 $c=1$ 的情形说明先手仍握主动权。板书按 $c=0,1,2$ 三行列表、每行填“先手胜条件”，学生能自己发现只有第一行取反。最后提醒：把“谁被迫取最后一颗”的主动权当作分析对象，是所有 misère 题的通用切入点。
\end{teach}
\begin{theorem}[反 Nim 判定]\label{thm:antinim}
设 $S=a_1\oplus\cdots\oplus a_n$，$c$ 为 $a$ 中大于 $1$ 的数的个数。若 $\max_i a_i\le1$，则先手必胜当且仅当 $S=0$；否则先手必胜当且仅当 $S\ne0$。也就是说，只有 $\max_i a_i\le1$ 的情况下胜负条件被取反，其余情况下胜负条件不变。
\end{theorem}
\begin{proof}
考虑对 $c$ 分类讨论。每一步只改变一堆，所以一次操作后 $c$ 至多变化 $1$，这保证了不同情形之间不会“串台”。
\begin{itemize}
\item $c=0$：所有堆都是 $1$，每步只能整堆取走，游戏恰好进行 $c_0$（总堆数）步，取最后一颗石子的人反而输，故先手胜当且仅当 $c_0$ 为偶数，即 $S=c_0\bmod 2=0$。这正是唯一一处胜负条件被真正取反的地方。
\item $c=1$：考虑策略窃取，先尝试把最大值变成 $1$，不行就变成 $0$。具体地，设唯一的大堆有 $A>1$ 个石子、其余 $1$ 堆有 $t$ 个。把 $A$ 变成 $1$ 后，留给后手的是 $t+1$ 个“$1$ 堆”；把 $A$ 变成 $0$ 后，留给后手的是 $t$ 个“$1$ 堆”。$t$ 与 $t+1$ 一奇一偶，而全“$1$ 堆”的局面先手胜当且仅当堆数为偶数（见 $c=0$），所以两种做法必有一种能把必败局面交给后手。这里的策略窃取很关键，它消除了规则取反带来的影响：先手通过选择“把 $A$ 停在 $1$ 还是 $0$”，强行把游戏导回 $c=0$ 的已解情形。
\item $c\ge2$：无论如何操作，操作后都有 $c\ge1$。因此直接按照 Nim 游戏的策略操作即可：$S\ne0$ 时按 Nim 的证明 1 把某一堆取到 $a_i\oplus S$，把 $S=0$ 交给对手。需要确认这一手不会把局面送成全“$1$ 堆”的陷阱：Nim 的归零手只动一堆，动完后大堆数仍 $\ge c-1\ge1$；即使恰剩一堆大于 $1$（$c$ 变成 $1$），由 $c=1$ 情形的分析，$c=1$ 时 $S$ 不可能为 $0$，而我们现在交出的局面满足 $S=0$，矛盾。所以对手永远接不到“全 $1$ 堆”的便宜局面，只能按正常 Nim 的节奏输掉。
\end{itemize}
配合石子总数递减的终止性，三种情况分别给出了先手的必胜应对（或 $S=0$ 时后手的维持策略），结论得证。
\end{proof}

反 Nim 给我们的教训是：\textbf{规则反转只影响“全是 $1$”这种极端局面}，遇到“无法行动者胜”的题目，先检查有没有让对手没法拿到便宜终局的手段。

\section{Nim-k 游戏}
\begin{problem}{Nim-k 游戏}
有 $n$ 堆石子，第 $i$ 堆有 $a_i$ 个，每次可以选择至少一堆、至多 $k$ 堆，每堆取走至少一个石子，不能行动者负。判定初始状态。
\end{problem}

普通 Nim 每步只许动一堆，Nim-k 放宽到至多 $k$ 堆，可以想见判定条件要从“异或”（模 $2$ 的按位和）升级为“模 $k+1$ 的按位和”——因为每一步在同一个二进制位上，最多能把 $k$ 个 $1$ 同时清零。

\paragraph{第一猜想及其错误.} 考虑推广 Nim 游戏的结论，猜想：设 $S$ 表示 $a_{1\ldots n}$ 的 $k+1$ 进制异或和（把 $a_i$ 写成 $k+1$ 进制后逐位相加、每一位模 $k+1$ 不进位），则先手必胜当且仅当 $S\ne0$。很遗憾，这个结论是错误的，因为一个 P 态可能到达另一个 P 态：比如说 $k=2$，初始有两堆石子，大小分别为 $1$ 和 $2$，此时 $S=0$（$1+2\equiv0\pmod 3$）。如果 $A$ 把两堆石子都取完，则操作后仍然有 $S=0$，矛盾。症结在于：$k+1$ 进制下“一次操作”并不是“只改一位”——把 $2$ 堆整堆取走相当于把它从 $2$ 改成 $0$，在 $k+1$ 进制下低位的 $2$ 直接消失，总和的每一位都被打乱，导致我们连“P 态走不到 P 态”都保证不了。

\begin{teach}[猜想--反驳--修复的教学节奏]
Nim-k 最适合演示“数学是如何犯错的”：先让全班对第一猜想（$k+1$ 进制异或）举手投票；再给出反例 $(1,2)$（$k=2$ 时 $S=0$ 却一步取完），请学生指出病根——$k+1$ 进制下一次操作会打乱每一位；最后引出修复（按二进制位、模 $k+1$ 计数）。这套节奏之后可复用到“广义 Nim-k 为什么做不了”：让学生自己检查两段式证明的哪一段断掉。
\end{teach}
\paragraph{修复.} 把 $a_{1\ldots n}$ 写成二进制形式，设 $S$ 表示把这些 01 串看成 $k+1$ 进制数后的异或和（即 $S$ 的第 $j$ 位等于“第 $j$ 个二进制位上为 $1$ 的堆数模 $k+1$”），则先手必胜当且仅当 $S\ne0$。\footnote{原稿此处写作“$S=0$”，按上下文（证明 1 是从 $S\ne0$ 调整到 $S=0$）应为 $S\ne0$，整理时已订正。}修复的要点是把按位运算放在二进制位上进行：一次操作至多影响 $k$ 堆，所以每个数位上至多变化 $k$，正好够不到 $k+1$——这就是把模数取成 $k+1$ 的原因。仍用上面的反例验证：$k=2$，两堆 $1,2$，二进制串为 $01$ 与 $10$，逐位模 $3$ 相加得 $S$ 的两位数位为 $1,1$，$S\ne0$，判定为 N 态；事实也确是如此——$A$ 一步取完两堆即可获胜。

\begin{theorem}[Nim-k 判定]\label{thm:nimk}
设 $S$ 的第 $j$ 位等于“$a_1,\ldots,a_n$ 的二进制第 $j$ 位为 $1$ 的堆数模 $k+1$”，则先手必胜当且仅当 $S\ne0$。
\end{theorem}
\begin{proof}
仍分两个部分，这里的 P 态即 $S=0$ 的局面。
\end{proof}

\begin{proof}[证明 1（$S\ne0$ 的后继中有 P 态）]
从高到低考虑 $S$ 的每一位，设当前考虑到了第 $i$ 位，$(S)_i=j$，已经操作过 $cur$ 堆石子，这些堆中石子数的高位都已经严格减小了，因此低位可以任意修改。那么，考虑把 $(S)_i$ 的 $k+1$ 种取值排成一个环，则如果只用 $cur$ 个堆进行调整，可以得到的是包含 $j$ 的一段弧。如果这段弧就包含 $0$，那直接调整即可。否则，设这段弧是一段区间 $[l,r]$（$1\le l\le j\le r\le k$）。考虑先用 $cur$ 个堆把 $(S)_i$ 调整到 $l$。现在，这 $cur$ 堆石子的数量的第 $i$ 位都为 $0$。又因为 $(S)_i=l$，所以至少存在 $l$ 堆石子的数量的第 $i$ 位为 $1$，再操作其中的 $l$ 堆即可。操作后 $cur$ 会变成 $r$，因为 $r\le k$ 所以操作合法。
\end{proof}

\paragraph{拆开读.} 我们的目标是：从高位到低位逐位把 $S$ 的所有数位清零，全程维护“已经动过的堆（$cur$ 堆）的高位已严格变小”。为什么要这个性质？因为一堆只要高位已经变小，其低位就可以自由改动（无论如何改都不影响已处理好的高位），等于获得了“编辑许可”；没动过的堆高位必须是 $0$，否则它的高位早就该被处理。现在看第 $i$ 位：设它当前的数字和为 $j$，$cur$ 堆在第 $i$ 位可以各自任意取 $0/1$，所以能调出的总和覆盖环上含 $j$ 的一段连续弧（长为 $cur+1$）。若弧内含 $0$，直接调整，第 $i$ 位清零，进入下一位。若弧不含 $0$，说明它是一段不环绕的整数区间 $[l,r]$，$l\ge1$：先把 $cur$ 堆的第 $i$ 位全部调成 $0$（总和降到 $l$），此时数字和为 $l$ 意味着“第 $i$ 位为 $1$ 的堆”恰有 $l$ 个（模 $k+1$ 的意义下至少 $l$ 个），而这些堆都不是 $cur$ 中的堆（$cur$ 的堆第 $i$ 位已是 $0$），它们的高位全为 $0$、可以安全改动；于是再操作其中 $l$ 堆，把它们的第 $i$ 位从 $1$ 改成 $0$，第 $i$ 位总和归零。动过的堆数累计为 $cur+l=r\le k$，操作合法，且这 $l$ 堆从此获得“低位自由”，可以继续处理下一位。如此逐位推进，最终得到 $S=0$ 的局面。

\begin{proof}[证明 2（P 态的后继全是 N 态）]
操作一次后 $S$ 一定会发生变化。
\end{proof}
\paragraph{补全.} 一次操作改变 $c$ 堆（$1\le c\le k$）。取这些堆中发生变化的最高的二进制位 $p$：第 $p$ 位以上所有堆都没有变化；而在第 $p$ 位，凡是变化的堆，其这一位必然是从 $1$ 变成 $0$——否则这堆的值不会变小。因此 $S$ 的第 $p$ 位总和变化了 $-c$，而 $1\le c\le k$ 使得 $-c\not\equiv0\pmod{k+1}$，故 $S$ 必变。特别地，$S=0$ 的局面走一步后 $S\ne0$，P 态的后继全是 N 态。这个“看最高变化位”的论证手法也适用于普通 Nim（$k=1$、模 $2$），是理解后面许多证明的钥匙。

\section{反 Nim-k 游戏}
\begin{problem}{反 Nim-k 游戏}
有 $n$ 堆石子，第 $i$ 堆有 $a_i$ 个，每次可以选择至少一堆、至多 $k$ 堆，每堆取走至少一个石子，\textbf{不能行动者胜}。判定初始状态。
\end{problem}

\begin{theorem}[反 Nim-k 判定]\label{thm:antinimk}
采用定理~\ref{thm:nimk} 中的 $S$。若 $\max_i a_i\le1$，则先手必胜当且仅当 $S\ne1$；否则先手必胜当且仅当 $S\ne0$。
\end{theorem}
\begin{proof}[证明概要]
把反 Nim 游戏和 Nim-k 游戏拼起来，思路与反 Nim 完全平行：先数出大于 $1$ 的堆数 $c$，再分两类看。
\begin{itemize}
\item $c=0$（全是 $1$ 堆）：每步至多取走 $k$ 个“$1$ 堆”，等价于两个人轮流从一堆 $c_0$ 个石子中拿走 $1\sim k$ 个、谁拿最后一个谁输的取石子游戏，先手负当且仅当 $c_0\equiv1\pmod{k+1}$。而全 $1$ 局面下 $S$ 的数位只在最低位有一个非零数字 $c_0\bmod(k+1)$，故这正是“$S=1$”，即先手胜当且仅当 $S\ne1$。
\item $c\ge1$：只要场上还有大堆，就按 Nim-k 的正常策略把 $S$ 归零；反 Nim 一节的经验在这里同样适用——“规则反转”只会在终局附近作祟，而只要有大堆存在，取最后一堆的主动权就握在能维持 $S=0$ 的一方手里。
\end{itemize}
\end{proof}

\begin{pitfall}[自造变体先对拍]
反 Nim-k 这个推广结论的正确性并不平凡。Nim-k 的归零操作一次可能同时动多堆，反 Nim 中“一步只动一堆所以大堆数不减”的论证不再成立，归零后局面可能恰好变成全 $1$ 的尴尬形状，因此“拼接”的细节比反 Nim 微妙得多（讲义作者亦未展开完整证明）。竞赛中遇到这类自造变体，稳妥的做法是先写个暴力对拍验证公式，再放心使用；附录的自测程序对有限范围做了验证。
\end{pitfall}

\chapter{Sprague--Grundy 理论}
\section{SG 函数}
SG 函数做的事情，是把任意公平组合游戏“折算”成 Nim 游戏：它给每个状态赋一个非负整数，称为该状态的 \textbf{SG 值}，表示这个状态“相当于”一个多大的 Nim 堆，从而把 Nim 的胜负结论搬到千变万化的游戏上。当然，这种折算是有代价的，可能会损失部分信息。

具体来说，首先，每个公平组合游戏都可以表示成一个 DAG，游戏规则是：初始时某个节点上放置着一枚棋子，每次操作可以把棋子移动到任意一个出点上，不能移动者负。为什么可以这样表示？回顾 1.3 节：公平组合游戏的状态与合法操作天然构成有向图，而“双方决策空间相同”意味着这个图对两名玩家一视同仁，棋子“属于谁”无关紧要。接下来，考虑对每个节点 $u$，把“棋子在 $u$ 上”的游戏状态对应到一个大小为 $\SG(u)$ 的 Nim 堆上——堆的大小就是这个局面的 SG 值，用来衡量它的“战斗值”。

\begin{teach}[怎样想到 SG 值？]
不要一上来就写 $\operatorname{mex}$。先让学生对“每次取 $1$ 或 $2$ 颗”手算胜负表（$n=0$ 到 $6$），发现 $0,1,2$ 周期；再问“这个周期数字在数什么”，引出“相当于多大的 Nim 堆”；$\operatorname{mex}$ 只是把直觉写成算子。课后追问：把“取 $1$ 或 $2$”改成“取 $1$ 或 $3$”，周期还在吗？（目的：打掉机械记忆，强制回到递推本身。）
\end{teach}
\begin{definition}[SG 函数与 mex]\label{def:sg}
设 $\Np(u)=\{v_1,v_2,\ldots,v_k\}$ 为 $u$ 的后继集合，$\mex(X)$ 表示非负整数集合 $X$ 中未出现的最小非负整数。定义
\[
 g(u)=\mex\{g(v):v\in\Np(u)\},
\]
函数值 $g(u)$ 称为状态 $u$ 的 \textbf{SG 值}（文献中也叫 Grundy 值）。
\end{definition}

先把这一节的来龙去脉说清：整套理论由 Sprague 与 Grundy 在 1930 年代各自独立建立，故得名 Sprague--Grundy 理论；把赋值看成定义在状态集上的函数，就得到定义~\ref{def:sg} 的 SG 函数。注意大小写：Nim 是游戏专名，必须大写（Nim 游戏、Nim 堆、Nim 和），SG 是 Sprague--Grundy 的缩写。做题时看到“SG 值”“Grundy 值”两个词，心里想着同一个东西即可：给公平组合游戏的每个状态赋一个非负整数，值为 $0$ 的状态恰是 P 态；多个独立游戏合成和游戏时，总值为各值之异或（SG 定理，见下）。

\begin{example}[热身：每次取 $1$ 或 $2$ 颗石子]
设大小为 $n$ 的堆的 SG 值为 $f_n$，则 $f_0=0$（无后继），$f_1=\mex\{f_0\}=1$，$f_2=\mex\{f_0,f_1\}=2$，$f_3=\mex\{f_1,f_2\}=0$，$f_4=\mex\{f_2,f_3\}=1$，依此类推周期为 $3$，而“$f_n=0$ 当且仅当 $3\mid n$”恰对应众所周知的结论（$3$ 的倍数时后手必胜）。
\end{example}

$\mex$ 的定义保证了两件事：第一，$0,1,\ldots,g(u)-1$ 这些值都出现在 $u$ 的后继中，即 $u$ 存在对应大小为 $0\ldots(g(u)-1)$ 的 Nim 堆的后继——好比 Nim 中可以把一堆从 $g(u)$ 取到任意更小值；第二，$u$ 不存在对应大小为 $g(u)$ 的 Nim 堆的后继——Nim 中一步不可能把一堆保持原样。这两条恰好复刻了 Nim 游戏的核心性质：每个 SG 值为 $a$ 的状态，其“能到达的后继的 SG 值集合”包含着通往 $0$ 的路（当 $a>0$），却不含 $a$ 自身。

\begin{theorem}[SG 值刻画 P 态]\label{thm:sgzero}
在正常规则（无法操作者输）下，$u$ 是 P 态当且仅当 $g(u)=0$。
\end{theorem}
\begin{proof}[证明（归纳）]
$g(u)=0$ 时后继 SG 值都不为 $0$（全是 N 态）；$g(u)>0$ 时存在 SG 值为 $0$ 的后继（N 态的定义成立）。
\end{proof}

损失的信息在哪里？不能保证 $u$ 不存在对应大小\textbf{大于} $g(u)$ 的 Nim 堆的后继。也就是说，一个“弱”状态也可能有“强”后继，只是这个后继本身是 N 态、理性玩家不会主动走进去而已。在只分胜负的正常规则下这无伤大雅，但一旦规则变化（例如反 Nim 的“不能动者胜”、或一步可以同时动多个子游戏），这些隐藏的后继就可能颠覆结论——这正是后面“广义反 Nim”“广义 Nim-k”两节讨论的麻烦来源。

\begin{teach}[mex 与强后继的课堂演示]
现场画一个小 DAG：给某点安排后继 SG 值 $\{0,1,7\}$，让学生算 $\operatorname{mex}$（得 $2$），再问“$7$ 这个强后继会不会坏事”。正常规则不会（没人走过去），同步和与反规则会——这一演示为广义反 Nim 与同步和两节埋下伏笔，比事后补救有效得多。
\end{teach}
\begin{pitfall}[SG 值不排除“强后继”]
$mex$ 只管“最小的未出现值”。$g(u)=2$ 的状态完全可以有 $g=7$ 的后继：正常规则下没人会走过去，可一旦规则反转或一步可动多个子游戏，“强后继”就会被利用。判断一个新规则能否沿用 SG 理论，先问一句：证明里哪里用到了“后继的 SG 值都不超过当前值”？其实哪里都没用到——这正是隐患所在。
\end{pitfall}

\section{广义 Nim 游戏}
\begin{problem}{广义 Nim 游戏}
有 $n$ 个独立游戏 $G_{1\ldots n}$，第 $i$ 个的 SG 值为 $a_i$，每次可以选择一个游戏移动一步，不能行动者负。判定初始状态。
\end{problem}

这就是 Nim 游戏的“角色扮演版”：把 Nim 里的每堆石子换成一个任意的公平组合游戏，把“取石子”换成“在这个游戏里走一步”，SG 值扮演堆大小的角色。

\begin{corollary}[广义 Nim 判定]
设 $S=a_1\oplus a_2\oplus\cdots\oplus a_n$。广义 Nim 游戏中先手必胜，当且仅当 $S\ne0$。
\end{corollary}
\begin{proof}[证明思路]
与 Nim 相同的两段式：第一段与 Nim 的证明 1 逐字相同（把“取到 $a_i\oplus S$”理解成“把第 $i$ 个游戏走到 SG 值为 $a_i\oplus S$ 的后继”，而这样的后继由 SG 值的定义保证存在，因为 $a_i\oplus S<a_i$）；第二段因为一步只改变一个游戏、其 SG 值必变，异或和随之改变，同样与 Nim 的证明 2 相同。
\end{proof}

有了广义 Nim 游戏，读者应该已经能猜到下一步：如何证明“$n$ 个游戏之和的 SG 值等于各游戏 SG 值的异或”？这正是下一节“Nim 和”的内容。

\section{Nim 和：SG 定理}
即 Sprague--Grundy 定理，它是组合博弈论的顶梁柱：\textbf{任意有限公平组合游戏的和游戏，其 SG 值等于各游戏 SG 值的异或}。这句话回答了三个问题：为什么 SG 值恰好是个“数”而不是更复杂的对象？为什么异或运算无处不在？以及（对做题人最实际的）为什么可以把大游戏拆成小游戏分别算 SG 值再异或。

先考虑两个独立游戏 $G_1,G_2$ 的和 $G_1+G_2$：玩家轮流选择其中一个游戏并操作一次。$\SG(G_1+G_2)$ 是多少？

\begin{theorem}[SG 定理]\label{thm:sgtheorem}
$\SG(G_1+G_2)=\SG(G_1)\oplus\SG(G_2)$。
\end{theorem}
\begin{proof}
设 $\SG(G_1)=n$，$\SG(G_2)=m$。首先 $G_1+G_2$ 的后继状态中不含 $n\oplus m$，否则不妨设操作了 $G_1$，则 $G_1$ 的后继状态含 $n$，与 SG 值的定义矛盾。

下面证明 $\forall x<n\oplus m$，存在 $G_1+G_2$ 的后继状态为 $x$。即证 $G_1$ 存在后继状态 $m\oplus x$，或 $G_2$ 存在后继状态 $n\oplus x$。即证 $n>m\oplus x$ 或 $m>n\oplus x$。

设 $k=\topbit(x\oplus n\oplus m)$，则 $(x)_k=0$，$(n\oplus m)_k=1$。不妨设 $(n)_k=0$，$(m)_k=1$。分情况讨论：
\begin{itemize}
\item 如果 $\topbit(n\oplus m)>k$：
\begin{itemize}
\item 如果 $n>m$，则 $n>m\oplus x$；
\item 如果 $n<m$，则 $m>n\oplus x$。
\end{itemize}
\item 如果 $\topbit(n\oplus m)=k$：有 $m>n\oplus x$。
\end{itemize}
\end{proof}

\paragraph{直觉压缩.} 第一段是“合游戏的后继 SG 值不含 $n\oplus m$”，因为一步只能动一个游戏，而 SG 值的定义禁止“原地踏步”；第二段是“小于 $n\oplus m$ 的值全都可达”，它依赖二进制的一个基本事实——对任意 $x<n\oplus m$，异或和的最高差异位 $k=\topbit(x\oplus n\oplus m)$ 上 $n,m$ 中必有一方的这一位是 $1$（另一方是 $0$），而这一位为 $1$ 的一方的 SG 值较大，可以把它“调”到需要的值。两段合起来，后继能取到的 SG 值集合覆盖了 $0,1,\ldots,(n\oplus m)-1$ 的全部值、却不含 $n\oplus m$，由 $\mex$ 的定义（只关心最小的未出现值，更大的值即使出现也不影响结论）即得 $\SG(G_1+G_2)=n\oplus m$。归纳到 $n$ 个游戏的和即得广义 Nim 游戏的结论，这也是下一节“验证”的根据。

SG 定理解决了“分治合”三步中的“合”。那么我们只需要考虑如何“分”和“治”：先把大的游戏分成独立的小游戏，再分别解决（对每个小游戏、每个状态求 SG 值）。这个范式是例题 [ABC278G]、[CF1439E]、[HNOI2014] 等题目共同的骨架。

\section{验证广义 Nim 游戏的结论}
考虑 $n$ 个游戏的和：由 SG 定理，和的 SG 值等于各分量 SG 值的异或，即 $S$；由 SG 值的性质（定理~\ref{thm:sgzero} 的“$g=0$ 恰为 P 态”），整个和游戏初始状态是 N 态当且仅当 $S\ne0$。这个结论与广义 Nim 游戏的结论相同——两条路殊途同归，也算是对 SG 定理的一次“体检”。

\section{广义反 Nim 游戏}
\begin{problem}{广义反 Nim 游戏}
有 $n$ 个独立游戏 $G_{1\ldots n}$，第 $i$ 个的 SG 值为 $a_i$，每次可以选择一个游戏移动一步，当 $a_1=a_2=\cdots=a_n=0$ 时当前执步者胜。
\end{problem}

先解释这个定义。若按字面把反 Nim 推广为“无法行动者胜”，那么在和游戏里“无法行动”意味着每个游戏都走到了自己的终点；把“终点态”集体标注成胜利并不优雅，因为一个游戏的终点（SG 值为 $0$）本身是个正常可到达的状态，把它视为“输”会让所有分析都变得反直觉。于是定义改为：当所有分量都归零（$a_1=\cdots=a_n=0$）时，当前执步者胜。这与反 Nim 的“无法行动者胜”一脉相承（在纯 Nim 的情形下，全零恰是无法行动），但它把“轮到自己时全零”变成了赢局而非输局，效果是：谁把局面做成全零，谁就输了——因为紧接着轮到的对手会立刻获胜。

广义反 Nim 游戏的胜负判定条件与反 Nim 游戏的直接推广稍有不同。这样定义的好处是 SG 值为 $0$ 的状态一定是 N 态——在正常规则下“所有游戏都已归零”是输局，这里它却是赢局——但代价是没法再用 SG 定理来验证结论了：SG 定理本身仍然是正确的，但是合并两个游戏不能递归到结果相同的子问题。说得更直白些：SG 定理的递归证明依赖“子游戏的和”与“子游戏本身”之间的同构关系，而这里的胜负规则取决于“是否所有分量同时为零”这一全局性质，把 $G_1+G_2$ 切成两半后，“某一半全零”并不等于“某一半结束”，递归就此断裂。因此本题的验证只能回到反 Nim 的分类讨论思路，而不能指望把每个游戏化简成一个数。

\begin{corollary}[广义反 Nim 判定]
设 $S=a_1\oplus\cdots\oplus a_n$。若 $\max_i a_i\le1$，则先手必胜当且仅当 $S=0$；否则先手必胜当且仅当 $S\ne0$。
\end{corollary}

直觉与反 Nim 如出一辙：$a_i\le1$ 时每个游戏至多还能走一步，胜负退化为“谁先被迫把局面做成全零”；只要某个 $a_i\ge2$，先手就可以把局面控制在异或非零的安全区，把“做成全零”的最后一手留给对方。

\begin{pitfall}[$\max a_i\le1$ 的真实含义]
这里 $\max_i a_i\le1$ 的含义与石子大小无关，而是每个独立游戏“剩余步数的上限”至多一步。SG 值小的游戏也可能还剩很多步可走（回想“强后继”），所以判定的适用对象是“每个分量都临近终点”的局面，不是数值小的局面。
\end{pitfall}

\section{广义 Nim-k 游戏？}
感觉做不了。因为 $u$ 可能存在对应大小大于 $\SG(u)$ 的 Nim 堆的后继，所以对于一个 P 态，操作大于一堆石子后可能到达另一个 P 态。展开说：Nim-k 一步可以同时动至多 $k$ 个游戏，而 SG 值对应关系只能保证“单游戏一步内不出现 SG 值相同的后继”；当多个游戏同时变动时，各自的隐藏后继（那些 SG 值更大的后继）可能互相补偿，使总和不变。反 Nim 一节尚且能靠“一步只动一堆”的性质维持证明，Nim-k 连这个性质都没有，两段式证明的第二段直接失效，故本讲义不收录它的结论——这本身也是一个提醒：\textbf{推广结论前先检查证明的每个关键步骤在推广后是否依然成立}。

\section{小测验：同步和}
\begin{problem}{小测验}
有 $n$ 个独立游戏 $G_{1\ldots n}$，每次对所有可以移动的游戏\textbf{同时}移动一步，不能移动者负。判定初始状态。
\end{problem}

\hint 这题的“和”不再是 SG 意义下的和（那里每次只动一个游戏），而是\textbf{同步和}：每轮所有还没结束的游戏都要被迫走一步。它更像“每局游戏各自按自己的时钟赛跑，谁先在自己的局里无路可走谁输”的并行竞赛，因此胜负不再由 SG 值决定，而由“谁的游戏先耗完”决定。

\sol 观察到对每个 N 态的独立游戏，$A$ 按照必胜策略操作不劣：如果 $A$ 选择走到一个 N 态点，则即使这个游戏结束得最晚，也是 $B$ 赢。

所以，$A$ 的最优策略是：对每个 N 态的独立游戏，按照赢得最慢的必胜策略操作；对每个 P 态的独立游戏，按照输得最快的策略操作。直觉是：N 态的游戏是 $A$ 的“资产”，资产要尽可能长久地续命（把结束时刻拖到其他游戏之后）；P 态的游戏是 $A$ 的“负债”，负债要尽早爆掉（把“轮到 $B$ 时这局已经结束”的局面尽快交给对方）。双方在所有游戏上同时行动，收益互不干扰，所以最优策略就是对每个游戏独立取最优，再比较“哪个游戏最后结束”。

对每个独立游戏的每个状态，DP 预处理出它是 P 态还是 N 态，以及将来还要走多少步。如果步数最长的独立游戏是 N 态则 $A$ 胜，否则 $B$ 胜。请读者自行验证这个结论与“每个游戏独立地按最优策略走到尽头”是一致的。

\section{分层例题与解析}

以下例题的共同骨架是“分、治、合”：识别独立子游戏，对每个子游戏求 SG 值，用 SG 定理异或结算；个别题先判断“能不能合”，不能合时退回直接分类。

\begin{problem}{[ABC278G] Generalized Subtraction Game}
本题为交互题。

给定整数 $N,L,R$。你将与评测系统进行如下游戏：

有 $N$ 张编号为 $1$ 到 $N$ 的卡片放在场上。

先手和后手轮流进行如下操作：选择一组整数 $(x,y)$，满足 $1\le x\le N$，$L\le y\le R$，并且编号为 $x,x+1,\ldots,x+y-1$ 的 $y$ 张卡片都还在场上，然后将这些卡片从场上移除。

不能进行操作的一方判负，另一方获胜。

你需要选择先手或后手。然后，在你选择的回合与评测系统进行游戏，并取得胜利。

数据范围：$N\le2000$。
\end{problem}

\hint 取走一段连续卡片就把局面切成左右两段的和游戏，先写 SG 递推；直接 DP 是 $O(N^3)$，再观察“取走正中间一段后镜像”可以覆盖 $R>L$ 的全部情形，只剩 $L=R$ 需要 $O(N^2)$ 的 DP。

\sol \paragraph{分析.} 一次操作是取走一段连续的卡片，把一段连续区间切成左右两段互不相干的区间——这正是“分治合”的“分”：整局游戏是若干连续区间段的和游戏，所以先直接上 SG 值。设 $f_n$ 表示连续 $n$ 张卡片（一段完整的区间）对应的 SG 值，则移除一个区间后剩下左右两段，由 SG 定理，新局面的 SG 值是两段 SG 值的异或。把 $(x,y)$ 理解为被移除区间的左右端点，则
\[
 f_n=\mex_{L\le y-x+1\le R}\bigl(f_{x-1}\oplus f_{n-y}\bigr),
\]
其中下标范围的意思是枚举所有长度在 $[L,R]$ 内的可移除子区间 $[x,y]$。但是直接 DP 复杂度是 $O(N^{3})$ 的（枚举 $n$、左端点 $x$、长度 $y-x+1$ 共三维），无法通过。

\paragraph{考虑下模仿棋.} 考虑先手视角。镜像策略的适用条件是：先手第一步能移走正中间的一段，使剩下的左右两段长度相等且结构完全对称，之后无论 B 在某一侧怎么取连续区间，A 都在另一侧的对称位置取镜像区间。因为两侧被取走的区间始终成对，对称性永不破坏，B 能取的区间其镜像一定还在，于是 A 永不卡壳、必胜。难点只在于第一步是否存在：当 $r>l$ 时可取长度至少有两种且连续，一定凑得出与 $n$ 同奇偶的长度 $y\in[l,\min(r,n)]$（前提是 $l\le n$，即场上确实存在可移除的区间；若 $n<l$，双方一步都走不了，初始就是 P 态）。取中间 $y$ 张，剩下左右各 $(n-y)/2$ 张，两侧严格对称。所以当 $r>l$ 且 $l\le n$ 时，如果 $n$ 是偶数，就移除中间的连续偶数张卡片，否则移除中间的连续奇数张卡片。之后，不断操作 B 操作的位置的对称位置即可。因此初始状态是 N 态（对评测的每个 $N$ 都选择先手并照此执行）。否则（$l=r$，可移长度唯一），奇偶凑不齐，镜像不可用，只能老老实实 DP 出 $f_N$ 再判胜负——此时每层只需枚举左端点，DP 的复杂度是 $O(N^{2})$ 的，可以通过（$N\le2000$）。

\src{https://atcoder.jp/contests/abc278/tasks/abc278_g}

\begin{problem}{[CF1439E] Cheat and Win}
考虑一个 $[0,10^{9}]\times[0,10^{9}]$ 的网格。如果 $x\ \mathrm{and}\ y=0$，则称格子 $(x,y)$ 是“好格子”。我们以所有好格子为顶点构建一张图，并在所有相邻的好格子之间连边。可以证明，这张图是一棵树。钦定 $(0,0)$ 是根节点。

AB 两人博弈。最初，部分好格子为黑色，其余为白色。每一回合，玩家选择一个黑色好格子及其若干祖先（可以选 $0$ 个祖先），并将这些格子的颜色反转（白变黑，黑变白）。如果某一玩家无法操作（即所有好格子均为白色），则判负。可以证明，这个游戏一定会在有限步内结束。

最初，所有格子都是白色。你会得到 $m$ 对格子，对于每一对格子，将它们之间的简单路径上的所有格子染成黑色。注意，这里是直接染黑，不是反转颜色。

B 想要获胜，于是决定作弊。他可以在游戏开始前多次进行如下操作：选择一个格子，将它到树根的路径上的所有顶点颜色反转。求 B 至少需要做多少次作弊操作。

数据范围：$m\le10^{5}$。
\end{problem}

\hint 四层结构：分形树与深度公式；把“翻转路径”等价改写成“独立棋子”；单枚棋子的 SG 值是 $2^{\dep_u}-1$；把初始 SG 值的二进制里 $1$ 的连续段数修到零。

\sol \paragraph{这题的结构.} 整道题的骨架可以拆成四层：(1) 好格子构成一棵分形树，深度即 $x+y$；(2) “翻转某点到根路径”的博弈通过异或对合转化成“独立棋子”的公平组合游戏；(3) 用树上差分求出初始局面的 SG 值；(4) B 的作弊次数等于把该 SG 值修到 $0$ 所需的翻转次数。下面逐层展开。

\paragraph{为什么图是一棵树？} 观察到满足 $x\ \mathrm{and}\ y=0$ 的 $(x,y)$ 构成一个分形：设 $G_k$ 表示 $[0,2^{k})\times[0,2^{k})$ 中的好格子的集合，初始时 $G_0=\{(0,0)\}$。依次考虑二进制下的每一位，因为 $x,y$ 的第 $k$ 位不能同时为 $1$，所以 $G_{k+1}$ 由 $G_k$ 本身、$G_k$ 向下平移 $2^{k}$ 个单位长度、$G_k$ 向右平移 $2^{k}$ 个单位长度三个形状相同的部分组成。不难发现这是一棵树，且每个格子的父亲要么是它上方的格子，要么是左侧的格子（这两个格子不可能同时是好的）。基于这个刻画，容易观察到 $(x,y)$ 的深度是 $x+y$：每走一步恰好让 $x+y$ 减少 $1$（向上或向左），而从 $(0,0)$ 到 $(x,y)$ 也恰好要走 $x+y$ 步。深度公式是整个“虚树 + 差分”方案的基石：它让我们不需要显式建树就能定位任意格子。

\paragraph{考虑给定一棵树和每个节点的初始颜色，如何判断谁会赢.} 黑点之间是不独立的，不好处理。因为异或比较容易造对合，考虑这样一个转化：初始时在每个黑点处放一枚棋子，每次可以选择一个点 $u$，移除 $u$ 上的一枚棋子，并对 $u$ 的每个祖先 $v$，选择是否在 $v$ 上放一枚棋子，不能操作者负。显然这个问题的结果与原问题是相同的，但是每枚棋子之间是独立的了。为什么独立了？因为一次操作只涉及“以 $u$ 为根的祖先链”这条路径，不同棋子只通过“同一条链上的加减”耦合，而“在 $v$ 上选择放或不放”正好对应颜色翻转的线性性（异或），把每条链上的操作分解开，就得到一堆彼此独立的单棋子游戏——这是异或对合把“集合翻转”变成“单个对象操作”的典型例子。

\paragraph{求单枚棋子的 SG 值.} 考虑点 $u$ 上的棋子对应的 SG 值，不难观察到是 $2^{\dep u}-1$（单个棋子每次只能“上行”并任选是否在祖先处卸子，其可达集合的结构恰好给出 $2^{\dep u}-1$ 这个值，可用对深度的归纳严格验证），所以 $u$ 父链上每个点的 SG 值的异或和就等于
\[
 2^{0}\oplus2^{1}\oplus\cdots\oplus2^{\dep u-1}=2^{\dep u}-1,
\]
（不同深度的 $2$ 的幂互不重叠，异或等于求和）。所以，如果能求出初始状态对应的 SG 值，设其二进制表示为 $\sum_{i\ge0}a_i2^{i}$，考虑差分，则最少的作弊次数为 $\sum_{i\ge0}\iva{a_i\ne a_{i+1}}$（约定 $a_{i+1}$ 在最高位之后视为 $0$；这个式子统计的其实是二进制串中“$1$ 的连续段”的个数——一次作弊操作恰好能清除一段连续低位的组合效应，因此段数就是下界且可达到，详细的位运算推导此处从略，读者可结合小例子验证）。结合格子 $(x,y)$ 的深度为 $x+y$，只需要建出 $m$ 条路径端点的虚树，然后树上差分。

为了模拟建虚树的过程，我们需要支持两个操作：比较两个节点的 DFS 序；求两个节点的 LCA。先规定一个 DFS 序：先向下递归再向右递归，则这两个问题都可以通过对分形结构分治解决。

\cost 时间复杂度 $O(m\log m)$。总结：本题把“博弈”外衣层层剥掉之后，剩下的其实是一个数据结构题——建虚树、树上差分、按深度定位，博弈含量集中在“每枚棋子 SG 值 $=2^{\dep}-1$”这一步推导上。

\src{https://codeforces.com/problemset/problem/1439/E}

\begin{problem}{Gra}
有 $m$ 个格子排成一行，初始时有 $n$ 个格子 $a_1,\ldots,a_n$ 里有棋子。每次可以选择一枚棋子，将它向右移动到第一个没有棋子的格子中，不能移动者负。判定初始状态。
\end{problem}

\hint 把“棋子”翻译成“空格之间的间隔”：空格是墙，墙与墙之间夹着的棋子是台阶上的石子，这是阶梯 Nim 的伪装。

\sol \paragraph{转化为阶梯 Nim.} 这题的妙处在于把“棋子”翻译成“空格之间的间隔”。把每个空格看成一面“墙”，棋子按空格自然分段：从右往左数第 $i$ 个空格与第 $i+1$ 个空格之间夹着的棋子，就看成一堆“第 $i$ 级台阶上的石子”。一次操作“把一枚棋子向右移到第一个空格”，效果是把某一段最靠左的棋子挪到段尾的空位——反映到台阶上，正好是某个台阶的石子数减一、相邻台阶的石子数加一（具体方向取决于编号，但搬运结构完全相同）。考虑以空格为阶梯分界线，从右往左第 $i$ 个空格左侧、第 $i+1$ 个空格右侧的每枚棋子对应第 $i$ 级台阶上的一个石子，转化为阶梯 Nim：于是胜负判定直接套用定理~\ref{thm:stairnim} 台阶游戏的结论——只有某个固定奇偶性的台阶（按上述编号的“奇数级”）参与异或，异或和为 $0$ 时后手必胜。

这个转化与下一题 [SDOI2019] 移动金币互为镜像：那里金币只能向左挪、这里棋子只能向右跳，左右翻转后两者结构一致。建议读者先拿小例子（如 $m=5$、三枚棋子）手工对拍，体会“空格分段”的计数方式，再回头理解为什么只有奇偶性固定的段参与判定。

\begin{problem}{[SDOI2019] 移动金币}
一个 $1\times n$ 的棋盘上最初摆放有 $m$ 枚金币。其中每一枚金币占据了一个独立的格子，任意一个格子内最多只有一枚金币。

AB 两人博弈。每次可以选择一枚金币，并将其向左移动任意多格，且至少移动一格。金币不能被移出棋盘，也不能越过其它金币。不能行动者负。

求有多少种初始状态是 N 态。答案对 $10^{9}+9$ 取模。

数据范围：$n\le1.5\times10^{5}$ 且 $m\le50$。
\end{problem}

\hint 先做判定：又是阶梯 Nim，观察对象是金币之间的空隙；判定化为“奇数项空隙异或和是否为 $0$”后，问题变成带异或约束的计数，用数位 DP。

\sol \paragraph{先做判定：又是阶梯 Nim.} 金币只会左移，我们把“空隙”当作石子来观察。把金币之间的空格以及最左端金币左侧的空格都数出来：设 $c_i$ 表示从右往左第 $i$ 个金币左侧、第 $i+1$ 个金币右侧的空格数（$i=1$ 对应最右金币与次右金币之间；$i=m$ 对应最左金币与棋盘左边界之间；最右金币右侧的空格不算——它们永远只增不减，是游戏的“垃圾场”）。若某金币左移 $d$ 格，它左侧的空隙减少 $d$、右侧的空隙增加 $d$：在从右往左的编号下，这就是“第 $i$ 级台阶上的 $d$ 个石子被搬到第 $i-1$ 级”，与台阶游戏的操作完全同构（最右金币右侧的空隙扮演“地面”第 $0$ 级，只进不出）。考虑以金币为阶梯分界线，从右往左第 $i$ 个金币左侧、第 $i+1$ 个金币右侧的每个空格对应第 $i$ 级台阶上的一个石子，转化为阶梯 Nim。由定理~\ref{thm:stairnim}：只有奇数级台阶参与异或，故初始状态是 N 态当且仅当第 $1,3,5,\ldots$ 项空隙的异或和不为 $0$（终局时所有金币挤在左端、各级台阶全空，恰好对应异或和为 $0$ 的 P 态，两端的刻画闭合）。

\paragraph{再变成计数问题.} 金币位置 $(p_1<\cdots<p_m)$ 与空隙数列 $(c_1,\ldots,c_m)$ 一一对应，且 $c_1+\cdots+c_m=p_m-m\le n-m$（棋盘右侧还可能剩空格）。于是问题转化为：求有多少个长为 $m$、和不超过 $n-m$ 的数列（即 $m$ 个非负整数，和不超过 $n-m$），满足第奇数项的异或和为 $0$——这类“少变量 + 异或约束 + 总量约束”的组合问题用数位 DP：按二进制位从高到低枚举，状态记录“奇数项的异或是否已归零的进位情况”以及“总和是否还贴着上界”，其中 $m\le50$ 个小变量的按位贡献可以用组合数批量统计。时间复杂度 $O(m^{3}\log n)$（对 $10^{9}+9$ 取模）。注意要求的是 N 态数量，而“第奇数项异或和为 $0$”计数的是 P 态，最后用方案总数减去它即可。

\src{https://www.luogu.com.cn/problem/P5363}

\begin{problem}{[HNOI2014] 江南乐}
小 A 是一个名副其实的狂热的回合制游戏玩家。在获得了许多回合制游戏的世界级奖项之后，小 A 有一天突然想起了他小时候在江南玩过的一个回合制游戏。

游戏的规则是这样的，首先给定一个数 $F$，然后游戏系统会产生 $T$ 组游戏。每一组游戏包含 $N$ 堆石子，小 A 和他的对手轮流操作。每次操作时，操作者先选定一个不小于 $2$ 的正整数 $M$（$M$ 是操作者自行选定的，而且每次操作时可不一样，但是 $M$ 不能大于选中的堆中的石子数量），然后将任意一堆数量不小于 $F$ 的石子分成 $M$ 堆，并且满足这 $M$ 堆石子中石子数最多的一堆至多比石子数最少的一堆多 $1$（即分的尽量平均，事实上按照这样的分石子方法，选定 $M$ 和一堆石子后，它分出来的状态是固定的）。当一个玩家不能操作的时候，也就是当每一堆石子的数量都严格小于 $F$ 时，他就输掉。补充：先手从 $N$ 堆石子中选择一堆数量不小于 $F$ 的石子分成 $M$ 堆后，此时共有 $N+M-1$ 堆石子，接下来小 A 从这 $N+M-1$ 堆石子中选择一堆数量不小于 $F$ 的石子，依此类推。

小 A 从小就是个有风度的男生，他邀请他的对手作为先手。小 A 现在想要知道，面对给定的一组游戏，而且他的对手也和他一样聪明绝顶的话，究竟谁能够获得胜利？

数据范围：$T\le100$，$N\le100$，$F\le10^{5}$，每堆石子数量 $\le10^{5}$。
\end{problem}

\hint “把一堆分成若干堆”是分裂型操作，分裂后各堆独立，是标准的分治合：算单堆 SG 值再异或；转移里的 $M$ 用整除分块压缩，相同 SG 值的偶数个副本互相抵消后每块只剩两个候选。

\sol \paragraph{分治合 + 整除分块的组合拳.} “把一堆分成若干堆”是典型的分裂型操作，但分裂后各堆互不相干、继续分裂时彼此独立，所以每一堆又是独立的子游戏，整局就是一堆子游戏的和——先算单堆 SG 值，再异或所有堆。观察到每一堆石子是独立的，考虑计算每一堆石子的 SG 值。设大小为 $n$ 的堆的 SG 值为 $f_n$。则当 $n<F$ 时，$f_n=0$（太小不能分，也不能再动，终局即 P 态）。否则，枚举 $2\le m\le n$，可以把大小为 $n$ 的堆划分成 $m$ 个大小极差不超过 $1$ 的堆：设 $n=qm+r$，则分成了 $r$ 个大小为 $q+1$ 的堆，$m-r$ 个大小为 $q$ 的堆。因为这些堆又是独立的，所以后继的 SG 值等于各独立游戏 SG 值的异或（SG 定理；注意是异或而不是算术和）。这里可以立刻省一半力气：相同 SG 值的偶数个副本互相异或抵消——$r$ 个 $f_{q+1}$ 只有当 $r$ 为奇数时才贡献 $f_{q+1}$，$m-r$ 个 $f_q$ 同理，所以后继的 SG 值等于
\[
 \bigl(\iva{2\nmid r}\,f_{q+1}\bigr)\oplus\bigl(\iva{2\nmid m-r}\,f_q\bigr).
\]
$f_n$ 就等于所有后继的 SG 值的 $\mex$。

剩下的问题是如何不枚举全部 $O(n)$ 个 $m$。观察到 $r=n-qm$，所以固定 $n,q$，$r$ 的奇偶性只和 $m$ 的奇偶性有关——$q$ 相同时（$m$ 落在一个整除块内），$r$ 的奇偶性由 $m$ 的奇偶性唯一决定，而每个后继 SG 值只依赖 $r$ 与 $m-r$ 的奇偶性，于是整块内所有 $m$ 产生的候选 SG 值至多只有 $2$ 种（对应 $m$ 的奇偶）。考虑整除分块，则每一块只需要考虑至多两个后继，把 $O(n)$ 个 $m$ 压成 $O(\sqrt n)$ 块。

\cost 时间复杂度 $O(V\sqrt V)$（$V\le10^{5}$），每组游戏的总判定再异或各堆的 $f$ 即可。

\begin{teach}[给学生的实验任务]
读完移动金币后让学生完成三件事：(1) 用 $n\le20$ 写暴力计数，验证“奇数项异或为零”的 P 态刻画；(2) 用数位 DP 重算并对比；(3) 把 $m$ 改成 $m+1$，看公式哪一项变化。目的：让“先判定、再计数、最后对拍”成为肌肉记忆。对拍不过时，先查编号方向（从右往左），再查两端（最右金币右侧算不算）。
\end{teach}
\chapter{拓展：非公平组合游戏与超现实数}
\section{公平这个前提用在哪里}
前面的 Sprague--Grundy 理论能成立，靠的是 1.3 节的公平性前提：同一局面下，双方的决策空间相同。这个前提一旦撤掉——国际象棋里白方只能动白子，轮到白方与轮到黑方时能走的棋完全是两本账——整个框架都会塌掉：
\begin{itemize}
\item \textbf{P/N 二分塌了。}1.3 节说过，非公平游戏的局面可能有四种结果：先手必胜、后手必胜、玩家甲必胜、玩家乙必胜。造一个最简单的例子：桌上有三颗石子，两颗只属于甲、一颗只属于乙，轮到谁谁从自己的石子堆里取走任意正数颗（自己的石子取完了就不能动，判负）。甲先手时先取 $1$ 颗留 $1$ 颗，乙只能拿光自己的 $1$ 颗，甲再拿光自己的，乙无路可走；乙先手时乙拿光自己的 $1$ 颗，甲拿光自己的 $2$ 颗，乙依然无路可走。结论是甲必胜，与先后手无关——“你是谁”比“谁先走”更本质，P/N 两个标签根本不够用。
\item \textbf{SG 函数塌了。}“给局面赋一个非负整数”的前提是双方对这个局面拥有同等选择权；一旦棋子认主人，同一个局面在甲、乙眼里是两本不同的账，一个数装不下两个人的视角。
\end{itemize}
所以非公平游戏需要更精细的语言。本节把 Conway 的整套框架补上：它在数学上完整地包含公平游戏这个特例（第 5 章的内容全部可以看作它的退化情形），这也是原稿提到超现实数的原因。路线是：先给游戏一个“双视角记号”，再定义负、和与比较三种运算，最后把“能比大小”的游戏称为数，并看看数如何“做算术”。

\section{双视角记号}
\begin{definition}[游戏的双视角记号]
把一个游戏写成左右两个选项集合的对
\[
 G=\{G^L\mid G^R\},
\]
其中 $G^L$ 是玩家甲（约定称 Left，记 L）走一步能到达的局面集合，$G^R$ 是玩家乙（Right，记 R）的。
\end{definition}
与全篇一致，约定游戏在有限步内结束、无法行动者负。公平游戏恰好是 $G^L=G^R$ 的特例——所以 Nim 游戏与 SG 值全都住在这个记号里。记号本身平凡，不平凡的是在全体游戏上定义三种运算：
\begin{itemize}
\item \textbf{负}：$-G$ 是把 $G$ 的左右选项对调（再逐个取负），含义是“两人角色互换”。$G+(-G)$ 一定是后手必胜：先手在 $G$ 里走什么，后手就在 $-G$ 的镜像位置回应什么，直到先手无路可走——这是“下模仿棋”在记号层面的化身。
\item \textbf{和}：$G+H$ 表示两局并列，轮到谁走、谁挑其中一局走一步；一局结束后只剩另一局。和满足交换律与结合律，配上负运算后“全体游戏”构成一个群：这是后面一切算术的合法性来源，也是全篇“分治合”中“合”的最终形态。
\item \textbf{比较}：$G=H$ 当且仅当 $G-H$ 是后手必胜；$G>H$ 当且仅当 $G-H$ 无论谁先手都是 L 胜；$G<H$ 对称；$G$ 与 $H$ 不可比（记 $G\parallel H$）当且仅当 $G-H$ 是先手必胜。
\end{itemize}
注意“谁先手”由外部局面决定，但“差是后手胜”是游戏的内在性质，与下棋的人无关——这保证了用“值”做算术的客观性。

\section{四个最小的游戏}
从最小的开始看，每个游戏后面的括号里都写着它到底是谁的优势：
\begin{itemize}
\item $0=\{\mid\}$：双方都没有合法操作，先手直接输——零优势；
\item $1=\{0\mid\}$：只有 L 能走一步，且只能走到 $0$——好比桌上有一颗只属于 L 的石子。R 先手时无路可走；L 先手时拿走石子，R 依然无路可走。无论谁先手都是 L 胜，故 $1>0$；
\item $-1=\{\mid 0\}$：把 $1$ 的左右对调，R 必胜，$-1<0$；注意 $-(-1)=1$，而 $1+(-1)$ 见下一条；
\item $1+(-1)=0$：L、R 各有一颗自己的石子。先手拿走自己那颗，后手拿走自己那颗，先手便无路可走——后手胜，恰好对应“$1$ 与 $-1$ 相加抵消”；一般地 $G+(-G)$ 都是后手胜。
\end{itemize}

看起来像在绕圈子，但立刻有收获：$1$ 可以读作“L 领先一步”，$-1$ 是“R 领先一步”，$1+(-1)$ 是“领先被抵消”。把这个直觉推广开：L 独占 $a$ 颗石子（每步可取任意正数颗、只能 L 取）的游戏值就是整数 $a$（把“独占 $a$ 颗”与数 $a$ 相减，差是后手胜，可归纳验证）；若 L、R 各独占 $a$、$b$ 颗，整局等于 $a+(-b)=a-b$，胜负只看差值。

\begin{example}[验证 $a=2,b=1$]
L 先手时先取 $1$ 颗留 $1$ 颗，R 只能拿光自己的 $1$ 颗，L 再拿光自己的，R 无路可走；R 先手时 R 拿光自己的 $1$ 颗，L 接着拿光自己的 $2$ 颗，R 依然无路。两种先后手都是 L 胜，与值 $2-1=1>0$ 一致。
\end{example}

注意这盘棋里没有共享堆、没有异或，博弈内容与 Nim 完全不同，胜负却被“差一个整数”干净地描述——这是“把局面算成数”的第一个成功案例，也是“游戏是一个数”这句话的含义。

\section{四种结果与值的符号}
值的符号与胜负结果的对应可以一次说清：
\begin{itemize}
\item $G>0$（正游戏）：L 必胜，先手后手都不怕；
\item $G<0$（负游戏）：R 必胜；
\item $G=0$（零游戏）：后手必胜；
\item $G\parallel 0$（与 $0$ 模糊）：先手必胜。
\end{itemize}
四行正好把 1.3 节非公平游戏的四种结果（先手必胜、后手必胜、甲必胜、乙必胜）按“值”重新分类。回头看公平游戏的退化：SG 值为 $0$ 的局面是后手胜，即 $G=0$；SG 值非零的局面是先手胜，即 $G\parallel0$——公平游戏从不出现 $G>0$ 或 $G<0$。这从根上解释了为什么公平游戏只需要问“是否为 $0$”：双方的选项完全对称，谁也积累不出“专属优势”，能积累的只有“模糊的先手优势”，于是 Sprague--Grundy 理论只关心一个比特的信息；非公平游戏则不然，L 可以拥有实打实的“正资产”。

\section{数：冷游戏的算术}
游戏世界里存在“不是数”的游戏（下一节的 $\ast$ 就是），Conway 把数定义为如下特殊类型的游戏。

\begin{definition}[数]
左右选项满足“每个左选项都小于每个右选项”的游戏称为\textbf{数}。
\end{definition}

$0=\{\mid\}$ 是第一个数；此后由简到繁依次造出整数（$1=\{0\mid\}$、$2=\{1\mid\}$、$\ldots$）、分数（$\tfrac12=\{0\mid 1\}$：L 的选项是 $0$、R 的选项是 $1$，值恰好卡在两者之间，即“比 $0$ 好、比 $1$ 差”的半目优势）、乃至无穷与无穷小（$\omega=\{0,1,2,\ldots\mid\}$）。全体数构成一个可以加减乘除的全序结构，这就是\textbf{超现实数}。数之间没有模糊性——任意两个数都能比出大小，这与“棋子认主人”造成的混乱正相反，因此取值是数的游戏也叫\textbf{冷游戏}：双方都没有抢着先手的动机，局面的价值像温度计读数一样稳定，先手可能反而吃亏。

冷游戏的好处是可以直接做算术：几个数游戏并成一局，总值就是各数的和；和为正则 L 占优、为负则 R 占优、为零则后手胜。

\begin{example}[$\tfrac12+\tfrac12=1$]
两局 $\{0\mid 1\}$ 并玩：L 先手时把任一局走到 $0$，剩下 $0+\tfrac12=\tfrac12>0$，R 无论怎么走都救不回来；R 先手时把任一局走到 $1$，剩下 $1+\tfrac12=\tfrac32>0$，L 更是稳赢。两个“半目”合成“一目”，算术与博弈严丝合缝。
\end{example}

这就是“超现实数是研究非公平组合游戏的有力工具”的技术内容：把局面翻译成值，把博弈翻译成算术。

\section{模糊游戏：star 与 SG 值}
再看那个“不是数”的游戏：
\[
 \ast=\{0\mid 0\},
\]
双方都只有一步可走，先手走到 $0$ 即胜。$\ast$ 与 $0$ 不可比：任何一方先手都能赢，谁也不能声称自己“更大”——它正是 SG 值为 $1$ 的公平游戏，即大小为 $1$ 的 Nim 堆。更有意思的是 $\ast+\ast=0$：两局 $\ast$ 并玩，先手把其中一局走到 $0$，后手把另一局走到 $0$，先手无路可走——后手胜。对照 SG 值：$1\oplus1=0$，分毫不差。一般地，SG 值为 $n$ 的公平游戏都记作 $\ast_n$（大小为 $n$ 的 Nim 堆），且
\[
 \ast_n+\ast_m=\ast_{(n\oplus m)},
\]
SG 定理（Nim 和）在这套记号下浓缩成一行：公平游戏相加，值按异或合成。异或为什么无处不在，答案就在这里：不是异或神奇，而是公平游戏在“游戏群”里的结构恰好是 $\ast$ 家族的加法——$\ast_n+\ast_n=\ast_{(n\oplus n)}=\ast_0=0$，任何公平游戏与自己相加都是后手胜，与“SG 值异或自己得 $0$”完全一致。

\begin{pitfall}[口径与边界]
这里有一个需要厘清的口径问题。如果把“超现实数”理解成 Conway 的整套“游戏即数”体系，那么 SG 值确实是其中定义清晰的一个大家族——原稿把 SG 值视为超现实数的一个特例，正是取这个广义口径；若取狭义（只指那些能比大小的数），SG 值不是数——$\ast$ 不是 $0$，$\ast_2$ 也不是实数 $2$。“公平游戏”与“数”是这套体系里两个不同的子类：前者左右选项相同、典型地与 $0$ 模糊，后者左小右大、彼此全能排出大小。两种口径都能自洽；关键只是别把两者混为一谈——$\ast$ 的 SG 值是 $1$，但 $\ast$ 本身并不等于数 $1$。

这套框架的边界也要说清楚。它默认正常规则（无法行动者负）与和游戏的封闭性：一旦规则反转（misère 系，如反 Nim 游戏），“差是后手胜即相等”这套公理不再成立，这正是反 Nim 的结论需要另起炉灶、“广义反 Nim”无法再用 SG 定理验证的深层原因；而“一步同时动多个子游戏”之类的变体同样会离开这套框架（见“广义 Nim-k 游戏？”一节的讨论）。另外，对一般的模糊游戏求“值”并不容易：两个游戏比大小在计算上可以非常困难，竞赛题几乎只会碰到冷游戏退化成整数的情形（即独占堆那种例子）。因此本节的实际用途有三：看懂公平游戏为什么只需要 SG 值；遇到“这不是公平游戏”的题目时能立刻认出 SG 不适用；以及在极少数构造题里，用这套记号做小规模的手工推演。
\end{pitfall}

\begin{teach}[拓展板块的课堂取舍]
建议把第 6 章安排为一到两次选修课，不必全讲。必带内容：双视角记号与四个最小游戏、单纯性法则加竹竿算值（$\tfrac12$、$\tfrac14$ 两个例子足够）、$\uparrow$ 的前两条性质与例 1、例 5。可弃内容：温度的任何严格理论、大棋盘 Domineering、口径之争的细节。教具：围棋子当蓝红边、骨牌当场玩 $2\times2$ Domineering，比幻灯片有效。
\end{teach}
\section{拓展：数从游戏中诞生——单纯性法则}
6.5 节把 $\tfrac12=\{0\mid 1\}$ 当作事实抛了出来：凭什么 $\{0\mid 1\}$ “就等于” $\tfrac12$？本节给出这套造数法的完整规则。它出自 Berlekamp、Conway、Guy 的《Winning Ways for Your Mathematical Plays》第一章（另见 Conway《On Numbers and Games》），是超现实数体系的“施工规范”。

\begin{definition}[出生日]
从 $\{\mid\}$（即 $0$）出发，每用已有的游戏构造一次新的数，就说新数“晚出生一天”：$0$ 出生于第 $0$ 天，$1=\{0\mid\}$ 与 $-1=\{\mid 0\}$ 出生于第 $1$ 天，$\{0\mid 1\}$ 出生于第 $2$ 天，依此类推。
\end{definition}

\begin{theorem}[单纯性法则]\label{thm:simplicity}
设每个左选项与每个右选项都是已构造出的数，且每个左选项严格小于每个右选项，则 $\{L\mid R\}$ 的值是\textbf{严格大于所有左选项、严格小于所有右选项的那个出生最早}的数。
\end{theorem}

严格证明见《Winning Ways》第二章与《On Numbers and Games》；本书把它当“施工规范”使用，重点是算对值。“相等”仍按 6.2 节的定义检验：两个游戏相等，当且仅当它们的差是后手必胜。

\begin{example}[二进分数的诞生]\label{ex:dyadic}
$\{0\mid 1\}$：严格介于 $0$ 与 $1$ 之间、出生最早的数是 $\tfrac12$（第 $2$ 天），故 $\{0\mid 1\}=\tfrac12$，这正是 6.5 节直接验证过的等式。类似地
\[
 \{1\mid 2\}=\tfrac32,\qquad \{0\mid \tfrac12\}=\tfrac14,\qquad \{-1\mid 0\}=-\tfrac12 .
\]
按此规范造出的数先是全体二进分数（分母为 $2$ 的幂），很晚之后才轮到 $\tfrac13$、$\pi$ 乃至无穷与无穷小——超现实数的整个数轴是“从游戏里长出来”的。
\end{example}

\paragraph{伐木游戏（Hackenbush）竹竿.} 单纯性法则最经典的应用对象是\textbf{蓝红伐木游戏}：图由若干条边接在地面（或接在别的边上）构成，Left 只能砍蓝边、Right 只能砍红边；一条边被砍掉时，它上方悬空的整个部分一起倒下；无法行动者负。一根由 $n$ 条蓝边叠成的竖直竹竿，Left 砍最底一条就让 Right 无事可做，值为 $\{n-1\mid\}=n$——这正是 6.3 节“L 独占 $n$ 颗石子”画成图的样子。混色竹竿才是新意。\textbf{蓝底红梢}（自下而上：蓝、红）：Left 砍蓝底，整根倒掉，得 $0$；Right 砍红梢，剩下一条蓝边，值 $1$。于是
\[
 \text{蓝红}=\{0\mid 1\}=\tfrac12 .
\]
再接一条红（自下而上：蓝、红、红）：Left 仍只有砍底一条路，得 $0$；Right 砍中间的红剩值 $1$、砍顶端的红剩 $\tfrac12$，Right 取对自己更小的 $\tfrac12$，故
\[
 \text{蓝红红}=\{0\mid \tfrac12\}=\tfrac14 .
\]
对称地，红底蓝梢的值为 $-\tfrac12$。竹竿每向上多接一节，值的绝对值就折半——二进分数就这样长在竹竿上，6.5 节“两个半目合成一目”的围棋直觉由此获得精确模型。

\section{拓展：开关、热博弈与温度}
若 $x\ge y$ 且 $x,y$ 都是数，则 $\{x\mid y\}$ 不满足“左选项 $<$ 右选项”，因此\textbf{不是数}，称为\textbf{开关}（switch）。开关是“热”的：先手总能从自己的选项里拿到好处。

\begin{example}[开关的两端]\label{ex:switch}
下一节 Domineering 的 $2\times2$ 棋盘值为 $\{1\mid -1\}$：Left 先手走成 $1$，Right 先手走成 $-1$，谁先谁赢、与 $0$ 模糊，记作 $\pm1$。一般地 $\{3\mid -1\}$ 读作“Left 先手可拿到 $3$、Right 先手可拿到 $-1$”的局面。
\end{example}

开关虽然不是数，实践中常用两条粗略刻画（选学内容，不作严格定理）：\textbf{均值} $\tfrac{x+y}{2}$ 大致描述“双方轮流各走一步后平均谁赚”；\textbf{温度} $\tfrac{x-y}{2}$ 大致描述“先手权在这一点上值多少”。由此得到一条实用的启发式——\textbf{先玩最热的分支}：和游戏里有多个分支时，先手应当先把先手权花在温度最高的分支上。

\begin{example}[先玩最热]\label{ex:hotstrat}
$\{3\mid -1\}+\tfrac12$：Left 先手走开关，得 $3+\tfrac12=\tfrac72$，胜；Right 先手也走开关，得 $-1+\tfrac12=-\tfrac12$，胜。双方的最优首着都是开关而不是 $\tfrac12$——先手权花在温度高的地方才不浪费。注意该局面与 $0$ 模糊：均值 $\tfrac{3-1}{2}+ \tfrac12=\tfrac32>0$ 只说明“长期轮流平分下 Left 净赚”，并不决定单局胜负；这正是开关不能当数用的原因。
\end{example}

\begin{center}
\begin{tabular}{cccl}
\toprule
开关 & 均值 & 温度 & 一句话画像\\
\midrule
$\{1\mid -1\}$ & $0$ & $1$ & 纯先手权，谁先谁赢（Domineering $2\times2$）\\
$\{2\mid -2\}$ & $0$ & $2$ & 更烫的纯先手权（兑现支票 $X_2$）\\
$\{3\mid -1\}$ & $1$ & $2$ & 均值偏向 Left 的争夺点\\
\bottomrule
\end{tabular}
\end{center}

再次强调口径：均值与温度是\textbf{启发式坐标}，不是相等关系；两个均值相同的开关未必相等。它们的价值在于给“先玩哪个分支”排序，以及快速估计“这步先手权值几目”。

\begin{teach}[微小值的课堂演示]
黑板两列对照：“$\ast$ 是先手权，与 $0$ 模糊；$\uparrow$ 是额外的一丝优势，严格大于 $0$”。请两位学生实际对玩 $\uparrow+\ast$ 与 $\uparrow+\downarrow$，记录胜负后回到两段式分析。课堂追问：为什么 $\ast+\ast=0$ 而 $\uparrow+\ast$ 不等于 $0$？（前者是两个先手权互相抵消，后者两个分支里只有一个提供先手权。）学生能答对这一问，拓展课的目的就达到了。
\end{teach}
\section{拓展：微小值 up、down 与 up-star}
数的世界之外还有一群“比 $0$ 精细、又比任何正数都小”的微小值，它们刻画的是先手权以外的第二层优势。

\begin{definition}[up 与 down]
\[
 \uparrow=\{0\mid \ast\},\qquad \downarrow=\{\ast\mid 0\}=-\uparrow .
\]
\end{definition}

以下三件事实都可以用 6.3 节的“分谁先手”手工验证（完整验证见本节例 3）：
\begin{itemize}
\item $\uparrow>0$：无论谁先手都是 Left 胜——$\uparrow$ 是货真价实的正优势，只是很小；
\item $\uparrow\parallel\ast$：微小值与先手权互为模糊，谁也不能压倒谁；
\item $\uparrow+\downarrow=0$：镜像策略的再现——Right 在 $\downarrow$ 中的每一步恰好回应 Left 在 $\uparrow$ 中的那一步。
\end{itemize}

\begin{theorem}[微小性]\label{thm:up}
$\uparrow>0$，且对任意正数 $x$ 都有 $0<\uparrow<x$；$\downarrow$ 对称地满足 $-x<\downarrow<0$。
\end{theorem}
\begin{proof}[证明思路]
对 $x$ 的出生日归纳：正数 $x$ 的左选项中必有非负数，逐一与 $\uparrow$ 的选项比较即可；完整归纳见《Winning Ways》。博弈含义是：$\uparrow$ 是“一丝真实优势”，它永远输给哪怕最小的正数优势，却永远压过 $0$。围棋贴目里“半目”的精确定价用的正是这类微小值。
\end{proof}

\section{拓展：两个经典游戏——Domineering 与蓝红伐木}
\begin{problem}{Domineering（纵横多米诺）}
在方格棋盘上，Left 每轮竖放一张 $1\times 2$ 骨牌，Right 每轮横放一张 $2\times 1$ 骨牌，格子不能重叠，先无法放置者负。求给定棋盘的值。
\end{problem}

它在 1.3 节的意义下显然非公平：同一张空棋盘，轮到 Left 与轮到 Right 时能放的骨牌方向不同。小棋盘可以直接按定义算值：
\begin{itemize}
\item $1\times n$ 单行棋盘：Left 永远无处落子。$1\times2$：Right 放满即空，值 $\{\mid 0\}=-1$；$1\times3$：Right 放完必剩 $1\times1$，右选项为 $0$，值 $\{\mid 0\}=-1$；$1\times4$：Right 的右选项为放中间（断成两个 $1\times1$，值 $0$）与放一端（剩 $1\times2$，值 $-1$），删去被支配的 $0$ 得 $\{\mid -1\}=-2$。归纳可得
\[
 v(1\times n)=-\bigl\lfloor n/2\bigr\rfloor .
\]
\item 单列 $2\times1$：Left 竖放占满即空，Right 无处可放，值 $\{0\mid\}=1$。
\item $2\times2$：Left 竖放剩单列 $2\times1$（值 $1$），Right 横放剩单行 $1\times2$（值 $-1$），故
\[
 v(2\times2)=\{1\mid -1\}=\pm1 ,
\]
是一枚开关：谁先手谁赢。
\end{itemize}
把前几个值排成一行：$v(1\times1)=0$（谁都不能放），$v(1\times2)=-1$，$v(1\times3)=-1$，$v(1\times4)=-2$，$v(1\times5)=-2$，$v(1\times6)=-3$，恰是 $-\lfloor n/2\rfloor$。更大的棋盘（如 $3\times n$、$8\times8$ 整盘）值要靠计算机枚举，《Winning Ways》里有成套分析；本书只取最小情形演示方法。

\begin{problem}{蓝红伐木游戏（Blue--Red Hackenbush）}
棋盘改为一般图形：边接在地面或别的边上，Left 只砍蓝边、Right 只砍红边，被砍边的悬空部分整体倒下，无法行动者负。求给定图形的值。
\end{problem}

竹竿只是最简单的图形。《Winning Ways》的一个基础结论是：\textbf{任意蓝红伐木局面的值都是数}——蓝红世界里没有模糊性，这正是它适合给超现实数“接生”的原因。若加入双方都能砍的\textbf{绿边}，模糊性立刻回来：单独一条绿边 $\{0\mid 0\}=\ast$，又回到第 5 章的世界。蓝、红、绿三色版伐木游戏因此可以看成“数的世界 + 模糊游戏”的完整舞台，也是检验 6.4 节四种结果分类的最佳沙盘。

\begin{teach}[值表与判断题组的课堂用法]
「拓展小结：一张值表」的十行表建议做成卡片：正面是游戏，背面是值与结果类，课上抽认。例 9、例 10 的判断题组适合四人小组抢答：每组先独立判断真伪，再逐线验证——重点盯住 (1)(3)：一份 $\uparrow$ 压不住先手权、两份才行，学生自己发现这一点，比讲十遍微小性定理都管用。常见误区预告：不少学生会把 $\uparrow+\ast$ 当成正数（“微小值嘛”），例 10(3) 专治这一条。
\end{teach}
\paragraph{绿竿值多少.} 单条绿边就是 $\ast$。两节绿竿：Left 砍底，整竿倒下，Right 面对空局面直接负；Right 对称地先手也胜。其左右选项同为 $\{0,\ast\}$，记作 $\{0,\ast\mid 0,\ast\}$，与 $0$ 模糊。更有意思的是绿与数的对峙：$1+\ast$ 中 Left 先手把 $\ast$ 走成 $0$，Right 面对 $1$ 无事可做；Right 先手同样只能把 $\ast$ 走成 $0$，Left 再把 $1$ 走成 $0$。两种先后手都是 Left 胜——\textbf{任何数都压得过先手权}（“数与模糊游戏比较”的通例，见《Winning Ways》）。同理 $\tfrac14+\uparrow+\downarrow>0$ 是定理~\ref{thm:up} 微小性的直接应用。于是绿边的加入让伐木游戏变成“数的部分 + 模糊的部分”的混合体：数的部分按算术结算，模糊部分用第 5 章的工具处理——这正是三色伐木被视为“连接第 5、6 两章之桥”的原因。

\begin{pitfall}[开关不是数]
$\pm1=\{1\mid -1\}$ 不满足“左选项 $<$ 右选项”，不能参与“比大小”意义上的算术；比较热博弈时用“均值”只是启发式，不是相等关系（例~\ref{ex:hotstrat} 中均值 $\tfrac32>0$ 的局面 Right 先手照样赢）。若一道题只关心胜负而出现了开关，先把开关当成独立分支按先后手讨论，别急着套数轴。
\end{pitfall}

\section{拓展小结：一张值表}
把本章出现过的值集中成一张表，复习时先盖住后两列，自己说出每个游戏的值与结果类。

\begin{center}
\begin{tabular}{llll}
\toprule
游戏 & 值 & 结果类 & 一句话\\
\midrule
$\{\mid\}$（空局面） & $0$ & 零游戏，后手胜 & 无事可做\\
$\{0\mid\}$（L 一颗石子） & $1$ & 正数，L 必胜 & L 领先一步\\
$\{\mid 0\}$（R 一颗石子） & $-1$ & 负数，R 必胜 & R 领先一步\\
$\{0\mid 1\}$（蓝红竹） & $\tfrac12$ & 正数，L 必胜 & 半目优势\\
$\{0\mid \tfrac12\}$（蓝红红） & $\tfrac14$ & 正数，L 必胜 & 四分之一目\\
$\{1\mid -1\}$（Domineering $2\times2$） & $\pm1$ & 模糊，先手胜 & 开关，不是数\\
$\{0\mid 0\}$（单条绿边） & $\ast$ & 模糊，先手胜 & 先手权\\
$\{0\mid \ast\}$（up） & $\uparrow$ & 正，L 必胜 & 一丝真实优势\\
$\{\ast\mid 0\}$（down） & $\downarrow$ & 负，R 必胜 & $-\uparrow$\\
$\uparrow+\downarrow$ & $0$ & 零游戏，后手胜 & 镜像抵消\\
\bottomrule
\end{tabular}
\end{center}

读表三句话：值为\textbf{数}的局面是“冷”的，直接做加减法、比大小就能判定胜负；值为\textbf{模糊}的局面必须分“谁先手”讨论；\textbf{微小值}介于两者之间——比 $0$ 强、比任何正数弱，而且它本身不是数（$\uparrow$ 的右选项 $\ast$ 不是数，不满足数的构造规则）。这张表也是 6.4 节“四种结果”分类的完整对照：正数与负数对应 L、R 的“专属优势”，零与模糊对应先后手的攻防。

\section{分层例题与解析}
以下例题把本章拓展的三件工具——单纯性法则、开关、微小值——各练一遍；解析默认读者已读过本章相应小节。

\begin{problem}{例 1：混色竹竿求值}
求蓝红红、红蓝两根竹竿的值，并判定局面“蓝红红 $+$ 红蓝”（两根并玩）谁必胜。
\end{problem}

\sol 蓝红红：Left 砍蓝底得 $0$；Right 砍中间红剩 $1$、砍顶端红剩 $\tfrac12$，取 $\tfrac12$，故值为 $\{0\mid\tfrac12\}=\tfrac14$。红蓝：Right 砍红底得 $0$；Left 砍蓝梢剩 $-1$，故值为 $\{-1\mid 0\}=-\tfrac12$。两根并玩，值的和为 $\tfrac14+(-\tfrac12)=-\tfrac14<0$：数相加仍是数，负数意味着无论谁先手都是 Right 胜。验证：Left 先手无论动哪根，Right 都按数的策略应对即可；这体现了“冷游戏可以直接做算术”。

\begin{problem}{例 2：Domineering 小局拼合}
求 $1\times4$、$2\times1$ 与 $2\times2$ 的值（直接引用上一节的结论亦可），并判定两个局面谁必胜：(a) $2\times1+1\times4$；(b) $2\times2+1\times2$。
\end{problem}

\sol 值分别为 $1$、$-2$、$\pm1$。(a) $1+(-2)=-1<0$，是数，Right 无论先后手都胜。(b) $\pm1+(-1)$：Left 先手把开关走成 $1$，得 $1+(-1)=0$，Right 无路可走，Left 胜；Right 先手把开关走成 $-1$，得 $-2$，Left 无路可走，Right 胜。与 $0$ 模糊，先手必胜。注意 (b) 的两个分支“一个是开关、一个是负数”，不能把开关当 $0$ 或 $\pm1$ 随意代入数轴相加，必须按先后手分情况——这是热博弈与冷游戏在操作上的分水岭。

\begin{problem}{例 3：验证 $\uparrow$ 的三条性质}
验证 $\uparrow>0$、$\uparrow\parallel\ast$、$\uparrow+\downarrow=0$。
\end{problem}

\sol \textbf{$\uparrow>0$：}Left 先手，把 $\uparrow$ 走成 $0$，Right 无路可走；Right 先手，只能把 $\uparrow$ 走成 $\ast$，Left 把 $\ast$ 走成 $0$，Right 仍无路可走。两种先后手 Left 都胜。\textbf{$\uparrow+\downarrow=0$：}Right 的每一步都能镜像回应——Left 动 $\uparrow$ 的某支，Right 就在 $\downarrow$ 的对应位置回应；Left 动 $\downarrow$ 的左选项（到 $\ast$），Right 把新出现的 $\ast$ 走成 $0$，局面回到 $\uparrow$。逐线验证可得后手必胜，与“互为相反游戏”一致。\textbf{$\uparrow\parallel\ast$：}考虑 $\uparrow+\ast$。Left 先手：走 $\ast\to0$，留 $\uparrow$ 给 Right，Right 走成 $\ast$，Left 走成 $0$，Right 负。Right 先手：走 $\uparrow\to\ast$，留 $\ast+\ast$ 给 Left，Left 把一支 $\ast$ 走成 $0$，Right 把另一支走成 $0$，Left 无路可走，Right 胜。先手必胜，故与 $0$ 模糊。三条性质合起来的直观：$\uparrow$ 是“压过 $0$ 却压不过先手权”的微小正优势。

\begin{problem}{例 4：兑现支票游戏}
定义一列游戏 $X_0=0$，$X_{n}=\{X_{n-1}+1\mid X_{n-1}-1\}$。求 $X_1,X_2,X_3$ 的值，并判定局面 $X_2+\tfrac12$ 与 $0$ 的关系。
\end{problem}

\sol $X_1=\{1\mid -1\}=\pm1$，即 6.8 节见过的开关。$X_2$ 的左选项是 $X_1+1$：与 $0$ 模糊的游戏加 $1$ 后 Left 先手拿到 $2$，故左选项值为 $2$；对称地右选项值为 $-2$，即 $X_2=\{2\mid -2\}=\pm2$。归纳即得 $X_n=\pm n$：这游戏像一张每天翻倍违约金的支票，谁先动手谁按票面拿钱。$X_2+\tfrac12$：Left 先手走开关得 $2+\tfrac12=\tfrac52$，胜；Right 先手走开关得 $-2+\tfrac12=-\tfrac32$，胜——与 $0$ 模糊，且双方的最优首着都是开关（先玩最热）。均值 $\tfrac{0+1}{2}=\tfrac12>0$ 只说明轮流平分下 Left 长期净赚，不改变单局谁先谁赢；开关不是数，这正是原因。

\begin{problem}{例 5：竹竿拼合与镜像}
求局面“蓝 $+$ 红”（两根单色竹竿并玩）与局面“红蓝 $+$ 蓝红红”的值，并判定胜负。
\end{problem}

\sol 蓝竿值 $1$、红竿值 $-1$，和为 $0$：后手必胜，且策略就是 6.2 节的下模仿棋——Right 在红竿上的每次砍截回应 Left 在蓝竿上的对应砍截。红蓝值 $-\tfrac12$、蓝红红值 $\tfrac14$，和为 $-\tfrac14<0$：是数，Right 无论先后手都胜。两题都不需要逐线分析——冷游戏直接做算术，这正是“值为数”的省力之处；也只有值为数时才能这么做。

\begin{problem}{例 6：Domineering 拼合与均值的边界}
求 $2\times1+2\times1$ 与 $2\times2+2\times1$ 的值（各分量值见本章 Domineering 一节），并判定胜负。
\end{problem}

\sol $2\times1$ 值 $1$，两列并玩得 $1+1=2>0$：Left 大胜，与先后手无关。$2\times2+2\times1$ 值为 $\pm1+1$：Left 先手把开关走成 $1$，得 $2$，胜；Right 先手把开关走成 $-1$，得 $0$，Left 面对全空局面无路可走，Right 胜。与 $0$ 模糊。注意 $\pm1$ 的均值为 $0$、整体均值 $1>0$，但 Right 先手照样赢——均值可以指导“先玩哪个分支”（先玩最热），不能代替对先后手的逐一分析。

\begin{teach}[ch6 例题节讲评建议]
例 4 的支票与例 7 的开关加法建议板书时并排放：$\pm1$、$\pm2$、$\pm3$ 一列写下去，学生自然会问“开关族是否自成体系”——这正是《Winning Ways》开关与温度一章的入口，讲多深由课时决定。例 5 与例 9(d) 只用加减法，适合请基础较弱的同学上台完成。例 10 的 (1)(3) 各留十分钟逐线验证：宁可少讲一题，不能跳过验证环节。
\end{teach}

\begin{problem}{例 9：数与模糊的对峙}
判定四个局面与 $0$ 的关系：(a) $\ast+1$；(b) $\ast+\ast$；(c) $\tfrac14+\uparrow+\downarrow$；(d) 蓝红红 $+$ 蓝红 $+$ 红（三根竹竿并玩）。
\end{problem}

\hint 数的部分与模糊的部分分开结算；(b) 想想两个先手权；(d) 全是数，直接做算术。

\sol (a) $>0$：Left 先手走 $\ast\to0$，Right 面对 $1$ 无路可走；Right 先手同样只能走进 $1$，随后 Left 收官。数压过先手权。(b) $=0$：先手在一支 $\ast$ 上走到 $0$，后手在另一支上走到 $0$，先手遂无路可走——两个先手权互相抵消。(c) 用 $\uparrow+\downarrow=0$ 化简得 $\tfrac14>0$，L 必胜；若去掉 $\downarrow$，结论不变，但理由换成微小性 $\uparrow<\tfrac14$。(d) 值为 $\tfrac14+\tfrac12+(-1)=-\tfrac14<0$，全是数，Right 无论先后手都胜。(a)(b)(c)(d) 恰好演示了拓展板块的完整工作流：先认出哪些分量是数并直接相加，剩下的模糊部分再按先后手逐线分析。

\begin{problem}{例 10（自拟）：判断题组}
判断真伪并说明理由：
(1) $\uparrow+\uparrow+\ast>0$；
(2) $\tfrac12+\ast$ 与 $0$ 模糊；
(3) $\uparrow+\ast$ 与 $0$ 模糊；
(4) $\uparrow+\uparrow+\downarrow$ 的值为 $\uparrow$；
(5) 蓝红红 $+\ast$ 与 $0$ 模糊。
\end{problem}

\hint 每题先分清“数 / 模糊 / 微小”三种成分；(1)(3) 要分先后手逐线验证，别背结论。

\sol (1) \textbf{真}：Left 先手先消 $\ast$；Right 砍任何一支 $\uparrow$ 都留下 $\uparrow+\ast$，Left 消 $\ast$ 剩 $\uparrow$，Right 走成 $\ast$，Left 收官，Right 负。Right 先手走 $\uparrow\to\ast$，得 $\uparrow+\ast+\ast=\uparrow$，Left 照样赢。(2) \textbf{伪}：数压过先手权，$\tfrac12+\ast>0$。(3) \textbf{真}：Left 先手消 $\ast$ 剩 $\uparrow$，Right 走成 $\ast$，Left 收官，Right 负；Right 先手走 $\uparrow\to\ast$ 得 $\ast+\ast=0$，Left 无路可走，Right 胜——先手必胜，模糊。它与 (1) 的差别值得细品：一份 $\uparrow$ 压不住先手权，两份才行。(4) \textbf{真}：$\uparrow+\downarrow=0$，故 $\uparrow+\uparrow+\downarrow=\uparrow+(\uparrow+\downarrow)=\uparrow$。(5) \textbf{伪}：蓝红红值 $\tfrac14$ 是数，$\tfrac14+\ast>0$。五题连读的教训：(2)(5) 靠“数压过模糊”直接结算，(3)(4) 提醒微小值仍是模糊家族的近亲——加法虽好，逐线验证不可省。

\begin{problem}{例 11（自拟）：换色竹竿}
把三节竹竿自下而上接成“蓝、红、蓝”，求其值并判定结果类。
\end{problem}

\hint 逐个颜色考虑“砍这里会剩下什么”，再按左右选项的支配关系取极值。

\sol Left 的全部选项：砍顶节剩蓝红（$\tfrac12$）；砍中间红节，顶节随之倒下，剩底节蓝（$1$）；砍底节得 $0$——取最大，左选项为 $1$。Right 的同样三个选项取最小，右选项为 $0$。故值为 $\{1\mid 0\}$：左选项大于右选项，\textbf{不是数}，且与 $0$ 模糊——L 先手砍红节留 $1$，R 无路可走；R 先手砍底节留 $0$，L 无路可走，先手必胜。同一根竹竿，只把颜色顺序从“蓝红红”换成“蓝红蓝”，就从冷游戏变成了先手权游戏：Right 的最优砍法从“砍到最小正值”变成“直接砍断命根”，颜色结构决定了谁能抢到先手。这也解释了为什么“任意蓝红局面都是数”里“蓝红”二字一个都不能少。

\begin{problem}{例 7：开关的加法}
求 $\pm2+\pm1$ 的值（其中 $\pm k=\{k\mid -k\}$），并判定局面与 $0$ 的关系。
\end{problem}

\sol 逐项相加：Left 先手时把 $\pm2$ 走成 $2$，再处理 $\pm1$ 最好可拿到 $2+1=3$；Right 先手对称地拿到 $-3$。归纳可证两个“纯开关”之和仍是纯开关：$\pm2+\pm1=\{3\mid -3\}=\pm3$。与 $0$ 模糊——先手赢 $3$、后手输 $3$，温度高达 $3$，双方的最优首着都应花在这枚最烫的开关上。开关加开关仍是开关，这是开关族最好的性质；一旦混入 $\ast$ 或 $\uparrow$ 这样的非数非开关项，就必须回到逐分支分析。

\begin{problem}{例 8：结果类分类}
对下列游戏，指出它与 $0$ 的关系（大于、小于、等于、模糊），并说明哪些是数：$0$；$\tfrac14$；$\pm1$；$\ast$；$\uparrow$。
\end{problem}

\sol $\tfrac14>0$，是数；$0=0$，是数；$\pm1\parallel0$、$\ast\parallel0$，都不是数（$\pm1$ 不满足“左选项小于右选项”，$\ast$ 的左右选项都含非零游戏）；$\uparrow>0$，但 $\uparrow$ 不是数——它的右选项 $\ast$ 不是数，不满足数的构造规则。五者中三个与 $0$ 模糊或为正，对应“先手胜”与“L 必胜”两种结果。分类经验：\textbf{先问是不是数}（左小右大且全是数），是数则直接比大小；不是数再看是否正负分明（如 $\uparrow$）；都不是才是模糊，需要按先后手分析。本题的五个值恰好覆盖值表的五行，建议对照上一节的小结表逐行复述。



\chapter{纳什均衡与零和博弈}
\section{纳什均衡}
静态非合作博弈问题的局部最优策略称为纳什均衡。

\begin{definition}[纳什均衡]
一个策略组合是\textbf{纳什均衡}，若其中任何一个玩家单方面改变策略，都不能提高他的收益（局部最优）。
\end{definition}

先解释“局部”二字：纳什均衡只要求“没有人想单方面反悔”，并不保证整体收益最大，甚至不保证均衡是唯一的；它刻画的是理性博弈中一个“稳定”的拍板状态——给定别人的选择，我已经做了我能做的最好的事。

“策略”指的是什么？最简单的想法是：每名玩家的策略就是固定一种决策（\textbf{纯策略}）。然而，按照这个定义，纳什均衡不一定存在。

\begin{example}[石头剪刀布没有纯策略均衡]
胜者收益 $+1$，败者收益 $-1$，平局双方收益均为 $0$。收益矩阵（行是 $A$ 的出拳，列是 $B$ 的出拳，格子里是 $A$ 的收益；$B$ 的收益恰为其相反数）为
\[
 W=\begin{pmatrix}0&1&-1\\-1&0&1\\1&-1&0\end{pmatrix}.
\]
任取非胜的一方，固定另一方的策略，这一方一定可以通过反悔取胜，所以不存在（纯策略）纳什均衡。换言之：在纯策略的世界里，石头剪刀布没有“稳定的拍板状态”，谁先亮拳谁吃亏。
\end{example}

根据生活经验，石头剪刀布博弈的最优策略是随机出拳，而不是一直出同样的拳。这就引出了一种更复杂的策略。

\begin{definition}[纯策略与混合策略]
固定一种决策的策略称为\textbf{纯策略}；为每种决策分配一个使用概率、每次依概率随机决策的策略称为\textbf{混合策略}。
\end{definition}

在混合策略下，纳什均衡比较的是玩家反悔前后的\textbf{期望}收益——决策随机化之后，“收益”不再是确定值，比较的基准自然变成期望。可以验证石头剪刀布中唯一的均衡是双方各以 $\tfrac13$ 的概率出三种拳：此时无论对方出什么，自己的期望收益都是 $0$，任何偏离（例如偏向石头）都会被对方用“专门出布”惩罚到负期望，因此没有人愿意单方面改变。

\begin{theorem}[纳什定理]\label{thm:nash}
混合策略纳什均衡一定存在。
\end{theorem}

纳什定理告诉我们：把策略空间放宽到概率分布后，“稳定拍板点”总是存在；这一定理是 1950 年纳什用角谷不动点定理证明的，也是他获得诺贝尔经济学奖的工作之一。

纯策略纳什均衡太简单了（枚举决策组合即可判定），下文以混合策略纳什均衡为主线；唯一需要两者并存的场合是双人零和博弈——那里纯策略均衡对应鞍点、混合均衡给出唯一的值，见后面“零和博弈”与“Minimax 定理”两小节。

纳什均衡有什么用？在 OI 中几乎没用。理由如下：
\begin{itemize}
\item 因为纳什均衡仅仅是一个局部最优策略，所以对于同一个博弈问题，可能存在多个毫不相关、无法横向对比的纳什均衡。题目若要求“求出最优策略”，需要的是全局最优；而纳什均衡连“哪个均衡更好”都说不清。
\item 更糟糕的是，求解纳什均衡的复杂度很高。\footnote{参考资料：\url{https://people.cs.pitt.edu/~kirk/CS1699Fall2014/lect4.pdf}。相关经典结论：判定是否存在多个纳什均衡是 NP-Complete 的（Gilboa--Zemel, 1989）；求任意一组纳什均衡是 PPAD-Complete 的（Chen--Deng, 2006）。}
\begin{itemize}
\item 判定是否存在两个纳什均衡是 NP-Complete 的：3-SAT 可以规约到双人匿名博弈的情况。
\item 求任意一组纳什均衡是 PPAD-Complete 的：求布劳威尔不动点可以规约到双人博弈的情况。
\end{itemize}
\end{itemize}
所以，在 OI 中可能有用的纳什均衡问题只有：双人零和博弈。下面先补一点零和博弈本身的性质，Minimax 定理紧随其后。

\section{零和博弈}
零和的意思是收益总和恒定——在全篇的约定里通常就是 $0$：一方所得即另一方所失，$B$ 的收益恰好是 $A$ 的收益的相反数。于是整场博弈只需要一张表就够了：$A$ 的收益矩阵 $W$，$B$ 的全部心思就是最小化它。注意 1.2 节把“非合作（零和）”列为组合游戏的特征，这里讨论的则是静态的两人零和博弈，视角不同，但“利益完全对立”的内核是同一个。

\begin{example}[猜硬币（Matching Pennies）]\label{ex:pennies}
两人同时各亮一枚硬币，同面则 $A$ 赢 $1$ 元、异面则 $A$ 输 $1$ 元，收益矩阵（A 视角）为
\[
 W=\begin{pmatrix}1&-1\\-1&1\end{pmatrix}.
\]
这个矩阵没有纯策略纳什均衡：$A$ 出正面时 $B$ 想出反面，$A$ 出反面时 $B$ 想出正面，任何“固定亮法”都会被人针对。所以双方只能随机：各以 $\tfrac12$ 的概率出正反，均衡值 $0$。矩阵是反对称的（$W_{i,j}=-W_{j,i}$，决策集相同），对称性保证了没人能占便宜、均衡值只能是 $0$；石头剪刀布（上一节开头那张 $3\times3$ 矩阵）就是它的三决策版本。
\end{example}

这里附带一个常用技巧：给收益矩阵整体加一个常数或乘一个正数，所有均衡策略都不变，只是博弈值随之平移或缩放——因此分析时可以把矩阵平移成非负、放大成整数，怎么方便怎么来（这也是下一小节线性规划写法能成立的伏笔）。

如果纯策略下就有“互相锁死”的格子，事情会简单得多。设 $A$ 能保底的收益是 $\max_i\min_j W_{i,j}$（A 先挑行，B 再挑最差列），$B$ 能把 $A$ 压到的上限是 $\min_j\max_i W_{i,j}$（B 先挑列，A 再挑最好行），显然恒有 $\max\min\le\min\max$。

\begin{definition}[鞍点]
若存在格子 $(i^*,j^*)$ 同时实现 $\max_i\min_j W_{i,j}$ 与 $\min_j\max_i W_{i,j}$，则称 $(i^*,j^*)$ 为\textbf{鞍点}。
\end{definition}

若鞍点存在，双方直接选 $i^*,j^*$ 就是纯策略纳什均衡。

\begin{example}[鞍点]
\[
 W=\begin{pmatrix}1&1\\2&0\end{pmatrix}
\]
中，$\max_i\min_j W_{i,j}=1$（A 选第一行，无论 B 怎么选都保底 $1$），$\min_j\max_i W_{i,j}=1$（B 选第二列，无论 A 怎么选都最多给 $1$），格子 $(1,2)$ 就是鞍点。鞍点的直观形象是“行最小、列最大”：A 不敢换行（换了可能更差），B 不敢换列（换了 A 可能更好），谁先偏离谁吃亏。反过来看猜硬币，$\max\min=-1$ 而 $\min\max=1$，空隙里全是“想反悔”的冲动，纯策略自然稳不住。
\end{example}

零和博弈最可爱的地方在于它有唯一的值：Minimax 定理保证 $\max\min=\min\max$，于是无论纯策略还是混合策略，所有纳什均衡都给出同一个期望收益 $v$，不存在一般博弈里“多个均衡、无从比较”的尴尬。这也正是前面说“OI 里有用的纳什均衡只剩双人零和”的底气：题目要的“最优策略下的期望收益”在零和博弈里是个良定义的数，换成非零和博弈，连“最优”都说不清。鞍点存在就直接用纯策略，没有鞍点（如猜硬币）就上混合策略——具体怎么解，看下面的 Minimax 定理与图解法。

\section{Minimax 定理}
\begin{problem}{Minimax 问题}
AB 两人博弈，A 有 $n$ 种决策，B 有 $m$ 种决策。如果 A 选择决策 $i$，B 选择决策 $j$，则游戏的分数是 $W_{i,j}$。$W$ 称为游戏的收益矩阵。A 要最大化期望分数，B 要最小化期望分数，求游戏的期望分数。
\end{problem}

\begin{theorem}[Minimax 定理]\label{thm:minimax}
这个问题的答案是唯一的。
\end{theorem}

先感受一下为什么“唯一”不平凡：B 想最小化分数，但他不知道 A 会怎么随机；A 想最大化分数，也不知道 B 怎么随机。两人都在跟“对方的未知策略”博弈。Minimax 定理说的是：这种相互猜测最终收敛到一个确定值——存在混合策略 $(p,q)$ 与常数 $v$，使得 A 无论怎么走，在 $q$ 下的期望至多 $v$；B 无论怎么走，在 $p$ 下的期望至少 $v$，而 $v$ 就是博弈的“公平价格”。注意这里 $v$ 同时是 A 的保底与 B 的天花板，两者相等是定理的全部内容。

\paragraph{A 的角度（下界）.} A 可以选择一个概率分布 $p_{1\ldots n}$，来最大化在 B 的任意策略下游戏的期望分数。具体地说，A 要最大化“最坏情况”：无论 B 选哪个纯策略 $j$，期望分数 $\sum_i p_iW_{i,j}$ 都必须 $\ge ans$。这个最坏情况下的最大值是答案的一个下界——A 至少能保证这么多。

\paragraph{B 的角度（上界）.} 对称地，B 可以选择一个概率分布 $q_{1\ldots m}$，来最小化在 A 的任意策略下游戏的期望分数：无论 A 选哪个纯策略 $i$，期望分数 $\sum_j q_jW_{i,j}$ 都必须 $\le ans'$。这是答案的一个上界——B 最多让 A 拿到这么多。

只要证明上界等于下界，就能证明原问题的答案唯一。若下界 $<$ 上界，两人各自的“保底”之间有空隙，说明存在相互占优的余地，均衡不存在；定理断言空隙不存在。

考虑 A 的角度。问题可以表示成线性规划的形式：设答案为 $ans$。
\[
 \begin{aligned}
 \max\quad & ans\\
 \text{s.t.}\quad & \sum_{i=1}^{n}p_i\le1\\
 & \forall j\in[1,m],\ ans-\sum_{i=1}^{n}p_iW_{i,j}\le0
 \end{aligned}
\]

\begin{pitfall}[$\sum p_i\le1$ 还是 $\sum p_i=1$]
严格地说，对一般的实收益矩阵，约束应写成 $\sum_i p_i=1$（概率质量必须恰好用完）；这里写 $\le1$ 相当于允许 A“弃权”部分概率质量，只有当收益矩阵非负（或先整体平移成非负）时两种写法才给出同一最优值。整体给矩阵加常数只平移博弈值、不改变均衡策略，乘正数同理——所以遇到一般矩阵，先平移成非负即可放心使用这套写法。OI 中的双人零和题（胜率、金钱收益）矩阵天然非负，更无此顾虑。
\end{pitfall}

考虑其对偶形式：
\[
 \begin{aligned}
 \min\quad & ans'\\
 \text{s.t.}\quad & \sum_{j=1}^{m}q_j\ge1\\
 & \forall i\in[1,n],\ ans'-\sum_{j=1}^{m}q_jW_{i,j}\ge0
 \end{aligned}
\]
观察到这就是 B 的答案对应的线性规划！证毕。这里的“证毕”其实调用了线性规划的\textbf{强对偶定理}：对偶问题的最优值等于原问题的最优值。而恰好 B 的“最小化最坏情况”写出来就是 A 的线性规划的对偶，两个最优值相等，上下界闭合，Minimax 定理得证。从理论层面看，Minimax 定理是“LP 强对偶”在博弈论里的一个化身。

使用求解线性规划的算法也可以构造出一组策略：解出上述 LP（或其任何等价形式）后，最优解 $p,q$ 本身就是双方的均衡混合策略。不过 OI 中一般用不到通用的 LP 求解器，实战的出路在下一小节。

\begin{teach}[图解法的画图检查表]
\begin{enumerate}
\item 两条直线是否一升一降？同向则均衡在端点，不要硬解交点。
\item 交点 $p^*$ 是否落在 $(0,1)$ 内？出界退回端点比较。
\item 交点纵坐标是否就是博弈值 $v$？提醒学生 $v$ 是期望收益，不是概率。
\item 反解 $q^*$ 时用“B 在交点处两列期望相等”，别再用 A 的直线。
\end{enumerate}
\end{teach}
\section{能不能不用线性规划？}
对于一般情况，不行。对于一般题目，可以。

实战中遇到的双人零和博弈，一般每人都只有两种决策。这种情况下可以使用\textbf{图解法}：设 A 使用第一种决策的概率为 $p$（则第二种决策的概率为 $1-p$）。观察到对于 B 的每种决策，当 B 只使用这种决策时，A 的期望分数是一个关于 $p$ 的一次函数：
\[
 f_j(p)=pW_{1,j}+(1-p)W_{2,j}.
\]
画出这些一次函数的图像，容易通过分类讨论求出纳什均衡。具体做法如下：
\begin{itemize}
\item 对每个 $j$ 画直线 $y=f_j(p)$（$p\in[0,1]$）。A 的保证收益是“最坏情况”，即所有直线中的\textbf{下包络} $\min_j f_j(p)$——B 一定会挑最有利于他的那条线来对付 A。
\item 纳什均衡中的 $p$ 就是下包络的最高点 $p^*=\arg\max_{p\in[0,1]}\min_j f_j(p)$，均衡值 $v=\min_j f_j(p^*)$。
\item 最高点只可能出现在两种地方：两条直线的交点，或区间端点 $p=0,1$。若 $f_1,f_2$ 的斜率符号相反（B 的两种决策对 A 的威胁一增一减），下包络呈“倒 V”形或“V”形，最大值点自然在两线交点处，解方程 $f_1(p)=f_2(p)$ 即可；否则退化到端点，检查哪个端点收益更大。
\end{itemize}

\begin{center}
\begin{tikzpicture}[scale=3.2,>=Stealth]
\draw[->] (-0.05,0) -- (1.15,0) node[right] {$p$};
\draw[->] (0,-0.05) -- (0,1.0);
\draw (1,-0.02) node[below] {$1$} -- (1,0.02);
\draw (0.015,1) node[left] {$y$} -- (-0.015,1);
\draw[thick,ink!70] (0,0.9) -- (1,0.2) node[right] {\small $f_1(p)$};
\draw[thick,accent] (0,0.3) -- (1,0.8) node[right] {\small $f_2(p)$};
\draw[very thick,warn] (0,0.3) -- (0.7,0.55) -- (1,0.2);
\filldraw[warn] (0.7,0.55) circle (0.6pt) node[above=2pt] {\small $(p^*,v)$};
\end{tikzpicture}\\[2pt]
{\small 图解法：实心折线是两条期望直线的下包络，其最高点 $(p^*,v)$ 即 A 的均衡混合概率与博弈值。}
\end{center}

当交点 $p^*\in(0,1)$ 时，B 的均衡混合策略 $q$ 也可以用对称方法求出，或者用“在均衡点处 B 两种纯决策的期望相等”这一条件反解。例题 [CF98E] Help Shrek and Donkey 和 [THUPC 2023 Pre] 欺诈游戏的解法都会落到这个流程上。

如果遇到了决策数很多的双人博弈，可以先想想能不能转化成每人只有两种决策的情况。这一步往往才是题目的真正难点：例如 [THUPC 2023 Pre] 欺诈游戏看似每方有 $n$ 种决策，其实可以用递推把“是否选择 $n$”提炼成 $2$ 种决策的博弈。

如果想不到，可以考虑使用一个网络上广为流传的乱搞做法：不妨设 $n\le m$。设 $w_j$ 表示如果 B 只使用决策 $j$，A 的期望分数。则如果 $w$ 不两两相等，那么直觉上可以“移多补少”，调整得到一个更优的方案。考虑直接列方程高斯消元：根据 $w$ 两两相等可以列出 $m-1$ 个方程，此外我们还有 $\sum p_i=1$，所以方程数量一定足够。为什么均衡时 $w$ 要两两相等？因为在均衡点处 B 对他用到的每个纯策略应当“无所谓”（期望相同），否则 B 会放弃收益差的那个策略；这个条件叫\textbf{无差异原理}，是手动解小型博弈的利器。因为题目往往有特殊性质，所以这样做很可能是对的。

根据 Itst 的说法，当 $n=m$ 时，如果从 A 的角度出发，方程的解满足 $p_i\ge0$，且从 B 的角度出发，方程的解满足 $q_i\ge0$，那解方程就是正确的。我们不会证明，但是这个判定条件似乎对解题并无帮助。

\begin{teach}[数值算例的课堂流程]
完整算例建议按五步走并让学生参与每一步：先猜 $v$ 的符号（矩阵略微偏 A，多数学生能猜对）；再画两条直线找交点；然后用“保底/天花板”双向验算；最后用无差异原理反解 $q^*$，并特意指出本例 $p^*=q^*$ 是巧合。高频错误：把无差异对象接反——$p^*$ 让 B 无差异，$q^*$ 让 A 无差异。练习 1 至练习 8 可留作分层作业，练习 4 的闭式公式与练习 7、练习 8 的数值要求推导而不许背。
\end{teach}
\section{图解法完整算例}
抽象流程说完了，本节把一个具体的 $2\times2$ 博弈从头算到尾，作为图解法的标准作业流程。\footnote{延伸资料：洛谷专栏《浅谈纳什均衡》（\url{https://www.luogu.me/article/pjhj07pz}），其中强调的“每个参与人的混合策略都使其余参与人任何纯策略的期望收益相等”正是本节反复使用的无差异原理；文中关于匿名博弈纳什均衡近似求解的讨论也值得一读。}

\begin{problem}{完整算例}
双方各两种决策，收益矩阵（A 视角）为
\[
 W=\begin{pmatrix}2&-1\\-1&1\end{pmatrix},
\]
求均衡混合策略与博弈值 $v$。
\end{problem}

\sol \paragraph{第一步：排除纯策略均衡（找鞍点）.} 行最小值为 $-1,-1$，故 $\max\min=-1$；列最大值为 $2,1$，故 $\min\max=1$。两者不等，无鞍点，必须混合。

\paragraph{第二步：写两条期望直线.} 设 A 以概率 $p$ 选第一行，则 B 各纯策略下 A 的期望为
\[
 f_1(p)=2p-(1-p)=3p-1,\qquad f_2(p)=-p+(1-p)=1-2p .
\]
\paragraph{第三步：找下包络的最高点.} $f_1$ 递增、$f_2$ 递减，包络呈“倒 V”形，交点满足 $3p-1=1-2p$，解得
\[
 p^{*}=\frac25,\qquad v=f_1(p^{*})=\frac65-\frac15=\frac15 .
\]
A 以 $\tfrac25$ 混合时，无论 B 出哪一列，A 的期望都是 $\tfrac15$——这就是 A 的保底。

\paragraph{第四步：反解 B 的混合（无差异原理）.} 设 B 以概率 $q$ 选第一列。均衡点处 A 对两行应当无所谓：
\[
 2q-(1-q)=-q+(1-q)\ \Longrightarrow\ 3q-1=1-2q\ \Longrightarrow\ q^{*}=\frac25,
\]
此时 B 把 A 恰好压到 $v=\tfrac15$。本例中 $p^{*}=q^{*}$ 纯属巧合（矩阵接近对称），一般并不相等。

\paragraph{第五步：验算.} 若 A 偏离到 $p=0$，B 用第一列把 A 压到 $-1<v$；若 B 偏离到全用第二列，A 拿到 $1>v$。双方都无利可图，均衡成立。整个流程——排除鞍点、画包络、解交点、无差异反解、验算——对任何 $2\times n$ 图解法题目都是同一套。

\begin{pitfall}[端点均衡：当交点不在开区间内时]
下包络的最高点未必是两条直线的交点：若两线同向（B 的两种策略对 A 的威胁同增同减），最高点落在 $p=0$ 或 $p=1$，均衡退化为纯策略——此时通常存在被支配的纯策略，先删支配再解往往更省事。判据：交点 $p^{*}\notin[0,1]$ 或两线斜率同号时，直接比较两个端点的包络值。
\end{pitfall}

\begin{teach}[扑克模型的三步讲法]
先玩：用三张牌（$1,2,3$ 各一张）代替均匀分布、赌注固定一元，课堂玩几轮，学生会自己发现“弱牌全弃”会被偷注打穿。再算：带学生算 $a=1$ 的三个数 $c=\tfrac13$、$p=\tfrac23$、$v=-\tfrac19$，每算一个就问“它让谁无差异”。后点题：把诈唬量化为确定频率，正是 1928 年那篇论文前后的工作，Kuhn 与 Tucker 称之为“连外行都能欣赏的突破”。可留一问下课：“先手吃亏，为什么现实中常有人抢着先表态？”
\end{teach}
\section{拓展：冯·诺依曼 1928——Minimax、对偶与扑克}
1928 年，二十出头的冯·诺依曼发表《Zur Theorie der Gesellschaftsspiele》（论博弈论）\footnote{Mathematische Annalen 第 100 卷，295--320 页；1926 年 12 月在格丁根数学学会首报，\url{https://doi.org/10.1007/BF01448847}。}，第一次对一般二人零和博弈证明了定理~\ref{thm:minimax}。1953 年他在给弗雷歇的公开答复里写道：“就我所见，没有这条定理，就没有现代意义上的博弈论……在 Minimax 定理被证明之前，我以为没有什么值得发表。”（ Econometrica 上关于波雷尔批注的通信。）事实上波雷尔从 1921 年起就在一系列短文里提出混合策略的想法，并在极小的例子上猜测这条定理，但始终没能证明；弗雷歇 1953 年为波雷尔争优先权，冯·诺依曼以此作答。7.3 节的证明（线性规划对偶）在历史上出现得晚得多：1928 年的原始证明纯组合且相当繁难，后来经不动点方法（冯·诺依曼 1937、纳什 1950）与线性规划对偶（丹齐格 1951）被反复重证、日益透明。据丹齐格回忆，1947 年他向冯·诺依曼介绍线性规划时，对方听完当即回敬了一小时即兴讲演，讲的正是对偶性以及它与矩阵博弈的等价——7.4 节讨论的问题，源头就是这次会面。

论文前后，冯·诺依曼用刚诞生的 Minimax 理论做了一件至今仍令人叫绝的事：分析扑克。他把“诈唬”从心理花招变成了数学必然——Kuhn 与 Tucker 后来称之为“连外行都能欣赏的突破”。本节用本章的工具把这个模型完整解掉。

\paragraph{模型.} 规则：双方各下底注 $1$ 元；各从 $[0,1]$ 均匀抽一个隐藏数（A 拿 $x$，B 拿 $y$），数大者胜，各自只见自己的数；A 先行动：弃牌（输掉底注，收益 $-1$）或下注 $a$ 元；若 A 下注，B 或弃牌（A 赢得 B 的底注 $+1$），或跟注（补足 $a$ 元后摊牌，大者全收）。不许加注。这是零和博弈，但信息不完全——谁也不知道对方的数。

\paragraph{均衡.} 设 B 用阈值策略：$y\ge c$ 才跟注。A 拿 $x<c$ 的牌下注：B 以概率 $c$ 弃牌送钱、以概率 $1-c$ 跟注且必胜，收益为 $c-(1-c)(1+a)=c(2+a)-(1+a)$；弃牌收益 $-1$。两者相等的临界值即
\[
 c=\frac{a}{2+a}.
\]
$x\ge c$ 时下注收益为 $c+(1+a)(2x-c-1)$，随 $x$ 递增，故 A 全部价值下注；$x<c$ 的手牌在均衡处诈唬与弃牌无差异，下注频率由 B 的无差异条件定：B 拿 $y=c$ 的牌，摊牌胜率是 A 诈唬牌占下注范围的份额 $\tfrac{pc}{pc+1-c}$（$p$ 为诈唬频率），令其等于 $\tfrac{a}{2(1+a)}$，解出
\[
 p=\frac{2}{2+a}.
\]
三类收益加总——弃牌区每手 $-1$、诈唬区每手恰为 $-1$、价值下注区积分恰为 $c(1-c)$——博弈值为
\[
 v=-c+c(1-c)=-c^{2}=-\Bigl(\frac{a}{2+a}\Bigr)^{2},
\]
对 A 为负：被迫先行动、先亮出“我下注了”这个信息的一方吃亏。

\begin{center}
\begin{tabular}{cccc}
\toprule
赌注 $a$ & B 跟注阈值 $c$ & A 诈唬频率 $p$ & 博弈值 $v$\\
\midrule
$\tfrac12$ & $\tfrac15$ & $\tfrac45$ & $-\tfrac1{25}$\\
$1$ & $\tfrac13$ & $\tfrac23$ & $-\tfrac19$\\
$2$ & $\tfrac12$ & $\tfrac12$ & $-\tfrac14$\\
\bottomrule
\end{tabular}
\end{center}

\paragraph{为什么必须诈唬.} 假如 A 从不诈唬，只在好牌下注，B 的最优回应是“牌不够大一律弃”，A 的价值下注便一无所获；于是 A 必须用一部分垃圾牌冒充强牌，把 B 逼回赌桌，而诈唬频率恰好被调到 B 愿意跟注为止——7.5 节的无差异原理在这里两头各捏一刀：A 的阈值由自己的无差异定，A 的频率由 B 的无差异定。注意赌注越大，B 越谨慎（$c$ 升高）、A 诈唬越罕见（$p$ 下降）但区间越宽、先手劣势 $c^{2}$ 也越大。

\begin{pitfall}[诈唬不是越多越好]
学生常把“均衡要诈唬”读成“诈唬越多越好”。练习 8 给出反例：$a=1$ 时把频率提到 $1$，B 的阈值降到 $\tfrac14$，每手垃圾牌比弃牌多亏 $\tfrac14$ 元。反向误读同样要防——“均衡里有诈唬，所以我随时可以诈”：离开均衡频率，诈唬立刻从必需品变成送给对手的礼物。
\end{pitfall}

这对策略在 7.1 节的意义下恰是纳什均衡：单方面多诈、少诈、放宽或收紧阈值都只会更差。冯·诺依曼比纳什早二十二年，就在零和世界里把“各玩各的最优回应”走到了尽头；2017 年击败职业选手的扑克 AI Libratus，在算力层面延续的仍是“按均衡频率诈唬”这条 1928 年的思路。

\section{分层例题与解析}

以下例题的共性：真实博弈被压缩成一个 $2\times2$ 零和子博弈，再用本章的图解法求解；博弈论含量集中在“找到正确的 $2\times2$ 子博弈”这一层。


\begin{problem}{练习 1（自拟）：三线包络}
双方博弈，A 有两种决策、B 有三种决策，收益矩阵（A 视角）为
\[
 W=\begin{pmatrix}2&-1&1\\-1&1&1\end{pmatrix},
\]
求 A 的均衡混合策略、博弈值，并说明 B 的第三种决策在均衡中的地位。
\end{problem}

\hint 三条期望直线逐对求交点，找出真正构成下包络的两条；整体位于包络上方的直线对应的纯策略会被闲置。

\sol 设 A 以概率 $p$ 选第一行：$f_1(p)=3p-1$，$f_2(p)=1-2p$，$f_3(p)\equiv1$。$f_1,f_2$ 交于 $p=\tfrac25$、值 $\tfrac15$；而 $f_3\equiv1$ 在整个 $[0,1]$ 上都高于 $f_1,f_2$ 的包络（包络最高点只有 $\tfrac15$），第三列对 B 毫无吸引力。于是均衡为
\[
 p^{*}=\frac25,\qquad v=\frac15,\qquad q^{*}=\Bigl(\frac25,\ \frac35,\ 0\Bigr).
\]
无差异核对：$2q_1-q_2=-q_1+q_2$ 且 $q_2=1-q_1$ 给出 $q_1=\tfrac25$，此时 A 对两行的期望同为 $\tfrac15$。教学点：图解法画出的直线若整体在包络上方，它对应的纯策略在均衡中概率为零——B 若用它，反而把更高的期望送给了 A。

\begin{problem}{练习 2（自拟）：非对称石头剪刀布}
收益规则：平局 $0$；A 赢常规一局得 $1$、输一局得 $-1$；唯独“石头砸剪刀”这一手 A 赢得 $2$（输时仍为 $-2$）。收益矩阵（行 $A\in$\{石头,布,剪刀\}，列 $B$ 同序）为
\[
 W=\begin{pmatrix}0&-1&2\\1&0&-1\\-2&1&0\end{pmatrix},
\]
求双方均衡混合策略与博弈值。
\end{problem}

\hint 矩阵反对称（$W_{i,j}=-W_{j,i}$），均衡值必为 $0$；用无差异原理列两个方程再加归一化即可。

\sol 反对称性使任何相同的混合策略相互对局时双方期望均为 $0$，故 $v=0$。设 A 出石头、布、剪刀的概率为 $p_1,p_2,p_3$。令 B 三种纯策略下 A 的期望全为 $0$（无差异原理）：
\[
 -p_2+2p_3=0,\qquad p_1-p_3=0,\qquad -2p_1+p_2=0,
\]
三式同解：$p_1=p_3$、$p_2=2p_3$，加归一化得 $(p_1,p_2,p_3)=\bigl(\tfrac14,\tfrac12,\tfrac14\bigr)$。由对称性 B 的均衡混合同为 $\bigl(\tfrac14,\tfrac12,\tfrac14\bigr)$。直观：石头的高回报被它“怕布”的软肋精确抵消，多出的那点概率全补在布上。

\begin{problem}{练习 3（自拟）：不对称猜硬币}
两人同时亮硬币，同面 A 赢 $1$ 元、异面 A 输 $1$ 元；但若 A 出反面而 B 出正面，A 输 $2$ 元。收益矩阵（行 \{正,反\}，列 \{正,反\}）为
\[
 W=\begin{pmatrix}1&-1\\-2&2\end{pmatrix},
\]
求均衡与博弈值。
\end{problem}

\hint 先查鞍点；再按“下包络最高点”解 $p$，用无差异原理反解 $q$。

\sol 行最小值 $-1,-2$，列最大值 $1,2$，无鞍点。$f_1(p)=p-2(1-p)=3p-2$，$f_2(p)=-p+2(1-p)=2-3p$，交点 $p^{*}=\tfrac23$，$v=0$。反解 B：A 对两行无差异 $q-(1-q)=-2q+2(1-q)$，得 $q^{*}=\tfrac12$。验算：A 以 $\tfrac23$ 出正面时，无论 B 怎么出，期望都是 $0$；B 若全出正面，A 改出全正面即可赚 $1>v$。有趣的是矩阵明显不对称，博弈值却仍是 $0$——“值”是一个数，未必带着矩阵的符号偏见。

\begin{problem}{练习 4：一般 $2\times2$ 混合均衡的闭式解}
设收益矩阵 $W=\begin{pmatrix}a&b\\ c&d\end{pmatrix}$ 且无鞍点，推导均衡混合概率与值的一般公式，并用它复算完整算例（$a=2,b=-1,c=-1,d=1$）与练习 3。
\end{problem}

\hint 对 B 的两列写无差异方程解 $p^{*}$；对 A 的两行写无差异方程解 $q^{*}$；代回任一行求 $v$。

\sol 令 B 两列下 A 的期望相等：$pa+(1-p)c=pb+(1-p)d$，整理得 $p(a-b)=(1-p)(d-c)$，故
\[
 p^{*}=\frac{d-c}{(a-b)+(d-c)},\qquad
 q^{*}=\frac{d-b}{(a-c)+(d-b)},
\]
$v$ 为把 $p^{*}$（或 $q^{*}$）代回任一条期望直线的值。复算：完整算例 $p^{*}=\tfrac{1-(-1)}{3+2}=\tfrac25$、$q^{*}=\tfrac{1-(-1)}{3+2}=\tfrac25$，$v=\tfrac15$；练习 3 的 $p^{*}=\tfrac{2-(-2)}{2+4}=\tfrac23$、$q^{*}=\tfrac{2-(-1)}{3+3}=\tfrac12$，均与前文一致。注意公式只在“无鞍点”前提下成立；有鞍点时分子分母可能同时为零，纯策略均衡就是答案。

\begin{problem}{练习 5（自拟）：鞍点检索}
收益矩阵（A 视角）为
\[
 W=\begin{pmatrix}3&2&4\\1&2&0\\5&2&1\end{pmatrix},
\]
判断是否存在纯策略纳什均衡（鞍点）；若有，给出双方的选择与博弈值。
\end{problem}

\hint 分别算 $\max_i\min_j W_{i,j}$ 与 $\min_j\max_i W_{i,j}$，相等即有鞍点。

\sol 行最小值为 $2,0,1$，故 $\max\min=2$；列最大值为 $5,2,4$，故 $\min\max=2$。两者相等，鞍点存在：格子 $(1,2)$ 同时是第一行的最小值与第二列的最大值。A 选第一行、B 选第二列，博弈值 $v=2$。此时双方都无偏离动机：A 换行也许局部多得，但 B 会立刻改列把它压回不高于 $2$；B 同理。纯策略均衡的判定成本只有 $O(nm)$，这正是“纯策略均衡太简单”的含义。

\begin{problem}{练习 6（自拟）：端点均衡与鞍点}
收益矩阵（A 视角）为
\[
 W=\begin{pmatrix}1&-2\\1&1\end{pmatrix},
\]
用图解法求均衡，并解释“最高点落在端点”意味着什么。
\end{problem}

\hint 写两条期望直线，找下包络最高点；端点均衡往往对应鞍点，回头用 $\max\min$ 与 $\min\max$ 验证。

\sol $f_1(p)\equiv1$（第一列两个收益都是 $1$），$f_2(p)=-2p+(1-p)=1-3p$。下包络为 $1-3p$，在 $p^{*}=0$ 处取得最大值 $v=1$——最高点落在端点，A 直接用纯策略（第二行）。验证：第二行两个收益都是 $1$，$\max\min=\min\max=1$，格子 $(2,1)$、$(2,2)$ 都是鞍点。教学点：端点均衡与鞍点是一回事的两种说法；图解法先画图、发现端点后退回鞍点检验，比硬套交点公式更稳。

\begin{problem}{练习 7（自拟）：小赌注的扑克}
取 $a=\tfrac12$，求 B 的跟注阈值 $c$、A 的诈唬频率 $p$ 与博弈值 $v$，并与 $a=1$ 比较：赌注变小后，B 更愿意跟注还是更谨慎？A 的诈唬更频繁还是更罕见？
\end{problem}

\hint 代入 $c=\tfrac{a}{2+a}$、$p=\tfrac{2}{2+a}$、$v=-\bigl(\tfrac{a}{2+a}\bigr)^{2}$。

\sol $c=\dfrac{1/2}{5/2}=\tfrac15$，$p=\dfrac{2}{5/2}=\tfrac45$，$v=-\tfrac1{25}$。与 $a=1$（$c=\tfrac13$、$p=\tfrac23$、$v=-\tfrac19$）相比：赌注减半后 B 的阈值从 $\tfrac13$ 降到 $\tfrac15$，跟得更松；A 的诈唬频率从 $\tfrac23$ 升到 $\tfrac45$，更频繁但区间 $[0,\tfrac15)$ 更窄；先手劣势从 $\tfrac19$ 缩到 $\tfrac1{25}$。小赌注的桌上人人放松警惕，先行动的代价也小。

\begin{problem}{练习 8（自拟）：诈唬过头的代价}
设 $a=1$，A 改为“所有手牌都下注”。求 B 的最优跟注阈值，并算出 A 的垃圾牌下注相对弃牌的净亏损，说明均衡频率 $\tfrac23$ 为什么不能再高。
\end{problem}

\hint 面对“全下注”，B 在哪条 $y$ 上跟注与弃牌无差异？

\sol A 全下注时 B 的摊牌胜率即 $P(x<y)=y$，无差异条件 $2(2y-1)=-1$ 给出 $y=\tfrac14<\tfrac13$：B 跟得更松。A 的垃圾牌下注收益为 $\tfrac14\cdot 1+\tfrac34\cdot(-2)=-\tfrac54$，比弃牌净亏 $\tfrac14=\tfrac{a^{2}}{2(1+a)}$——多出的每一手诈唬都被“松跟”的 B 收税。反过来若诈唬不足，B 收紧，A 的价值下注失色。无差异原理把频率钉死在 $\tfrac23$：它正是“多一手都亏、少一手都亏”的平衡点。
\begin{problem}{[CF98E] Help Shrek and Donkey}
史瑞克和驴决定玩一种叫做 YAGame 的牌类游戏。规则非常简单：最开始，史瑞克手中有 $m$ 张牌，驴手中有 $n$ 张牌（两位玩家都看不到对方手中的牌），还有一张牌正面朝下地放在桌子上，两位玩家同样也看不到。因此，游戏开始时共有 $m+n+1$ 张牌。此外，玩家们知道整副牌由哪些牌组成，也清楚自己手中有哪些牌（但他们不知道哪一张牌在桌子上，以及对方手中有哪些牌）。游戏由史瑞克先手，接下来两位玩家轮流行动。每轮中，某位玩家可以进行以下两种操作之一：
\begin{itemize}
\item 尝试猜测桌子上的那张牌。如果他猜对了，游戏立即结束，他获胜。如果猜错了，游戏同样立即结束，但对方获胜。
\item 指定整副牌中的任意一张牌。如果对方手中有这张牌，则必须亮出来并将其移出本局（这张牌不再参与后续游戏）。如果对方没有该牌，则直接说明。
\end{itemize}
如果两位玩家均采取最优策略，求史瑞克获胜的概率。

数据范围：$n,m\le1000$。
\end{problem}

\hint 设 $f_{n,m}$ 为 A（先手，$n$ 张牌）对 B（$m$ 张牌）的胜率，递归会互换攻守，需要“$1-f$”补集写法；决策二的分支里藏着一个 $2\times2$ 静态子博弈，用图解法在两直线交点处解均衡。

\sol \paragraph{考虑递归子问题.} 设 $f_{n,m}$ 表示 A（先手）有 $n$ 张牌，B 有 $m$ 张牌，A 的胜率。显然，固定一方的手牌数，另一方的手牌数越大，胜算就越大，因为这一方更难猜中桌上的牌。递归会互换攻守（B 亮牌后轮到 B 行动），所以需要“$1-f$”这种补集写法：$f_{m-1,n}$ 是以 B 为先手的局面里 B 的胜率，A 的胜率就是 $1-f_{m-1,n}$。注意每一层递归两个手牌数之和都在减少（$(m-1,n)$ 与 $(m,n-1)$），状态图无环，按 $n+m$ 升序记忆化即可。

\paragraph{考虑 A 的策略.} 如果 A 使用了决策 1，那游戏会立即结束。所以，A 采取两个决策的混合策略是无意义的（因为不会相互影响），所以只需要分别计算胜率并取较大值。使用决策 1 的胜率显然是
\[
 \frac{1}{m+1}
\]
——A 手里 $n$ 张牌以外还有 $m+1$ 张他没见过（B 的 $m$ 张 + 桌上 $1$ 张），均匀随机猜一张才是最优（猜自己见过的牌必错），猜中的概率是 $\tfrac1{m+1}$。考虑计算使用决策 2 的胜率。

如果 A 使用决策 2，一个简单的想法是，首先，A 为了获得更多的信息，应当问一张他没有的牌。又因为对 A 来说，他没有的牌都是相同的，所以 A 应当均匀随机问一张他没有的牌。如果 B 有这张牌，那么递归到 $1-f_{m-1,n}$（B 被迫亮牌弃牌，且轮到 B 行动），否则答案就是这张牌——注意“否则”有两种可能：牌在 A 自己手里，或就在桌上。如果答案就是这张牌，则 B 坐以待毙显然是不优的，B 应该直接猜这张牌就是答案——B 猜中则立即获胜。但是这就引出 A 的另一个策略：故意问一张他有的牌，误导 B 自爆。于是，这就形成了一个双人静态博弈，它只出现在“B 手里没有 A 问的那张牌”的情形：A 的决策：问他有的牌、问他没有的牌；如果 B 没有 A 问的牌，则 B 的决策：直接猜这张牌、认为这张牌属于 A 而忽略。收益矩阵如下（行 = A 的两种问法，列 = B 的两种应对，格子里是 A 的胜率）：
\[
 \begin{pmatrix}
 1 & 1-f_{m,n-1}\\[4pt]
 \dfrac{m}{m+1}\Bigl(1-f_{m-1,n}\Bigr) & \dfrac{1}{m+1}+\dfrac{m}{m+1}\Bigl(1-f_{m-1,n}\Bigr)
 \end{pmatrix}
\]
把四个格子逐个翻译一遍：
\begin{itemize}
\item （问自己有的牌，B 猜）：A 问的牌在自己手里，桌上必不是它，B 猜了必错，游戏结束 A 胜，胜率 $1$。
\item （问自己有的牌，B 忽略）：B 获得了“这张牌属于 A”的信息，等价于从 A 的未知牌里划掉一张——未知牌少一张的递归局面 $1-f_{m,n-1}$。注意这里的便宜在于 A 相当于“白赚”了一次提问而不损失任何牌。
\item （问自己没有的牌，B 猜）：被问的牌在 B 手里（概率 $\tfrac{m}{m+1}$）时 B 被迫亮牌，递归到 $1-f_{m-1,n}$；在桌上（概率 $\tfrac1{m+1}$）时 B 一猜即中，A 胜率为 $0$。
\item （问自己没有的牌，B 忽略）：在 B 手里的情形同上（$\tfrac{m}{m+1}$ 概率递归）；在桌上的情形 B 放过了唯一一次猜中机会，A 反而由此确定桌上牌并赢下游戏（$\tfrac1{m+1}$ 的概率胜率 $1$）。
\end{itemize}
矩阵里两行给出的两条“期望直线”方向相反（B 的两个应对分别对 A 的“诚实问”与“诈问”有利），且交点在 $[0,1]$ 内，所以均衡在交点处取得。考虑图解法，观察到两个决策的直线方向相反，且交点在 $[0,1]$ 内，所以求交点即可——交点处 B 对两种应对无差异，解一次方程得到混合概率与 $f_{n,m}$。每个状态的矩阵是 $2\times2$、解方程 $O(1)$，状态共 $O(nm)$ 个。

\cost 时间复杂度 $O(n^{2})$（$n,m\le1000$）。本题是“图解法 + 概率递推”的教科书组合：牌类游戏的不确定性被压缩成一个 $2\times2$ 收益矩阵，博弈论含量在“找到正确的 $2\times2$ 子博弈”这一层。

\src{https://codeforces.com/problemset/problem/98/E}

\begin{problem}{[THUPC 2023 Pre] 欺诈游戏}
在《LIAR GAME》中，小 E 看到了一个有趣的游戏。

这个游戏名叫《走私游戏》。游戏规则大概是这样的：一名玩家扮演走私者，一名玩家扮演检察官。走私者可以将 $x$ 日元（$x$ 为 $[0,n]$ 内的整数，由走私者决定）秘密放入箱子中，而检察官需要猜测箱子中的金额。假设检察官猜了 $y$（$y$ 也必须是整数）。如果 $x=y$，则走私失败，走私者一分钱也拿不到。如果 $x>y$，则走私成功，走私者可以从检察官那里拿走 $x$ 日元。如果 $x<y$，则走私失败，但是由于冤枉检察官需要赔付给走私者 $y/2$ 日元。游戏分有限回合进行，双方轮流做走私者和检察官。

可以证明，最优情况下每个回合走私者会采用同一种策略，检察官也会采用同一种策略。

小 E 想知道在一个回合中，双方的最优策略分别是什么，构造方案。答案对 $998244353$ 取模。

数据范围：$1\le n\le4\times10^{5}$。
\end{problem}

\hint 看似每人 $n$ 种决策，其实可以分步决策：先把“最大金额 $n$”的攻防单独拎出来，剩下的自动退化为 $n-1$ 的子问题；每层解一个 $2\times2$ 零和博弈。

\sol 称 A 为走私者，B 为检察官。设 $\Val(n)$ 为“金额上限是 $n$”时一回合的均衡值（A 的期望收益）。本题看起来每人有 $n$ 种决策（A 选走私金额 $x$，B 猜金额 $y$），但实际上可以分步决策：AB 分别决策自己是否选择 $n$——先把“最大金额”的攻防单独拎出来，剩下的部分自动退化成上限 $n-1$ 的子问题。四个分支的收益如下（A 的收益）：
\begin{itemize}
\item 如果 AB 同时选择 $n$（$x=n$，$y=n$），则 $x=y$，走私失败，收益为 $0$；
\item 如果 A 选 $n$ 而 B 不选 $n$（$y<n$），则 $x=n>y$，走私成功拿走 $n$，收益为 $n$；
\item 如果 A 不选 $n$ 而 B 选 $n$，则 $x<n=y$，属“冤枉”情形，B 赔付 $y/2=n/2$——注意只要 $x<n$ 收益恒为 $n/2$，与 $x$ 具体取多少无关；
\item 如果 AB 都不选 $n$，双方都只在 $[0,n-1]$ 内决策，则收益等于 $n-1$ 的问题的答案 $\Val(n-1)$，可以递归求解。
\end{itemize}
于是每层只需解一个 $2\times2$ 的零和博弈，收益矩阵是
\[
 \begin{pmatrix}0&n\\ n/2&\Val(n-1)\end{pmatrix}
\]
（行：A 选 $n$／不选 $n$；列：B 猜 $n$／不猜 $n$）。边界 $\Val(0)=0$（$x=y=0$ 只能平局收场）。这矩阵没有支配关系（第一列第二行更大、第二列第一行更大），均衡在混合策略内部：设 A 选 $n$ 的概率为 $p$、B 猜 $n$ 的概率为 $q$，A 的期望是
\[
 E(p,q)=p\bigl(0\cdot q+n(1-q)\bigr)+(1-p)\Bigl(\frac n2 q+\Val(n-1)(1-q)\Bigr),
\]
它是 $p,q$ 的双线性函数。用图解法：固定 $p$，B 会挑使期望更小的端点 $q\in\{0,1\}$；均衡的 $p^{*}$ 让 B 在两个端点之间无差异，反过来 $q^{*}$ 让 A 无差异——各解一次一元一次方程，$O(1)$ 求出一层的 $p^{*},q^{*}$ 与 $\Val(n)$。解完后用图解法求解纳什均衡即可。因为每一层的解只依赖上一层，且要求“构造方案”，自顶向下把每层选 $n$ 的概率存下来、按均匀随机实现即可（答案对 $998244353$ 取模，概率用逆元表示）。

\cost 时间复杂度 $O(n)$。本题演示了“压缩决策”这一招：$n$ 种决策的博弈若能找到某种“分层的最大元素”结构，就可以层层剥开化成 $2\times2$ 子博弈的递归。

\src{https://www.luogu.com.cn/problem/P9142}

\begin{teach}[课堂追问：纳什均衡还剩下什么？]
结课时抛出三个判断题抢答：(1) 零和博弈的两个均衡可能给出不同期望收益吗？（不可能，Minimax。）(2) 一般博弈求纳什均衡是多项式时间的吗？（PPAD-Complete。）(3) 广义 Nim-k 为什么没有结论？（强后继互相补偿。）三个都答对，说明学生带走的不只是套路，还有边界感。
\end{teach}

% !TeX program = xelatex
% ============================================================================
%  第四部分（续）  概率论
%  第 8 章   样本空间与随机事件
%  第 9 章   随机变量、期望与方差
%  第 10 章  常用概率模型与随机化算法
% ============================================================================

\chapter{样本空间与随机事件}
\begin{chapterintro}[章节导读]
概率论的第一道门槛不是公式，而是\textbf{把随机现象写成集合}。一旦写清了“所有可能结果是什么”“我们关心其中哪些”，古典概型就退化成一个计数问题——于是第三部分的排列组合、容斥、递推全部可以直接搬过来用。本章先建立这套语言，再处理条件概率与独立性这两个真正容易出错的地方。
\end{chapterintro}

\begin{tip}[概率论的统一口诀]
\textbf{先写样本空间，再写事件，最后才写公式。}凡是概率题做错，十有八九是第一二步含糊：或者漏了某些结果，或者把“顺序算不算不同”搞混了。这两件事与第三部分第 1 章的要求完全一致。
\end{tip}

\section{样本空间与事件}
\begin{definition}[样本空间]
一个随机试验的\textbf{所有可能结果}组成的集合 $\Omega$ 叫样本空间，$\Omega$ 中的元素叫样本点。$\Omega$ 的子集叫\textbf{事件}；只含一个样本点的事件叫基本事件。
\end{definition}

\begin{definition}[古典概型]
若样本空间 $\Omega$ 有限，且每个样本点等可能发生，则对任意事件 $A$，
\[
\Pp(A)=\frac{|A|}{|\Omega|}.
\]
\end{definition}

\begin{teach}[古典概型的全部难点都在 $|A|$ 与 $|\Omega|$]
分子分母必须用\textbf{同一套}计数口径。最常见的错误是分母按“有序”数（例如把两枚骰子看成 $36$ 个有序结果），分子却按“无序”数（例如把“点数之和为 $7$”按 $3$ 种组合数）。\textbf{先约定口径，再分别数分子分母。}
\end{teach}

\section{事件的运算与概率公理}
\begin{definition}[事件的运算]
\begin{itemize}
\item $A\cup B$：$A$ 或 $B$ 发生；
\item $A\cap B$：$A$ 且 $B$ 发生，简写 $AB$；
\item $\overline A=\Omega\setminus A$：$A$ 不发生；
\item $A\setminus B=A\cap\overline B$：$A$ 发生且 $B$ 不发生；
\item $A,B$ \textbf{互斥}：$A\cap B=\varnothing$。
\end{itemize}
\end{definition}

\begin{theorem}[概率的三条公理]
\begin{enumerate}
\item $\Pp(A)\ge0$；
\item $\Pp(\Omega)=1$；
\item 若 $A_1,A_2,\ldots$ 两两互斥，则 $\Pp\left(\bigcup_i A_i\right)=\sum_i\Pp(A_i)$。
\end{enumerate}
由此立刻得到：
\[
\Pp(\overline A)=1-\Pp(A),\qquad
\Pp(A\cup B)=\Pp(A)+\Pp(B)-\Pp(A\cap B).
\]
\end{theorem}
\begin{proof}
第二条：$\Omega\cup\varnothing=\Omega$ 且两者互斥，故 $\Pp(\Omega)=\Pp(\Omega)+\Pp(\varnothing)$，得 $\Pp(\varnothing)=0$。再由 $A\cup\overline A=\Omega$ 互斥，$\Pp(A)+\Pp(\overline A)=1$。

第三条是容斥原理的概率版：把 $A\cup B$ 拆成 $A$ 与 $B\setminus A$ 两个互斥部分，$\Pp(A\cup B)=\Pp(A)+\Pp(B\setminus A)$，而 $\Pp(B)=\Pp(B\setminus A)+\Pp(A\cap B)$，代入即得。
\end{proof}

\begin{pitfall}[“或”仍然不是“异或”]
$\Pp(A\cup B)=\Pp(A)+\Pp(B)-\Pp(A\cap B)$ 中减去的正是“两者都发生”的部分。如果题目中的 $A,B$ 已经互斥，这一项为 $0$；\textbf{不互斥时千万别忘了减}。这与第三部分第 3 章的容斥原理是同一个道理。
\end{pitfall}

\section{条件概率}
\begin{definition}[条件概率]
设 $\Pp(B)>0$，在 $B$ 发生的条件下 $A$ 发生的概率为
\[
\Pp(A\mid B)=\frac{\Pp(A\cap B)}{\Pp(B)}.
\]
直观上：样本空间被“缩小”到 $B$，然后数 $A\cap B$ 占 $B$ 的比例。
\end{definition}

\begin{theorem}[乘法公式]
\[
\Pp(A\cap B)=\Pp(B)\Pp(A\mid B)=\Pp(A)\Pp(B\mid A).
\]
更一般地，
\[
\Pp(A_1\cap\cdots\cap A_n)=\Pp(A_1)\Pp(A_2\mid A_1)\cdots\Pp(A_n\mid A_1\cap\cdots\cap A_{n-1}).
\]
\end{theorem}

\begin{theorem}[全概率公式]
若 $B_1,\ldots,B_n$ 两两互斥且 $\bigcup_i B_i=\Omega$（叫一个\textbf{划分}），则对任意事件 $A$，
\[
\Pp(A)=\sum_{i=1}^{n}\Pp(B_i)\Pp(A\mid B_i).
\]
\end{theorem}
\begin{proof}
$A=A\cap\Omega=A\cap\bigcup_i B_i=\bigcup_i(A\cap B_i)$，右端各项两两互斥（因为 $B_i$ 互斥），由公理 3 与乘法公式即得。
\end{proof}

\begin{theorem}[贝叶斯公式]
在同样的记号下，
\[
\Pp(B_j\mid A)=\frac{\Pp(B_j)\Pp(A\mid B_j)}{\sum_i\Pp(B_i)\Pp(A\mid B_i)}.
\]
\end{theorem}
\begin{proof}
分子由乘法公式为 $\Pp(A\cap B_j)=\Pp(A)\Pp(B_j\mid A)$；分母是全概率公式给出的 $\Pp(A)$。相除即得。
\end{proof}

\begin{teach}[贝叶斯公式的“先验—似然—证据”读法]
$\Pp(B_j)$ 是\textbf{先验}（看到证据前认为 $B_j$ 成立的概率），$\Pp(A\mid B_j)$ 是\textbf{似然}（在 $B_j$ 成立时看到 $A$ 的概率），分母 $\Pp(A)$ 是\textbf{证据}（看到 $A$ 的总概率）。读题时先把三者分别指出来，再套公式，比直接背公式可靠得多。
\end{teach}

\section{独立性}
\begin{definition}[事件的独立性]
若 $\Pp(A\cap B)=\Pp(A)\Pp(B)$，则称 $A$ 与 $B$ \textbf{相互独立}。等价地（当 $\Pp(B)>0$），$\Pp(A\mid B)=\Pp(A)$。
\end{definition}

\begin{pitfall}[互斥 $\ne$ 独立]
两个概率都为正的事件若互斥，则 $\Pp(A\cap B)=0\ne\Pp(A)\Pp(B)$，\textbf{一定不独立}。互斥说“不能同时发生”，独立说“互不影响”，二者方向相反。这是概率论里最常见的一处概念混淆。
\end{pitfall}

\begin{definition}[多个事件的独立性]
事件 $A_1,\ldots,A_n$ 相互独立，若对任意 $1\le i_1<\cdots<i_k\le n$ 都有
\[
\Pp(A_{i_1}\cap\cdots\cap A_{i_k})=\Pp(A_{i_1})\cdots\Pp(A_{i_k}).
\]
\textbf{两两独立不蕴含相互独立}——验证时不能只检查 $k=2$。
\end{definition}

\section{分层例题与解析}
\begin{problem}{例 1：古典概型}
掷两枚均匀的骰子，求点数之和为 $7$ 的概率。
\end{problem}
\hint 样本空间按“有序对”数，共 $36$ 个；再数有利结果。

\sol $\Omega=\{(i,j):1\le i,j\le6\}$，$|\Omega|=36$。和为 $7$ 的有序对为 $(1,6),(2,5),(3,4),(4,3),(5,2),(6,1)$，共 $6$ 个。故
\[
\Pp(A)=\frac6{36}=\frac16.
\]

\begin{problem}{例 2：容斥与“或”}
一副 $52$ 张牌中随机抽一张，求它是红桃或 $K$ 的概率。
\end{problem}
\hint $|\text{红桃}|=13$、$|K|=4$、红桃 $K$ 有 $1$ 张。

\sol
\[
\Pp=\frac{13+4-1}{52}=\frac{16}{52}=\frac4{13}.
\]

\begin{problem}{例 3：条件概率}
一个家庭有两个孩子，已知其中至少有一个是女孩，求两个孩子都是女孩的概率。
\end{problem}
\hint 样本空间是 $\{BB,BG,GB,GG\}$（$B$ 男、$G$ 女，顺序表示出生先后）；条件“至少一个女孩”排除了 $BB$。

\sol 设 $A$ 为“都是女孩”，$B$ 为“至少一个女孩”。$\Pp(B)=3/4$，$\Pp(A\cap B)=\Pp(A)=1/4$，故
\[
\Pp(A\mid B)=\frac{1/4}{3/4}=\frac13.
\]
\textbf{不是 $1/2$。}若把条件改成“较大的那个是女孩”（即先出生的是女孩），答案才是 $1/2$。两个条件的区别是本题的全部内容。

\begin{teach}[“已知有一个女孩”为什么答案会变]
这道题在概率论教学史上很有名。要点在于：\textbf{信息是如何产生的}。如果这句话来自“我随机告诉你其中一个孩子的性别，结果是女孩”，答案还是 $1/2$；如果来自“我检查了所有孩子，发现至少一个是女孩”，答案是 $1/3$。让学生把两种情形的样本空间分别写出来，比争辩更有用。
\end{teach}

\begin{problem}{例 4：全概率与贝叶斯}
某疾病在人群中的患病率是 $0.1\%$。一种检测的灵敏度是 $99\%$（患病者检测为阳性），特异度是 $99\%$（健康者检测为阴性）。若某人检测为阳性，他真正患病的概率是多少？
\end{problem}
\hint 用全概率公式算 $\Pp(+)$，再用贝叶斯。

\sol 设 $D$ 为患病，$+$ 为阳性。$\Pp(D)=0.001$，$\Pp(+\mid D)=0.99$，$\Pp(+\mid\overline D)=0.01$。由全概率公式
\[
\Pp(+)=0.001\cdot0.99+0.999\cdot0.01=0.00099+0.00999=0.01098.
\]
由贝叶斯公式
\[
\Pp(D\mid +)=\frac{0.001\cdot0.99}{0.01098}=\frac{0.00099}{0.01098}\approx0.0902.
\]
约 $9\%$。\textbf{即使检测准确率高达 $99\%$，阳性者真正患病的概率也不到十分之一}——因为健康人群实在太大。

\begin{pitfall}[基础率越低，越不能只看准确率]
上题的结论对直觉冲击很大，值得反复讲。要点：阳性者中绝大多数是“健康人被误判”的，因为 $\Pp(\overline D)\approx1$ 而误判率 $0.01$ 已经足够大。这提示：\textbf{任何检测类问题都必须先问基础率。}
\end{pitfall}

\begin{problem}{统一练习：生日问题}
一个班有 $23$ 人，求至少两人同生日的概率（假设 $365$ 天等可能，忽略闰年）。
\end{problem}
\hint 算补事件“所有人生日互不相同”。

\sol
\[
\Pp(\text{无人同日})=\frac{365\cdot364\cdot\cdots\cdot(365-22)}{365^{23}}
=\prod_{k=0}^{22}\left(1-\frac{k}{365}\right)\approx0.4927.
\]
故所求概率为
\[
1-0.4927\approx0.5073.
\]

\begin{teach}[为什么 $23$ 人就超过一半]
把 $23$ 人两两配对，共有 $\dbinom{23}{2}=253$ 对，每对同生日的概率约 $1/365$。虽然这些事件不独立，但“期望对数”约为 $253/365\approx0.69$，已经足以让概率过半。这个估计解释了数量级，严格的概率要用上面的乘积。\textbf{“数对数”是概率题里最常见的降维手段。}
\end{teach}

\begin{problem}{统一练习：独立与互斥的辨析}
判断下列说法是否正确，并说明理由：
\begin{enumerate}[label=(\alph*)]
\item 若 $A,B$ 互斥，则 $\Pp(A\cup B)=\Pp(A)+\Pp(B)$。
\item 若 $A,B$ 独立，则 $\Pp(A\cup B)=\Pp(A)+\Pp(B)$。
\item 若 $\Pp(A\mid B)=\Pp(A)$，则 $A,B$ 独立。
\item 若 $A,B$ 互斥且 $\Pp(A),\Pp(B)>0$，则 $A,B$ 独立。
\end{enumerate}
\end{problem}
\hint 逐条对照定义。

\sol (a) 正确，这是公理 3 的直接应用。(b) 错误，独立时应为 $\Pp(A\cup B)=\Pp(A)+\Pp(B)-\Pp(A)\Pp(B)$。(c) 正确（需 $\Pp(B)>0$），这是独立的等价定义。(d) 错误，互斥且概率为正的事件必不独立。

\chapter{随机变量、期望与方差}
\begin{chapterintro}[章节导读]
把随机事件的“是/否”变成一个数，就得到随机变量。有了随机变量，概率问题的大部分内容可以归结为计算两个数：\textbf{期望}（平均值）与\textbf{方差}（波动）。本章最重要的是\textbf{期望的线性性}——它不要求任何独立性，是概率工具里性价比最高的一条性质。
\end{chapterintro}

\section{随机变量与分布}
\begin{definition}[随机变量]
从样本空间 $\Omega$ 到 $\R$ 的函数 $X:\Omega\to\R$ 叫随机变量。若 $X$ 只取有限或可数多个值，叫\textbf{离散随机变量}。
\end{definition}

\begin{definition}[分布列与期望]
设离散随机变量 $X$ 取值 $x_1,x_2,\ldots$，且 $\Pp(X=x_i)=p_i$（$\sum_i p_i=1$），则
\[
\E X=\sum_i x_ip_i
\]
叫 $X$ 的\textbf{数学期望}（均值）。
\end{definition}

\begin{theorem}[期望的线性性]
对任意随机变量 $X,Y$（不要求独立）与常数 $a,b$，
\[
\E[aX+bY]=a\E X+b\E Y.
\]
\end{theorem}
\begin{proof}
只需对有限个取值的情形验证。设 $X$ 取值 $x_i$、$Y$ 取值 $y_j$，联合概率 $p_{ij}=\Pp(X=x_i,Y=y_j)$，则
\[
\E[X+Y]=\sum_{i,j}(x_i+y_j)p_{ij}
=\sum_{i,j}x_ip_{ij}+\sum_{i,j}y_jp_{ij}
=\sum_i x_i\Pp(X=x_i)+\sum_j y_j\Pp(Y=y_j)=\E X+\E Y.
\]
\textbf{证明中没有用到任何独立性假设}——这是线性性如此强大的原因。
\end{proof}

\begin{teach}[指示变量法]
若 $X$ 是“满足条件 $i$ 的对象个数”，令 $I_i=\iva{\text{对象 }i\text{ 满足条件}}$，则 $X=\sum_i I_i$，而 $\E I_i=\Pp(\text{对象 }i\text{ 满足条件})$。于是
\[
\E X=\sum_i\Pp(\text{对象 }i\text{ 满足条件}).
\]
\textbf{把“数个数”变成“求单个对象的概率再求和”}，这是处理计数型随机变量最通用的办法。
\end{teach}

\begin{definition}[方差与标准差]
\[
\Var X=\E[(X-\E X)^2]=\E[X^2]-(\E X)^2,\qquad
\sigma_X=\sqrt{\Var X}.
\]
\end{definition}
\begin{theorem}[方差的性质]
$\Var(aX+b)=a^2\Var X$。若 $X,Y$ \textbf{独立}，则 $\Var(X+Y)=\Var X+\Var Y$。
\end{theorem}
\begin{proof}
$\Var(aX+b)=\E[(aX+b-a\E X-b)^2]=a^2\E[(X-\E X)^2]=a^2\Var X$。独立时
\[
\E[(X+Y)^2]=\E[X^2]+2\E[XY]+\E[Y^2]=\E[X^2]+2\E X\E Y+\E[Y^2],
\]
减去 $(\E X+\E Y)^2$ 即得。
\end{proof}

\begin{pitfall}[方差的线性性需要独立性]
期望的线性性无条件成立，\textbf{方差的加法需要独立}。若 $X=Y$，则 $\Var(X+Y)=\Var(2X)=4\Var X$ 而不是 $2\Var X$。这是初学者最常犯的错之一。
\end{pitfall}

\section{常用离散分布}
\begin{center}
\begin{tabular}{llll}
\toprule
分布 & 分布列 & 期望 & 方差\\
\midrule
$0$-$1$ 分布（参数 $p$） & $\Pp(X=1)=p$ & $p$ & $p(1-p)$\\
二项分布 $B(n,p)$ & $\Pp(X=k)=\dbinom nkp^k(1-p)^{n-k}$ & $np$ & $np(1-p)$\\
几何分布（参数 $p$） & $\Pp(X=k)=(1-p)^{k-1}p$ & $1/p$ & $(1-p)/p^2$\\
泊松分布（参数 $\lambda$） & $\Pp(X=k)=\dfrac{\lambda^k\mathrm e^{-\lambda}}{k!}$ & $\lambda$ & $\lambda$\\
\midrule
均匀分布（$1,\ldots,n$） & $\Pp(X=k)=1/n$ & $(n+1)/2$ & $(n^2-1)/12$\\
\bottomrule
\end{tabular}
\end{center}

\begin{proof}[二项分布期望的推导（两种方法）]
\textbf{方法一（定义）}：
\[
\E X=\sum_{k=0}^{n}k\binom nkp^k(1-p)^{n-k}
=np\sum_{k=1}^{n}\binom{n-1}{k-1}p^{k-1}(1-p)^{n-k}=np.
\]
\textbf{方法二（线性性）}：把 $X$ 写成 $n$ 个独立的 $0$-$1$ 变量之和 $X=\sum_i I_i$，其中 $\E I_i=p$。由线性性 $\E X=np$。

第二种方法只需两行，而且\textbf{揭示了二项分布的本质：$n$ 次独立重复试验的成功次数}。方差的推导同理：$\Var X=\sum_i\Var I_i=np(1-p)$。
\end{proof}

\begin{proof}[几何分布期望的推导]
设 $\E X=\mu$。第一次试验若失败（概率 $1-p$），已经用掉一次且要重新开始，故
\[
\mu=1+(1-p)\mu\quad\Longrightarrow\quad \mu=\frac1p.
\]
这是\textbf{首步分析}的标准写法：按第一步的结果分类，得到关于 $\mu$ 的方程。
\end{proof}

\begin{teach}[“首步分析”是概率递推的万能钥匙]
凡是一个过程每一步形态相同、状态有限，就可以对每个状态 $s$ 设 $f(s)=\E[\text{从 }s\text{ 到结束的代价}]$，按第一步分类列出
\[
f(s)=\sum_t p_{s,t}\bigl(c_{s,t}+f(t)\bigr).
\]
第四章博弈论里的期望值分析、第 9 章的“掷骰子走路”都是这个式子。\textbf{先设未知函数，再按首步分类}，比硬算分布列可靠。
\end{teach}

\section{两个不等式}
\begin{theorem}[马尔可夫不等式]
设 $X\ge0$，则对任意 $a>0$，
\[
\Pp(X\ge a)\le\frac{\E X}{a}.
\]
\end{theorem}
\begin{proof}
由 $X\ge a\iva{X\ge a}$（因为 $X\ge0$），取期望得 $\E X\ge a\Pp(X\ge a)$。
\end{proof}

\begin{theorem}[切比雪夫不等式]
对任意随机变量 $X$ 与 $a>0$，
\[
\Pp\bigl(|X-\E X|\ge a\bigr)\le\frac{\Var X}{a^2}.
\]
\end{theorem}
\begin{proof}
对非负随机变量 $Y=(X-\E X)^2$ 与阈值 $a^2$ 用马尔可夫不等式：
\[
\Pp\bigl((X-\E X)^2\ge a^2\bigr)\le\frac{\E[(X-\E X)^2]}{a^2}=\frac{\Var X}{a^2}.
\]
\end{proof}

\begin{tip}[两个不等式怎么用]
它们给出的都是\textbf{上界}，通常很松，但足以证明“偏离平均值的概率很小”。在算法分析中，常用切比雪夫不等式把“期望 $O(n)$”加强为“以高概率 $O(n)$”。第二部分第 13 章的 Miller--Rabin 与 Pollard--Rho 都要靠这类论证。
\end{tip}

\section{分层例题与解析}
\begin{problem}{例 1：期望的线性性}
同时掷 $n$ 枚均匀骰子，求点数之和的期望。
\end{problem}
\hint 把总和写成 $n$ 个单枚骰子点数之和。

\sol 每枚骰子的期望为 $(1+2+\cdots+6)/6=7/2$。由线性性，总和期望为 $7n/2$。

\begin{problem}{例 2：指示变量法}
随机排列 $n$ 个数，求不动点（满足 $a_i=i$ 的位置）个数的期望。
\end{problem}
\hint 令 $I_i=\iva{a_i=i}$；$\Pp(a_i=i)=1/n$。

\sol $\E[\text{不动点数}]=\sum_{i=1}^{n}\Pp(a_i=i)=n\cdot\dfrac1n=1$。\textbf{与 $n$ 无关}：无论排列多长，平均总有一个不动点。注意我们完全不需要知道不动点个数的分布。

\begin{problem}{例 3：首步分析}
从原点出发，每次等概率向右走 $1$ 格或向左走 $1$ 格。求首次到达 $+3$ 的期望步数。
\end{problem}
\hint 一维对称随机游走首次到达 $+3$ 的期望步数是无穷大；先证明这件事。

\sol 设 $f(k)$ 为从 $k$ 出发首次到达 $3$ 的期望步数。按第一步分类：
\[
f(k)=1+\frac12 f(k+1)+\frac12 f(k-1),\qquad f(3)=0.
\]
整理得二阶差分方程
\[
f(k+1)-2f(k)+f(k-1)=-2.
\]
它的一般解是 $f(k)=-k^2+ak+b$——一条开口向下的抛物线。但 $f(k)$ 是期望步数，必须 $f(k)\ge0$，而当 $k\to-\infty$ 时 $-k^2+ak+b\to-\infty$，矛盾。

所以	extbf{不存在满足条件的非负解}，即期望时间为无穷大。事实上对称随机游走必然到达 $+3$（一维随机游走常返，见第 10 章定理），但期望时间是 $\infty$。

\begin{pitfall}[“必然发生”不等于“期望有限”]
一维对称随机游走以概率 $1$ 回到原点，但期望回归时间是 $\infty$。初学者常把“迟早会发生”误当成“平均很快发生”。\textbf{概率为 $1$ 只说明时间随机变量几乎必然有限，不说明它的期望有限。}
\end{pitfall}

\begin{problem}{例 4：优惠券收集}
一包零食附赠 $n$ 种卡片之一（等概率）。要集齐全部 $n$ 种，期望要买多少包？
\end{problem}
\hint 分阶段：已经收集了 $k$ 种时，再获得一种新卡片的期望包数是 $n/(n-k)$。

\sol 设 $T$ 为总包数，$T=T_0+T_1+\cdots+T_{n-1}$，其中 $T_k$ 是“从 $k$ 种到 $k+1$ 种”所需的包数。此时每包获得新卡片的概率为 $(n-k)/n$，由几何分布 $\E T_k=n/(n-k)$。于是
\[
\E T=\sum_{k=0}^{n-1}\frac{n}{n-k}=n\sum_{j=1}^{n}\frac1j=nH_n.
\]
其中 $H_n$ 是第 $n$ 个调和数，$H_n=\ln n+\gamma+O(1/n)$。所以集齐 $n$ 种卡片大约需要 $\mathrm e\,n\ln n$ 量级的时间。

\begin{teach}[“分阶段”是期望递推的第二种套路]
第 9 章讲过按\textbf{首步}分类，这里按\textbf{阶段}分类：把一个复杂过程切成若干段，每段内部形态相同、期望好算，再用线性性相加。优惠券收集、随机打乱的比较次数、快速选择的比较次数，都靠这一招。
\end{teach}

\begin{problem}{例 5：切比雪夫不等式}
掷 $n$ 枚均匀骰子，用切比雪夫不等式估计点数之和偏离期望超过 $10\sqrt n$ 的概率上界。
\end{problem}
\hint 单枚骰子点数 $X_i$：$\E X_i=7/2$、$\Var X_i=(1^2+\cdots+6^2)/6-(7/2)^2=91/6-49/4=35/12$。

\sol 总和 $S=\sum_iX_i$，$\E S=7n/2$，$\Var S=n\cdot\dfrac{35}{12}$。由切比雪夫不等式，
\[
\Pp\bigl(|S-\E S|\ge10\sqrt n\bigr)\le\frac{35n/12}{100n}=\frac{35}{1200}=\frac7{240}\approx0.0292.
\]
注意这个上界\textbf{与 $n$ 无关}——它虽然正确，但比较粗糙；真实的偏差概率随 $n$ 增大迅速下降（可用中心极限定理给出更精确的估计）。

\begin{problem}{统一练习：把“或”变成“至少一个”}
$n$ 个人各自独立地把自己的帽子随机放进一个帽子堆再随机取回一顶。求没有人拿到自己帽子的概率（$n=5$ 时计算数值）。
\end{problem}
\hint 用容斥；这正是第三部分第 3 章的错位排列。

\sol 设 $A_i$ 为“第 $i$ 个人拿到自己的帽子”。$\Pp(A_i)=1/n$，更一般地 $\Pp(A_{i_1}\cap\cdots\cap A_{i_k})=(n-k)!/n!$。由容斥，
\[
\Pp\left(\bigcap_i\overline{A_i}\right)
=\sum_{k=0}^{n}(-1)^k\binom nk\frac{(n-k)!}{n!}
=\sum_{k=0}^{n}\frac{(-1)^k}{k!}.
\]
$n=5$ 时通分到 $120$：
\[
1-1+\frac12-\frac16+\frac1{24}-\frac1{120}
=\frac{120-120+60-20+5-1}{120}=\frac{44}{120}=\frac{11}{30}\approx0.3667.
\]
当 $n\to\infty$ 时趋于 $\mathrm e^{-1}\approx0.3679$。

\begin{teach}[概率与计数的接口]
本题的容斥过程与第三部分第 3 章推导错位排列数 $D_n$ 完全一样，只是在那里数“排列个数”，在这里数“概率”。\textbf{古典概型下，概率问题就是计数问题除以 $|\Omega|$。}这条对应关系值得反复强调。
\end{teach}

\chapter{常用概率模型与随机化算法}
\begin{chapterintro}[章节导读]
本章把前两章的工具用在两个方向：一是算法设计中的\textbf{随机化}——用随机选择把最坏情况变成期望情况；二是数论与概率的交叉——随机大整数的素性、随机游走的回归。这两块内容在竞赛中出现频率很高，而且都与第二部分数论直接衔接。
\end{chapterintro}

\section{两类随机化算法}
\begin{definition}[拉斯维加斯算法]
总是给出正确答案，但运行时间是随机变量。例如随机快速排序、随机快速选择。
\end{definition}
\begin{definition}[蒙特卡洛算法]
运行时间确定，但有小概率给出错误答案。例如 Miller--Rabin 素性测试、Pollard--Rho 因数分解。
\end{definition}

\begin{teach}[两类算法的区别要讲透]
拉斯维加斯算法“可能慢，但不会错”；蒙特卡洛算法“快，但可能错”。重复运行 $k$ 次可以把蒙特卡洛算法的错误率降到 $p^k$——这就是第二部分第 12 章 Miller--Rabin 用若干个底数把错误率压到忽略不计的原因。
\end{teach}

\section{随机打乱}
\begin{theorem}[Fisher--Yates 洗牌]
从后往前遍历 $i=n,n-1,\ldots,2$，每次在 $[1,i]$ 中等概率选一个 $j$ 并交换 $a_i,a_j$。该算法产生全部 $n!$ 个排列的概率相等。
\end{theorem}
\begin{proof}[归纳]
记 $P_n$ 为算法对 $n$ 个元素产生均匀分布。$n=1$ 显然。假设 $P_{n-1}$ 成立。第一步在 $[1,n]$ 中等概率选 $j$ 并交换 $a_n,a_j$，故元素 $x$ 落到位置 $n$ 的概率为 $1/n$，与 $x$ 无关。给定 $x$ 在位置 $n$ 后，剩下 $n-1$ 个元素的排列由同一算法在 $n-1$ 个元素上产生，由归纳假设是均匀的。于是每个完整排列的概率为
\[
\frac1n\cdot\frac1{(n-1)!}=\frac1{n!}.
\]
\end{proof}

\begin{pitfall}[“每个元素随机交换一次”不是均匀洗牌]
对 $i=1,\ldots,n$ 都执行“与随机位置交换”会产生 $n^n$ 个等概率的结果，但 $n^n$ 不是 $n!$ 的倍数时无法均分到排列上，分布\textbf{不均匀}。错误写法是竞赛里最常出现的随机化 bug 之一。
\end{pitfall}

\section{随机快速选择}
\begin{theorem}[随机快速选择的期望复杂度]
在 $n$ 个元素中随机选主元求第 $k$ 小，期望比较次数为 $O(n)$。
\end{theorem}
\begin{proof}[按主元排名分类]
设 $C(n)$ 为在 $n$ 个元素上的期望比较次数，$C(0)=C(1)=0$。若主元的排名为 $k$，则划分为左右两部分，规模为 $k-1$ 与 $n-k$，比较次数恰为 $n-1$。于是
\[
C(n)\le n+\frac1n\sum_{k=1}^{n}C\bigl(\max(k,n-k)\bigr).
\]
把 $k$ 与 $n-k$ 配对（$k=1,\ldots,n$ 时 $\max(k,n-k)$ 的取值在 $\lceil n/2\rceil$ 到 $n-1$ 之间，且每个值至多出现两次）：
\[
\frac1n\sum_{k=1}^{n}C\bigl(\max(k,n-k)\bigr)
\le\frac{2}{n}\sum_{j=\lceil n/2\rceil}^{n-1}C(j).
\]
$C$ 单调不减，且项数为 $n/2$、最大下标 $n-1\le\dfrac{3n}{4}$（$n\ge4$），故
\[
C(n)\le n+\frac{2}{n}\cdot\frac n2\cdot C\!\left(\frac{3n}{4}\right)
=n+C\!\left(\frac{3n}{4}\right).
\]
设 $C(n)\le cn$（$n\ge4$），则
\[
C(n)\le n+c\cdot\frac{3n}{4}=\left(1+\frac{3c}{4}\right)n\le cn
\quad\Longleftrightarrow\quad c\ge4.
\]
取 $c=4$，归纳得 $C(n)\le4n$。更精细的分析可把常数降到约 $3.39$；本估计已足以证明线性。
\end{proof}

\begin{teach}[“好主元”的期望试验次数]
每次随机选到中间一半的概率是 $1/2$，所以得到好主元的期望试验次数是 $2$（几何分布，参数 $1/2$）。把“每次都要重新随机”写成“期望 $2$ 次”，是这个证明里唯一的技巧，其余都是代数。
\end{teach}

\section{随机游走}
\begin{theorem}[一维随机游走常返]
在一维整数格 $\Z$ 上，从原点出发每步等概率向 $\pm1$ 移动，则以概率 $1$ 回到原点。
\end{theorem}
\begin{proof}[用组合计数]
回到原点的充要条件是走偶数步 $2n$，其中恰有 $n$ 步向右。设 $u_{2n}$ 为“第 $2n$ 步首次回到原点”的概率，$p_{2n}=\dbinom{2n}{n}2^{-2n}$ 为“第 $2n$ 步在原点”的概率。由首步分解
\[
p_{2n}=\sum_{k=1}^{n}u_{2k}p_{2n-2k},\qquad p_0=1.
\]
设 $P(z)=\sum_n p_{2n}z^n$、$U(z)=\sum_n u_{2n}z^n$，则 $P(z)=1+P(z)U(z)$，即 $U(z)=1-1/P(z)$。于是
\[
\sum_{n\ge0}u_{2n}=U(1)=1-\frac1{P(1)}.
\]
而
\[
P(1)=\sum_n\binom{2n}{n}2^{-2n}=\frac1{\sqrt{1-4\cdot(1/4)}}=\infty
\]
（二项式系数生成函数 $\sum\binom{2n}{n}x^n=(1-4x)^{-1/2}$ 在 $x=1/4$ 处发散）。故 $\sum u_{2n}=1$，即以概率 $1$ 回到原点。
\end{proof}

\begin{tip}[维度的分界]
一维与二维简单随机游走\textbf{常返}（必然回到原点），三维及更高维\textbf{瞬态}（有正概率永不回来）。这个结论的证明要用到 $d$ 维的返回概率渐近 $p_{2n}\asymp n^{-d/2}$：$\sum n^{-d/2}$ 当 $d\le2$ 发散、当 $d\ge3$ 收敛。这个事实在第五部分处理生成函数时会给出一个干净的推导。
\end{tip}

\section{与数论的交叉}
\subsection{随机大整数的素性}
由素数定理，$x$ 附近的随机整数是素数的概率约为 $1/\ln x$。这给出一个实用的搜索策略：随机取奇数、过小素数试除、再做 Miller--Rabin。第二部分第 12 章会给出完整流程。

\subsection{二次剩余与“是否为平方”}
在模奇素数 $p$ 下，$\dfrac{p-1}{2}$ 个非零剩余都是二次剩余，所以“随机非零数是二次剩余”的概率是 $1/2$。第二部分第 6 章用这一点分析 Tonelli--Shanks 的期望表现。

\subsection{随机化 gcd 与 Pollard--Rho}
Pollard--Rho 的期望复杂度 $O(n^{1/4})$ 依赖于生日悖论：随机函数 $f$ 的迭代序列出现碰撞的期望时间是 $O(\sqrt{n})$。这与本章开头生日问题的“数对数”估计是同一件事。

\begin{problem}{统一练习：生日悖论的精确形式}
在 $1,\ldots,n$ 中独立随机取 $k$ 个数（可重复），求至少两个数相等的概率；并由此估计 Pollard--Rho 需要多少次迭代才能出现碰撞。
\end{problem}
\hint 算补事件；用 $1-x\le\mathrm e^{-x}$ 估计。

\sol 全部不同的概率为
\[
\prod_{i=0}^{k-1}\left(1-\frac in\right)
\le\exp\left(-\sum_{i=0}^{k-1}\frac in\right)
=\exp\left(-\frac{k(k-1)}{2n}\right).
\]
故碰撞概率至少为
\[
1-\exp\left(-\frac{k(k-1)}{2n}\right).
\]
取 $k=\sqrt{2n}$ 得碰撞概率至少 $1-\mathrm e^{-1}\approx0.63$，取 $k=2\sqrt n$ 则几乎必然碰撞。Pollard--Rho 正是利用这一点：在 $O(\sqrt p)$ 次迭代内以高概率找到模 $p$ 的因子。

\begin{teach}[把生日问题背下来是值得的]
“$k\approx\sqrt n$ 时碰撞概率过半”这个数量级估计，在哈希、随机游走、因数分解、字符串匹配里反复出现。\textbf{用 $1-x\le\mathrm e^{-x}$ 把乘积变成指数}，是得到这个估计的标准动作。
\end{teach}

\section{综合例题}
\begin{problem}{综合 1：期望 DP}
掷一枚均匀骰子，从格子 $0$ 出发，点数即前进格数。若到达或超过格子 $10$ 就停止。求期望掷骰次数。
\end{problem}
\hint 设 $f(i)$ 为从格子 $i$ 出发的期望次数；$i\ge10$ 时 $f(i)=0$。

\sol 设 $f(i)$ 为从格子 $i$ 出发的期望投掷次数。越界即停，故
\[
f(i)=0\ (i\ge10),\qquad
f(i)=1+\frac16\sum_{d=1}^{6}f(i+d)\ (0\le i<10).
\]
从 $i=9$ 往 $i=0$ 倒推（下表保留六位小数）：
\[
\begin{array}{c|cccccccccc}
i&0&1&2&3&4&5&6&7&8&9\\\hline
f(i)&3.3237&3.0433&2.7752&2.5216&2.1614&1.8526&1.5880&1.3611&1.1667&1.0000
\end{array}
\]
故所求期望为 $f(0)\approx3.32$ 次。

（当 $3\le i\le9$ 时有简洁的闭式 $f(i)=\left(\dfrac76\right)^{9-i}$，例如 $f(9)=1$、$f(6)=\dfrac{343}{216}$、$f(3)=\dfrac{117649}{46656}$；$i\le2$ 时该式失效。）

\begin{teach}[期望 DP 的三步]
\textbf{设对状态、写对转移、注意边界。}本题的状态是“当前格子”，转移是“六种点数各六分之一”，边界是“到达或超过 $10$ 就停”。三步都对，剩下的只是倒推填表。写代码时把 $f(\min(10,i+d))$ 的截断写清楚，是最容易出错的地方。
\end{teach}

倒推填表的参考实现：
\begin{lstlisting}[language=C++]
double f[20];
for (int i = 9; i >= 0; --i) {
    double s = 0;
    for (int d = 1; d <= 6; ++d) s += f[min(10, i + d)];
    f[i] = 1 + s / 6;          // f[10..15] 初值即为 0
}
// 答案 f[0]，约为 3.3237
\end{lstlisting}

\begin{problem}{综合 2：随机化与确定性}
一个数组有 $n$ 个元素，其中恰好有一个“特殊”元素。你每次可以随机选一个元素检查。求找到它的期望检查次数。若允许先随机打乱再顺序检查，期望次数是多少？若已知特殊元素在数组的后半段呢？
\end{problem}
\hint 分别用几何分布与均匀分布的期望。

\sol 随机检查：每次命中概率 $1/n$，期望 $n$ 次。随机打乱后顺序检查等价于在均匀随机的位置中找特殊元素，期望位置是 $(n+1)/2$。已知在后半段：均匀分布在 $n/2+1,\ldots,n$，期望为
\[
\frac{(n/2+1)+n}{2}=\frac{3n}{4}+\frac12.
\]
\textbf{信息能换时间}：知道得越多，期望检查次数越少。

\begin{problem}{综合 3：概率与数论}
设 $p$ 为奇素数，$a$ 为模 $p$ 的非零二次剩余。随机选 $x\in\{1,\ldots,p-1\}$，求 $x^2\equiv a\pmod p$ 的解的个数；并由此说明为什么 Tonelli--Shanks 需要“找非平凡平方根”。
\end{problem}
\hint 模 $p$ 的方程 $x^2=a$ 恰有两个解 $x$ 与 $p-x$。

\sol 恰有 $2$ 个解。所以在 $p-1$ 个非零 $x$ 中随机取，命中一个解的概率是 $2/(p-1)$——太小，不能靠随机试。Tonelli--Shanks 的做法是反过来：\textbf{随机取 $x$，计算 $x^2$，若它不是 $a$ 就得到一个非平凡平方根}（即 $x^2$ 的某个平方根 $y$ 满足 $y\not\equiv\pm x$），从而可以逐步降阶。这正是第二部分第 6 章的内容。

\begin{problem}{综合 4：把方差用起来}
一个哈希表用 $n$ 个桶、$m$ 个键（$m=n$），每个键独立均匀地选一个桶。设 $X$ 为最大桶长的上界估计：用马尔可夫不等式或切比雪夫不等式，估计“某个桶至少有 $10$ 个键”的概率上界。
\end{problem}
\hint 先算单个桶的键数 $B_i\sim B(m,1/n)$，再用联合界（union bound）。

\sol 固定一个桶，$B_i\sim B(n,1/n)$，$\E B_i=1$、$\Var B_i=1-\dfrac1n$。由切比雪夫不等式
\[
\Pp(B_i\ge10)\le\Pp(|B_i-1|\ge9)\le\frac{1-1/n}{81}\approx0.0123.
\]
由联合界，存在某个桶键数 $\ge10$ 的概率至多
\[
n\cdot\frac{1-1/n}{81}\approx\frac{n}{81},
\]
当 $n>81$ 时这个上界超过 $1$，\textbf{说明切比雪夫给出的界太松}。这是不等式使用的常态：它们给出量级，不给出精确值。精确分析要用泊松化（键数近似 $\mathrm{Pois}(1)$）。

\begin{teach}[第四部分的结课检查]
让学生独立完成三道小题：
\begin{enumerate}
\item 证明：随机排列 $n$ 个数的逆序对数的期望是 $\dfrac{n(n-1)}4$（提示：用指示变量，每对贡献 $1/2$）。
\item 从 $1$ 到 $n$ 中等概率取一个数，重复取直到出现重复，求期望取多少次（提示：分阶段，答案是 $\dfrac{n(n+1)}{2(n-1)}$ 量级；先算 $n$ 小的情形）。
\item 一个随机化算法每次运行成功概率为 $p$。要使命失败概率小于 $10^{-6}$，需要独立运行多少次？用 $p=1/2$ 计算。
\end{enumerate}
三题全对即可进入第五部分。第三题的答案是 $k\ge20$（因为 $2^{-20}\approx9.5\times10^{-7}$），算不出就回到第 10 章“重复运行压低错误率”那一节。
\end{teach}


% ============================================================================
%  第五部分  多项式与生成函数
% ============================================================================
\BeginPart{12}{3}{V}{第五部分}{多项式与生成函数}
% !TeX program = xelatex
% ============================================================================
%  第五部分  多项式与生成函数
%
%  本部分由使用者另行提供，此处仅占位。保留以下内容的目的：
%    1. 让全书结构完整、可独立编译（XeLaTeX 两遍即可得到含第五部分的目录）；
%    2. 为后续补写标出章节编号、记号与交叉引用的落点。
%
%  补写时请注意：
%    - 章号从“第 1 章”开始（\\BeginPart 已把 chapter 计数器归零）；
%    - 超链接目标使用 \\theHchapter = <partno>.<chapter>，与前面四部分不会冲突；
%    - 记号沿用全书速查表：\\Z \\N \\Pp \\eps \\one \\id \\ord \\lcm \\rad \\iva
%     \\floor \\ceil，概率部分另有 \\E \\Var；若需新的宏，请加在 main.tex
%     的“简写命令”一节，不要在各部分文件里重复定义。
%    - 前置知识交叉引用：第三部分（组合恒等式、递推、格路）、第四部分
%     （概率生成函数、随机游走）、第二部分（模意义下的多项式与 NTT）。
% ============================================================================

\chapter{多项式与生成函数（占位）}
\begin{chapterintro}[章节导读]
本部分计划覆盖三块内容：\textbf{多项式运算与快速傅里叶变换}、\textbf{形式幂级数与普通生成函数}、\textbf{指数生成函数与指数公式}。它们是第三部分组合恒等式与递推的“机械化”版本，也是第一部分复数单位根的自然应用场所。此处仅保留骨架，正文由使用者提供。
\end{chapterintro}

\section{计划目录}
\begin{enumerate}
\item 多项式的表示与基本运算（系数形式、点值形式、乘法逆元）
\item 快速傅里叶变换与数论变换（NTT）
\item 形式幂级数：加法、乘法、逆、开方、对数与指数
\item 普通生成函数与线性递推的闭式（Binet 公式、有理生成函数）
\item 指数生成函数与指数公式、Bell 数与第二类 Stirling 数
\item 多项式与组合计数的综合应用
\end{enumerate}

\section{与全书的衔接点}
\begin{itemize}
\item \textbf{承接第一部分第 6 章}：单位根与复数乘法在那里引入，本部分用它构造 DFT。
\item \textbf{承接第三部分第 2、4 章}：二项式定理、范德蒙与朱世杰恒等式都是生成函数系数的特例；斐波那契、卡特兰数列的闭式由有理／代数生成函数给出。
\item \textbf{承接第三部分第 6 章}：格路计数的反射原理与生成函数是同一问题的两种语言。
\item \textbf{承接第四部分第 9、10 章}：概率生成函数 $G_X(z)=\E[z^X]$ 把离散分布的卷积变成多项式乘法；生日悖论的碰撞概率也可由生成函数重新推导。
\item \textbf{承接第二部分}：模素数意义下的 NTT、以及“多项式同余”与“整数同余”的类比。
\end{itemize}

\begin{tip}[给补写者的一句话]
生成函数的全部力量来自一件事：\textbf{把“数数”翻译成“多项式的系数”}。凡是第三部分里用分类讨论得到的递推，都可以在这里用一次乘法或一次求逆重新得到。补写时建议先复用第三部分的例题，再推广到新场景。
\end{tip}


% ============================================================================
%  附录
% ============================================================================
% 附录章号改用 A--D；同时改写 \theHchapter，否则附录 A 的锚点名会与
% 第五部分第 1 章相同（都是 5.1），hyperref 会报重复目标。
\appendix
\renewcommand{\theHchapter}{\Alph{chapter}}
\renewcommand{\theHsection}{\theHchapter.\arabic{section}}
\renewcommand{\theHsubsection}{\theHsection.\arabic{subsection}}
\renewcommand{\theHdefinition}{\theHchapter.\arabic{definition}}
\renewcommand{\theHequation}{\theHchapter.\arabic{equation}}
\renewcommand{\theHtable}{\theHchapter.\arabic{table}}
\renewcommand{\theHfigure}{\theHchapter.\arabic{figure}}
\ctexset{chapter={format=\raggedright,
  nameformat=\Large\bfseries\color{ink},
  titleformat=\Huge\bfseries\color{ink},
  name={附录,},aftername=\par,number=\Alph{chapter}}}
% !TeX program = xelatex
% 本文件由 assemble.py 自 OI-Docs 原稿抽取生成，请勿手改（会被覆盖）。
% 源文件：Number Theory/number_theory_teacher.tex
% 源行号：3762-4149 / 4152-4350 / 4353-4639（数论）+ 1642-1728（博弈）
% 说明：已删去独立成册所需的导言区与页码重置；章节正文与交叉引用原样保留。

% ==========================================================================
% 附录：A--D。章号由 main.tex 中的 \appendix 切换为 A、B、C、D。
% ==========================================================================

\chapter{分层训练与参考解析}
\begin{chapterintro}[章节导读]
前十五章建立的是一条从基本数论工具一路走到高级专题的主线；从这里开始，正文不再继续引入新的知识体系，而是转向训练与教学资源。附录 A 把前面学过的方法重新组合进分层练习中，重点从“会不会公式”转向“能不能选对工具、检查条件并写出完整推导”。
\end{chapterintro}

本章是新增训练，不替代前文原题。学生版只保留题目；教师版在每题后给出结果、关键步骤与评价要点。建议先独立完成，再检查条件和推导，而不是仅比对最后一个数。

\section{A 组：整除、通解与基础建模}
\begin{problem}{A1：从 gcd 回代到不定方程}
求 $\gcd(414,662)$，给出 $414x+662y=6$ 的一组整数解与通解。
\end{problem}

\begin{pitfall}[讲解：先看这题到底在做什么]
先把 gcd 的回代结果写成一个线性组合，再整体乘到右端。这里练的是“exgcd 的输出怎么变成方程特解”，而不是重新证明裴蜀定理。
\end{pitfall}

\sol gcd 为 $2$。欧几里得回代得 $2=8\cdot414-5\cdot662$，乘三后得到 $(x_0,y_0)=(24,-15)$。通解为
\[
 x=24+331t,\qquad y=-15-207t,\quad t\in\Z.
\]

\begin{tip}[三档练习的用法]
A 组只用到第 1--3 章，做完应该能独立写出通解区间；B 组对应第 4--6 章，重点是\textbf{先列周期再套公式}；C、D 组要求你识别工具而不是回忆公式。建议每题限时 $15$ 分钟，写不出来就读解析的第一句，仍然卡住就先跳到下一组，回头再战。
\end{tip}

评价时应分别检查：整除条件、缩放方向、通解步长；只写一个特解不能算完成。


\begin{problem}{A2：正整数解的计数}
求 $14x+21y=147$ 的全部正整数解、解数以及两个变量的极值。
\end{problem}

\begin{pitfall}[讲解：先看这题到底在做什么]
先写出全部整数解，再把 $x>0,y>0$ 各自产生的条件变成同一个参数 $t$ 的上下界。这样“数多少组正整数解”就完全变成了整数区间计数。
\end{pitfall}

\sol 化为 $2x+3y=21$。取 $x=3t,y=7-2t$，约束给 $1\le t\le3$，所以解为 $(3,5),(6,3),(9,1)$。解数 $3$，$x$ 最小 $3$、最大 $9$，$y$ 最小 $1$、最大 $5$。强调两个最小值不来自同一组解。


\begin{problem}{A3：约分会改变模数}
解 $18x\equiv30\pmod{42}$，并在区间 $[0,42)$ 中列出全部解。
\end{problem}

\begin{pitfall}[讲解：先看这题到底在做什么]
同余式两边约掉公共因子以后，模数也必须同步除掉这个因子。先得到约分后的最小模数，再列出全部同余解；只写一个代表元会漏掉区间里的其他解。
\end{pitfall}

\sol 除 gcd $6$ 得 $3x\equiv5\pmod7$，故 $x\equiv4\pmod7$。区间内解为 $4,11,18,25,32,39$，共有六个。只给 $4$ 是漏解；把约分后模数仍写成 $42$ 是错误。


\begin{problem}{A4：CRT 与下界}
求满足 $x\equiv1\pmod4$、$x\equiv3\pmod6$ 且 $x\ge40$ 的最小整数。
\end{problem}

\begin{pitfall}[讲解：先看这题到底在做什么]
CRT 负责找到所有满足同余条件的数，即一个同余类；题目还有“至少多大”的额外限制，所以最后还要从这个同余类中找第一个不小于下界的代表元。
\end{pitfall}

\sol 令 $x=1+4t$，$4t\equiv2\pmod6$，故 $t\equiv2\pmod3$，合并为 $x\equiv9\pmod{12}$。不小于 $40$ 的首项为 $45$。直接输出最小非负代表元 $9$ 忽略了大小约束。


\begin{problem}{A5：一个质因数分解对应多种函数}
计算 $\varphi(72),\mu(72),\tau(72)$ 与 $\sum_{d\mid72}\varphi(d)$。
\end{problem}

\begin{pitfall}[讲解：先看这题到底在做什么]
把同一个质因数分解同时喂给 $\varphi,\mu,	au$，可以很好地训练“各函数到底看分解的什么信息”：$\varphi$ 看质因子及其幂次，$\mu$ 只关心是否出现平方因子，$	au$ 直接把各指数加一后相乘。
\end{pitfall}

\sol $72=2^3\cdot3^2$，所以 $\varphi(72)=72(1-1/2)(1-1/3)=24$，$\mu(72)=0$，$\tau(72)=(3+1)(2+1)=12$。最后一项直接由欧拉恒等式等于 $72$，无需枚举十二个约数。评价重点是区分“有平方因子”与“不同素因子个数”。


\begin{problem}{A6：单次逆元与批量逆元}
求 $37^{-1}\pmod{101}$；再说明模 $17$ 同时求 $2,3,5$ 的逆元时怎样只做一次一般求逆。
\end{problem}

\begin{pitfall}[讲解：先看这题到底在做什么]
先用一次 exgcd 得到基准逆元，再观察前缀积之间的乘积逆元关系。批量法的节省之处不在于“exgcd 更快”，而在于原本很多次求逆被一次逆元和若干乘法替代。
\end{pitfall}

\sol $1=11\cdot101-30\cdot37$，故逆元为 $71$。批量问题先算前缀积 $1,2,6,13$，求 $13^{-1}=4$，再逆推得到 $5^{-1}=7,3^{-1}=6,2^{-1}=9$。每一步应用的是乘积逆元关系，而不是将三个独立 exgcd 称为“批量”。


\begin{problem}{A7：环的连通分量}
在 $18$ 个点的环上每次加 $12$，有几个循环？从点 $3$ 出发能否到点 $4$？列出点 $3$ 所在循环。
\end{problem}

\begin{pitfall}[讲解：先看这题到底在做什么]
每走一步就是给编号加 $k$ 并模 $n$。能走到哪些点只由 $\gcd(n,k)$ 决定；因此先算 gcd 就能直接知道环被分成多少个分量，而不是逐点模拟。
\end{pitfall}

\sol gcd 为 $6$，所以有六个循环，每个长三。每次保持模 $6$ 余数，$3$ 与 $4$ 不在同一循环。对应循环为 $3\to15\to9\to3$。可以进一步追问：到某个点的最少步数与“是否同环”不是同一个问题。


\begin{problem}{A8：无整数解与无正整数解}
比较 $6x+10y=7$ 与 $6x+10y=2$。第二个方程中，单独可取得的最小正 $x$、最小正 $y$ 分别是多少？它们能否同时取得？
\end{problem}

\begin{pitfall}[讲解：先看这题到底在做什么]
先区分两个问题：$d
mid c$ 时连整数解都没有；$d\mid c$ 时整数解存在，但正整数限制可能让参数区间为空。即使没有同时为正的解，最小正 $x$ 和最小正 $y$ 仍可能分别存在。
\end{pitfall}

\sol 前者因为 $2\nmid7$ 而无整数解。后者有解 $(-3,2)$，通解 $x=-3+5t,y=2-3t$，没有同时为正的参数。最小正 $x=2$（对应 $y=-1$），最小正 $y=2$（对应 $x=-3$），不能同时取得。这正是 P5656 特殊输出分支的含义。


\section{B 组：指数、赋值与组合数}
\begin{problem}{B1：欧拉指数与真实周期}
求 $7^{123456789}\bmod40$，并分别使用欧拉函数与阶解释降幂依据。
\end{problem}

\begin{pitfall}[讲解：先看这题到底在做什么]
欧拉定理给出的 $\varphi(m)$ 只是一个一定有效的指数，不一定是最小周期。先通过小次幂找到真实的阶，再决定指数应该对什么数取模，能避免把“周期”和“可用上界”混为一谈。
\end{pitfall}

\sol $\gcd(7,40)=1$，$\varphi(40)=16$，指数模 $16$ 为 $5$；而 $7^2\equiv9,7^4\equiv1$，真实阶为 $4$。指数模 $4$ 为一，答案 $7$。欧拉定理给出一个周期，未必给最小周期。


\begin{problem}{B2：非互质大指数}
求 $6^{10^{100}}\bmod72$，并说明不计算任何大指数即可得到答案的原因。
\end{problem}

\begin{pitfall}[讲解：先看这题到底在做什么]
这里不要机械套欧拉降幂。直接看各素因子的赋值什么时候达到模数所需的指数，一旦整个数已经被模数整除，答案就立即是 $0$。这是非互质底数时最直接的判定方式。
\end{pitfall}

\sol $72=2^3\cdot3^2$，$6^e$ 中两个素数的赋值均为 $e$。当 $e\ge3$ 即同时被两个素数幂整除，故答案为零。此处直接按赋值判零，比构造欧拉链更简单。不要把所有巨大指数问题都机械套用递归降幂。


\begin{problem}{B3：先列周期，再写离散对数通解}
求 $\ord_{17}(4)$，并解 $4^x\equiv13\pmod{17}$。
\end{problem}

\begin{pitfall}[讲解：先看这题到底在做什么]
先把 $4$ 的幂循环写出来，阶自然就出现了。找到最小解以后，所有解都只是在这个周期上平移，因此通解就是“最小解 + 阶的倍数”。
\end{pitfall}

\sol $4^2\equiv16=-1,4^4\equiv1$，所以阶为四。周期为 $1,4,16,13$，最小非负解为三，全部非负解为 $x=3+4t$（$t\ge0$）。这里模 $17$ 的阶不是 $\varphi(17)=16$。


\begin{problem}{B4：离散对数的有限前缀}
分别求 $6^x\equiv12\pmod{18}$ 与 $6^x\equiv0\pmod{18}$ 的最小非负解。
\end{problem}

\begin{pitfall}[讲解：先看这题到底在做什么]
底数与模数不互质时，序列可能先快速坍缩，再进入一个短周期。此时不能定义一个正常的乘法阶，因此要先观察前缀行为，再判断是否存在目标值。
\end{pitfall}

\sol 序列为 $1,6,0,0,\ldots$，第一式无解，第二式最小解为二。这一序列不是从指数零开始的纯周期；也不存在 $\ord_{18}(6)$。学生若写“指数对阶取模”，应先让其指出阶为何不存在。


\begin{problem}{B5：奇素数 LTE}
求 $\nu_3(10^{81}-1)$，并列出所使用的全部前提。
\end{problem}

\begin{pitfall}[讲解：先看这题到底在做什么]
LTE 不是只背最后一个公式。讲解时先逐条检查素数、整除关系和指数条件，再把赋值直接代入。学生能明确说出“为什么这条 LTE 能用”，比记住结果更重要。
\end{pitfall}

\sol 素数三为奇数，$3\mid10-1$，$3\nmid10$、$3\nmid1$，指数为正。故结果为 $\nu_3(9)+\nu_3(81)=2+4=6$。满分答案必须写出这些条件，而不是只写六。


\begin{problem}{B6：二的赋值}
求 $\nu_2(5^{12}-1)$，用一般偶指数形式与 $4\mid x-y$ 的简化形式各算一次。
\end{problem}

\begin{pitfall}[讲解：先看这题到底在做什么]
模 $2$ 的 LTE 与奇素数版本不一样，最稳妥的做法是先写通式，再看额外条件能否简化。最后两种形式得到相同结果时，可以把它当作一次条件自检。
\end{pitfall}

\sol 一般式为 $\nu_2(4)+\nu_2(6)+\nu_2(12)-1=2+1+2-1=4$。由于 $4\mid5-1$，简化式为 $\nu_2(4)+\nu_2(12)=4$。两种形式一致是一次很好的条件检查。


\begin{problem}{B7：阶乘赋值与组合数赋值}
求 $\nu_3(100!)$ 与 $\nu_2\binom{100}{50}$。
\end{problem}

\begin{pitfall}[讲解：先看这题到底在做什么]
先把阶乘中的 $p$ 因子数量分别算出来，再做减法得到组合数的赋值。Legendre 公式和二进制数位和公式是两种不同视角，适合让学生验证同一个结果。
\end{pitfall}

\sol 第一项为 $33+11+3+1=48$。第二项为 $\nu_2(100!)-2\nu_2(50!)=97-2\cdot47=3$。也可用二进制数位和：$100$、$50$ 都有三个一，所以 $3+3-3=3$。结果说明该组合数能被八整除，不能被十六整除。


\begin{problem}{B8：一次 Lucas 计算}
求 $\binom{20}{6}\bmod7$。若把下标改成八，结果是否仍非零？
\end{problem}

\begin{pitfall}[讲解：先看这题到底在做什么]
Lucas 定理的关键是把 $n,k$ 写成 $p$ 进制，然后逐位取组合数。只要出现某一位的下标超过上标，对应因子就是 $0$；没有越界时才继续相乘。
\end{pitfall}

\sol $20=(26)_7,6=(06)_7$，余数为 $\binom20\binom66=1$。$8=(11)_7$，余数为 $\binom21\binom61=12\equiv5$，仍非零。这两题都没有位上越界；改成 $7$ 进制某位超过上标的数才会自动为零，不能仅凭下标“变大”判断。


\begin{problem}{B9：合数模数不能直接除阶乘}
求 $\binom{12}{4}\bmod18$，要求按模二和模九分别解释。
\end{problem}

\begin{pitfall}[讲解：先看这题到底在做什么]
组合数取模时最危险的地方就是把普通除法机械搬进模环。先分别处理素数幂中的 $p$ 因子和单位部分，再用 CRT 合并；尤其要明确“分母没有逆元”时不能直接除。
\end{pitfall}

\sol 模二：$\nu_2=10-3-7=0$，组合数为奇数，余数一。模九：$\nu_3=5-1-2=2$，至少含 $3^2$，余数零。CRT 合并“奇数且为九的倍数”，模十八余九。实际 $\binom{12}{4}=495$。不允许把 $4!\bmod18=6$ 的“逆元”代入，因为它根本不存在。


\begin{problem}{B10：辨认计数对象}
(1) P7488 的三角形模型中 $n=4,m=2$ 的组数是多少？(2) 五件不同礼物分别给两位有编号的人两件、一件，有多少方案？
\end{problem}

\begin{pitfall}[讲解：先看这题到底在做什么]
两道题表面都是“选一些对象”，但排列是否算不同、接收者是否有身份，决定了计数公式。先明确对象是什么、哪些选择不能重复，再决定是否需要乘或除阶乘。
\end{pitfall}

\sol (1) $\binom42^2\cdot2!=72$，也可由 $R(3,2)+(8-2)R(3,1)=18+6\cdot9=72$。 (2) $\binom52\binom31=30$。第一题三角形组无序，第二题接收者有身份，但两题最终都通过恰好计数避免重复；不能见到“组”就随意除阶乘。


\section{C 组：反演与亚线性求和}
\begin{problem}{C1：手算一次杜教筛}
计算 $S_\varphi(12)$ 和 $S_\mu(12)$，并写出前者按商分块后的递归式。
\end{problem}

\begin{pitfall}[讲解：先看这题到底在做什么]
先把递推式中的每一个块看成一个整体，而不是逐项展开。$\floor{n/j}$ 相同的连续 $j$ 会共享同一个状态，所以真正要算的是块长和对应的递归状态；手算一次后再看代码会清楚很多。
\end{pitfall}

\sol
\[
S_\varphi(12)=78-S_\varphi(6)-S_\varphi(4)-S_\varphi(3)-2S_\varphi(2)-6S_\varphi(1)=46.
\]
$S_\mu(10)=-1,\mu(11)=-1,\mu(12)=0$，所以 $S_\mu(12)=-2$。本题检查的是块长与从 $j=2$ 开始，漏掉这两处任意一处都会出错。


\begin{problem}{C2：矩形上的 gcd 与互质计数}
求 $\sum_{i\le4,j\le6}\gcd(i,j)$，以及其中互质有序数对的个数。
\end{problem}

\begin{pitfall}[讲解：先看这题到底在做什么]
这两种计数都能拆成“公共因子 $d$”，但权重分别来自 $\varphi$ 和 $\mu$。先把计数对象说清楚，再写权重，能避免把两个看起来相似的公式混在一起。
\end{pitfall}

\sol gcd 和为 $24+6+4+2=36$；互质数对为 $24-6-2=16$。前者权重是 $\varphi(d)$，后者是 $\mu(d)$，两式商因子相同，不能把权重混用。也不能在矩形问题中将两个商都写成 $\floor{4/d}$。


\begin{problem}{C3：gcd 的平方应该使用哪个系数？}
推导 $\sum_{i,j\le n}\gcd(i,j)^2$ 的一维求和式，给出系数在素数幂处的值，并计算 $n=3$ 的答案。
\end{problem}

\begin{pitfall}[讲解：先看这题到底在做什么]
平方 gcd 对应的是 Jordan 型卷积系数，而不是把 $\varphi$ 直接平方。最稳妥的解释是从 $\id_2=J_2*\one$ 出发，再用同样的公共因子计数套路。
\end{pitfall}

\sol 令 $J_2=\id_2*\mu$，则 $\id_2=J_2*\one$，所以答案是
\[
 \sum_{d\le n}J_2(d)\floor{n/d}^2,\qquad
 J_2(p^e)=p^{2e}-p^{2(e-1)}.
\]
$n=3$ 时为 $9\cdot1+1\cdot3+1\cdot8=20$。系数不是 $\varphi(d)^2$，因为卷积不保持这种逐点平方关系。


\begin{problem}{C4：PN 思想的最简单计数应用}
求不超过三十的平方自由数个数，要求不逐个判定。
\end{problem}

\begin{pitfall}[讲解：先看这题到底在做什么]
这个例子只是在数没有平方因子的整数，但它完整体现了 PN 的套路：先做局部卷积消去主要贡献，再留下平方数支撑，最后只枚举 $d\le\sqrt n$。
\end{pitfall}

\sol 使用 $\sum_{d\le\sqrt{30}}\mu(d)\floor{30/d^2}$，得到 $30-7-3+0-1=19$。讲解时可让学生说明为什么含多个平方因子的数会通过莫比乌斯系数正确容斥。


\begin{problem}{C5：PN 卷积商的交叉项}
设积性函数满足 $f(p^k)=p^{2k}+[k\ge2]$。取 $g(n)=n^2$ 并令 $f=g*h$，求 $h(p),h(p^2),h(p^3)$，设计 $S_f(n)$ 的稀疏求和式。
\end{problem}

\begin{pitfall}[讲解：先看这题到底在做什么]
卷积商在 $p^3$ 处会混入低次项的贡献，所以不能直接写成“原函数减基准函数”。先按素数幂递推把前面的项扣干净，再看剩余支撑为什么仍然稀疏。
\end{pitfall}

\sol $h(p)=0$，$h(p^2)=1$，而
\[
 h(p^3)=f(p^3)-g(p^3)-g(p)h(p^2)=1-p^2.
\]
故不能对所有幂次都写成 $f(p^k)-g(p^k)$。有
\[
 S_f(n)=\sum_{d\in\mathrm{PN},d\le n}h(d)\,
 \frac{q(q+1)(2q+1)}6,\quad q=\floor{n/d}.
\]
在局部系数及枚举准备完成后，只需 $O(\sqrt n)$ 个支撑项；若取模且分母不可逆，要整数约分后再取模。


\begin{problem}{C6：验证 P4240 的简化系数}
计算 $f(6),f(12)$，其中 $f(k)=\sum_{d\mid k}d\mu(k/d)/\varphi(d)$；再用 $G$ 形式算 $\sum_{i,j\le3}\varphi(ij)$。
\end{problem}

\begin{pitfall}[讲解：先看这题到底在做什么]
这一题适合专门检查“局部系数到底有没有算对”。先在整数或有理数层面验证几个小 $k$，确认系数结构正确，再把它们放回模意义下的实现。
\end{pitfall}

\sol $6$ 平方自由，$f(6)=1/\varphi(6)=1/2$；$12$ 含平方因子，$f(12)=0$。后一个求和中，$k=1,2,3$ 的贡献为 $16,1,2$，合计 $19$。分数可先在有理数中理解，再在题目模数中用逆元表示。


\section{D 组：证明辨析与综合拓展}
\begin{problem}{D1：非平凡平方根为何能分解？}
列出 $x^2\equiv1\pmod{15}$ 的所有解。用其中一个非平凡解找到十五的两个素因子。
\end{problem}

\begin{pitfall}[讲解：先看这题到底在做什么]
CRT 告诉我们模素因子分别取 $1$ 或 $-1$，组合起来可以得到非平凡平方根。只要它既不是 $1$ 也不是 $-1$，那么与 $V-1,V+1$ 的 gcd 就会分别捞出不同因子。
\end{pitfall}

\sol CRT 分别在模三、模五选择 $\pm1$，得 $1,4,11,14$。选 $V=4$，则 $\gcd(V-1,15)=3$，$\gcd(V+1,15)=5$。若选 $V=14=-1$，后一 gcd 等于十五，不会分裂；这正是随机分解步骤为何排除 $-1$。


\begin{problem}{D2：相同阶不代表同一个子群}
比较模十五的 $2,7$ 的阶。是否存在 $a\ge0$ 使 $2^a\equiv7\pmod{15}$？说明 P5605 中原根条件不可删除。
\end{problem}

\begin{pitfall}[讲解：先看这题到底在做什么]
“阶相同”和“生成同一个子群”是两件事。先列出两个元素的幂集合，再直接比较集合是否相同，就能看到为什么在非循环单位群里不能只靠阶判断子群包含。
\end{pitfall}

\sol 两者阶均为四，但 $2$ 的幂集合为 $\{1,2,4,8\}$，$7$ 的幂集合为 $\{1,7,4,13\}$，后者并不属于前者生成的子群。模十五的单位群不循环；阶的整除关系不能替代子群包含关系。


\begin{problem}{D3：一步边也可能产生多个路径余数}
两个点仅通过一条边权二的无向边相连，模数为七。允许反复走边，从一端到另一端能得到哪些权值余数？
\end{problem}

\begin{pitfall}[讲解：先看这题到底在做什么]
先把路径权值写成一个几何级数，再研究 $2^h$ 在允许的奇数 $h$ 上能出现哪些余数。这里必须同时关注图上的路径结构和数值状态，不能只看其中一个。
\end{pitfall}

\sol 路径长度 $h$ 必须为奇数，原权值为 $2(1+2+\cdots+2^{h-1})=2(2^h-1)$。$2^h$ 在奇数 $h$ 上可取 $2,1,4$，所以余数为 $2,0,6$。图的拓扑连通性与数值状态连通性必须同时考虑。


\begin{problem}{D4：二维卷积的一次完全消去}
取 $S(n)=\tau(n)$。证明
\[
 \tau(xy)=\sum_{d\mid\gcd(x,y)}\mu(d)\tau(x/d)\tau(y/d),
\]
并写出 $\sum_{i,j\le n}\tau(ij)$ 的一维表达式。
\end{problem}

\begin{pitfall}[讲解：先看这题到底在做什么]
这题真正要看的不是最后那一行公式，而是局部生成函数里的 $1-XY$。它意味着所有坐标轴上的主要贡献都被消掉，只剩对角线上的稀疏支撑，于是二维求和又降成一维。
\end{pitfall}

\sol 固定素数，$S(p^k)=k+1$。二元局部级数为
\[
 \sum_{a,b\ge0}(a+b+1)X^aY^b
 =\frac{1-XY}{(1-X)^2(1-Y)^2}.
\]
分母对应 $\tau(x)\tau(y)$，剩余 $1-XY$ 表示卷积商只在局部 $(p,p)$ 为 $-1$，因此全局只在 $x=y=d$ 为 $\mu(d)$。答案为
\[
 \sum_{d\le n}\mu(d)S_\tau(\floor{n/d})^2.
\]
$n=3$ 时为 $25-1-1=23$。这展示了“算法二”在某些函数上把所有非对角残差完全消掉，而不是只得到一个较稀疏的支撑。


\begin{problem}{D5：从最大素因子开始还原}
已知 $M=1080=N\varphi(N)$ 且解存在，求 $N$，并解释为什么不是直接取 $\sqrt M$。
\end{problem}

\begin{pitfall}[讲解：先看这题到底在做什么]
已知 $N\varphi(N)$ 后，可以从最大的素因子开始拆，因为 $N$ 和 $\varphi(N)$ 中各素因子的指数关系很受限制。每还原出一个素因子，就更新剩余量，再继续处理下一层。
\end{pitfall}

\sol $1080=2^3\cdot3^3\cdot5$，最大素因子五的指数为一，对应 $N$ 中一个五；除去 $5\cdot4=20$ 后剩五十四。三的指数为三，对应 $N$ 中两个三；除去 $3^3\cdot2=54$ 后结束。所以 $N=45$，$\varphi(45)=24$，乘积为 $1080$。$\varphi(N)$ 一般不等于 $N$，平方根只能提供量级，不能给答案。


\begin{problem}{D6：构造题还需要全局去重}
当 $k=2$，从种子 $(1,2,6)$ 写出两个孩子。为什么它们都不能作为该种子之后的第二组输出？沿 BFS 继续找一组可用输出。
\end{problem}

\begin{pitfall}[讲解：先看这题到底在做什么]
构造出一个满足方程的三元组只是局部正确；题目通常还要求不同组之间不能重复使用整数。因此搜索状态必须同时记录“这个组是否合法”和“这些数是否已经在全局出现过”。
\end{pitfall}

\sol 两个孩子是 $(2,6,15)$、$(1,6,12)$，都与种子共享已使用整数。不能因为“三元组不相同”就输出。按正文孩子顺序 BFS，继续可找到 $(15,40,104)$；它与 $1,2,6$ 无交，且
\[
15^2+40^2+104^2=12641=2(15\cdot40+40\cdot104+15\cdot104)+1.
\]
把“单组合法”和“所有组之间合法”分别验证，是构造题的通用习惯。





\chapter{核心模板与接口}
\begin{chapterintro}[章节导读]
附录 A 训练的是分析与推导，接下来需要把这些结论真正落到程序接口上。本附录集中整理高频数论模板，并把输入前提、溢出边界、返回值以及空间成本写在代码旁边，方便把数学证明转化为可靠实现。
\end{chapterintro}

\section{模板使用说明}

\begin{tip}[附录代码怎么读]
先读注释里的 Contract（前提条件：模数范围、是否互质、输入上限），再读函数体。凡是本书在正文里写\textbf{必须先检查}的地方，模板里都有对应的 \texttt{assert} 或返回值；提交时要按题面删掉 \texttt{assert} 以免被卡常数。
\end{tip}

完整 C++17 实现在资料包 \texttt{code/number\_theory.hpp}。本附录仅摘录关键部分，依赖同文件中的类型别名、\texttt{exgcd} 与相关函数；摘录不是可单独编译的提交文件。测试程序为 \texttt{code/selftest.cpp}。

\begin{longtable}{p{0.28\textwidth}p{0.61\textwidth}}
\toprule
模板 & 限制\\\midrule
\texttt{mul\_mod / pow\_mod} & 无符号 $64$ 位整数，模数必须大于零，乘积使用无符号 $128$ 位。\\
\texttt{inverse\_mod} & 不可逆时返回空 optional；模数一返回零是程序约定，不用于通常单位群讨论。\\
\texttt{exgcd} & 参数非负且不全为零；返回 gcd 及其一组整数线性组合系数。\\
\texttt{merge} & 输入及最终模数必须放入正的有符号 $64$ 位。合并最小公倍数超界时抛异常，不静默溢出。\\
\texttt{bsgs / exbsgs} & 返回最小非负解或空值；哈希表空间为平方根级，不能因为数值能放入 $64$ 位就对 $10^{18}$ 直接开表。\\
\texttt{PrimePowerBinom} & 素数参数需预先保证；默认最多建 $2\times10^6$ 个量级的表项。支持非负 $n,k$，$k>n$ 返回零。\\
\texttt{is\_prime / factor} & 判素适用于完整无符号 $64$ 位；分解使用可复现的默认随机种子，允许调用方修改，不是密码学实现。\\
\texttt{order} & 需要 $a$ 与 $m$ 互质，并提供正确的 $\varphi(m)$ 及其有序素因子列表。\\
\bottomrule
\end{longtable}

\section{安全模乘、快速幂与逆元}
\begin{lstlisting}
inline u64 mul_mod(u64 a, u64 b, u64 m) {
    assert(m > 0);
    return (u128)a * b % m;
}
inline u64 pow_mod(u64 a, u64 e, u64 m) {
    assert(m > 0);
    u64 r = 1 % m; a %= m;
    for (; e; e >>= 1, a = mul_mod(a, a, m))
        if (e & 1) r = mul_mod(r, a, m);
    return r;
}
inline std::optional<u64> inverse_mod(u64 a, u64 m) {
    assert(m > 0);
    // The residue for modulus 1 is returned as a programming convention.
    if (m == 1) return 0;
    i128 x, y;
    if (exgcd(a % m, m, x, y) != 1) return std::nullopt;
    return (u64)norm(x, m);
}
\end{lstlisting}

\section{最小非负离散对数}
\begin{lstlisting}
inline std::optional<u64> bsgs(u64 a, u64 b, u64 m) {
    assert(m > 0);
    if (m == 1) return 0;
    a %= m; b %= m;
    if (std::gcd(a,m) != 1) return std::nullopt;
    if (b == 1) return 0;
    if (std::gcd(b,m) != 1) return std::nullopt;
    u64 B = (u64)std::sqrt((long double)m);
    while ((u128)B*B < m) ++B;
    while (B && (u128)(B-1)*(B-1) >= m) --B;
    // Caller is responsible for the O(sqrt(m)) memory budget.
    std::unordered_map<u64,u64> baby;
    baby.reserve((std::size_t)(B*2+1));
    u64 v = 1;
    for (u64 j=0;j<B;++j) {
        baby.emplace(v,j); // Keep the earliest/smallest j.
        v = mul_mod(v,a,m);
    }
    u64 step = *inverse_mod(v,m), giant = b;
    for (u64 i=0;i<=B;++i) {
        auto it = baby.find(giant);
        if (it != baby.end()) {
            u128 x = (u128)i*B+it->second;
            if (x < m) return (u64)x;
        }
        giant = mul_mod(giant,step,m);
    }
    return std::nullopt;
}
inline std::optional<u64> exbsgs(u64 a, u64 b, u64 m) {
    assert(m > 0);
    if (m == 1) return 0;
    a %= m; b %= m;
    if (b == 1) return 0;
    if (a == 0) return b == 0 ? std::optional<u64>(1) : std::nullopt;
    u64 count = 0, k = 1;
    while (true) {
        u64 d = std::gcd(a,m);
        if (d == 1) break;
        if (b == k) return count;
        if (b % d) return std::nullopt;
        b /= d; m /= d; ++count;
        if (m == 1) return count;
        k = mul_mod(k,a/d,m);
    }
    auto inv = inverse_mod(k,m);
    assert(inv.has_value());
    auto tail = bsgs(a,mul_mod(b,*inv,m),m);
    if (!tail) return std::nullopt;
    return *tail + count;
}
\end{lstlisting}

\section{素数幂模数下的组合数}
\begin{lstlisting}
struct PrimePowerBinom {
    u64 p, q = 1;
    unsigned alpha;
    std::vector<u64> pref;
    // Contract: p is prime. Table size must be affordable.
    PrimePowerBinom(u64 prime, unsigned exponent,
                    u64 table_limit=2000000):p(prime),alpha(exponent) {
        assert(p>=2 && alpha>=1);
        for (unsigned i=0;i<alpha;++i) {
            if (q>table_limit/p) throw std::length_error("prime-power table too large");
            q*=p;
        }
        pref.assign(q+1,1);
        for (u64 i=1;i<=q;++i)
            pref[i]=(i%p) ? mul_mod(pref[i-1],i,q) : pref[i-1];
        assert(pref[q]==1 || pref[q]==q-1);
    }
    u64 valuation(u64 n) const {
        u64 e=0;
        while (n) { n/=p; e+=n; }
        return e;
    }
    u64 unit(u64 n) const {
        u64 r=1;
        while (n) {
            r=mul_mod(r,pref[n%q],q);
            if ((n/q)&1) r=mul_mod(r,pref[q],q);
            n/=p;
        }
        return r;
    }
    u64 choose(u64 n,u64 k) const {
        if (k>n) return 0;
        u64 e=valuation(n)-valuation(k)-valuation(n-k);
        if (e>=alpha) return 0;
        u64 r=unit(n);
        r=mul_mod(r,*inverse_mod(unit(k),q),q);
        r=mul_mod(r,*inverse_mod(unit(n-k),q),q);
        return mul_mod(r,pow_mod(p,e,q),q);
    }
};
\end{lstlisting}

\section{完整无符号 64 位的 Miller--Rabin}
\begin{lstlisting}
inline bool is_prime(u64 n) {
    if (n<2) return false;
    for (u64 p:{2ULL,3ULL,5ULL,7ULL,11ULL,13ULL,
                17ULL,19ULL,23ULL,29ULL,31ULL,37ULL})
        if (n%p==0) return n==p;
    u64 d=n-1; unsigned s=0;
    while (!(d&1)) { d>>=1; ++s; }
    for (u64 base:{2ULL,325ULL,9375ULL,28178ULL,
                   450775ULL,9780504ULL,1795265022ULL}) {
        u64 a=base%n;
        if (!a) continue;
        u64 x=pow_mod(a,d,n);
        if (x==1 || x==n-1) continue;
        bool pass=false;
        for (unsigned j=1;j<s;++j) {
            x=mul_mod(x,x,n);
            if (x==n-1) { pass=true; break; }
            if (x==1) break;
        }
        if (!pass) return false;
    }
    return true;
}
\end{lstlisting}

\section{怎样复现验证}

\begin{tip}[怎么做：对拍的最小可行版本]
这四条命令是\textbf{对拍}的最小可行版本：一份正确但慢的实现 + 一份你写的实现 + 随机数据。竞赛里 90\% 的\textbf{WA}能被它抓出来，包括取模负号、最小解、边界 $n=1$ 这类。
\end{tip}

在资料目录中运行：
\begin{lstlisting}[language=bash]
g++ -std=c++17 -O2 -Wall -Wextra -fsanitize=undefined \
    code/selftest.cpp -o nt_selftest
./nt_selftest
python validation/verify_math.py
python validation/check_happy.py
python code/vieta_construct.py --verify
\end{lstlisting}
C++ 测试程序以 Boost 的大整数作为组合数的独立对照；模板本身无需 Boost。Python 的数学验证仅依赖标准库。

\begin{tip}[提醒：Boost 只出现在自测里]
Boost 只在自测里当\textbf{独立裁判}（用 \texttt{cpp\_int} 精确算组合数再取模）。提交的代码只依赖 \texttt{number\_theory.hpp}。这个区分很重要：对拍用的实现要\textbf{笨且显然正确}，被测的实现要\textbf{快且接口干净}。
\end{tip}



\begin{tip}[代码与验证脚本的位置]
本附录摘录的模板，完整 C++17 实现在资料包的 \texttt{Number Theory/code/number\_theory.hpp}；
测试程序为 \texttt{Number Theory/code/selftest.cpp}；数学验证脚本为
\texttt{Number Theory/validation/verify\_math.py}、\texttt{check\_happy.py}
与 \texttt{Number Theory/code/vieta\_construct.py}。
这些文件相对本书源码目录的上一级，编译本书本身不需要它们。
\end{tip}

\chapter{教师使用：课堂设计、追问与评量}
\begin{chapterintro}[章节导读]
附录 A 解决“学生怎样练”，附录 B 解决“程序怎样写”；最后的附录 C 则把这两部分重新放回课堂。这里不新增一套数学主线，而是给出课程安排、追问、评分与验证框架，用来说明同一份数论材料怎样分别服务于讲授、训练和复习。
\end{chapterintro}

\section{如何安排一轮而不是一节“数论课”}
建议把本资料看作课程资源库，而不是一次讲完的幻灯片。每讲至少留下一个可验证的学习产出：一段具有输入输出契约的代码、一份没有跳步的推导、一组能击穿错误解法的反例。下面给出六个 $90$ 分钟课堂包；高级章节可另设专题阅读。

\begin{longtable}{p{0.15\textwidth}p{0.38\textwidth}p{0.37\textwidth}}
\toprule
阶段 & 通常的时间分配 & 教师的观察对象\\\midrule
诊断与情境 & $0$--$10$ 分钟：前置小测、具体难题 & 学生卡在数学概念，还是只忘了实现？\\
提出结构 & $10$--$30$ 分钟：换元、等价条件、不变量 & 能否说出“为什么这样设”？\\
证明与手算 & $30$--$50$ 分钟：核心证明、小数值验算 & 每一步是否依赖互质、素数、正性？\\
引导练习 & $50$--$70$ 分钟：同伴推导与纠错 & 先让学生指出错误，再补正确算法。\\
实现与迁移 & $70$--$85$ 分钟：代码接口、变式 & 能否从结论转换为边界安全的实现？\\
离堂检查 & $85$--$90$ 分钟：一道短题、一句解释 & 不只收答案，应收关键条件与理由。\\
\bottomrule
\end{longtable}

\section{课堂包一：exgcd 从“能解”到“求最小解”}
\paragraph{学习目标.} 学生能区分整数解与正整数解；能从一组特解生成全部解；能把两个正性约束转成同一个整数参数区间。
\paragraph{课前诊断.} 让学生回答 $6x+10y=7$ 为什么无解、$6x+10y=2$ 为什么可能有整数解却没有正整数解。不必先出现代码。
\paragraph{板书主线.}
\[
\text{公约数集合不变}\ \longrightarrow\ au+bv=d
\ \longrightarrow\ (x_0,y_0)=\frac cd(u,v)
\ \longrightarrow\ t\text{ 的区间}.
\]
\paragraph{引导顺序.}
先用 $252,198$ 手算 gcd；把余数倒代为整数线性组合。让学生自己发现乘 $c/d$，再用两组解相减证明 $\Delta x$ 必须是 $b/d$ 的倍数。最后讲 P5656，不要从通解公式直接开始背诵。
\paragraph{预设错误与回应.}
学生说“有解当且仅当 $c\mid d$”时，要求把 $a=6,b=10,c=32$ 代入判断，再找出 $(2,2)$；学生说“求到一个负数解说明无解”时，要求展示参数平移；学生把 ceil 写成 C++ 除法时，使用 $-5/3$ 进行对照。
\paragraph{同伴活动.} 一人负责给出一个特解，另一人只利用通解在区间内数解，再交换检查。两人使用不同特解时，最终正整数解集合应相同；这可以检验“特解选择不影响答案”的理解。
\paragraph{离堂题与作业.} 离堂做 A4；作业做 A1、A2、A8，再实现 P5656 的数学核心。评阅时要求附至少一个“有整数解但无正整数解”的测试。

\section{课堂包二：指数取模的三个不同问题}
\paragraph{学习目标.} 区分“算 $a^E$ 的余数”“求指数 $x$”“求最小周期”三种问题；知道欧拉函数是可用周期但不一定最小；知道非互质指数不能只保留余数。
\paragraph{课前诊断.} 写出模八的 $2^x$ 前五项，模七的 $2^x$ 前七项。请学生观察一个有尾巴、一个从开头就循环。
\paragraph{板书分栏.}
\begin{center}
\begin{tabular}{lll}
\toprule
问题 & 核心工具 & 不能丢的信息\\\midrule
计算幂余数 & 快速幂、扩展欧拉 & 指数是否达到阈值\\
求最小指数 & BSGS/exBSGS & 零指数、最小命中方式\\
求周期 & 阶与素因子试除 & 单位条件、正整数最小值\\\bottomrule
\end{tabular}
\end{center}
\paragraph{证明重点.} 欧拉定理只讲清一次“乘法置换剩余系”，不反复背费马与欧拉两套相同证明。扩展欧拉在素数幂上分类，把“不互质”变成“足够多素因子时归零”。
\paragraph{引导练习.} 先计算 $2^4\bmod8$，再对照指数只取模四的错误结果；接着让学生设计二元组接口 $(E\bmod q,E\ge q)$。随后手算 $2^x\equiv5\pmod{13}$，画出 baby/giant 的匹配表。
\paragraph{预设错误与回应.} “两个连续塔高余数相同就稳定”需要完整状态论证；“固定模数就能复用 BSGS 表”还要检查底数；“求阶时 BSGS 返回零”说明忘记阶要求最小\emph{正}指数。
\paragraph{离堂题与作业.} 做 B3、B4，要求先写能否使用阶的判断。提高作业为 P6736：只实现截断塔和模塔两个函数，再解释它们为什么返回不同信息，暂不要求最优时限提交。

\section{课堂包三：组合数中不能取逆的部分去哪了？}
\paragraph{学习目标.} 将组合数分为素数指数与单位部分；能解释 Lucas 只适合素数模数；能独立完成一个小合数模数的组合数计算。
\paragraph{情境.} 给出 $\binom{10}{4}\bmod12$，先允许学生尝试阶乘逆元法。等有人发现分母不能求逆时，再提出“并非不能除，而是不能在取模后直接除”。
\paragraph{板书主线.}
\[
 n!\ \longrightarrow\ p^{\nu_p(n!)}U_p(n)
 \quad\longrightarrow\ p^e\cdot\text{单位分式}
 \quad\longrightarrow\text{各素数幂余数}\quad\longrightarrow\mathrm{CRT}.
\]
\paragraph{证明顺序.} 先用倍数计数讲 Legendre；再用数位和讲 Kummer；最后讲 Lucas 的生成函数证明。扩展 Lucas 不是在原 Lucas 公式中把素数改成合数，而是更换为素数幂单位阶乘路线。
\paragraph{操作活动.} 将 $1,\ldots,10$ 中每个数写成二的幂乘奇数，排成两行。让学生亲自看见“删掉偶数”与“除去二的因子”得到不同乘积，然后推 $U_p(n)$ 的递归。
\paragraph{预设错误与回应.} $e\ge\alpha$ 可直接判零，但 $e<\alpha$ 时单位部分才决定剩余答案。单独的 $k!\bmod M=0$ 并不能说明组合数为零。若模数为一，整个计算都可提前结束。
\paragraph{离堂题与作业.} 离堂做 B9；基础作业做 B7、B8，提高作业做 P2183 和 P7488。要求附上一个最终余数为零但计数对象确实存在的例子。

\section{课堂包四：反演是为了分离变量}
\paragraph{学习目标.} 学生能将 $C(\gcd(i,j))$ 写成约数贡献和；能明确交换求和后每个因子的上界；能从反演过渡到前缀和递归。
\paragraph{情境.} 先画 $3\times3$ 的 gcd 矩阵并算总和十二。提出：“直接处理每个点，需要平方次；能否让每个公约数统一贡献？”
\paragraph{板书主线.}
\[
\gcd(i,j)=\sum_{d\mid i,d\mid j}\varphi(d)
\quad\Longrightarrow\quad
\sum_d\varphi(d)\floor{n/d}\floor{m/d}.
\]
再把 $\varphi$ 改为 $D=C*\mu$，得到一般加权公式。最后从 $\varphi*\one=\id$ 推导杜教筛，使同一卷积关系既用于拆点值，也用于拆前缀。
\paragraph{同伴活动.} 每组在一张纸上只写求和指标的条件，不写函数值：$1\le i\le n,1\le j\le m,d\mid i,d\mid j$。另一组把它改成先枚举 $d$ 的形式。这样可以暴露上界与整除条件的丢失，而不是让错误藏在长公式中。
\paragraph{预设错误与回应.} 把 gcd 拆成约数和后不能继续乘原 gcd；把 $i=da,j=db$ 后要同步修改两边上界；分块时两个商变化不同，必须取共同右端点。$S_\mu(n)$ 为零不代表它未被计算。
\paragraph{离堂题与作业.} 做 C2、C3；上机先写朴素 gcd 双循环和反演公式对拍，再实现记忆化杜教筛。提高学生可分析 P4240 的变长表项数。

\section{课堂包五：PN 与二维积性函数的“消主项”}
\paragraph{学习目标.} 理解为什么素数处匹配会让卷积商的支撑稀疏；理解二维坐标轴与对角线是两个不同层次的主要贡献；能正确叠加预处理和查询成本。
\paragraph{前置要求.} 必须已经熟练进行卷积的素数幂递推，知道形式幂级数常数项为一时可逆。若学生尚不能算出 C5 的三个局部系数，不宜直接讲 UOJ 885。
\paragraph{板书主线.}
\[
 f(p)=g(p)\Rightarrow h(p)=0\Rightarrow\text{PN 支撑};
\]
\[
 f(p^a,1)=g(p^a,1)\Rightarrow\text{去掉坐标轴};
\]
\[
 h=g_3*h'\Rightarrow h'(p^k,p^k)=0\Rightarrow\text{再去掉对角线}.
\]
\paragraph{引导练习.} 先做平方自由数计数，看到 $1-z^2$；再做 D4，看到 $1-XY$。这两个可完全手算的局部级数，是一般稀疏支撑论证的最好入口。
\paragraph{证明分层.} 基础提高班证明 PN 只有 $O(\sqrt n)$ 个；进一步证明同质因子对数为 $O(n)$；专题班再用收敛的局部乘积证明 $O((AB)^{1/3})$。不要同时把三层计数压缩成一句“显然很少”。
\paragraph{离堂题与作业.} 要求学生写出一个函数匹配素数处但 $h(p^3)\ne f(p^3)-g(p^3)$ 的例子，并指出总复杂度中哪项是前缀表预处理。研究作业为算法三中 $\rho<2/3$ 的收敛阈值，不直接要求完整最优实现。

\section{课堂包六：概率算法的答案为什么可信？}
\paragraph{学习目标.} 区分单边错误与零错误重试；知道概率的随机对象是底数而不是输入；能通过非平凡平方根解释判素和分解的联系。
\paragraph{情境.} 先用 $561$ 展示费马测试的失败，再展示 $263,166,67,1$ 的平方链。让学生自己指出 $67$ 为什么在素数模数下不可能出现。
\paragraph{板书对照.}
\begin{center}
\begin{tabular}{lll}
\toprule
工具 & 成功得到什么 & 失败时意味着什么\\\midrule
Miller--Rabin 见证 & 确定合数 & 未见证不自动证明为素数\\
Pollard--Rho 真因子 & 可直接验算的因子 & 更换参数重试\\
已知群阶倍数分裂 & 非平凡平方根 & 若全为 $\pm1$，重新随机\\\bottomrule
\end{tabular}
\end{center}
\paragraph{课堂实验.} 固定一个合数，枚举所有小底数，统计哪些会通过；再换合数，观察比例。必须先固定输入再讨论底数比例。让学生将输出因子乘回原数，理解可验证结果与概率判定结果的不同。
\paragraph{预设错误与回应.} “做了十次就绝对正确”需要明确随机误差；“这些七个底数适用任意大整数”要指出 $64$ 位范围；“得到 gcd 等于 $n$ 就递归 $n$”会无限循环；“发现 $-1$ 就得到真因子”用 gcd 立即反驳。
\paragraph{离堂题与作业.} 做 D1、D5。上机只对受控大小整数实现判素与分解，记录固定种子以便复现。进阶阅读 QOJ 1860 时，要求专门证明条件于 $a^T=1$ 的子群分布，而不是偷用 $\varphi(m)\mid T$。

\section{一次完整作业的评分框架}
\begin{longtable}{p{0.19\textwidth}p{0.12\textwidth}p{0.58\textwidth}}
\toprule
项目 & 建议分值 & 应检查的内容\\\midrule
建模与条件 & $20$ & 变量范围、互质／素数条件、是否允许零指数、对象是否有序。\\
正确性推导 & $30$ & 必要性与充分性、无漏解、无重复计数、不变量、递归终止。\\
算法与实现 & $20$ & 数据结构与推导一致，表的合法索引，模数一与负余数。\\
复杂度 & $15$ & 分解、预处理、每问及空间分开写；随机与最坏界不混淆。\\
验证与反例 & $15$ & 独立朴素程序、小范围穷举、边界、输出回代、失败案例。\\\bottomrule
\end{longtable}
建议允许学生用一个准确的反例修复先前不成立的结论，并对此给分。教师版的目标是建立可审查的推理，不是让学生把所有长公式背到最后。

\section{上机作业的最低测试清单}
\begin{enumerate}
\item \textbf{exgcd/CRT：} 不相容模数、一个模数整除另一个、余数为负、模数一、只有下界不限制余数、最小公倍数超出类型。
\item \textbf{离散对数：} 最小解为零、底数零、目标零、gcd 不为一、剥离后模数一、同余类重复导致的非最小命中。
\item \textbf{组合数：} $k=0,n$，$k>n$，$n=0$，$M=1$，模八的单位块乘积为一，赋值刚好等于模数指数。
\item \textbf{函数求和：} $n=1$，$S_\mu(n)=0$，两个商区间错位，多询问缓存，矩形两边差距极大。
\item \textbf{判素分解：} $0,1,2$，偶数、素数幂、Carmichael 数、强伪素数、接近 $2^{64}$ 的数、批量 gcd 等于 $n$ 的回退。
\item \textbf{构造：} 每组回代、组内大小范围、跨组全局去重、所有允许参数的规模验证。
\end{enumerate}

% \section{本次实际运行的检查}
% \begin{itemize}
% \item C++ 自测共执行 $1\,781\,581$ 次检查，覆盖带符号除法、逆元、最小离散对数、CRT、溢出拒绝、素数幂阶乘、组合数、判素、分解与求阶；启用了未定义行为检查器。
% \item 小范围 exgcd 正解区间与 P8255 构造实例共 $32\,165$ 项；LTE 有 $1088$ 项；P4240 覆盖 $40\times40$ 个矩形。
% \item 两阶段筛的 $720$ 个输入逐一核对了全部商集状态；二维卷积的两次分解用十二组系数与直接双循环对照。
% \item CF516E 与逐日模拟比较了 $14\,400$ 个小输入；AGC031F 的 $180$ 张小图对全部端点与余数询问比较完整状态图。
% \item WC2020 的 $62$ 个案例穷举子集覆盖，含两个原题样例；P7488 的递推与闭式在 $n\le12$ 上核对。
% \item P1400 覆盖 $k=2,\ldots,1000$，每个参数构造 $1000$ 组，验算方程、位数与全部 $3000$ 个数互异；逐参数 CSV 随附。
% \end{itemize}
% 这些记录说明具体实现和有限算例通过了检查，\textbf{不表示所有例题都已经在原评测平台提交满分，也不能替代无限范围的数学证明}。原题完整输入输出、极限常数、平台编译器与数据范围仍应在正式提交前核对。

% \chapter{主要勘误与严谨性补充}
% 下表记录合并时的实质性修正，既包括错误，也包括缺失条件与复杂度表述不完整。不是所有列出的原结论都错误：例如 Miller--Rabin 的 $1/2$ 界是可用弱界，但原证明需要重写；经验性构造也应与一般定理区分。

% \begingroup\small
% \begin{longtable}{p{0.055\textwidth}p{0.20\textwidth}p{0.645\textwidth}}
% \toprule
% 编号 & 位置 & 正确表述／本次处理\\\midrule
% \endfirsthead
% \toprule 编号 & 位置 & 正确表述／本次处理\\\midrule
% \endhead
% 1 & 第一稿：难度定位 & 基础整除可入门，但原根、二维积性函数及随机分解不能定位为小学二年级内容；改为按前置知识分层。\\
% 2 & 第一稿：P5656 & 整数解存在当且仅当 $\gcd(a,b)\mid c$，不是 $c\mid\gcd(a,b)$。\\
% 3 & 第一稿：P5656 & 从 $ax+by=d$ 的特解转到右边为 $c$ 时乘 $c/d$，不是 $d/c$。\\
% 4 & 第一稿：P5656 & 正整数解还需参数区间；有整数解不等于有正整数解；两个单独最小正值不一定能同时取得。\\
% 5 & 第一稿：P8255 & 应为 $d^2=\gcd(z/x,x^2)$，不是 $\gcd(z/x,a^2)$；还需 $x\mid z$ 与完全平方判定。\\
% 6 & 第一稿：欧拉展开 & $\gcd(i,n)$ 展开后约数项应为 $\varphi(d)$；保留原 gcd 会重复计数。\\
% 7 & 第一稿：CRT & 构造应为 $\sum_i a_i(M/m_i)((M/m_i)^{-1}\bmod m_i)$，不能把 $M/m_i$ 误写为 $m_i$。\\
% 8 & 两稿：CRT & 标准构造需模数两两互质；唯一性是模总乘积的唯一同余类。非互质时先检查相容条件，再取最小公倍数。\\
% 9 & 第一稿：屠龙勇士 & 选剑是“不高于生命值”的最大值，即 $\le$，不是严格小于。\\
% 10 & 第一稿：屠龙勇士 & 除同余式外还需攻击伤害下界 $x\ge\ceil{a_i/c_i}$；$a_i\le p_i$ 不是所有测试点的统一条件。\\
% 11 & 第一稿：LTE & $n=2^e u$ 去掉奇数部分后应为 $x^{2^e}-y^{2^e}$，不是 $x^{2^n}-y^{2^n}$。\\
% 12 & 第一稿：高阶箭号 & $3\uparrow^33$ 是高度 $3^{27}$ 的三的幂塔，不等于仅四层的 $3^{3^{27}}$。有限枚举需证明稳定高度。\\
% 13 & 第二稿：分配律 & $(a+b)c\equiv ac+bc$，不是 $ac+ab$。\\
% 14 & 第二稿：同余除法 & 同除公共因子后模数同步变为 $m/d$；等式中的商仍为原整数 $k$，不要求 $k/d$ 是整数。\\
% 15 & 第二稿：欧拉定理 & 模数必须为正，通常在 $m\ge2$ 讨论；“非负整数模数”会错误包含零。\\
% 16 & 第二稿：裴蜀证明 & 用“整数线性组合构成的子群”代替含糊的实向量空间说法，避免张成系数域不同导致误解。\\
% 17 & 第二稿：根数界 & 因式定理与非零根差的可逆性是在 $\mathbb F_p$ 中使用；合数模数可有超出次数的根数。\\
% 18 & 两稿：Wilson & 一般单位乘积中的符号按群结构分类；模二时 $1$ 与 $-1$ 是相同余数。\\
% 19 & 两稿：去因子阶乘 & $U_p(n)$ 是对每个数除尽素数因子后再相乘，不是直接删除所有该素数的倍数。\\
% 20 & 第二稿：扩展 Lucas & 查询成本含 $\log_p n$ 层递归与逆元成本，不能仅写成模数的函数；预处理要另计模数分解。\\
% 21 & 第一稿：P2183 & 不可行时原题输出 Impossible；可行方案数取模为零时仍应输出零，二者不等价。\\
% 22 & 第一稿：P7488 & 腰长彼此不同、底长彼此不同；不是全部腰长与底长数值整体两两不同，等边三角形允许出现。\\
% 23 & 第二稿：BSGS & 一般互质模数的指数上界来自欧拉定理，不来自只适用于素数的费马小定理。\\
% 24 & 第二稿：多询问 BSGS & 表可直接复用需要底数与模数都固定；仅固定模数不充分。\\
% 25 & 第二稿：BSGS 最小解 & 需规定重复余数存储策略与块枚举顺序；随便返回一次碰撞只保证一组解，不保证最小。\\
% 26 & 第二稿：exBSGS & 每层先检查零指数候选，除 gcd 前需 $d\mid b$；模数一、底数零要单独处理。\\
% 27 & 第二稿：原根判定 & 枚举的是 $\varphi(m)$ 的不同素因子，且须先保证 $g$ 为单位；不是任意整数因子。\\
% 28 & 第二稿：求阶复杂度 & 计算／分解 $\varphi(m)$ 是预处理成本，不能藏在每问快速幂成本中；求最小正阶不能直接接受指数零。\\
% 29 & 第二稿：CF516E & 同一循环只保留最小编号源的支配结论不成立；正文给反例并改成带权多源环扫线。\\
% 30 & 第二稿：CF571E & 参数是非负整数，不是“正整数且可为零”；各素因子的指数必须共享同一参数。\\
% 31 & 第二稿：WC2020 & 给出固定子集的最小覆盖询问数及其期望计数式，避免与在线未知答案的策略期望混淆。\\
% 32 & 第二稿：常见函数 & $\mu(p^k)=-[k=1]$；原稿已有误作 $I$ 的脚注，本次统一常一函数为 $\one$，约数个数为 $\tau$。\\
% 33 & 两稿：筛法与内存 & $O(N)$ 时间不等于小内存；$N=10^8$ 时多个 int 数组的空间可能不可接受，应按用途选择筛法。\\
% 34 & 第二稿：倍增估计 & 用端点代替一块各项是常数倍估计，应使用 $O,\Theta$ 或不等式，不是逐步精确相等。\\
% 35 & 第二稿：杜教筛 & 命中预处理范围应返回存储的前缀和，原稿的“返回 $\sqrt n$”是错误文字。补全阈值条件 $B\ge\sqrt n$。\\
% 36 & 第二稿：杜教筛 & $O(\sqrt n)$ 个状态，每个状态还有分块工作；未预处理总成本通常为 $O(n^{3/4})$，不是只数状态个数。\\
% 37 & 第二稿：PN 表示 & 不限制 $b$ 时 $a^2b^3$ 表示不唯一；要求 $b$ 平方自由可得唯一表示。计数上界不需要未经成立的双射。\\
% 38 & 第二稿：PN 总成本 & $O(\sqrt n)$ 是在相关前缀、局部系数与枚举准备完成后的支撑求和成本；完整算法还要加预处理。\\
% 39 & 第二稿：洲阁筛 & 第一阶段基函数必须能在非互质乘法中拆积，典型条件是完全积性；仅称积性不够。\\
% 40 & 第二稿：洲阁筛状态 & 明确排除整数一，原始前缀应减去一；素数幂枚举需有 $p^e\le v$ 上界；原地更新必须匹配旧层依赖。\\
% 41 & 第二稿：二维计数 & 从 $x,y\le n$ 放宽后应是 $xy\le n^2$，不是 $xy\le n$。\\
% 42 & 第二稿：二维支撑 & 对未出现的素数允许两边指数都为零；不能对所有素数都要求两边指数至少一。局部常数项必须为一。\\
% 43 & 第二稿：增长推广 & 应看第一个“非零”局部项而非“非负”项；本讲义明确归一化、非负整数系数与多项式增长条件。\\
% 44 & 第二稿：二维算法二 & 总成本必须加 $S,D$ 前缀预处理；算法三需要新的批量查询结构，不能把研究方向当成已完成实现。\\
% 45 & 第二稿：Carmichael & 费马性质针对 $\gcd(a,n)=1$ 的底数，而不是所有 $n\nmid a$ 的底数。\\
% 46 & 第二稿：Miller--Rabin & $1/2$ 可作为弱界；重新用局部循环群和 CRT 证明，并另给经典 $1/4$ 界。修正原证明中素数幂模数和指数的混写。\\
% 47 & 第二稿：一般分解 & 应说“没有已知的经典多项式时间算法”，不是断言不存在；复杂度变量是输入位数，并区别量子计算模型。\\
% 48 & 第二稿：Rho 两两比较 & $O(\sqrt p)$ 个状态若全部两两检查，会产生 $O(p)$ 次 gcd；判环才避免平方次检查。\\
% 49 & 第二稿：Rho 批处理 & $B\asymp\log n$ 时恢复项为 $O(\log^2n)$，不能不加条件略去；$n^{1/4}$ 是常用启发式量级而非无条件最坏界。\\
% 50 & 第二稿：$M=N\varphi(N)$ & $V=-1$ 时 $\gcd(V+1,m)=m$，必须找非平凡平方根。不能假设所有递归模数均满足 $\varphi(m)\mid M$，改用条件子群分析。\\
% 51 & 第二稿：韦达递降 & 另一根是整数但可能非正；应写清正整数解空间的停止条件。全局数字互异不同于三元组互异。\\
% 52 & 第二稿：P1400 范围 & 最多一百位指小于 $10^{100}$；经验性“128 位足够”不是普遍保证，平方验算还需更大范围。本次提供大整数实现与全参数有限验证。\\
% \bottomrule
% \end{longtable}
% \endgroup

% \section{排版与来源处理}
% 两稿统一为可用 XeLaTeX 编译的现代 UTF-8 中文文档，使用标准 $11$ 磅字号，避免不受支持的 $10.5$ 磅选项；统一定理、问题、解析、教学提示与警示样式。标题日期明确是本次整理日期，不冒充原讲稿日期；插图使用内嵌 TikZ，移除外部缺失图片依赖。


% \chapter{合并对应表、例题索引与延伸阅读}
% \section{知识点合并对应表}
% \begin{longtable}{p{0.34\textwidth}p{0.24\textwidth}p{0.31\textwidth}}
% \toprule
% 原稿内容 & 本书位置 & 整理方式\\\midrule
% 最大公约数、裴蜀、exgcd、步长循环 & 第 1 章 & 合并基本证明，补整数参数区间与边界。\\
% 同余、逆元、费马、欧拉、扩展欧拉 & 第 2 章 & 统一条件与记号，补截断指数接口与两种在线逆元方案。\\
% 中国剩余定理、扩展 CRT & 第 3 章 & 标准构造与逐式合并并列，补相容性、下界与溢出。\\
% 阶、原根、离散对数与高阶箭号 & 第 4 章 & 按周期到求指数重排，白泽教育归于此处而非 LTE。\\
% 第一稿的升幂引理 & 第 5 章 & 完整保留三类结论，重写证明，增加阶的提升例题。\\
% Wilson、阶乘、Legendre、Kummer、Lucas & 第 6 章 & 两稿合并，补单位阶乘与组合计数建模。\\
% 积性函数、卷积、线性筛、狄利克雷变换 & 第 7 章 & 统一 $\one,\mu,\tau$，区分完全积性与积性。\\
% 莫比乌斯／欧拉反演、gcd 求和 & 第 8 章 & 补通式、原题解析及多询问有效表存储。\\
% 倍增估计、块筛、杜教筛、PN & 第 9 章 & 补预处理／查询成本与局部支撑论证。\\
% 洲阁筛 & 第 10 章 & 写清两阶段状态、更新顺序与完全积性前提。\\
% 二维积性函数、UOJ 885 三类算法 & 第 11 章 & 补局部系数递推与更直接的支撑计数；最优改进列为研究拓展。\\
% 同余代数综合例题 & 第 12 章 & 补指数向量、环扫线、六状态图与期望计数解析。\\
% 素性判定、Miller--Rabin、Pollard--Rho & 第 13 章 & 区分概率界、确定性范围与启发式性能。\\
% 给定 $n,\varphi(n)$；构造 $n\varphi(n)$ & 第 14 章 & 补随机分裂的条件证明与选择性剪枝。\\
% 韦达定理、递降树、P1400 & 第 15 章 & 补正性边界、统一种子、全局去重与有限验证。\\
% 新增训练与教学资源 & 第 16 章及教师专章、附录 & 三十道分层练习、六个课堂包、模板与勘误。\\
% \bottomrule
% \end{longtable}

% \section{原有例题去重索引}
% \begin{longtable}{p{0.39\textwidth}p{0.12\textwidth}p{0.38\textwidth}}
% \toprule
% 题目 & 章节 & 教师版解析重点\\\midrule
% P5656 二元一次不定方程 & 1 & 通解、正解区间、无正解时的特殊输出。\\
% P8255 数学游戏 & 1 & 平方 gcd 的双向证明与解的唯一性。\\
% 五类乘法逆元问题 & 2 & 单次、线性、批量、减半递归与分块预处理。\\
% P4774 屠龙勇士 & 3 & 选剑、带系数同余、伤害下界、exCRT。\\
% P5605 小 A 与两位神仙 & 4 & 循环单位群、阶整除与非循环反例。\\
% P6736 白泽教育 & 4 & 幂塔截断、欧拉链稳定、三阶候选上界。\\
% 原稿两道赋值练习 & 5 & $7$-进差与二进偶指数 LTE。\\
% HDU 2973 YAPTCHA & 6 & 嵌套取整化为 Wilson 的判素示性函数。\\
% P2183 礼物 & 6 & 多项式系数与素数幂单位部分。\\
% P7488 黑色毛衣 & 6 & 阶梯棋盘递推与大参数合数取模。\\
% P2398 GCD SUM & 8 & 欧拉展开与交换求和。\\
% 一般加权 gcd 求和 & 8 & 任意数组的卷积差分与倍数后缀。\\
% P4240 毒瘤之神的考验 & 8 & $\mu^2/\varphi$、变长表与小／大参数分治。\\
% P4213 杜教筛 & 9 & 两个典型递归、记忆化与阈值分析。\\
% LOJ 6053 & 9 & 奇素数匹配 $\varphi$，单独处理素数二。\\
% P5325 & 10 & 素数贡献筛与按最小素因子补合数。\\
% UOJ 885 红场阅兵 & 11 & 坐标轴／对角线消去、支撑规模；算法三为研究拓展。\\
% CF516E & 12 & 带权多源环的前后缀最小值，修正原支配结论。\\
% CF571E & 12 & 两条指数向量射线的交集；大数据 Bonus 为研究拓展。\\
% AGC031F & 12 & 可逆状态、平移子群、每顶点六状态。\\
% P6730 WC2020 猜数游戏 & 12 & 阶分组、非单位短幂链、分量贡献期望。\\
% QOJ 4920 & 14 & 已知群阶倍数的随机非平凡平方根分裂。\\
% QOJ 1860 & 14 & 唯一性、选择性分解与条件子群分析。\\
% P1400 Easy Equation & 15 & 韦达跳跃、BFS、全局互异、大整数与全参数验证。\\
% \bottomrule
% \end{longtable}

% \section{来源与引用说明}
% \begin{enumerate}
% \item 原稿 A 为用户提供的《数论》文本；原稿 B 为用户提供、署名周雨衡、日期为 2025 年 9 月 26 日的同名讲义整理稿。本次不是对原作者原文的逐字复制，而是合并、重排、校订及新增教学阐释。
% \item 各原题的公开题面链接列在对应例题后，用于复查定义与输入输出。原稿提供的两个博客链接在本次访问时要求阅读密码，因此未把它们当成已阅读内容，也未尝试访问受保护正文；CF571E 与 AGC031F 的讲解根据数学推导及可公开核对的材料补充。
% \item AGC031F 的六状态路线可参考公开讲解：\url{https://www.luogu.com.cn/article/hz11gwzl}。本书补充了平移等价的条件、更新形式和商模数说明，并用完整小状态图对拍。
% \item UOJ 885 的延伸阅读为周康阳的专题文章：\url{https://www.cnblogs.com/zkyJuruo/p/18204106}。本书对算法一、二给出独立整理的卷积式与收敛支撑证明，算法三明确保留为专题研究方向。
% \item 组合数素数幂模数下更深入的算法可参考原稿链接：\url{https://faqak.github.io/misc/binom-mod-pe/}。本书采用适用于素数幂模数较小时的可解释预处理方案，不宣称已覆盖所有最优算法。
% \item 判素与分解可进一步阅读：\url{https://cp-algorithms.com/algebra/primality_tests.html} 与 \url{https://cp-algorithms.com/algebra/factorization.html}。固定底数保证、随机底数误差和 Rho 的启发式性能需分别理解。
% \end{enumerate}

% \section{后续修订建议}
% 对于新的题目或更大的数据范围，先重新写出函数接口、模数条件与资源上界，再复用本书的数学结构。尤其是多次询问、任意模数、高精度输入与最小解要求，经常改变一个本来正确模板的适用范围。任何新增推导都建议配一份独立朴素实现或最小反例，不把一次通过测试当成无条件证明。


% --------------------------------------------------------------------------
% 附录 D：来自《博弈论》教师版的“核心判定模板与接口”。
% --------------------------------------------------------------------------
\chapter{核心判定模板与接口}
\section{模板使用说明}
本附录摘录全书反复出现的判定框架，完整 C++17 实现见资料包 \texttt{code/selftest.cpp}；摘录是教学模板，不是可直接提交的完整程序。所有模板共用一个前提：游戏在有限步内结束、无法行动者负（三分类模板额外支持平局）。自测程序对有限范围的随机与穷举数据，把每个公式与暴力搜索逐点比对。

\begin{lstlisting}[language=C++]
// mex：集合中未出现的最小非负整数。
int mex(const std::vector<int>& s) {
    std::vector<char> seen(s.size() + 1, 0);
    for (int x : s)
        if (0 <= x && x < (int)seen.size()) seen[x] = 1;
    int m = 0;
    while (seen[m]) ++m;
    return m;
}
\end{lstlisting}

\begin{lstlisting}[language=C++]
// 模板一：DAG 上的记忆化 SG（隐式转移版）。
// 约定：move(u) 枚举 u 的全部后继；图无环（游戏有限步结束）。
// g(u) == 0 恰为 P 态（定理 5.x）。
struct SGTable {
    std::function<void(int, const std::function<void(int)>&)> move;
    std::vector<int> sg;
    std::vector<char> done;
    SGTable(int n, decltype(move) m) : move(m), sg(n, 0), done(n, 0) {}
    int dfs(int u) {
        if (done[u]) return sg[u];
        done[u] = 1;
        std::vector<int> nxt;
        move(u, [&](int v) { nxt.push_back(dfs(v)); });
        return sg[u] = mex(nxt);
    }
};
\end{lstlisting}

\begin{lstlisting}[language=C++]
// 模板二：一般状态图（允许有环）上的三分类。
// 返回 col：0 = P 态，1 = N 态，2 = 平局（过河卒框架）。
// u 是 N 态 <=> 存在 P 态后继；u 是 P 态 <=> 全部后继都是 N 态。
std::vector<int> solve_game(int n,
                            const std::vector<std::pair<int,int>>& edges) {
    std::vector<std::vector<int>> g(n), rg(n);
    for (auto [u, v] : edges) { g[u].push_back(v); rg[v].push_back(u); }
    std::vector<int> col(n, 2), und(n);
    std::vector<char> done(n, 0);
    std::queue<int> q;
    for (int u = 0; u < n; ++u) {
        und[u] = (int)g[u].size();
        if (und[u] == 0) { col[u] = 0; done[u] = 1; q.push(u); }
    }
    while (!q.empty()) {
        int u = q.front(); q.pop();
        for (int p : rg[u]) {
            if (done[p]) continue;
            if (col[u] == 0) { col[p] = 1; done[p] = 1; q.push(p); }
            else if (--und[p] == 0) { col[p] = 0; done[p] = 1; q.push(p); }
        }
    }
    return col;  // 仍为 2 的点：双方都不离开未定子图 => 平局。
}
\end{lstlisting}

\begin{lstlisting}[language=C++]
// 模板三：阶梯 Nim 判定（定理 2.x）。
// 输入 a[1..n]：第 i 级台阶上的石子数；a[0] 弃用（地面）。
// 每次把第 i 级的至少一颗石子移到第 i-1 级，不能行动者负。
// 先手必胜 <=> 奇数级异或和不为 0（只有奇数级参与！）
bool stair_nim_wins(const std::vector<long long>& a) {
    long long s = 0;
    for (int i = 1; i < (int)a.size(); i += 2) s ^= a[i];
    return s != 0;
}
\end{lstlisting}

Nim、反 Nim、Nim-k 与反 Nim-k 的判定公式各只有几行，不再单独摘录：它们在自测程序中以函数形式出现，并直接与暴力搜索对拍。

\begin{teach}[备课验证]
例题 [CF98E] 与 [THUPC 2023 Pre] 的 $2\times2$ 矩阵建议课前各解一遍数值，确认交点在 $0$ 到 $1$ 之间；本章练习 1--6 的数值答案（$p^{*}=\tfrac25,v=\tfrac15$；$(\tfrac14,\tfrac12,\tfrac14)$；$p^{*}=\tfrac23,q^{*}=\tfrac12$；闭式公式；$v=2$；$p^{*}=0,v=1$；扑克 $a=1$：$c=\tfrac13,p=\tfrac23,v=-\tfrac19$，$a=\tfrac12$：$\tfrac15,\tfrac45,-\tfrac1{25}$）应课前亲算一遍。附录自测程序应在课前编译运行一次，正常输出为一万二千余项检查全部通过。若你修改了任何判定模板，先跑自测再进课堂。
\end{teach}
\section{怎样复现验证}
在资料目录中运行：
\begin{lstlisting}[language=bash]
g++ -std=c++17 -O2 -Wall -Wextra code/selftest.cpp -o gt_selftest
./gt_selftest
\end{lstlisting}
自测程序完成五组有限范围验证：(1) “每次取 $1$ 或 $2$ 颗”的 SG 值周期 $3$；(2) Nim、反 Nim、Nim-k（$k=1,2$）与反 Nim-k（$k=1,2$）的判定公式对全部 $4$ 堆以内、每堆不超过 $5$ 颗的局面与暴力博弈搜索一致；(3) 三分类模板与不动点迭代在随机带环图上一致；(4) 阶梯 Nim 判定与暴力搜索一致；(5) 记忆化 SG 与显式转移表一致。暴力验证依赖的是“状态数小到可以穷举”这一事实，与正文证明互为印证而非替代；更大的参数范围仍以正文证明为准。


\begin{tip}[代码与验证脚本的位置]
本附录摘录的判定模板，完整 C++17 实现在资料包的 \texttt{Game Theory/code/selftest.cpp}。
自测程序对有限范围的随机与穷举数据，把每个公式与暴力搜索逐点比对；
它相对本书源码目录的上一级，编译本书本身不需要它。
\end{tip}


\end{document}
